\documentclass{article}
\usepackage{amsmath, amssymb, amsthm}
\usepackage{hyperref}
\usepackage{geometry}
\usepackage{titlesec}
\geometry{margin=1in}

\title{Erd\H{o}s Problems --- Open (Unsolved)\\[0.5em]\large A Collection of 622 Problems}
\date{Generated: 2026-06-09}
\author{Source: \url{https://www.erdosproblems.com}}

\begin{document}
\maketitle
\tableofcontents
\newpage


%% ============================================================
\section{Problem \#1}
\label{prob:1}

\noindent\textbf{Status:} OPEN \quad|\quad \textbf{Prize:} \$500 \quad|\quad \textbf{Updated:} 2025-08-31

\noindent\textbf{Tags:} number theory, additive combinatorics

\medskip

\noindent\textbf{Statement:}

If $A\subseteq \{1,\ldots,N\}$ with $\lvert A\rvert=n$ is such that the subset sums $\sum_{a\in S}a$ are distinct for all $S\subseteq A$ then\[N \gg 2^{n}.\]

\noindent\rule{\linewidth}{0.4pt}


%% ============================================================
\section{Problem \#3}
\label{prob:3}

\noindent\textbf{Status:} OPEN \quad|\quad \textbf{Prize:} \$5000 \quad|\quad \textbf{Updated:} 2025-08-31

\noindent\textbf{Tags:} number theory, additive combinatorics, arithmetic progressions

\medskip

\noindent\textbf{Statement:}

If $A\subseteq \mathbb{N}$ has $\sum_{n\in A}\frac{1}{n}=\infty$ then must $A$ contain arbitrarily long arithmetic progressions?

\noindent\rule{\linewidth}{0.4pt}


%% ============================================================
\section{Problem \#5}
\label{prob:5}

\noindent\textbf{Status:} OPEN \quad|\quad \textbf{Prize:} none \quad|\quad \textbf{Updated:} 2025-08-31

\noindent\textbf{Tags:} number theory, primes

\medskip

\noindent\textbf{Statement:}

Let $C\geq 0$. Is there an infinite sequence of $n_i$ such that\[\lim_{i\to \infty}\frac{p_{n_i+1}-p_{n_i}}{\log n_i}=C?\]

\noindent\rule{\linewidth}{0.4pt}


%% ============================================================
\section{Problem \#9}
\label{prob:9}

\noindent\textbf{Status:} OPEN \quad|\quad \textbf{Prize:} none \quad|\quad \textbf{Updated:} 2025-08-31

\noindent\textbf{Tags:} number theory, additive basis, primes

\medskip

\noindent\textbf{Statement:}

Let $A$ be the set of all odd integers $\geq 1$ not of the form $p+2^{k}+2^l$ (where $k,l\geq 0$ and $p$ is prime). Is the upper density of $A$ positive?

\noindent\rule{\linewidth}{0.4pt}


%% ============================================================
\section{Problem \#10}
\label{prob:10}

\noindent\textbf{Status:} OPEN \quad|\quad \textbf{Prize:} none \quad|\quad \textbf{Updated:} 2025-08-31

\noindent\textbf{Tags:} number theory, additive basis, primes

\medskip

\noindent\textbf{Statement:}

Is there some $k$ such that every large integer is the sum of a prime and at most $k$ powers of 2?

\noindent\rule{\linewidth}{0.4pt}


%% ============================================================
\section{Problem \#11}
\label{prob:11}

\noindent\textbf{Status:} OPEN \quad|\quad \textbf{Prize:} none \quad|\quad \textbf{Updated:} 2026-03-14

\noindent\textbf{Tags:} number theory, additive basis

\medskip

\noindent\textbf{Statement:}

Is every large odd integer $n$ the sum of a squarefree number and a power of 2?

\noindent\rule{\linewidth}{0.4pt}


%% ============================================================
\section{Problem \#12}
\label{prob:12}

\noindent\textbf{Status:} OPEN \quad|\quad \textbf{Prize:} none \quad|\quad \textbf{Updated:} 2025-08-31

\noindent\textbf{Tags:} number theory

\medskip

\noindent\textbf{Statement:}

Let $A$ be an infinite set such that there are no distinct $a,b,c\in A$ such that $a\mid (b+c)$ and $b,c&#62;a$. Is there such an $A$ with\[\liminf \frac{\lvert A\cap\{1,\ldots,N\}\rvert}{N^{1/2}}&#62;0?\]Does there exist some absolute constant $c&#62;0$ such that there are always infinitely many $N$ with\[\lvert A\cap\{1,\ldots,N\}\rvert&#60;N^{1-c}?\]Is it true that\[\sum_{n\in A}\frac{1}{n}&#60;\infty?\]

\noindent\rule{\linewidth}{0.4pt}


%% ============================================================
\section{Problem \#14}
\label{prob:14}

\noindent\textbf{Status:} OPEN \quad|\quad \textbf{Prize:} none \quad|\quad \textbf{Updated:} 2025-08-31

\noindent\textbf{Tags:} number theory, sidon sets, additive combinatorics

\medskip

\noindent\textbf{Statement:}

Let $A\subseteq \mathbb{N}$. Let $B\subseteq \mathbb{N}$ be the set of integers which are representable in exactly one way as the sum of two elements from $A$. Is it true that for all $\epsilon&#62;0$ and large $N$\[\lvert \{1,\ldots,N\}\backslash B\rvert \gg_\epsilon N^{1/2-\epsilon}?\]Is it possible that\[\lvert \{1,\ldots,N\}\backslash B\rvert =o(N^{1/2})?\]

\noindent\rule{\linewidth}{0.4pt}


%% ============================================================
\section{Problem \#15}
\label{prob:15}

\noindent\textbf{Status:} OPEN \quad|\quad \textbf{Prize:} none \quad|\quad \textbf{Updated:} 2025-08-31

\noindent\textbf{Tags:} number theory, primes

\medskip

\noindent\textbf{Statement:}

Is it true that\[\sum_{n=1}^\infty(-1)^n\frac{n}{p_n}\]converges, where $p_n$ is the sequence of primes?

\noindent\rule{\linewidth}{0.4pt}


%% ============================================================
\section{Problem \#17}
\label{prob:17}

\noindent\textbf{Status:} OPEN \quad|\quad \textbf{Prize:} none \quad|\quad \textbf{Updated:} 2025-08-31

\noindent\textbf{Tags:} number theory, primes

\medskip
\noindent\textit{cluster primes}


\noindent\textbf{Statement:}

Are there infinitely many primes $p$ such that every even number $n\leq p-3$ can be written as a difference of primes $n=q_1-q_2$ where $q_1,q_2\leq p$?

\noindent\rule{\linewidth}{0.4pt}


%% ============================================================
\section{Problem \#18}
\label{prob:18}

\noindent\textbf{Status:} OPEN \quad|\quad \textbf{Prize:} none \quad|\quad \textbf{Updated:} 2025-08-31

\noindent\textbf{Tags:} number theory, divisors, factorials

\medskip
\noindent\textit{practical numbers}


\noindent\textbf{Statement:}

We call $m$ practical if every integer $1\leq n&#60;m$ is the sum of distinct divisors of $m$. If $m$ is practical then let $h(m)$ be such that $h(m)$ many divisors always suffice. Are there infinitely many practical $m$ such that\[h(m) &#60; (\log\log m)^{O(1)}?\]Is it true that $h(n!)&#60;n^{o(1)}$? Or perhaps even $h(n!)&#60;(\log n)^{O(1)}$?

\noindent\rule{\linewidth}{0.4pt}


%% ============================================================
\section{Problem \#20}
\label{prob:20}

\noindent\textbf{Status:} OPEN \quad|\quad \textbf{Prize:} \$1000 \quad|\quad \textbf{Updated:} 2025-08-31

\noindent\textbf{Tags:} combinatorics

\medskip
\noindent\textit{sunflower conjecture}


\noindent\textbf{Statement:}

Let $f(n,k)$ be minimal such that every family $\mathcal{F}$ of $n$-uniform sets with $\lvert \mathcal{F}\rvert \geq f(n,k)$ contains a $k$-sunflower. Is it true that\[f(n,k) &lt; c_k^n\]for some constant $c_k&gt;0$?

\noindent\rule{\linewidth}{0.4pt}


%% ============================================================
\section{Problem \#25}
\label{prob:25}

\noindent\textbf{Status:} OPEN \quad|\quad \textbf{Prize:} none \quad|\quad \textbf{Updated:} 2025-08-31

\noindent\textbf{Tags:} number theory

\medskip

\noindent\textbf{Statement:}

Let $1\leq n_1&#60;n_2&#60;\cdots$ be an arbitrary sequence of integers, each with an associated residue class $a_i\pmod{n_i}$. Let $A$ be the set of integers $n$ such that for every $i$ either $n&#60;n_i$ or $n\not\equiv a_i\pmod{n_i}$. Must the logarithmic density of $A$ exist?

\noindent\rule{\linewidth}{0.4pt}


%% ============================================================
\section{Problem \#28}
\label{prob:28}

\noindent\textbf{Status:} OPEN \quad|\quad \textbf{Prize:} \$500 \quad|\quad \textbf{Updated:} 2025-08-31

\noindent\textbf{Tags:} number theory, additive basis

\medskip

\noindent\textbf{Statement:}

If $A\subseteq \mathbb{N}$ is such that $A+A$ contains all but finitely many integers then $\limsup 1_A\ast 1_A(n)=\infty$.

\noindent\rule{\linewidth}{0.4pt}


%% ============================================================
\section{Problem \#30}
\label{prob:30}

\noindent\textbf{Status:} OPEN \quad|\quad \textbf{Prize:} \$1000 \quad|\quad \textbf{Updated:} 2025-08-31

\noindent\textbf{Tags:} number theory, sidon sets, additive combinatorics

\medskip

\noindent\textbf{Statement:}

Let $h(N)$ be the maximum size of a Sidon set in $\{1,\ldots,N\}$. Is it true that, for every $\epsilon&#62;0$,\[h(N) = N^{1/2}+O_\epsilon(N^\epsilon)?\]

\noindent\rule{\linewidth}{0.4pt}


%% ============================================================
\section{Problem \#32}
\label{prob:32}

\noindent\textbf{Status:} OPEN \quad|\quad \textbf{Prize:} none \quad|\quad \textbf{Updated:} 2025-08-31

\noindent\textbf{Tags:} number theory, additive basis

\medskip

\noindent\textbf{Statement:}

Is there a set $A\subset\mathbb{N}$ such that\[\lvert A\cap\{1,\ldots,N\}\rvert = o((\log N)^2)\]and such that every large integer can be written as $p+a$ for some prime $p$ and $a\in A$? Can the bound $O(\log N)$ be achieved? Must such an $A$ satisfy\[\liminf \frac{\lvert A\cap\{1,\ldots,N\}\rvert}{\log N}&#62; 1?\]

\noindent\rule{\linewidth}{0.4pt}


%% ============================================================
\section{Problem \#33}
\label{prob:33}

\noindent\textbf{Status:} OPEN \quad|\quad \textbf{Prize:} none \quad|\quad \textbf{Updated:} 2025-08-31

\noindent\textbf{Tags:} number theory, additive basis

\medskip

\noindent\textbf{Statement:}

Let $A\subset\mathbb{N}$ be such that every large integer can be written as $n^2+a$ for some $a\in A$ and $n\geq 0$. What is the smallest possible value of\[\limsup \frac{\lvert A\cap\{1,\ldots,N\}\rvert}{N^{1/2}}?\]Is\[\liminf \frac{\lvert A\cap\{1,\ldots,N\}\rvert}{N^{1/2}}&#62;1?\]

\noindent\rule{\linewidth}{0.4pt}


%% ============================================================
\section{Problem \#36}
\label{prob:36}

\noindent\textbf{Status:} OPEN \quad|\quad \textbf{Prize:} none \quad|\quad \textbf{Updated:} 2025-08-31

\noindent\textbf{Tags:} number theory, additive combinatorics

\medskip
\noindent\textit{minimum overlap problem}


\noindent\textbf{Statement:}

Find the optimal constant $c&#62;0$ such that the following holds. For all sufficiently large $N$, if $A\sqcup B=\{1,\ldots,2N\}$ is a partition into two equal parts, so that $\lvert A\rvert=\lvert B\rvert=N$, then there is some $x$ such that the number of solutions to $a-b=x$ with $a\in A$ and $b\in B$ is at least $cN$.

\noindent\rule{\linewidth}{0.4pt}


%% ============================================================
\section{Problem \#39}
\label{prob:39}

\noindent\textbf{Status:} OPEN \quad|\quad \textbf{Prize:} \$500 \quad|\quad \textbf{Updated:} 2025-08-31

\noindent\textbf{Tags:} number theory, sidon sets, additive combinatorics

\medskip

\noindent\textbf{Statement:}

Is there an infinite Sidon set $A\subset \mathbb{N}$ such that\[\lvert A\cap \{1\ldots,N\}\rvert \gg_\epsilon N^{1/2-\epsilon}\]for all $\epsilon&#62;0$?

\noindent\rule{\linewidth}{0.4pt}


%% ============================================================
\section{Problem \#40}
\label{prob:40}

\noindent\textbf{Status:} OPEN \quad|\quad \textbf{Prize:} \$500 \quad|\quad \textbf{Updated:} 2025-08-31

\noindent\textbf{Tags:} number theory, additive basis

\medskip

\noindent\textbf{Statement:}

For what functions $g(N)\to \infty$ is it true that\[\lvert A\cap \{1,\ldots,N\}\rvert \gg \frac{N^{1/2}}{g(N)}\]implies $\limsup 1_A\ast 1_A(n)=\infty$?

\noindent\rule{\linewidth}{0.4pt}


%% ============================================================
\section{Problem \#41}
\label{prob:41}

\noindent\textbf{Status:} OPEN \quad|\quad \textbf{Prize:} \$500 \quad|\quad \textbf{Updated:} 2025-08-31

\noindent\textbf{Tags:} number theory, sidon sets, additive combinatorics

\medskip

\noindent\textbf{Statement:}

Let $A\subset\mathbb{N}$ be an infinite set such that the triple sums $a+b+c$ are all distinct for $a,b,c\in A$ (aside from the trivial coincidences). Is it true that\[\liminf \frac{\lvert A\cap \{1,\ldots,N\}\rvert}{N^{1/3}}=0?\]

\noindent\rule{\linewidth}{0.4pt}


%% ============================================================
\section{Problem \#44}
\label{prob:44}

\noindent\textbf{Status:} OPEN \quad|\quad \textbf{Prize:} none \quad|\quad \textbf{Updated:} 2025-08-31

\noindent\textbf{Tags:} number theory, sidon sets, additive combinatorics

\medskip

\noindent\textbf{Statement:}

Let $N\geq 1$ and $A\subset \{1,\ldots,N\}$ be a Sidon set. Is it true that, for any $\epsilon&#62;0$, there exist $M$ and $B\subset \{N+1,\ldots,M\}$ (which may depend on $N,A,\epsilon$) such that $A\cup B\subset \{1,\ldots,M\}$ is a Sidon set of size at least $(1-\epsilon)M^{1/2}$?

\noindent\rule{\linewidth}{0.4pt}


%% ============================================================
\section{Problem \#50}
\label{prob:50}

\noindent\textbf{Status:} OPEN \quad|\quad \textbf{Prize:} \$250 \quad|\quad \textbf{Updated:} 2025-08-31

\noindent\textbf{Tags:} number theory

\medskip

\noindent\textbf{Statement:}

Schoenberg proved that for every $c\in [0,1]$ the density of\[\{ n\in \mathbb{N} : \phi(n)&#60;cn\}\]exists. Let this density be denoted by $f(c)$. Is it true that there are no $x$ such that $f'(x)$ exists and is positive?

\noindent\rule{\linewidth}{0.4pt}


%% ============================================================
\section{Problem \#51}
\label{prob:51}

\noindent\textbf{Status:} OPEN \quad|\quad \textbf{Prize:} none \quad|\quad \textbf{Updated:} 2025-08-31

\noindent\textbf{Tags:} number theory

\medskip

\noindent\textbf{Statement:}

Is there an infinite set $A\subset \mathbb{N}$ such that for every $a\in A$ there is an integer $n$ such that $\phi(n)=a$, and yet if $n_a$ is the smallest such integer then $n_a/a\to \infty$ as $a\to\infty$?

\noindent\rule{\linewidth}{0.4pt}


%% ============================================================
\section{Problem \#52}
\label{prob:52}

\noindent\textbf{Status:} OPEN \quad|\quad \textbf{Prize:} \$250 \quad|\quad \textbf{Updated:} 2026-05-28

\noindent\textbf{Tags:} number theory, additive combinatorics

\medskip
\noindent\textit{sum-product problem}


\noindent\textbf{Statement:}

Let $A$ be a finite set of integers. Is it true that for every $\epsilon&#62;0$\[\max( \lvert A+A\rvert,\lvert AA\rvert)\gg_\epsilon \lvert A\rvert^{2-\epsilon}?\]

\noindent\rule{\linewidth}{0.4pt}


%% ============================================================
\section{Problem \#60}
\label{prob:60}

\noindent\textbf{Status:} OPEN \quad|\quad \textbf{Prize:} none \quad|\quad \textbf{Updated:} 2025-08-31

\noindent\textbf{Tags:} graph theory, cycles

\medskip

\noindent\textbf{Statement:}

Does every graph on $n$ vertices with $&#62;\mathrm{ex}(n;C_4)$ edges contain $\gg n^{1/2}$ many copies of $C_4$?

\noindent\rule{\linewidth}{0.4pt}


%% ============================================================
\section{Problem \#61}
\label{prob:61}

\noindent\textbf{Status:} OPEN \quad|\quad \textbf{Prize:} none \quad|\quad \textbf{Updated:} 2025-08-31

\noindent\textbf{Tags:} graph theory

\medskip

\noindent\textbf{Statement:}

For any graph $H$ is there some $c=c(H)&#62;0$ such that every graph $G$ on $n$ vertices that does not contain $H$ as an induced subgraph contains either a complete graph or independent set on $\geq n^c$ vertices?

\noindent\rule{\linewidth}{0.4pt}


%% ============================================================
\section{Problem \#62}
\label{prob:62}

\noindent\textbf{Status:} OPEN \quad|\quad \textbf{Prize:} none \quad|\quad \textbf{Updated:} 2025-08-31

\noindent\textbf{Tags:} graph theory

\medskip

\noindent\textbf{Statement:}

If $G_1,G_2$ are two graphs with chromatic number $\aleph_1$ then must there exist a graph $G$ whose chromatic number is $4$ (or even $\aleph_0$) which is a subgraph of both $G_1$ and $G_2$?

\noindent\rule{\linewidth}{0.4pt}


%% ============================================================
\section{Problem \#65}
\label{prob:65}

\noindent\textbf{Status:} OPEN \quad|\quad \textbf{Prize:} none \quad|\quad \textbf{Updated:} 2025-08-31

\noindent\textbf{Tags:} graph theory, cycles

\medskip

\noindent\textbf{Statement:}

Let $G$ be a graph with $n$ vertices and $kn$ edges, and $a_1&#60;a_2&#60;\cdots $ be the lengths of cycles in $G$. Is it true that\[\sum\frac{1}{a_i}\gg \log k?\]Is the sum $\sum\frac{1}{a_i}$ minimised when $G$ is a complete bipartite graph?

\noindent\rule{\linewidth}{0.4pt}


%% ============================================================
\section{Problem \#66}
\label{prob:66}

\noindent\textbf{Status:} OPEN \quad|\quad \textbf{Prize:} \$500 \quad|\quad \textbf{Updated:} 2025-08-31

\noindent\textbf{Tags:} number theory, additive basis

\medskip

\noindent\textbf{Statement:}

Is there $A\subseteq \mathbb{N}$ such that\[\lim_{n\to \infty}\frac{1_A\ast 1_A(n)}{\log n}\]exists and is $\neq 0$?

\noindent\rule{\linewidth}{0.4pt}


%% ============================================================
\section{Problem \#68}
\label{prob:68}

\noindent\textbf{Status:} OPEN \quad|\quad \textbf{Prize:} none \quad|\quad \textbf{Updated:} 2025-08-31

\noindent\textbf{Tags:} number theory, irrationality

\medskip

\noindent\textbf{Statement:}

Is\[\sum_{n\geq 2}\frac{1}{n!-1}\]irrational?

\noindent\rule{\linewidth}{0.4pt}


%% ============================================================
\section{Problem \#70}
\label{prob:70}

\noindent\textbf{Status:} OPEN \quad|\quad \textbf{Prize:} none \quad|\quad \textbf{Updated:} 2025-08-31

\noindent\textbf{Tags:} graph theory, ramsey theory, set theory

\medskip

\noindent\textbf{Statement:}

Let $\mathfrak{c}$ be the ordinal of the real numbers, $\beta$ be any countable ordinal, and $2\leq n&#60;\omega$. Is it true that $\mathfrak{c}\to (\beta, n)_2^3$?

\noindent\rule{\linewidth}{0.4pt}


%% ============================================================
\section{Problem \#74}
\label{prob:74}

\noindent\textbf{Status:} OPEN \quad|\quad \textbf{Prize:} \$500 \quad|\quad \textbf{Updated:} 2025-08-31

\noindent\textbf{Tags:} graph theory, chromatic number, cycles

\medskip

\noindent\textbf{Statement:}

Let $f(n)\to \infty$ (possibly very slowly). Is there a graph of infinite chromatic number such that every finite subgraph on $n$ vertices can be made bipartite by deleting at most $f(n)$ edges?

\noindent\rule{\linewidth}{0.4pt}


%% ============================================================
\section{Problem \#75}
\label{prob:75}

\noindent\textbf{Status:} OPEN \quad|\quad \textbf{Prize:} none \quad|\quad \textbf{Updated:} 2025-08-31

\noindent\textbf{Tags:} graph theory, chromatic number

\medskip

\noindent\textbf{Statement:}

Is there a graph of chromatic number $\aleph_1$ with $\aleph_1$ vertices such that for all $\epsilon&#62;0$ if $n$ is sufficiently large and $H$ is a subgraph on $n$ vertices then $H$ contains an independent set of size $&#62;n^{1-\epsilon}$? What about an independent set of size $\gg n$?

\noindent\rule{\linewidth}{0.4pt}


%% ============================================================
\section{Problem \#77}
\label{prob:77}

\noindent\textbf{Status:} OPEN \quad|\quad \textbf{Prize:} \$250 \quad|\quad \textbf{Updated:} 2025-08-31

\noindent\textbf{Tags:} graph theory, ramsey theory

\medskip

\noindent\textbf{Statement:}

If $R(k)$ is the Ramsey number for $K_k$, the minimal $n$ such that every $2$-colouring of the edges of $K_n$ contains a monochromatic copy of $K_k$, then find the value of\[\lim_{k\to \infty}R(k)^{1/k}.\]

\noindent\rule{\linewidth}{0.4pt}


%% ============================================================
\section{Problem \#78}
\label{prob:78}

\noindent\textbf{Status:} OPEN \quad|\quad \textbf{Prize:} \$100 \quad|\quad \textbf{Updated:} 2025-08-31

\noindent\textbf{Tags:} graph theory, ramsey theory

\medskip

\noindent\textbf{Statement:}

Let $R(k)$ be the Ramsey number for $K_k$, the minimal $n$ such that every $2$-colouring of the edges of $K_n$ contains a monochromatic copy of $K_k$. Give a constructive proof that $R(k)&#62;C^k$ for some constant $C&#62;1$.

\noindent\rule{\linewidth}{0.4pt}


%% ============================================================
\section{Problem \#80}
\label{prob:80}

\noindent\textbf{Status:} OPEN \quad|\quad \textbf{Prize:} none \quad|\quad \textbf{Updated:} 2025-08-31

\noindent\textbf{Tags:} graph theory, ramsey theory

\medskip

\noindent\textbf{Statement:}

Let $c&#62;0$ and let $f_c(n)$ be the maximal $m$ such that every graph $G$ with $n$ vertices and at least $cn^2$ edges, where each edge is contained in at least one triangle, must contain a book of size $m$, that is, an edge shared by at least $m$ different triangles. Estimate $f_c(n)$. In particular, is it true that $f_c(n)&#62;n^{\epsilon}$ for some $\epsilon&#62;0$? Or $f_c(n)\gg \log n$?

\noindent\rule{\linewidth}{0.4pt}


%% ============================================================
\section{Problem \#81}
\label{prob:81}

\noindent\textbf{Status:} OPEN \quad|\quad \textbf{Prize:} none \quad|\quad \textbf{Updated:} 2025-08-31

\noindent\textbf{Tags:} graph theory

\medskip

\noindent\textbf{Statement:}

Let $G$ be a chordal graph on $n$ vertices - that is, $G$ has no induced cycles of length greater than $3$. Can the edges of $G$ be partitioned into $n^2/6+O(n)$ many cliques?

\noindent\rule{\linewidth}{0.4pt}


%% ============================================================
\section{Problem \#82}
\label{prob:82}

\noindent\textbf{Status:} OPEN \quad|\quad \textbf{Prize:} none \quad|\quad \textbf{Updated:} 2025-08-31

\noindent\textbf{Tags:} graph theory

\medskip

\noindent\textbf{Statement:}

Let $F(n)$ be maximal such that every graph on $n$ vertices contains a regular induced subgraph on at least $F(n)$ vertices. Prove that $F(n)/\log n\to \infty$.

\noindent\rule{\linewidth}{0.4pt}


%% ============================================================
\section{Problem \#84}
\label{prob:84}

\noindent\textbf{Status:} OPEN \quad|\quad \textbf{Prize:} none \quad|\quad \textbf{Updated:} 2025-08-31

\noindent\textbf{Tags:} graph theory, cycles

\medskip

\noindent\textbf{Statement:}

The cycle set of a graph $G$ on $n$ vertices is a set $A\subseteq \{3,\ldots,n\}$ such that there is a cycle in $G$ of length $\ell$ if and only if $\ell \in A$. Let $f(n)$ count the number of possible such $A$. Prove that $f(n)=o(2^n)$. Prove that $f(n)/2^{n/2}\to \infty$.

\noindent\rule{\linewidth}{0.4pt}


%% ============================================================
\section{Problem \#85}
\label{prob:85}

\noindent\textbf{Status:} OPEN \quad|\quad \textbf{Prize:} none \quad|\quad \textbf{Updated:} 2026-03-14

\noindent\textbf{Tags:} graph theory

\medskip

\noindent\textbf{Statement:}

Let $n\geq 4$ and $f(n)$ be minimal such that every graph on $n$ vertices with minimal degree $\geq f(n)$ contains a $C_4$. Is it true that, for all large $n$, $f(n+1)\geq f(n)$?

\noindent\rule{\linewidth}{0.4pt}


%% ============================================================
\section{Problem \#86}
\label{prob:86}

\noindent\textbf{Status:} OPEN \quad|\quad \textbf{Prize:} \$100 \quad|\quad \textbf{Updated:} 2025-08-31

\noindent\textbf{Tags:} graph theory

\medskip

\noindent\textbf{Statement:}

Let $Q_n$ be the $n$-dimensional hypercube graph (so that $Q_n$ has $2^n$ vertices and $n2^{n-1}$ edges). Is it true that every subgraph of $Q_n$ with\[\geq \left(\frac{1}{2}+o(1)\right)n2^{n-1}\]many edges contains a $C_4$?

\noindent\rule{\linewidth}{0.4pt}


%% ============================================================
\section{Problem \#87}
\label{prob:87}

\noindent\textbf{Status:} OPEN \quad|\quad \textbf{Prize:} none \quad|\quad \textbf{Updated:} 2025-08-31

\noindent\textbf{Tags:} graph theory, ramsey theory

\medskip

\noindent\textbf{Statement:}

Let $\epsilon &#62;0$. Is it true that, if $k$ is sufficiently large, then\[R(G)&#62;(1-\epsilon)^kR(k)\]for every graph $G$ with chromatic number $\chi(G)=k$? Even stronger, is there some $c&#62;0$ such that, for all large $k$, $R(G)&#62;cR(k)$ for every graph $G$ with chromatic number $\chi(G)=k$?

\noindent\rule{\linewidth}{0.4pt}


%% ============================================================
\section{Problem \#89}
\label{prob:89}

\noindent\textbf{Status:} OPEN \quad|\quad \textbf{Prize:} \$500 \quad|\quad \textbf{Updated:} 2025-08-31

\noindent\textbf{Tags:} geometry, distances

\medskip
\noindent\textit{Erdős distance problem}


\noindent\textbf{Statement:}

Does every set of $n$ distinct points in $\mathbb{R}^2$ determine $\gg n/\sqrt{\log n}$ many distinct distances?

\noindent\rule{\linewidth}{0.4pt}


%% ============================================================
\section{Problem \#91}
\label{prob:91}

\noindent\textbf{Status:} OPEN \quad|\quad \textbf{Prize:} none \quad|\quad \textbf{Updated:} 2025-08-31

\noindent\textbf{Tags:} geometry, distances

\medskip

\noindent\textbf{Statement:}

Let $n$ be a sufficiently large integer. Suppose $A\subset \mathbb{R}^2$ has $\lvert A\rvert=n$ and minimises the number of distinct distances between points in $A$. Prove that there are at least two (and probably many) such $A$ which are non-similar.

\noindent\rule{\linewidth}{0.4pt}


%% ============================================================
\section{Problem \#96}
\label{prob:96}

\noindent\textbf{Status:} OPEN \quad|\quad \textbf{Prize:} none \quad|\quad \textbf{Updated:} 2025-08-31

\noindent\textbf{Tags:} geometry, distances, convex

\medskip

\noindent\textbf{Statement:}

If $n$ points in $\mathbb{R}^2$ form a convex polygon then there are $O(n)$ many pairs which are distance $1$ apart.

\noindent\rule{\linewidth}{0.4pt}


%% ============================================================
\section{Problem \#98}
\label{prob:98}

\noindent\textbf{Status:} OPEN \quad|\quad \textbf{Prize:} none \quad|\quad \textbf{Updated:} 2025-08-31

\noindent\textbf{Tags:} geometry, distances

\medskip

\noindent\textbf{Statement:}

Let $h(n)$ be such that any $n$ points in $\mathbb{R}^2$, with no three on a line and no four on a circle, determine at least $h(n)$ distinct distances. Does $h(n)/n\to \infty$?

\noindent\rule{\linewidth}{0.4pt}


%% ============================================================
\section{Problem \#99}
\label{prob:99}

\noindent\textbf{Status:} OPEN \quad|\quad \textbf{Prize:} \$100 \quad|\quad \textbf{Updated:} 2025-08-31

\noindent\textbf{Tags:} geometry, distances

\medskip

\noindent\textbf{Statement:}

Let $A\subseteq\mathbb{R}^2$ be a set of $n$ points with minimum distance equal to 1, chosen to minimise the diameter of $A$. If $n$ is sufficiently large then must there be three points in $A$ which form an equilateral triangle of size 1?

\noindent\rule{\linewidth}{0.4pt}


%% ============================================================
\section{Problem \#100}
\label{prob:100}

\noindent\textbf{Status:} OPEN \quad|\quad \textbf{Prize:} none \quad|\quad \textbf{Updated:} 2025-08-31

\noindent\textbf{Tags:} geometry, distances

\medskip

\noindent\textbf{Statement:}

Let $A$ be a set of $n$ points in $\mathbb{R}^2$ such that all pairwise distances are at least $1$ and if two distinct distances differ then they differ by at least $1$. Is the diameter of $A$ $\gg n$?

\noindent\rule{\linewidth}{0.4pt}


%% ============================================================
\section{Problem \#101}
\label{prob:101}

\noindent\textbf{Status:} OPEN \quad|\quad \textbf{Prize:} \$100 \quad|\quad \textbf{Updated:} 2025-08-31

\noindent\textbf{Tags:} geometry

\medskip

\noindent\textbf{Statement:}

Given $n$ points in $\mathbb{R}^2$, no five of which are on a line, the number of lines containing four points is $o(n^2)$.

\noindent\rule{\linewidth}{0.4pt}


%% ============================================================
\section{Problem \#102}
\label{prob:102}

\noindent\textbf{Status:} OPEN \quad|\quad \textbf{Prize:} none \quad|\quad \textbf{Updated:} 2025-08-31

\noindent\textbf{Tags:} geometry

\medskip

\noindent\textbf{Statement:}

Let $c&#62;0$ and $h_c(n)$ be such that for any $n$ points in $\mathbb{R}^2$ such that there are $\geq cn^2$ lines each containing more than three points, there must be some line containing $h_c(n)$ many points. Estimate $h_c(n)$. Is it true that, for fixed $c&#62;0$, we have $h_c(n)\to \infty$?

\noindent\rule{\linewidth}{0.4pt}


%% ============================================================
\section{Problem \#103}
\label{prob:103}

\noindent\textbf{Status:} OPEN \quad|\quad \textbf{Prize:} none \quad|\quad \textbf{Updated:} 2025-08-31

\noindent\textbf{Tags:} geometry, distances

\medskip

\noindent\textbf{Statement:}

Let $h(n)$ count the number of incongruent sets of $n$ points in $\mathbb{R}^2$ which minimise the diameter subject to the constraint that $d(x,y)\geq 1$ for all points $x\neq y$. Is it true that $h(n)\to \infty$?

\noindent\rule{\linewidth}{0.4pt}


%% ============================================================
\section{Problem \#104}
\label{prob:104}

\noindent\textbf{Status:} OPEN \quad|\quad \textbf{Prize:} \$100 \quad|\quad \textbf{Updated:} 2025-08-31

\noindent\textbf{Tags:} geometry

\medskip

\noindent\textbf{Statement:}

Given $n$ points in $\mathbb{R}^2$ the number of distinct unit circles containing at least three points is $o(n^2)$.

\noindent\rule{\linewidth}{0.4pt}


%% ============================================================
\section{Problem \#108}
\label{prob:108}

\noindent\textbf{Status:} OPEN \quad|\quad \textbf{Prize:} none \quad|\quad \textbf{Updated:} 2025-08-31

\noindent\textbf{Tags:} graph theory, chromatic number, cycles

\medskip

\noindent\textbf{Statement:}

For every $r\geq 4$ and $k\geq 2$ is there some finite $f(k,r)$ such that every graph of chromatic number $\geq f(k,r)$ contains a subgraph of girth $\geq r$ and chromatic number $\geq k$?

\noindent\rule{\linewidth}{0.4pt}


%% ============================================================
\section{Problem \#111}
\label{prob:111}

\noindent\textbf{Status:} OPEN \quad|\quad \textbf{Prize:} none \quad|\quad \textbf{Updated:} 2025-08-31

\noindent\textbf{Tags:} graph theory, chromatic number, set theory

\medskip

\noindent\textbf{Statement:}

If $G$ is a graph let $h_G(n)$ be defined such that any subgraph of $G$ on $n$ vertices can be made bipartite after deleting at most $h_G(n)$ edges. What is the behaviour of $h_G(n)$? Is it true that $h_G(n)/n\to \infty$ for every graph $G$ with chromatic number $\aleph_1$?

\noindent\rule{\linewidth}{0.4pt}


%% ============================================================
\section{Problem \#112}
\label{prob:112}

\noindent\textbf{Status:} OPEN \quad|\quad \textbf{Prize:} none \quad|\quad \textbf{Updated:} 2025-08-31

\noindent\textbf{Tags:} graph theory, ramsey theory

\medskip

\noindent\textbf{Statement:}

Let $k=k(n,m)$ be minimal such that any directed graph on $k$ vertices must contain either an independent set of size $n$ or a transitive tournament of size $m$. Determine $k(n,m)$.

\noindent\rule{\linewidth}{0.4pt}


%% ============================================================
\section{Problem \#117}
\label{prob:117}

\noindent\textbf{Status:} OPEN \quad|\quad \textbf{Prize:} none \quad|\quad \textbf{Updated:} 2025-08-31

\noindent\textbf{Tags:} group theory

\medskip

\noindent\textbf{Statement:}

Let $h(n)$ be minimal such that any group $G$ with the property that any subset of $&#62;n$ elements contains some $x\neq y$ such that $xy=yx$ can be covered by at most $h(n)$ many Abelian subgroups. Estimate $h(n)$ as well as possible.

\noindent\rule{\linewidth}{0.4pt}


%% ============================================================
\section{Problem \#119}
\label{prob:119}

\noindent\textbf{Status:} OPEN \quad|\quad \textbf{Prize:} \$100 \quad|\quad \textbf{Updated:} 2025-08-31

\noindent\textbf{Tags:} analysis, polynomials

\medskip

\noindent\textbf{Statement:}

Let $z_i$ be an infinite sequence of complex numbers such that $\lvert z_i\rvert=1$ for all $i\geq 1$, and for $n\geq 1$ let\[p_n(z)=\prod_{i\leq n} (z-z_i).\]Let $M_n=\max_{\lvert z\rvert=1}\lvert p_n(z)\rvert$. Is it true that $\limsup M_n=\infty$? Is it true that there exists $c&#62;0$ such that for infinitely many $n$ we have $M_n &#62; n^c$? Is it true that there exists $c&#62;0$ such that, for all large $n$,\[\sum_{k\leq n}M_k &#62; n^{1+c}?\]

\noindent\rule{\linewidth}{0.4pt}


%% ============================================================
\section{Problem \#120}
\label{prob:120}

\noindent\textbf{Status:} OPEN \quad|\quad \textbf{Prize:} \$100 \quad|\quad \textbf{Updated:} 2025-08-31

\noindent\textbf{Tags:} combinatorics

\medskip
\noindent\textit{Erdős similarity problem}


\noindent\textbf{Statement:}

Let $A\subseteq\mathbb{R}$ be an infinite set. Must there be a set $E\subset \mathbb{R}$ of positive measure which does not contain any set of the shape $aA+b$ for some $a,b\in\mathbb{R}$ and $a\neq 0$?

\noindent\rule{\linewidth}{0.4pt}


%% ============================================================
\section{Problem \#122}
\label{prob:122}

\noindent\textbf{Status:} OPEN \quad|\quad \textbf{Prize:} none \quad|\quad \textbf{Updated:} 2025-08-31

\noindent\textbf{Tags:} number theory

\medskip

\noindent\textbf{Statement:}

For which number theoretic functions $f$ is it true that, for any $F(n)$ such that $F(n)/f(n)\to 0$ for almost all $n$, there are infinitely many $x$ such that\[\frac{\#\{ n\in \mathbb{N} : n+f(n)\in (x,x+F(x))\}}{F(x)}\to \infty?\]

\noindent\rule{\linewidth}{0.4pt}


%% ============================================================
\section{Problem \#123}
\label{prob:123}

\noindent\textbf{Status:} OPEN \quad|\quad \textbf{Prize:} \$250 \quad|\quad \textbf{Updated:} 2025-08-31

\noindent\textbf{Tags:} number theory

\medskip

\noindent\textbf{Statement:}

Let $a,b,c\geq 1$ be three integers which are pairwise coprime. Is every large integer the sum of distinct integers of the form $a^kb^lc^m$ ($k,l,m\geq 0$), none of which divide any other?

\noindent\rule{\linewidth}{0.4pt}


%% ============================================================
\section{Problem \#124}
\label{prob:124}

\noindent\textbf{Status:} OPEN \quad|\quad \textbf{Prize:} none \quad|\quad \textbf{Updated:} 2025-08-31

\noindent\textbf{Tags:} number theory, base representations

\medskip

\noindent\textbf{Statement:}

For any $d\geq 1$ and $k\geq 0$ let $P(d,k)$ be the set of integers which are the sum of distinct powers $d^i$ with $i\geq k$. Let $3\leq d_1&#60;d_2&#60;\cdots &#60;d_r$ be integers such that\[\sum_{1\leq i\leq r}\frac{1}{d_r-1}\geq 1.\]Can all sufficiently large integers be written as a sum of the shape $\sum_i c_ia_i$ where $c_i\in \{0,1\}$ and $a_i\in P(d_i,0)$? If we further have $\mathrm{gcd}(d_1,\ldots,d_r)=1$ then, for any $k\geq 1$, can all sufficiently large integers be written as a sum of the shape $\sum_i c_ia_i$ where $c_i\in \{0,1\}$ and $a_i\in P(d_i,k)$?

\noindent\rule{\linewidth}{0.4pt}


%% ============================================================
\section{Problem \#126}
\label{prob:126}

\noindent\textbf{Status:} OPEN \quad|\quad \textbf{Prize:} \$250 \quad|\quad \textbf{Updated:} 2025-08-31

\noindent\textbf{Tags:} number theory

\medskip

\noindent\textbf{Statement:}

Let $f(n)$ be maximal such that if $A\subseteq\mathbb{N}$ has $\lvert A\rvert=n$ then $\prod_{a\neq b\in A}(a+b)$ has at least $f(n)$ distinct prime factors. Is it true that $f(n)/\log n\to\infty$?

\noindent\rule{\linewidth}{0.4pt}


%% ============================================================
\section{Problem \#129}
\label{prob:129}

\noindent\textbf{Status:} OPEN \quad|\quad \textbf{Prize:} none \quad|\quad \textbf{Updated:} 2025-08-31

\noindent\textbf{Tags:} graph theory, ramsey theory

\medskip
\noindent\textit{ambiguous statement}


\noindent\textbf{Statement:}

Let $R(n;k,r)$ be the smallest $N$ such that if the edges of $K_N$ are $r$-coloured then there is a set of $n$ vertices which does not contain a copy of $K_k$ in at least one of the $r$ colours. Prove that there is a constant $C=C(r)&#62;1$ such that\[R(n;3,r) &#60; C^{\sqrt{n}}.\]

\noindent\rule{\linewidth}{0.4pt}


%% ============================================================
\section{Problem \#130}
\label{prob:130}

\noindent\textbf{Status:} OPEN \quad|\quad \textbf{Prize:} none \quad|\quad \textbf{Updated:} 2025-08-31

\noindent\textbf{Tags:} graph theory, chromatic number

\medskip

\noindent\textbf{Statement:}

Let $A\subset\mathbb{R}^2$ be an infinite set which contains no three points on a line and no four points on a circle. Consider the graph with vertices the points in $A$, where two vertices are joined by an edge if and only if they are an integer distance apart. How large can the chromatic number and clique number of this graph be? In particular, can the chromatic number be infinite?

\noindent\rule{\linewidth}{0.4pt}


%% ============================================================
\section{Problem \#131}
\label{prob:131}

\noindent\textbf{Status:} OPEN \quad|\quad \textbf{Prize:} none \quad|\quad \textbf{Updated:} 2025-08-31

\noindent\textbf{Tags:} number theory

\medskip

\noindent\textbf{Statement:}

Let $F(N)$ be the maximal size of $A\subseteq\{1,\ldots,N\}$ such that no $a\in A$ divides the sum of any distinct elements of $A\backslash\{a\}$. Estimate $F(N)$. In particular, is it true that\[F(N) &#62; N^{1/2-o(1)}?\]

\noindent\rule{\linewidth}{0.4pt}


%% ============================================================
\section{Problem \#132}
\label{prob:132}

\noindent\textbf{Status:} OPEN \quad|\quad \textbf{Prize:} \$100 \quad|\quad \textbf{Updated:} 2025-08-31

\noindent\textbf{Tags:} distances

\medskip

\noindent\textbf{Statement:}

Let $A\subset \mathbb{R}^2$ be a set of $n$ points. Must there be two distances which occur at least once but between at most $n$ pairs of points? Must the number of such distances $\to \infty$ as $n\to \infty$?

\noindent\rule{\linewidth}{0.4pt}


%% ============================================================
\section{Problem \#137}
\label{prob:137}

\noindent\textbf{Status:} OPEN \quad|\quad \textbf{Prize:} none \quad|\quad \textbf{Updated:} 2025-08-31

\noindent\textbf{Tags:} number theory

\medskip

\noindent\textbf{Statement:}

We say that $N$ is powerful if whenever $p\mid N$ we also have $p^2\mid N$. Let $k\geq 3$. Can the product of any $k$ consecutive positive integers ever be powerful?

\noindent\rule{\linewidth}{0.4pt}


%% ============================================================
\section{Problem \#138}
\label{prob:138}

\noindent\textbf{Status:} OPEN \quad|\quad \textbf{Prize:} \$500 \quad|\quad \textbf{Updated:} 2025-08-31

\noindent\textbf{Tags:} additive combinatorics

\medskip

\noindent\textbf{Statement:}

Let the van der Waerden number $W(k)$ be such that whenever $N\geq W(k)$ and $\{1,\ldots,N\}$ is $2$-coloured there must exist a monochromatic $k$-term arithmetic progression. Improve the bounds for $W(k)$ - for example, prove that $W(k)^{1/k}\to \infty$.

\noindent\rule{\linewidth}{0.4pt}


%% ============================================================
\section{Problem \#141}
\label{prob:141}

\noindent\textbf{Status:} OPEN \quad|\quad \textbf{Prize:} none \quad|\quad \textbf{Updated:} 2025-08-31

\noindent\textbf{Tags:} additive combinatorics, primes, arithmetic progressions

\medskip

\noindent\textbf{Statement:}

Let $k\geq 3$. Are there $k$ consecutive primes in arithmetic progression?

\noindent\rule{\linewidth}{0.4pt}


%% ============================================================
\section{Problem \#142}
\label{prob:142}

\noindent\textbf{Status:} OPEN \quad|\quad \textbf{Prize:} \$10000 \quad|\quad \textbf{Updated:} 2025-08-31

\noindent\textbf{Tags:} additive combinatorics, arithmetic progressions

\medskip

\noindent\textbf{Statement:}

Let $r_k(N)$ be the largest possible size of a subset of $\{1,\ldots,N\}$ that does not contain any non-trivial $k$-term arithmetic progression. Prove an asymptotic formula for $r_k(N)$.

\noindent\rule{\linewidth}{0.4pt}


%% ============================================================
\section{Problem \#143}
\label{prob:143}

\noindent\textbf{Status:} OPEN \quad|\quad \textbf{Prize:} \$500 \quad|\quad \textbf{Updated:} 2025-08-31

\noindent\textbf{Tags:} primitive sets

\medskip

\noindent\textbf{Statement:}

Let $A\subset (1,\infty)$ be a countably infinite set such that for all $x\neq y\in A$ and integers $k\geq 1$ we have\[ \lvert kx -y\rvert \geq 1.\]Does this imply that $A$ is sparse? In particular, does this imply that\[\sum_{x\in A}\frac{1}{x\log x}&#60;\infty\]or\[\sum_{\substack{x &#60;n\\ x\in A}}\frac{1}{x}=o(\log n)?\]

\noindent\rule{\linewidth}{0.4pt}


%% ============================================================
\section{Problem \#145}
\label{prob:145}

\noindent\textbf{Status:} OPEN \quad|\quad \textbf{Prize:} none \quad|\quad \textbf{Updated:} 2025-08-31

\noindent\textbf{Tags:} number theory

\medskip

\noindent\textbf{Statement:}

Let $s_1&#60;s_2&#60;\cdots$ be the sequence of squarefree numbers. Is it true that, for any $\alpha \geq 0$,\[\lim_{x\to \infty}\frac{1}{x}\sum_{s_n\leq x}(s_{n+1}-s_n)^\alpha\]exists?

\noindent\rule{\linewidth}{0.4pt}


%% ============================================================
\section{Problem \#146}
\label{prob:146}

\noindent\textbf{Status:} OPEN \quad|\quad \textbf{Prize:} \$500 \quad|\quad \textbf{Updated:} 2025-08-31

\noindent\textbf{Tags:} graph theory, turan number

\medskip

\noindent\textbf{Statement:}

If $H$ is bipartite and is $r$-degenerate, that is, every induced subgraph of $H$ has minimum degree $\leq r$, then\[\mathrm{ex}(n;H) \ll n^{2-1/r}.\]

\noindent\rule{\linewidth}{0.4pt}


%% ============================================================
\section{Problem \#148}
\label{prob:148}

\noindent\textbf{Status:} OPEN \quad|\quad \textbf{Prize:} none \quad|\quad \textbf{Updated:} 2025-08-31

\noindent\textbf{Tags:} number theory, unit fractions

\medskip

\noindent\textbf{Statement:}

Let $F(k)$ be the number of solutions to\[ 1= \frac{1}{n_1}+\cdots+\frac{1}{n_k},\]where $1\leq n_1&#60;\cdots&#60;n_k$ are distinct integers. Find good estimates for $F(k)$.

\noindent\rule{\linewidth}{0.4pt}


%% ============================================================
\section{Problem \#149}
\label{prob:149}

\noindent\textbf{Status:} OPEN \quad|\quad \textbf{Prize:} none \quad|\quad \textbf{Updated:} 2025-08-31

\noindent\textbf{Tags:} graph theory

\medskip

\noindent\textbf{Statement:}

The strong chromatic index of a graph $G$, denoted by $\mathrm{sq}(G)$, is the minimum $k$ such that the edges of $G$ can be partitioned into $k$ sets of 'strongly independent' edges, that is, such that the subgraph of $G$ induced by each set is the union of vertex-disjoint edges. Is it true that, for any graph $G$ with maximum degree $\Delta$,\[\mathrm{sq}(G)\leq\frac{5}{4}\Delta^2?\]

\noindent\rule{\linewidth}{0.4pt}


%% ============================================================
\section{Problem \#151}
\label{prob:151}

\noindent\textbf{Status:} OPEN \quad|\quad \textbf{Prize:} none \quad|\quad \textbf{Updated:} 2025-08-31

\noindent\textbf{Tags:} graph theory

\medskip

\noindent\textbf{Statement:}

For a graph $G$ let $\tau(G)$ denote the minimal number of vertices that include at least one from each maximal clique of $G$ on at least two vertices (sometimes called the clique transversal number). Let $H(n)$ be maximal such that every triangle-free graph on $n$ vertices contains an independent set on $H(n)$ vertices. If $G$ is a graph on $n$ vertices then is\[\tau(G)\leq n-H(n)?\]

\noindent\rule{\linewidth}{0.4pt}


%% ============================================================
\section{Problem \#153}
\label{prob:153}

\noindent\textbf{Status:} OPEN \quad|\quad \textbf{Prize:} none \quad|\quad \textbf{Updated:} 2025-08-31

\noindent\textbf{Tags:} sidon sets

\medskip

\noindent\textbf{Statement:}

Let $A$ be a finite Sidon set and $A+A=\{s_1&#60;\cdots&#60;s_t\}$. Is it true that\[\frac{1}{t}\sum_{1\leq i&#60;t}(s_{i+1}-s_i)^2 \to \infty\]as $\lvert A\rvert\to \infty$?

\noindent\rule{\linewidth}{0.4pt}


%% ============================================================
\section{Problem \#155}
\label{prob:155}

\noindent\textbf{Status:} OPEN \quad|\quad \textbf{Prize:} none \quad|\quad \textbf{Updated:} 2025-08-31

\noindent\textbf{Tags:} additive combinatorics, sidon sets

\medskip

\noindent\textbf{Statement:}

Let $F(N)$ be the size of the largest Sidon subset of $\{1,\ldots,N\}$. Is it true that for every $k\geq 1$ we have\[F(N+k)\leq F(N)+1\]for all sufficiently large $N$?

\noindent\rule{\linewidth}{0.4pt}


%% ============================================================
\section{Problem \#156}
\label{prob:156}

\noindent\textbf{Status:} OPEN \quad|\quad \textbf{Prize:} none \quad|\quad \textbf{Updated:} 2025-08-31

\noindent\textbf{Tags:} sidon sets

\medskip

\noindent\textbf{Statement:}

Does there exist a maximal Sidon set $A\subset \{1,\ldots,N\}$ of size $O(N^{1/3})$?

\noindent\rule{\linewidth}{0.4pt}


%% ============================================================
\section{Problem \#158}
\label{prob:158}

\noindent\textbf{Status:} OPEN \quad|\quad \textbf{Prize:} none \quad|\quad \textbf{Updated:} 2025-08-31

\noindent\textbf{Tags:} sidon sets

\medskip

\noindent\textbf{Statement:}

Let $A\subset \mathbb{N}$ be an infinite set such that, for any $n$, there are most $2$ solutions to $a+b=n$ with $a\leq b$. Must\[\liminf_{N\to\infty}\frac{\lvert A\cap \{1,\ldots,N\}\rvert}{N^{1/2}}=0?\]

\noindent\rule{\linewidth}{0.4pt}


%% ============================================================
\section{Problem \#159}
\label{prob:159}

\noindent\textbf{Status:} OPEN \quad|\quad \textbf{Prize:} none \quad|\quad \textbf{Updated:} 2025-08-31

\noindent\textbf{Tags:} graph theory, ramsey theory

\medskip

\noindent\textbf{Statement:}

There exists some constant $c&#62;0$ such that $$R(C_4,K_n) \ll n^{2-c}.$$

\noindent\rule{\linewidth}{0.4pt}


%% ============================================================
\section{Problem \#160}
\label{prob:160}

\noindent\textbf{Status:} OPEN \quad|\quad \textbf{Prize:} none \quad|\quad \textbf{Updated:} 2025-08-31

\noindent\textbf{Tags:} additive combinatorics, arithmetic progressions

\medskip

\noindent\textbf{Statement:}

Let $h(N)$ be the smallest $k$ such that $\{1,\ldots,N\}$ can be coloured with $k$ colours so that every four-term arithmetic progression must contain at least three distinct colours. Estimate $h(N)$.

\noindent\rule{\linewidth}{0.4pt}


%% ============================================================
\section{Problem \#161}
\label{prob:161}

\noindent\textbf{Status:} OPEN \quad|\quad \textbf{Prize:} \$500 \quad|\quad \textbf{Updated:} 2025-08-31

\noindent\textbf{Tags:} combinatorics, ramsey theory, discrepancy

\medskip

\noindent\textbf{Statement:}

Let $\alpha\in[0,1/2)$ and $n,t\geq 1$. Let $F^{(t)}(n,\alpha)$ be the smallest $m$ such that we can $2$-colour the edges of the complete $t$-uniform hypergraph on $n$ vertices such that if $X\subseteq [n]$ with $\lvert X\rvert \geq m$ then there are at least $\alpha \binom{\lvert X\rvert}{t}$ many $t$-subsets of $X$ of each colour. For fixed $n,t$ as we change $\alpha$ from $0$ to $1/2$ does $F^{(t)}(n,\alpha)$ increase continuously or are there jumps? Only one jump?

\noindent\rule{\linewidth}{0.4pt}


%% ============================================================
\section{Problem \#162}
\label{prob:162}

\noindent\textbf{Status:} OPEN \quad|\quad \textbf{Prize:} none \quad|\quad \textbf{Updated:} 2025-08-31

\noindent\textbf{Tags:} combinatorics, ramsey theory, discrepancy

\medskip

\noindent\textbf{Statement:}

Let $\alpha&#62;0$ and $n\geq 1$. Let $F(n,\alpha)$ be the largest $k$ such that there exists some 2-colouring of the edges of $K_n$ in which any induced subgraph $H$ on at least $k$ vertices contains more than $\alpha\binom{\lvert H\rvert}{2}$ many edges of each colour. Prove that for every fixed $0\leq \alpha \leq 1/2$, as $n\to\infty$,\[F(n,\alpha)\sim c_\alpha \log n\]for some constant $c_\alpha$.

\noindent\rule{\linewidth}{0.4pt}


%% ============================================================
\section{Problem \#165}
\label{prob:165}

\noindent\textbf{Status:} OPEN \quad|\quad \textbf{Prize:} \$250 \quad|\quad \textbf{Updated:} 2025-08-31

\noindent\textbf{Tags:} graph theory, ramsey theory

\medskip

\noindent\textbf{Statement:}

Give an asymptotic formula for $R(3,k)$.

\noindent\rule{\linewidth}{0.4pt}


%% ============================================================
\section{Problem \#168}
\label{prob:168}

\noindent\textbf{Status:} OPEN \quad|\quad \textbf{Prize:} none \quad|\quad \textbf{Updated:} 2025-08-31

\noindent\textbf{Tags:} additive combinatorics

\medskip

\noindent\textbf{Statement:}

Let $F(N)$ be the size of the largest subset of $\{1,\ldots,N\}$ which does not contain any set of the form $\{n,2n,3n\}$. What is\[ \lim_{N\to \infty}\frac{F(N)}{N}?\]Is this limit irrational?

\noindent\rule{\linewidth}{0.4pt}


%% ============================================================
\section{Problem \#169}
\label{prob:169}

\noindent\textbf{Status:} OPEN \quad|\quad \textbf{Prize:} none \quad|\quad \textbf{Updated:} 2025-08-31

\noindent\textbf{Tags:} additive combinatorics, arithmetic progressions

\medskip

\noindent\textbf{Statement:}

Let $k\geq 3$ and $f(k)$ be the supremum of $\sum_{n\in A}\frac{1}{n}$ as $A$ ranges over all sets of positive integers which do not contain a $k$-term arithmetic progression. Estimate $f(k)$. Is\[\lim_{k\to \infty}\frac{f(k)}{\log W(k)}=\infty\]where $W(k)$ is the van der Waerden number?

\noindent\rule{\linewidth}{0.4pt}


%% ============================================================
\section{Problem \#170}
\label{prob:170}

\noindent\textbf{Status:} OPEN \quad|\quad \textbf{Prize:} none \quad|\quad \textbf{Updated:} 2025-08-31

\noindent\textbf{Tags:} additive combinatorics

\medskip
\noindent\textit{sparse ruler problem}


\noindent\textbf{Statement:}

Let $F(N)$ be the smallest possible size of $A\subset \{0,1,\ldots,N\}$ such that $\{0,1,\ldots,N\}\subset A-A$. Find the value of\[\lim_{N\to \infty}\frac{F(N)}{N^{1/2}}.\]

\noindent\rule{\linewidth}{0.4pt}


%% ============================================================
\section{Problem \#172}
\label{prob:172}

\noindent\textbf{Status:} OPEN \quad|\quad \textbf{Prize:} none \quad|\quad \textbf{Updated:} 2025-08-31

\noindent\textbf{Tags:} additive combinatorics, ramsey theory

\medskip

\noindent\textbf{Statement:}

Is it true that in any finite colouring of $\mathbb{N}$ there exist arbitrarily large finite $A$ such that all sums and products of distinct elements in $A$ are the same colour?

\noindent\rule{\linewidth}{0.4pt}


%% ============================================================
\section{Problem \#173}
\label{prob:173}

\noindent\textbf{Status:} OPEN \quad|\quad \textbf{Prize:} none \quad|\quad \textbf{Updated:} 2025-08-31

\noindent\textbf{Tags:} geometry, ramsey theory

\medskip

\noindent\textbf{Statement:}

In any $2$-colouring of $\mathbb{R}^2$, for all but at most one triangle $T$, there is a monochromatic congruent copy of $T$.

\noindent\rule{\linewidth}{0.4pt}


%% ============================================================
\section{Problem \#174}
\label{prob:174}

\noindent\textbf{Status:} OPEN \quad|\quad \textbf{Prize:} none \quad|\quad \textbf{Updated:} 2025-08-31

\noindent\textbf{Tags:} geometry, ramsey theory

\medskip

\noindent\textbf{Statement:}

A finite set $A\subset \mathbb{R}^n$ is called Ramsey if, for any $k\geq 1$, there exists some $d=d(A,k)$ such that in any $k$-colouring of $\mathbb{R}^d$ there exists a monochromatic copy of $A$. Characterise the Ramsey sets in $\mathbb{R}^n$.

\noindent\rule{\linewidth}{0.4pt}


%% ============================================================
\section{Problem \#176}
\label{prob:176}

\noindent\textbf{Status:} OPEN \quad|\quad \textbf{Prize:} none \quad|\quad \textbf{Updated:} 2025-08-31

\noindent\textbf{Tags:} additive combinatorics, arithmetic progressions, discrepancy

\medskip

\noindent\textbf{Statement:}

Let $N(k,\ell)$ be the minimal $N$ such that for any $f:\{1,\ldots,N\}\to\{-1,1\}$ there must exist a $k$-term arithmetic progression $P$ such that\[ \left\lvert \sum_{n\in P}f(n)\right\rvert\geq \ell.\]Find good upper bounds for $N(k,\ell)$. Is it true that for any $c&#62;0$ there exists some $C&#62;1$ such that\[N(k,ck)\leq C^k?\]What about\[N(k,2)\leq C^k\]or\[N(k,\sqrt{k})\leq C^k?\]

\noindent\rule{\linewidth}{0.4pt}


%% ============================================================
\section{Problem \#177}
\label{prob:177}

\noindent\textbf{Status:} OPEN \quad|\quad \textbf{Prize:} none \quad|\quad \textbf{Updated:} 2025-08-31

\noindent\textbf{Tags:} discrepancy, arithmetic progressions

\medskip

\noindent\textbf{Statement:}

Find the smallest $h(d)$ such that the following holds. There exists a function $f:\mathbb{N}\to\{-1,1\}$ such that, for every $d\geq 1$,\[\max_{P_d}\left\lvert \sum_{n\in P_d}f(n)\right\rvert\leq h(d),\]where $P_d$ ranges over all finite arithmetic progressions with common difference $d$.

\noindent\rule{\linewidth}{0.4pt}


%% ============================================================
\section{Problem \#180}
\label{prob:180}

\noindent\textbf{Status:} OPEN \quad|\quad \textbf{Prize:} none \quad|\quad \textbf{Updated:} 2025-08-31

\noindent\textbf{Tags:} graph theory, turan number

\medskip

\noindent\textbf{Statement:}

If $\mathcal{F}$ is a finite set of finite graphs then $\mathrm{ex}(n;\mathcal{F})$ is the maximum number of edges a graph on $n$ vertices can have without containing any subgraphs from $\mathcal{F}$. Note that it is trivial that $\mathrm{ex}(n;\mathcal{F})\leq \mathrm{ex}(n;G)$ for every $G\in\mathcal{F}$. Is it true that, for every $\mathcal{F}$, there exists $G\in\mathcal{F}$ such that\[\mathrm{ex}(n;G)\ll_{\mathcal{F}}\mathrm{ex}(n;\mathcal{F})?\]

\noindent\rule{\linewidth}{0.4pt}


%% ============================================================
\section{Problem \#181}
\label{prob:181}

\noindent\textbf{Status:} OPEN \quad|\quad \textbf{Prize:} none \quad|\quad \textbf{Updated:} 2025-08-31

\noindent\textbf{Tags:} graph theory, ramsey theory

\medskip

\noindent\textbf{Statement:}

Let $Q_n$ be the $n$-dimensional hypercube graph (so that $Q_n$ has $2^n$ vertices and $n2^{n-1}$ edges). Prove that\[R(Q_n) \ll 2^n.\]

\noindent\rule{\linewidth}{0.4pt}


%% ============================================================
\section{Problem \#183}
\label{prob:183}

\noindent\textbf{Status:} OPEN \quad|\quad \textbf{Prize:} \$250 \quad|\quad \textbf{Updated:} 2025-08-31

\noindent\textbf{Tags:} graph theory, ramsey theory

\medskip

\noindent\textbf{Statement:}

Let $R(3;k)$ be the minimal $n$ such that if the edges of $K_n$ are coloured with $k$ colours then there must exist a monochromatic triangle. Determine\[\lim_{k\to \infty}R(3;k)^{1/k}.\]

\noindent\rule{\linewidth}{0.4pt}


%% ============================================================
\section{Problem \#184}
\label{prob:184}

\noindent\textbf{Status:} OPEN \quad|\quad \textbf{Prize:} none \quad|\quad \textbf{Updated:} 2025-08-31

\noindent\textbf{Tags:} graph theory, cycles

\medskip

\noindent\textbf{Statement:}

Any graph on $n$ vertices can be decomposed into $O(n)$ many edge-disjoint cycles and edges.

\noindent\rule{\linewidth}{0.4pt}


%% ============================================================
\section{Problem \#187}
\label{prob:187}

\noindent\textbf{Status:} OPEN \quad|\quad \textbf{Prize:} none \quad|\quad \textbf{Updated:} 2025-08-31

\noindent\textbf{Tags:} additive combinatorics, ramsey theory

\medskip

\noindent\textbf{Statement:}

Find the best function $f(d)$ such that, in any 2-colouring of the integers, at least one colour class contains an arithmetic progression with common difference $d$ of length $f(d)$ for infinitely many $d$.

\noindent\rule{\linewidth}{0.4pt}


%% ============================================================
\section{Problem \#188}
\label{prob:188}

\noindent\textbf{Status:} OPEN \quad|\quad \textbf{Prize:} none \quad|\quad \textbf{Updated:} 2025-08-31

\noindent\textbf{Tags:} geometry, ramsey theory

\medskip

\noindent\textbf{Statement:}

What is the smallest $k$ such that $\mathbb{R}^2$ can be red/blue coloured with no pair of red points unit distance apart, and no $k$-term arithmetic progression of blue points with distance $1$?

\noindent\rule{\linewidth}{0.4pt}


%% ============================================================
\section{Problem \#193}
\label{prob:193}

\noindent\textbf{Status:} OPEN \quad|\quad \textbf{Prize:} none \quad|\quad \textbf{Updated:} 2025-08-31

\noindent\textbf{Tags:} geometry

\medskip

\noindent\textbf{Statement:}

Let $S\subseteq \mathbb{Z}^3$ be a finite set and let $A=\{a_1,a_2,\ldots,\}\subset \mathbb{Z}^3$ be an infinite $S$-walk, so that $a_{i+1}-a_i\in S$ for all $i$. Must $A$ contain three collinear points?

\noindent\rule{\linewidth}{0.4pt}


%% ============================================================
\section{Problem \#195}
\label{prob:195}

\noindent\textbf{Status:} OPEN \quad|\quad \textbf{Prize:} none \quad|\quad \textbf{Updated:} 2025-08-31

\noindent\textbf{Tags:} arithmetic progressions

\medskip

\noindent\textbf{Statement:}

What is the largest $k$ such that in any permutation of $\mathbb{Z}$ there must exist a monotone $k$-term arithmetic progression $x_1&#60;\cdots&#60;x_k$?

\noindent\rule{\linewidth}{0.4pt}


%% ============================================================
\section{Problem \#196}
\label{prob:196}

\noindent\textbf{Status:} OPEN \quad|\quad \textbf{Prize:} none \quad|\quad \textbf{Updated:} 2025-08-31

\noindent\textbf{Tags:} arithmetic progressions

\medskip

\noindent\textbf{Statement:}

Must every permutation of $\mathbb{N}$ contain a monotone 4-term arithmetic progression? In other words, given a permutation $x$ of $\mathbb{N}$ must there be indices with either $i&lt;j&lt;k&lt;l$ or $i&gt;j&#62;k&#62;l$ such that $x_i,x_j,x_k,x_l$ are an arithmetic progression?

\noindent\rule{\linewidth}{0.4pt}


%% ============================================================
\section{Problem \#197}
\label{prob:197}

\noindent\textbf{Status:} OPEN \quad|\quad \textbf{Prize:} none \quad|\quad \textbf{Updated:} 2025-08-31

\noindent\textbf{Tags:} arithmetic progressions

\medskip

\noindent\textbf{Statement:}

Can $\mathbb{N}$ be partitioned into two sets, each of which can be permuted to avoid monotone 3-term arithmetic progressions?

\noindent\rule{\linewidth}{0.4pt}


%% ============================================================
\section{Problem \#200}
\label{prob:200}

\noindent\textbf{Status:} OPEN \quad|\quad \textbf{Prize:} none \quad|\quad \textbf{Updated:} 2025-08-31

\noindent\textbf{Tags:} primes, arithmetic progressions

\medskip

\noindent\textbf{Statement:}

Does the longest arithmetic progression of primes in $\{1,\ldots,N\}$ have length $o(\log N)$?

\noindent\rule{\linewidth}{0.4pt}


%% ============================================================
\section{Problem \#201}
\label{prob:201}

\noindent\textbf{Status:} OPEN \quad|\quad \textbf{Prize:} none \quad|\quad \textbf{Updated:} 2025-08-31

\noindent\textbf{Tags:} additive combinatorics, arithmetic progressions

\medskip

\noindent\textbf{Statement:}

Let $G_k(N)$ be such that any set of $N$ integers contains a subset of size at least $G_k(N)$ which does not contain a $k$-term arithmetic progression. Determine the size of $G_k(N)$. How does it relate to $R_k(N)$, the size of the largest subset of $\{1,\ldots,N\}$ without a $k$-term arithmetic progression? Is it true that\[\lim_{N\to \infty}\frac{R_3(N)}{G_3(N)}=1?\]

\noindent\rule{\linewidth}{0.4pt}


%% ============================================================
\section{Problem \#203}
\label{prob:203}

\noindent\textbf{Status:} OPEN \quad|\quad \textbf{Prize:} none \quad|\quad \textbf{Updated:} 2025-08-31

\noindent\textbf{Tags:} primes, covering systems

\medskip

\noindent\textbf{Statement:}

Is there an integer $m\geq 1$ with $(m,6)=1$ such that none of $2^k3^\ell m+1$ are prime, for any $k,\ell\geq 0$?

\noindent\rule{\linewidth}{0.4pt}


%% ============================================================
\section{Problem \#208}
\label{prob:208}

\noindent\textbf{Status:} OPEN \quad|\quad \textbf{Prize:} none \quad|\quad \textbf{Updated:} 2025-08-31

\noindent\textbf{Tags:} number theory

\medskip

\noindent\textbf{Statement:}

Let $s_1&lt;s_2&lt;\cdots$ be the sequence of squarefree numbers. Is it true that, for any $\epsilon&gt;0$ and large $n$,\[s_{n+1}-s_n \ll_\epsilon s_n^{\epsilon}?\]Is it true that\[s_{n+1}-s_n \leq (1+o(1))\frac{\pi^2}{6}\frac{\log s_n}{\log\log s_n}?\]

\noindent\rule{\linewidth}{0.4pt}


%% ============================================================
\section{Problem \#212}
\label{prob:212}

\noindent\textbf{Status:} OPEN \quad|\quad \textbf{Prize:} none \quad|\quad \textbf{Updated:} 2025-08-31

\noindent\textbf{Tags:} geometry, distances

\medskip

\noindent\textbf{Statement:}

Is there a dense subset of $\mathbb{R}^2$ such that all pairwise distances are rational?

\noindent\rule{\linewidth}{0.4pt}


%% ============================================================
\section{Problem \#213}
\label{prob:213}

\noindent\textbf{Status:} OPEN \quad|\quad \textbf{Prize:} none \quad|\quad \textbf{Updated:} 2025-08-31

\noindent\textbf{Tags:} geometry, distances

\medskip

\noindent\textbf{Statement:}

Let $n\geq 4$. Are there $n$ points in $\mathbb{R}^2$, no three on a line and no four on a circle, such that all pairwise distances are integers?

\noindent\rule{\linewidth}{0.4pt}


%% ============================================================
\section{Problem \#217}
\label{prob:217}

\noindent\textbf{Status:} OPEN \quad|\quad \textbf{Prize:} none \quad|\quad \textbf{Updated:} 2025-08-31

\noindent\textbf{Tags:} geometry, distances

\medskip

\noindent\textbf{Statement:}

For which $n$ are there $n$ points in $\mathbb{R}^2$, no three on a line and no four on a circle, which determine $n-1$ distinct distances and so that (in some ordering of the distances) the $i$th distance occurs $i$ times?

\noindent\rule{\linewidth}{0.4pt}


%% ============================================================
\section{Problem \#218}
\label{prob:218}

\noindent\textbf{Status:} OPEN \quad|\quad \textbf{Prize:} none \quad|\quad \textbf{Updated:} 2025-08-31

\noindent\textbf{Tags:} number theory, primes

\medskip

\noindent\textbf{Statement:}

Let $d_n=p_{n+1}-p_n$. The set of $n$ such that $d_{n+1}\geq d_n$ has density $1/2$, and similarly for $d_{n+1}\leq d_n$. Furthermore, there are infinitely many $n$ such that $d_{n+1}=d_n$.

\noindent\rule{\linewidth}{0.4pt}


%% ============================================================
\section{Problem \#222}
\label{prob:222}

\noindent\textbf{Status:} OPEN \quad|\quad \textbf{Prize:} none \quad|\quad \textbf{Updated:} 2025-08-31

\noindent\textbf{Tags:} number theory, squares

\medskip

\noindent\textbf{Statement:}

Let $n_1&#60;n_2&#60;\cdots$ be the sequence of integers which are the sum of two squares. Explore the behaviour of (i.e. find good upper and lower bounds for) the consecutive differences $n_{k+1}-n_k$.

\noindent\rule{\linewidth}{0.4pt}


%% ============================================================
\section{Problem \#233}
\label{prob:233}

\noindent\textbf{Status:} OPEN \quad|\quad \textbf{Prize:} none \quad|\quad \textbf{Updated:} 2025-08-31

\noindent\textbf{Tags:} number theory, primes

\medskip

\noindent\textbf{Statement:}

Let $d_n=p_{n+1}-p_n$, where $p_n$ is the $n$th prime. Prove that\[\sum_{1\leq n\leq N}d_n^2 \ll N(\log N)^2.\]

\noindent\rule{\linewidth}{0.4pt}


%% ============================================================
\section{Problem \#234}
\label{prob:234}

\noindent\textbf{Status:} OPEN \quad|\quad \textbf{Prize:} none \quad|\quad \textbf{Updated:} 2025-08-31

\noindent\textbf{Tags:} number theory, primes

\medskip

\noindent\textbf{Statement:}

For every $c\geq 0$ the density $f(c)$ of integers for which\[\frac{p_{n+1}-p_n}{\log n}&#60; c\]exists and is a continuous function of $c$.

\noindent\rule{\linewidth}{0.4pt}


%% ============================================================
\section{Problem \#236}
\label{prob:236}

\noindent\textbf{Status:} OPEN \quad|\quad \textbf{Prize:} none \quad|\quad \textbf{Updated:} 2025-08-31

\noindent\textbf{Tags:} number theory, primes

\medskip

\noindent\textbf{Statement:}

Let $f(n)$ count the number of solutions to $n=p+2^k$ for prime $p$ and $k\geq 0$. Is it true that $f(n)=o(\log n)$?

\noindent\rule{\linewidth}{0.4pt}


%% ============================================================
\section{Problem \#238}
\label{prob:238}

\noindent\textbf{Status:} OPEN \quad|\quad \textbf{Prize:} none \quad|\quad \textbf{Updated:} 2025-08-31

\noindent\textbf{Tags:} number theory, primes

\medskip

\noindent\textbf{Statement:}

Let $c_1,c_2&#62;0$. Is it true that, for any sufficiently large $x$, there exist more than $c_1\log x$ many consecutive primes $\leq x$ such that the difference between any two is $&#62;c_2$?

\noindent\rule{\linewidth}{0.4pt}


%% ============================================================
\section{Problem \#241}
\label{prob:241}

\noindent\textbf{Status:} OPEN \quad|\quad \textbf{Prize:} \$100 \quad|\quad \textbf{Updated:} 2025-08-31

\noindent\textbf{Tags:} additive combinatorics, sidon sets

\medskip

\noindent\textbf{Statement:}

Let $f(N)$ be the maximum size of $A\subseteq \{1,\ldots,N\}$ such that the sums $a+b+c$ with $a,b,c\in A$ are all distinct (aside from the trivial coincidences). Is it true that\[ f(N)\sim N^{1/3}?\]

\noindent\rule{\linewidth}{0.4pt}


%% ============================================================
\section{Problem \#243}
\label{prob:243}

\noindent\textbf{Status:} OPEN \quad|\quad \textbf{Prize:} none \quad|\quad \textbf{Updated:} 2025-08-31

\noindent\textbf{Tags:} number theory

\medskip

\noindent\textbf{Statement:}

Let $1\leq a_1&#60;a_2&#60;\cdots$ be a sequence of integers such that\[\lim_{n\to \infty}\frac{a_n}{a_{n-1}^2}=1\]and $\sum\frac{1}{a_n}\in \mathbb{Q}$. Then, for all sufficiently large $n\geq 1$,\[ a_n = a_{n-1}^2-a_{n-1}+1.\]

\noindent\rule{\linewidth}{0.4pt}


%% ============================================================
\section{Problem \#244}
\label{prob:244}

\noindent\textbf{Status:} OPEN \quad|\quad \textbf{Prize:} none \quad|\quad \textbf{Updated:} 2025-08-31

\noindent\textbf{Tags:} number theory, primes

\medskip

\noindent\textbf{Statement:}

Let $C&#62;1$. Does the set of integers of the form $p+\lfloor C^k\rfloor$, for some prime $p$ and $k\geq 0$, have density $&#62;0$?

\noindent\rule{\linewidth}{0.4pt}


%% ============================================================
\section{Problem \#247}
\label{prob:247}

\noindent\textbf{Status:} OPEN \quad|\quad \textbf{Prize:} none \quad|\quad \textbf{Updated:} 2025-08-31

\noindent\textbf{Tags:} number theory, irrationality

\medskip

\noindent\textbf{Statement:}

Let $1\leq a_1&#60;a_2&#60;\cdots$ be a sequence of integers such that\[\limsup \frac{a_n}{n}=\infty.\]Is\[\sum_{n=1}^\infty \frac{1}{2^{a_n}}\]transcendental?

\noindent\rule{\linewidth}{0.4pt}


%% ============================================================
\section{Problem \#249}
\label{prob:249}

\noindent\textbf{Status:} OPEN \quad|\quad \textbf{Prize:} none \quad|\quad \textbf{Updated:} 2025-08-31

\noindent\textbf{Tags:} number theory, irrationality

\medskip

\noindent\textbf{Statement:}

Is\[\sum_n \frac{\phi(n)}{2^n}\]irrational? Here $\phi$ is the Euler totient function .

\noindent\rule{\linewidth}{0.4pt}


%% ============================================================
\section{Problem \#251}
\label{prob:251}

\noindent\textbf{Status:} OPEN \quad|\quad \textbf{Prize:} none \quad|\quad \textbf{Updated:} 2025-08-31

\noindent\textbf{Tags:} number theory, irrationality

\medskip

\noindent\textbf{Statement:}

Is\[\sum \frac{p_n}{2^n}\]irrational? (Here $p_n$ is the $n$th prime.)

\noindent\rule{\linewidth}{0.4pt}


%% ============================================================
\section{Problem \#252}
\label{prob:252}

\noindent\textbf{Status:} OPEN \quad|\quad \textbf{Prize:} none \quad|\quad \textbf{Updated:} 2025-08-31

\noindent\textbf{Tags:} number theory, irrationality

\medskip

\noindent\textbf{Statement:}

Let $k\geq 1$ and $\sigma_k(n)=\sum_{d\mid n}d^k$. Is\[\sum \frac{\sigma_k(n)}{n!}\]irrational?

\noindent\rule{\linewidth}{0.4pt}


%% ============================================================
\section{Problem \#254}
\label{prob:254}

\noindent\textbf{Status:} OPEN \quad|\quad \textbf{Prize:} none \quad|\quad \textbf{Updated:} 2025-08-31

\noindent\textbf{Tags:} number theory

\medskip

\noindent\textbf{Statement:}

Let $A\subseteq \mathbb{N}$ be such that\[\lvert A\cap [1,2x]\rvert -\lvert A\cap [1,x]\rvert \to \infty\textrm{ as }x\to \infty\]and\[\sum_{n\in A} \{ \theta n\}=\infty\]for every $\theta\in (0,1)$, where $\{x\}$ is the distance of $x$ from the nearest integer. Then every sufficiently large integer is the sum of distinct elements of $A$.

\noindent\rule{\linewidth}{0.4pt}


%% ============================================================
\section{Problem \#256}
\label{prob:256}

\noindent\textbf{Status:} OPEN \quad|\quad \textbf{Prize:} none \quad|\quad \textbf{Updated:} 2025-08-31

\noindent\textbf{Tags:} analysis

\medskip

\noindent\textbf{Statement:}

Let $n\geq 1$ and $f(n)$ be maximal such that for any integers $1\leq a_1\leq \cdots \leq a_n$ we have\[\max_{\lvert z\rvert=1}\left\lvert \prod_{i}(1-z^{a_i})\right\rvert\geq f(n).\]Estimate $f(n)$ - in particular, is it true that there exists some constant $c&#62;0$ such that\[\log f(n) \gg n^c?\]

\noindent\rule{\linewidth}{0.4pt}


%% ============================================================
\section{Problem \#257}
\label{prob:257}

\noindent\textbf{Status:} OPEN \quad|\quad \textbf{Prize:} none \quad|\quad \textbf{Updated:} 2025-08-31

\noindent\textbf{Tags:} irrationality

\medskip

\noindent\textbf{Statement:}

Let $A\subseteq \mathbb{N}$ be an infinite set. Is\[\sum_{n\in A}\frac{1}{2^n-1}\]irrational?

\noindent\rule{\linewidth}{0.4pt}


%% ============================================================
\section{Problem \#260}
\label{prob:260}

\noindent\textbf{Status:} OPEN \quad|\quad \textbf{Prize:} none \quad|\quad \textbf{Updated:} 2025-08-31

\noindent\textbf{Tags:} irrationality

\medskip

\noindent\textbf{Statement:}

Let $a_1&#60;a_2&#60;\cdots$ be an increasing sequence such that $a_n/n\to \infty$. Is the sum\[\sum_n \frac{a_n}{2^{a_n}}\]irrational?

\noindent\rule{\linewidth}{0.4pt}


%% ============================================================
\section{Problem \#261}
\label{prob:261}

\noindent\textbf{Status:} OPEN \quad|\quad \textbf{Prize:} none \quad|\quad \textbf{Updated:} 2025-08-31

\noindent\textbf{Tags:} number theory

\medskip

\noindent\textbf{Statement:}

Are there infinitely many $n$ such that there exists some $t\geq 2$ and distinct integers $a_1,\ldots,a_t\geq 1$ such that\[\frac{n}{2^n}=\sum_{1\leq k\leq t}\frac{a_k}{2^{a_k}}?\]Is this true for all $n$? Is there a rational $x$ such that\[x = \sum_{k=1}^\infty \frac{a_k}{2^{a_k}}\]has at least $2^{\aleph_0}$ solutions?

\noindent\rule{\linewidth}{0.4pt}


%% ============================================================
\section{Problem \#263}
\label{prob:263}

\noindent\textbf{Status:} OPEN \quad|\quad \textbf{Prize:} none \quad|\quad \textbf{Updated:} 2025-08-31

\noindent\textbf{Tags:} irrationality

\medskip

\noindent\textbf{Statement:}

Let $a_n$ be an increasing sequence of positive integers such that for every sequence of positive integers $b_n$ with $b_n/a_n\to 1$ the sum\[\sum\frac{1}{b_n}\]is irrational. Is $a_n=2^{2^n}$ such a sequence? Must such a sequence satisfy $a_n^{1/n}\to \infty$?

\noindent\rule{\linewidth}{0.4pt}


%% ============================================================
\section{Problem \#264}
\label{prob:264}

\noindent\textbf{Status:} OPEN \quad|\quad \textbf{Prize:} none \quad|\quad \textbf{Updated:} 2025-08-31

\noindent\textbf{Tags:} irrationality

\medskip

\noindent\textbf{Statement:}

Let $a_n$ be a sequence of positive integers such that for every bounded sequence of integers $b_n$ (with $a_n+b_n\neq 0$ and $b_n\neq 0$ for all $n$) the sum\[\sum \frac{1}{a_n+b_n}\]is irrational. Are $a_n=2^n$ or $a_n=n!$ examples of such a sequence?

\noindent\rule{\linewidth}{0.4pt}


%% ============================================================
\section{Problem \#265}
\label{prob:265}

\noindent\textbf{Status:} OPEN \quad|\quad \textbf{Prize:} none \quad|\quad \textbf{Updated:} 2025-08-31

\noindent\textbf{Tags:} irrationality

\medskip

\noindent\textbf{Statement:}

Let $1\leq a_1&#60;a_2&#60;\cdots$ be an increasing sequence of integers. How fast can $a_n\to \infty$ grow if\[\sum\frac{1}{a_n}\quad\textrm{and}\quad\sum\frac{1}{a_n-1}\]are both rational?

\noindent\rule{\linewidth}{0.4pt}


%% ============================================================
\section{Problem \#267}
\label{prob:267}

\noindent\textbf{Status:} OPEN \quad|\quad \textbf{Prize:} none \quad|\quad \textbf{Updated:} 2025-08-31

\noindent\textbf{Tags:} irrationality

\medskip

\noindent\textbf{Statement:}

Let $F_1=F_2=1$ and $F_{n+1}=F_n+F_{n-1}$ be the Fibonacci sequence. Let $n_1&lt;n_2&lt;\cdots $ be an infinite sequence with $n_{k+1}/n_k \geq c&gt;1$. Must\[\sum_k\frac{1}{F_{n_k}}\]be irrational?

\noindent\rule{\linewidth}{0.4pt}


%% ============================================================
\section{Problem \#269}
\label{prob:269}

\noindent\textbf{Status:} OPEN \quad|\quad \textbf{Prize:} none \quad|\quad \textbf{Updated:} 2025-08-31

\noindent\textbf{Tags:} irrationality

\medskip

\noindent\textbf{Statement:}

Let $P$ be a finite set of primes with $\lvert P\rvert \geq 2$ and let $\{a_1&#60;a_2&#60;\cdots\}=\{ n\in \mathbb{N} : \textrm{if }p\mid n\textrm{ then }p\in P\}$. Is the sum\[\sum_{n=1}^\infty \frac{1}{[a_1,\ldots,a_n]},\]where $[a_1,\ldots,a_n]$ is the lowest common multiple of $a_1,\ldots,a_n$, irrational?

\noindent\rule{\linewidth}{0.4pt}


%% ============================================================
\section{Problem \#271}
\label{prob:271}

\noindent\textbf{Status:} OPEN \quad|\quad \textbf{Prize:} none \quad|\quad \textbf{Updated:} 2025-08-31

\noindent\textbf{Tags:} additive combinatorics, arithmetic progressions

\medskip
\noindent\textit{Stanley sequences}


\noindent\textbf{Statement:}

Let $A(n)=\{a_0&#60;a_1&#60;\cdots\}$ be the sequence defined by $a_0=0$ and $a_1=n$, and for $k\geq 1$ define $a_{k+1}$ as the least positive integer such that there is no three-term arithmetic progression in $\{a_0,\ldots,a_{k+1}\}$. Can the $a_k$ be explicitly determined? How fast do they grow?

\noindent\rule{\linewidth}{0.4pt}


%% ============================================================
\section{Problem \#272}
\label{prob:272}

\noindent\textbf{Status:} OPEN \quad|\quad \textbf{Prize:} none \quad|\quad \textbf{Updated:} 2025-08-31

\noindent\textbf{Tags:} additive combinatorics, arithmetic progressions

\medskip

\noindent\textbf{Statement:}

Let $N\geq 1$. What is the largest $t$ such that there are $A_1,\ldots,A_t\subseteq \{1,\ldots,N\}$ with $A_i\cap A_j$ a non-empty arithmetic progression for all $i\neq j$?

\noindent\rule{\linewidth}{0.4pt}


%% ============================================================
\section{Problem \#273}
\label{prob:273}

\noindent\textbf{Status:} OPEN \quad|\quad \textbf{Prize:} none \quad|\quad \textbf{Updated:} 2025-08-31

\noindent\textbf{Tags:} number theory, covering systems

\medskip

\noindent\textbf{Statement:}

Is there a covering system all of whose moduli are of the form $p-1$ for some primes $p\geq 5$?

\noindent\rule{\linewidth}{0.4pt}


%% ============================================================
\section{Problem \#274}
\label{prob:274}

\noindent\textbf{Status:} OPEN \quad|\quad \textbf{Prize:} none \quad|\quad \textbf{Updated:} 2025-08-31

\noindent\textbf{Tags:} group theory, covering systems

\medskip
\noindent\textit{Herzog-Schönheim conjecture}


\noindent\textbf{Statement:}

If $G$ is a group then can there exist an exact covering of $G$ by more than one cosets of different sizes? (i.e. each element is contained in exactly one of the cosets)

\noindent\rule{\linewidth}{0.4pt}


%% ============================================================
\section{Problem \#276}
\label{prob:276}

\noindent\textbf{Status:} OPEN \quad|\quad \textbf{Prize:} none \quad|\quad \textbf{Updated:} 2025-08-31

\noindent\textbf{Tags:} number theory, covering systems

\medskip

\noindent\textbf{Statement:}

Is there an infinite Lucas sequence $a_0,a_1,\ldots$ where $a_{n+2}=a_{n+1}+a_n$ for $n\geq 0$ such that all $a_k$ are composite, and yet no integer has a common factor with every term of the sequence?

\noindent\rule{\linewidth}{0.4pt}


%% ============================================================
\section{Problem \#278}
\label{prob:278}

\noindent\textbf{Status:} OPEN \quad|\quad \textbf{Prize:} none \quad|\quad \textbf{Updated:} 2025-08-31

\noindent\textbf{Tags:} number theory, covering systems

\medskip

\noindent\textbf{Statement:}

Let $A=\{n_1&#60;\cdots&#60;n_r\}$ be a finite set of positive integers. What is the maximum density of integers covered by a suitable choice of congruences $a_i\pmod{n_i}$? Is the minimum density achieved when all the $a_i$ are equal?

\noindent\rule{\linewidth}{0.4pt}


%% ============================================================
\section{Problem \#279}
\label{prob:279}

\noindent\textbf{Status:} OPEN \quad|\quad \textbf{Prize:} none \quad|\quad \textbf{Updated:} 2025-08-31

\noindent\textbf{Tags:} number theory, covering systems, primes

\medskip

\noindent\textbf{Statement:}

Let $k\geq 3$. Is there a choice of congruence classes $a_p\pmod{p}$ for every prime $p$ such that all sufficiently large integers can be written as $a_p+tp$ for some prime $p$ and integer $t\geq k$?

\noindent\rule{\linewidth}{0.4pt}


%% ============================================================
\section{Problem \#282}
\label{prob:282}

\noindent\textbf{Status:} OPEN \quad|\quad \textbf{Prize:} none \quad|\quad \textbf{Updated:} 2025-08-31

\noindent\textbf{Tags:} number theory, unit fractions

\medskip

\noindent\textbf{Statement:}

Let $A\subseteq \mathbb{N}$ be an infinite set and consider the following greedy algorithm for a rational $x\in (0,1)$: choose the minimal $n\in A$ such that $n\geq 1/x$ and repeat with $x$ replaced by $x-\frac{1}{n}$. If this terminates after finitely many steps then this produces a representation of $x$ as the sum of distinct unit fractions with denominators from $A$. Does this process always terminate if $x$ has odd denominator and $A$ is the set of odd numbers? More generally, for which pairs $x$ and $A$ does this process terminate?

\noindent\rule{\linewidth}{0.4pt}


%% ============================================================
\section{Problem \#288}
\label{prob:288}

\noindent\textbf{Status:} OPEN \quad|\quad \textbf{Prize:} none \quad|\quad \textbf{Updated:} 2025-08-31

\noindent\textbf{Tags:} number theory, unit fractions

\medskip

\noindent\textbf{Statement:}

Is it true that there are only finitely many pairs of intervals $I_1,I_2$ such that\[\sum_{n_1\in I_1}\frac{1}{n_1}+\sum_{n_2\in I_2}\frac{1}{n_2}\in \mathbb{N}?\]

\noindent\rule{\linewidth}{0.4pt}


%% ============================================================
\section{Problem \#289}
\label{prob:289}

\noindent\textbf{Status:} OPEN \quad|\quad \textbf{Prize:} none \quad|\quad \textbf{Updated:} 2025-08-31

\noindent\textbf{Tags:} number theory, unit fractions

\medskip

\noindent\textbf{Statement:}

Is it true that, for all sufficiently large $k$, there exist finite intervals $I_1,\ldots,I_k\subset \mathbb{N}$, distinct, not overlapping or adjacent, with $\lvert I_i\rvert \geq 2$ for $1\leq i\leq k$ such that\[1=\sum_{i=1}^k \sum_{n\in I_i}\frac{1}{n}?\]

\noindent\rule{\linewidth}{0.4pt}


%% ============================================================
\section{Problem \#291}
\label{prob:291}

\noindent\textbf{Status:} OPEN \quad|\quad \textbf{Prize:} none \quad|\quad \textbf{Updated:} 2025-08-31

\noindent\textbf{Tags:} number theory, unit fractions

\medskip

\noindent\textbf{Statement:}

Let $n\geq 1$ and define $L_n$ to be the least common multiple of $\{1,\ldots,n\}$ and $a_n$ by\[\sum_{1\leq k\leq n}\frac{1}{k}=\frac{a_n}{L_n}.\]Is it true that $(a_n,L_n)=1$ and $(a_n,L_n)&#62;1$ both occur for infinitely many $n$?

\noindent\rule{\linewidth}{0.4pt}


%% ============================================================
\section{Problem \#293}
\label{prob:293}

\noindent\textbf{Status:} OPEN \quad|\quad \textbf{Prize:} none \quad|\quad \textbf{Updated:} 2025-08-31

\noindent\textbf{Tags:} number theory, unit fractions

\medskip

\noindent\textbf{Statement:}

Let $k\geq 1$ and let $v(k)$ be the minimal integer which does not appear as some $n_i$ in a solution to\[1=\frac{1}{n_1}+\cdots+\frac{1}{n_k}\]with $1\leq n_1&#60;\cdots &#60;n_k$. Estimate the growth of $v(k)$.

\noindent\rule{\linewidth}{0.4pt}


%% ============================================================
\section{Problem \#295}
\label{prob:295}

\noindent\textbf{Status:} OPEN \quad|\quad \textbf{Prize:} none \quad|\quad \textbf{Updated:} 2025-08-31

\noindent\textbf{Tags:} number theory, unit fractions

\medskip

\noindent\textbf{Statement:}

Let $N\geq 1$ and let $k(N)$ denote the smallest $k$ such that there exist $N\leq n_1&#60;\cdots &#60;n_k$ with\[1=\frac{1}{n_1}+\cdots+\frac{1}{n_k}.\]Is it true that\[\lim_{N\to \infty} k(N)-(e-1)N=\infty?\]

\noindent\rule{\linewidth}{0.4pt}


%% ============================================================
\section{Problem \#301}
\label{prob:301}

\noindent\textbf{Status:} OPEN \quad|\quad \textbf{Prize:} none \quad|\quad \textbf{Updated:} 2025-08-31

\noindent\textbf{Tags:} number theory, unit fractions

\medskip

\noindent\textbf{Statement:}

Let $f(N)$ be the size of the largest $A\subseteq \{1,\ldots,N\}$ such that there are no solutions to\[\frac{1}{a}= \frac{1}{b_1}+\cdots+\frac{1}{b_k}\]with distinct $a,b_1,\ldots,b_k\in A$? Estimate $f(N)$. In particular, is it true that $f(N)=(\tfrac{1}{2}+o(1))N$?

\noindent\rule{\linewidth}{0.4pt}


%% ============================================================
\section{Problem \#302}
\label{prob:302}

\noindent\textbf{Status:} OPEN \quad|\quad \textbf{Prize:} none \quad|\quad \textbf{Updated:} 2025-08-31

\noindent\textbf{Tags:} number theory, unit fractions

\medskip

\noindent\textbf{Statement:}

Let $f(N)$ be the size of the largest $A\subseteq \{1,\ldots,N\}$ such that there are no solutions to\[\frac{1}{a}= \frac{1}{b}+\frac{1}{c}\]with distinct $a,b,c\in A$? Estimate $f(N)$. In particular, is $f(N)=(\tfrac{1}{2}+o(1))N$?

\noindent\rule{\linewidth}{0.4pt}


%% ============================================================
\section{Problem \#304}
\label{prob:304}

\noindent\textbf{Status:} OPEN \quad|\quad \textbf{Prize:} none \quad|\quad \textbf{Updated:} 2025-08-31

\noindent\textbf{Tags:} number theory, unit fractions

\medskip

\noindent\textbf{Statement:}

For integers $1\leq a&#60;b$ let $N(a,b)$ denote the minimal $k$ such that there exist integers $1&#60;n_1&#60;\cdots&#60;n_k$ with\[\frac{a}{b}=\frac{1}{n_1}+\cdots+\frac{1}{n_k}.\]Estimate $N(b)=\max_{1\leq a&#60;b}N(a,b)$. Is it true that $N(b) \ll \log\log b$?

\noindent\rule{\linewidth}{0.4pt}


%% ============================================================
\section{Problem \#306}
\label{prob:306}

\noindent\textbf{Status:} OPEN \quad|\quad \textbf{Prize:} none \quad|\quad \textbf{Updated:} 2025-08-31

\noindent\textbf{Tags:} number theory, unit fractions

\medskip

\noindent\textbf{Statement:}

Let $a/b\in \mathbb{Q}_{&#62;0}$ with $b$ squarefree. Are there integers $1&#60;n_1&#60;\cdots&#60;n_k$, each the product of two distinct primes, such that\[\frac{a}{b}=\frac{1}{n_1}+\cdots+\frac{1}{n_k}?\]

\noindent\rule{\linewidth}{0.4pt}


%% ============================================================
\section{Problem \#311}
\label{prob:311}

\noindent\textbf{Status:} OPEN \quad|\quad \textbf{Prize:} none \quad|\quad \textbf{Updated:} 2025-08-31

\noindent\textbf{Tags:} number theory, unit fractions

\medskip

\noindent\textbf{Statement:}

Let $\delta(N)$ be the minimal non-zero value of $\lvert 1-\sum_{n\in A}\frac{1}{n}\rvert$ as $A$ ranges over all subsets of $\{1,\ldots,N\}$. Is it true that\[\delta(N)=e^{-(c+o(1))N}\]for some constant $c\in (0,1)$?

\noindent\rule{\linewidth}{0.4pt}


%% ============================================================
\section{Problem \#312}
\label{prob:312}

\noindent\textbf{Status:} OPEN \quad|\quad \textbf{Prize:} none \quad|\quad \textbf{Updated:} 2025-08-31

\noindent\textbf{Tags:} number theory, unit fractions

\medskip

\noindent\textbf{Statement:}

Does there exist some $c&#62;0$ such that, for any $K&#62;1$, whenever $A$ is a sufficiently large finite multiset of positive integers with $\sum_{n\in A}\frac{1}{n}&#62;K$ there exists some $S\subseteq A$ such that\[1-e^{-cK} &#60; \sum_{n\in S}\frac{1}{n}\leq 1?\]

\noindent\rule{\linewidth}{0.4pt}


%% ============================================================
\section{Problem \#313}
\label{prob:313}

\noindent\textbf{Status:} OPEN \quad|\quad \textbf{Prize:} none \quad|\quad \textbf{Updated:} 2025-08-31

\noindent\textbf{Tags:} number theory, unit fractions

\medskip

\noindent\textbf{Statement:}

Are there infinitely many solutions to\[\frac{1}{p_1}+\cdots+\frac{1}{p_k}=1-\frac{1}{m},\]where $m\geq 2$ is an integer and $p_1&#60;\cdots&#60;p_k$ are distinct primes?

\noindent\rule{\linewidth}{0.4pt}


%% ============================================================
\section{Problem \#317}
\label{prob:317}

\noindent\textbf{Status:} OPEN \quad|\quad \textbf{Prize:} none \quad|\quad \textbf{Updated:} 2025-08-31

\noindent\textbf{Tags:} number theory, unit fractions

\medskip

\noindent\textbf{Statement:}

Is there some constant $c&#62;0$ such that for every $n\geq 1$ there exists some $\delta_k\in \{-1,0,1\}$ for $1\leq k\leq n$ with\[0&lt; \left\lvert \sum_{1\leq k\leq n}\frac{\delta_k}{k}\right\rvert &lt; \frac{c}{2^n}?\]Is it true that for sufficiently large $n$, for any $\delta_k\in \{-1,0,1\}$,\[\left\lvert \sum_{1\leq k\leq n}\frac{\delta_k}{k}\right\rvert &gt; \frac{1}{[1,\ldots,n]}\]whenever the left-hand side is not zero?

\noindent\rule{\linewidth}{0.4pt}


%% ============================================================
\section{Problem \#319}
\label{prob:319}

\noindent\textbf{Status:} OPEN \quad|\quad \textbf{Prize:} none \quad|\quad \textbf{Updated:} 2025-08-31

\noindent\textbf{Tags:} number theory, unit fractions

\medskip

\noindent\textbf{Statement:}

What is the size of the largest $A\subseteq \{1,\ldots,N\}$ such that there is a function $\delta:A\to \{-1,1\}$ such that\[\sum_{n\in A}\frac{\delta_n}{n}=0\]and\[\sum_{n\in A'}\frac{\delta_n}{n}\neq 0\]for all non-empty $A'\subsetneq A$?

\noindent\rule{\linewidth}{0.4pt}


%% ============================================================
\section{Problem \#320}
\label{prob:320}

\noindent\textbf{Status:} OPEN \quad|\quad \textbf{Prize:} none \quad|\quad \textbf{Updated:} 2025-08-31

\noindent\textbf{Tags:} number theory, unit fractions

\medskip

\noindent\textbf{Statement:}

Let $S(N)$ count the number of distinct sums of the form $\sum_{n\in A}\frac{1}{n}$ for $A\subseteq \{1,\ldots,N\}$. Estimate $S(N)$.

\noindent\rule{\linewidth}{0.4pt}


%% ============================================================
\section{Problem \#321}
\label{prob:321}

\noindent\textbf{Status:} OPEN \quad|\quad \textbf{Prize:} none \quad|\quad \textbf{Updated:} 2025-08-31

\noindent\textbf{Tags:} number theory, unit fractions

\medskip

\noindent\textbf{Statement:}

What is the size of the largest $A\subseteq \{1,\ldots,N\}$ such that all sums $\sum_{n\in S}\frac{1}{n}$ are distinct for $S\subseteq A$?

\noindent\rule{\linewidth}{0.4pt}


%% ============================================================
\section{Problem \#322}
\label{prob:322}

\noindent\textbf{Status:} OPEN \quad|\quad \textbf{Prize:} none \quad|\quad \textbf{Updated:} 2025-08-31

\noindent\textbf{Tags:} number theory, powers

\medskip

\noindent\textbf{Statement:}

Let $k\geq 3$ and $A\subset \mathbb{N}$ be the set of $k$th powers. What is the order of growth of $1_A^{(k)}(n)$, i.e. the number of representations of $n$ as the sum of $k$ many $k$th powers? Does there exist some $c&#62;0$ and infinitely many $n$ such that\[1_A^{(k)}(n) &#62;n^c?\]

\noindent\rule{\linewidth}{0.4pt}


%% ============================================================
\section{Problem \#323}
\label{prob:323}

\noindent\textbf{Status:} OPEN \quad|\quad \textbf{Prize:} none \quad|\quad \textbf{Updated:} 2025-08-31

\noindent\textbf{Tags:} number theory, powers

\medskip

\noindent\textbf{Statement:}

Let $1\leq m\leq k$ and $f_{k,m}(x)$ denote the number of integers $\leq x$ which are the sum of $m$ many nonnegative $k$th powers. Is it true that\[f_{k,k}(x) \gg_\epsilon x^{1-\epsilon}\]for all $\epsilon&#62;0$? Is it true that if $m&#60;k$ then\[f_{k,m}(x) \gg x^{m/k}\]for sufficiently large $x$?

\noindent\rule{\linewidth}{0.4pt}


%% ============================================================
\section{Problem \#324}
\label{prob:324}

\noindent\textbf{Status:} OPEN \quad|\quad \textbf{Prize:} none \quad|\quad \textbf{Updated:} 2025-08-31

\noindent\textbf{Tags:} number theory, powers

\medskip

\noindent\textbf{Statement:}

Does there exist a polynomial $f(x)\in\mathbb{Z}[x]$ such that all the sums $f(a)+f(b)$ with $a&#60;b$ nonnegative integers are distinct?

\noindent\rule{\linewidth}{0.4pt}


%% ============================================================
\section{Problem \#325}
\label{prob:325}

\noindent\textbf{Status:} OPEN \quad|\quad \textbf{Prize:} none \quad|\quad \textbf{Updated:} 2025-08-31

\noindent\textbf{Tags:} number theory, powers

\medskip

\noindent\textbf{Statement:}

Let $k\geq 3$ and $f_{k,3}(x)$ denote the number of integers $\leq x$ which are the sum of three nonnegative $k$th powers. Is it true that\[f_{k,3}(x) \gg x^{3/k}\]or even $\gg_\epsilon x^{3/k-\epsilon}$?

\noindent\rule{\linewidth}{0.4pt}


%% ============================================================
\section{Problem \#326}
\label{prob:326}

\noindent\textbf{Status:} OPEN \quad|\quad \textbf{Prize:} none \quad|\quad \textbf{Updated:} 2025-08-31

\noindent\textbf{Tags:} number theory, additive basis

\medskip

\noindent\textbf{Statement:}

Does there exist $A=\{a_1&#60;a_2&#60;\cdots\}\subset \mathbb{N}$ which is a minimal basis of order $2$ (i.e. every large integer is the sum of $2$ elements from $A$, and no proper subset of $A$ has this property), such that\[\lim_{k\to \infty}\frac{a_k}{k^2}=c\]for some $c\neq 0$?

\noindent\rule{\linewidth}{0.4pt}


%% ============================================================
\section{Problem \#327}
\label{prob:327}

\noindent\textbf{Status:} OPEN \quad|\quad \textbf{Prize:} none \quad|\quad \textbf{Updated:} 2025-08-31

\noindent\textbf{Tags:} number theory, unit fractions

\medskip

\noindent\textbf{Statement:}

Suppose $A\subseteq \{1,\ldots,N\}$ is such that if $a,b\in A$ and $a\neq b$ then $a+b\nmid ab$. Can $A$ be 'substantially more' than the odd numbers? What if $a,b\in A$ with $a\neq b$ implies $a+b\nmid 2ab$? Must $\lvert A\rvert=o(N)$?

\noindent\rule{\linewidth}{0.4pt}


%% ============================================================
\section{Problem \#329}
\label{prob:329}

\noindent\textbf{Status:} OPEN \quad|\quad \textbf{Prize:} none \quad|\quad \textbf{Updated:} 2025-08-31

\noindent\textbf{Tags:} number theory

\medskip

\noindent\textbf{Statement:}

Suppose $A\subseteq \mathbb{N}$ is a Sidon set. How large can\[\limsup_{N\to \infty}\frac{\lvert A\cap \{1,\ldots,N\}\rvert}{N^{1/2}}\]be?

\noindent\rule{\linewidth}{0.4pt}


%% ============================================================
\section{Problem \#332}
\label{prob:332}

\noindent\textbf{Status:} OPEN \quad|\quad \textbf{Prize:} none \quad|\quad \textbf{Updated:} 2025-08-31

\noindent\textbf{Tags:} number theory

\medskip

\noindent\textbf{Statement:}

Let $A\subseteq \mathbb{N}$ and $D(A)$ be the set of those numbers which occur infinitely often as $a_1-a_2$ with $a_1,a_2\in A$. What conditions on $A$ are sufficient to ensure $D(A)$ has bounded gaps?

\noindent\rule{\linewidth}{0.4pt}


%% ============================================================
\section{Problem \#334}
\label{prob:334}

\noindent\textbf{Status:} OPEN \quad|\quad \textbf{Prize:} none \quad|\quad \textbf{Updated:} 2025-08-31

\noindent\textbf{Tags:} number theory

\medskip

\noindent\textbf{Statement:}

Find the best function $f(n)$ such that every $n$ can be written as $n=a+b$ where both $a,b$ are $f(n)$-smooth (that is, are not divisible by any prime $p&#62;f(n)$.)

\noindent\rule{\linewidth}{0.4pt}


%% ============================================================
\section{Problem \#335}
\label{prob:335}

\noindent\textbf{Status:} OPEN \quad|\quad \textbf{Prize:} none \quad|\quad \textbf{Updated:} 2025-08-31

\noindent\textbf{Tags:} number theory, additive combinatorics

\medskip

\noindent\textbf{Statement:}

Let $d(A)$ denote the density of $A\subseteq \mathbb{N}$. Characterise those $A,B\subseteq \mathbb{N}$ with positive density such that\[d(A+B)=d(A)+d(B).\]

\noindent\rule{\linewidth}{0.4pt}


%% ============================================================
\section{Problem \#336}
\label{prob:336}

\noindent\textbf{Status:} OPEN \quad|\quad \textbf{Prize:} none \quad|\quad \textbf{Updated:} 2025-08-31

\noindent\textbf{Tags:} number theory, additive basis

\medskip

\noindent\textbf{Statement:}

For $r\geq 2$ let $h(r)$ be the maximal finite $k$ such that there exists a basis $A\subseteq \mathbb{N}$ of order $r$ (so every large integer is the sum of at most $r$ integers from $A$) and exact order $k$ (so every large integer is the sum of exactly $k$ integers from $A$). Find the value of\[\lim_r \frac{h(r)}{r^2}.\]

\noindent\rule{\linewidth}{0.4pt}


%% ============================================================
\section{Problem \#338}
\label{prob:338}

\noindent\textbf{Status:} OPEN \quad|\quad \textbf{Prize:} none \quad|\quad \textbf{Updated:} 2025-08-31

\noindent\textbf{Tags:} number theory, additive basis

\medskip

\noindent\textbf{Statement:}

The restricted order of a basis is the least integer $t$ (if it exists) such that every large integer is the sum of at most $t$ distinct summands from $A$. What are necessary and sufficient conditions that this exists? Can it be bounded (when it exists) in terms of the order of the basis? What are necessary and sufficient conditions that this is equal to the order of the basis?

\noindent\rule{\linewidth}{0.4pt}


%% ============================================================
\section{Problem \#340}
\label{prob:340}

\noindent\textbf{Status:} OPEN \quad|\quad \textbf{Prize:} none \quad|\quad \textbf{Updated:} 2025-08-31

\noindent\textbf{Tags:} number theory, additive combinatorics, sidon sets

\medskip

\noindent\textbf{Statement:}

Let $A=\{1,2,4,8,13,21,31,45,66,81,97,\ldots\}$ be the greedy Sidon sequence: we begin with $1$ and iteratively include the next smallest integer that preserves the Sidon property (i.e. there are no non-trivial solutions to $a+b=c+d$). What is the order of growth of $A$? Is it true that\[\lvert A\cap \{1,\ldots,N\}\rvert \gg N^{1/2-\epsilon}\]for all $\epsilon&#62;0$ and large $N$?

\noindent\rule{\linewidth}{0.4pt}


%% ============================================================
\section{Problem \#341}
\label{prob:341}

\noindent\textbf{Status:} OPEN \quad|\quad \textbf{Prize:} none \quad|\quad \textbf{Updated:} 2025-08-31

\noindent\textbf{Tags:} number theory

\medskip

\noindent\textbf{Statement:}

Let $A=\{a_1&#60;\cdots&#60;a_k\}$ be a finite set of positive integers and extend it to an infinite sequence $\overline{A}=\{a_1&#60;a_2&#60;\cdots \}$ by defining $a_{n+1}$ for $n\geq k$ to be the least integer exceeding $a_n$ which is not of the form $a_i+a_j$ with $i,j\leq n$. Is it true that the sequence of differences $a_{m+1}-a_m$ is eventually periodic?

\noindent\rule{\linewidth}{0.4pt}


%% ============================================================
\section{Problem \#342}
\label{prob:342}

\noindent\textbf{Status:} OPEN \quad|\quad \textbf{Prize:} none \quad|\quad \textbf{Updated:} 2025-08-31

\noindent\textbf{Tags:} number theory

\medskip

\noindent\textbf{Statement:}

With $a_1=1$ and $a_2=2$ let $a_{n+1}$ for $n\geq 2$ be the least integer $&#62;a_n$ which can be expressed uniquely as $a_i+a_j$ for $i&#60;j\leq n$. What can be said about this sequence? Do infinitely many pairs $a,a+2$ occur? Does this sequence eventually have periodic differences? Is the density $0$?

\noindent\rule{\linewidth}{0.4pt}


%% ============================================================
\section{Problem \#345}
\label{prob:345}

\noindent\textbf{Status:} OPEN \quad|\quad \textbf{Prize:} none \quad|\quad \textbf{Updated:} 2025-08-31

\noindent\textbf{Tags:} number theory, complete sequences

\medskip

\noindent\textbf{Statement:}

Let $A\subseteq \mathbb{N}$ be a complete sequence, and define the threshold of completeness $T(A)$ to be the least integer $m$ such that all $n\geq m$ are in\[P(A) = \left\{\sum_{n\in B}n : B\subseteq A\textrm{ finite }\right\}\](the existence of $T(A)$ is guaranteed by completeness). Is it true that there are infinitely many $k$ such that $T(n^k)&#62;T(n^{k+1})$?

\noindent\rule{\linewidth}{0.4pt}


%% ============================================================
\section{Problem \#346}
\label{prob:346}

\noindent\textbf{Status:} OPEN \quad|\quad \textbf{Prize:} none \quad|\quad \textbf{Updated:} 2025-08-31

\noindent\textbf{Tags:} number theory, complete sequences

\medskip

\noindent\textbf{Statement:}

Let $A=\{1\leq a_1&lt; a_2&lt;\cdots\}$ be a set of integers such that {UL} {LI} $A\backslash B$ is complete for any finite subset $B$ and {/LI} {LI} $A\backslash B$ is not complete for any infinite subset $B$.{/LI} {/UL} (Here 'complete' means all sufficiently large integers can be written as a sum of distinct members of the sequence.) Is it true that if $a_{n+1}/a_n \geq 1+\epsilon$ for some $\epsilon&gt;0$ and all $n$ then\[\lim_n \frac{a_{n+1}}{a_n}=\frac{1+\sqrt{5}}{2}?\]

\noindent\rule{\linewidth}{0.4pt}


%% ============================================================
\section{Problem \#348}
\label{prob:348}

\noindent\textbf{Status:} OPEN \quad|\quad \textbf{Prize:} none \quad|\quad \textbf{Updated:} 2025-08-31

\noindent\textbf{Tags:} number theory, complete sequences

\medskip

\noindent\textbf{Statement:}

For what values of $0\leq m&#60;n$ is there a complete sequence $A=\{a_1\leq a_2\leq \cdots\}$ of integers such that {UL} {LI} $A$ remains complete after removing any $m$ elements, but {/LI} {LI} $A$ is not complete after removing any $n$ elements? {/LI} {/UL}

\noindent\rule{\linewidth}{0.4pt}


%% ============================================================
\section{Problem \#349}
\label{prob:349}

\noindent\textbf{Status:} OPEN \quad|\quad \textbf{Prize:} none \quad|\quad \textbf{Updated:} 2025-08-31

\noindent\textbf{Tags:} number theory, complete sequences

\medskip

\noindent\textbf{Statement:}

For what values of $t,\alpha \in (0,\infty)$ is the sequence $\lfloor t\alpha^n\rfloor$ complete (that is, all sufficiently large integers are the sum of distinct integers of the form $\lfloor t\alpha^n\rfloor$)?

\noindent\rule{\linewidth}{0.4pt}


%% ============================================================
\section{Problem \#352}
\label{prob:352}

\noindent\textbf{Status:} OPEN \quad|\quad \textbf{Prize:} none \quad|\quad \textbf{Updated:} 2025-08-31

\noindent\textbf{Tags:} geometry

\medskip

\noindent\textbf{Statement:}

Is there some $c&#62;0$ such that every measurable $A\subseteq \mathbb{R}^2$ of measure $\geq c$ contains the vertices of a triangle of area 1?

\noindent\rule{\linewidth}{0.4pt}


%% ============================================================
\section{Problem \#354}
\label{prob:354}

\noindent\textbf{Status:} OPEN \quad|\quad \textbf{Prize:} none \quad|\quad \textbf{Updated:} 2025-08-31

\noindent\textbf{Tags:} number theory

\medskip

\noindent\textbf{Statement:}

Let $\alpha,\beta\in \mathbb{R}_{&#62;0}$ such that $\alpha/\beta$ is irrational. Is the multiset\[\{ \lfloor \alpha\rfloor,\lfloor 2\alpha\rfloor,\lfloor 4\alpha\rfloor,\ldots\}\cup \{ \lfloor \beta\rfloor,\lfloor 2\beta\rfloor,\lfloor 4\beta\rfloor,\ldots\}\]complete? That is, can all sufficiently large natural numbers $n$ be written as\[n=\sum_{s\in S}\lfloor 2^s\alpha\rfloor+\sum_{t\in T}\lfloor 2^t\beta\rfloor\]for some finite $S,T\subset \mathbb{N}$? What if $2$ is replaced by some $\gamma\in(1,2)$?

\noindent\rule{\linewidth}{0.4pt}


%% ============================================================
\section{Problem \#357}
\label{prob:357}

\noindent\textbf{Status:} OPEN \quad|\quad \textbf{Prize:} none \quad|\quad \textbf{Updated:} 2025-08-31

\noindent\textbf{Tags:} number theory

\medskip

\noindent\textbf{Statement:}

Let $1\leq a_1&#60;\cdots &#60;a_k\leq n$ be integers such that all sums of the shape $\sum_{u\leq i\leq v}a_i$ are distinct. Let $f(n)$ be the maximal such $k$. How does $f(n)$ grow? Is $f(n)=o(n)$?

\noindent\rule{\linewidth}{0.4pt}


%% ============================================================
\section{Problem \#359}
\label{prob:359}

\noindent\textbf{Status:} OPEN \quad|\quad \textbf{Prize:} none \quad|\quad \textbf{Updated:} 2025-08-31

\noindent\textbf{Tags:} number theory

\medskip
\noindent\textit{segmented numbers}


\noindent\textbf{Statement:}

Let $a_1&lt;a_2&lt;\cdots$ be an infinite sequence of integers such that $a_1=n$ and $a_{i+1}$ is the least integer which is not a sum of consecutive earlier $a_j$s. What can be said about the density of this sequence? In particular, in the case $n=1$, can one prove that $a_k/k\to \infty$ and $a_k/k^{1+c}\to 0$ for any $c&gt;0$?

\noindent\rule{\linewidth}{0.4pt}


%% ============================================================
\section{Problem \#361}
\label{prob:361}

\noindent\textbf{Status:} OPEN \quad|\quad \textbf{Prize:} none \quad|\quad \textbf{Updated:} 2025-08-31

\noindent\textbf{Tags:} number theory

\medskip

\noindent\textbf{Statement:}

Let $c&#62;0$ and $n$ be some large integer. What is the size of the largest $A\subseteq \{1,\ldots,\lfloor cn\rfloor\}$ such that $n$ is not a sum of a subset of $A$? Does this depend on $n$ in an irregular way?

\noindent\rule{\linewidth}{0.4pt}


%% ============================================================
\section{Problem \#365}
\label{prob:365}

\noindent\textbf{Status:} OPEN \quad|\quad \textbf{Prize:} none \quad|\quad \textbf{Updated:} 2025-09-20

\noindent\textbf{Tags:} number theory

\medskip

\noindent\textbf{Statement:}

Do all pairs of consecutive powerful numbers $n$ and $n+1$ come from solutions to Pell equations ? In other words, must either $n$ or $n+1$ be a square? Is the number of such $n\leq x$ bounded by $(\log x)^{O(1)}$?

\noindent\rule{\linewidth}{0.4pt}


%% ============================================================
\section{Problem \#367}
\label{prob:367}

\noindent\textbf{Status:} OPEN \quad|\quad \textbf{Prize:} none \quad|\quad \textbf{Updated:} 2025-08-31

\noindent\textbf{Tags:} number theory

\medskip

\noindent\textbf{Statement:}

Let $B_2(n)$ be the $2$-full part of $n$ (that is, $B_2(n)=n/n'$ where $n'$ is the product of all primes that divide $n$ exactly once). Is it true that, for every fixed $k\geq 1$,\[\prod_{n\leq m&#60;n+k}B_2(m) \ll n^{2+o(1)}?\]Or perhaps even $\ll_k n^2$?

\noindent\rule{\linewidth}{0.4pt}


%% ============================================================
\section{Problem \#368}
\label{prob:368}

\noindent\textbf{Status:} OPEN \quad|\quad \textbf{Prize:} none \quad|\quad \textbf{Updated:} 2025-08-31

\noindent\textbf{Tags:} number theory

\medskip

\noindent\textbf{Statement:}

How large is the largest prime factor of $n(n+1)$?

\noindent\rule{\linewidth}{0.4pt}


%% ============================================================
\section{Problem \#371}
\label{prob:371}

\noindent\textbf{Status:} OPEN \quad|\quad \textbf{Prize:} none \quad|\quad \textbf{Updated:} 2025-08-31

\noindent\textbf{Tags:} number theory

\medskip

\noindent\textbf{Statement:}

Let $P(n)$ denote the largest prime factor of $n$. Show that the set of $n$ with $P(n)&#60;P(n+1)$ has density $1/2$.

\noindent\rule{\linewidth}{0.4pt}


%% ============================================================
\section{Problem \#373}
\label{prob:373}

\noindent\textbf{Status:} OPEN \quad|\quad \textbf{Prize:} none \quad|\quad \textbf{Updated:} 2025-08-31

\noindent\textbf{Tags:} number theory, factorials

\medskip

\noindent\textbf{Statement:}

Show that the equation\[n! = a_1!a_2!\cdots a_k!,\]with $n-1&#62;a_1\geq a_2\geq \cdots \geq a_k\geq 2$, has only finitely many solutions.

\noindent\rule{\linewidth}{0.4pt}


%% ============================================================
\section{Problem \#374}
\label{prob:374}

\noindent\textbf{Status:} OPEN \quad|\quad \textbf{Prize:} none \quad|\quad \textbf{Updated:} 2025-08-31

\noindent\textbf{Tags:} number theory

\medskip

\noindent\textbf{Statement:}

For any $m\in \mathbb{N}$, let $F(m)$ be the minimal $k\geq 2$ (if it exists) such that there are $a_1&#60;\cdots &#60;a_k=m$ with $a_1!\cdots a_k!$ a square. Let $D_k=\{ m : F(m)=k\}$. What is the order of growth of $\lvert D_k\cap\{1,\ldots,n\}\rvert$ for $3\leq k\leq 6$? For example, is it true that $\lvert D_6\cap \{1,\ldots,n\}\rvert \gg n$?

\noindent\rule{\linewidth}{0.4pt}


%% ============================================================
\section{Problem \#376}
\label{prob:376}

\noindent\textbf{Status:} OPEN \quad|\quad \textbf{Prize:} none \quad|\quad \textbf{Updated:} 2025-08-31

\noindent\textbf{Tags:} number theory, binomial coefficients, base representations

\medskip

\noindent\textbf{Statement:}

Are there infinitely many $n$ such that $\binom{2n}{n}$ is coprime to $105$?

\noindent\rule{\linewidth}{0.4pt}


%% ============================================================
\section{Problem \#377}
\label{prob:377}

\noindent\textbf{Status:} OPEN \quad|\quad \textbf{Prize:} none \quad|\quad \textbf{Updated:} 2025-08-31

\noindent\textbf{Tags:} number theory, binomial coefficients

\medskip

\noindent\textbf{Statement:}

Is there some absolute constant $C&#62;0$ such that\[\sum_{p\leq n}1_{p\nmid \binom{2n}{n}}\frac{1}{p}\leq C\]for all $n$ (where the summation is restricted to primes $p\leq n$)?

\noindent\rule{\linewidth}{0.4pt}


%% ============================================================
\section{Problem \#382}
\label{prob:382}

\noindent\textbf{Status:} OPEN \quad|\quad \textbf{Prize:} none \quad|\quad \textbf{Updated:} 2025-08-31

\noindent\textbf{Tags:} number theory

\medskip

\noindent\textbf{Statement:}

Let $u\leq v$ be such that the largest prime dividing $\prod_{u\leq m\leq v}m$ appears with exponent at least $2$. Is it true that $v-u=v^{o(1)}$? Can $v-u$ be arbitrarily large?

\noindent\rule{\linewidth}{0.4pt}


%% ============================================================
\section{Problem \#383}
\label{prob:383}

\noindent\textbf{Status:} OPEN \quad|\quad \textbf{Prize:} none \quad|\quad \textbf{Updated:} 2025-08-31

\noindent\textbf{Tags:} number theory

\medskip

\noindent\textbf{Statement:}

Is it true that for every $k$ there are infinitely many primes $p$ such that the largest prime divisor of\[\prod_{0\leq i\leq k}(p^2+i)\]is $p$?

\noindent\rule{\linewidth}{0.4pt}


%% ============================================================
\section{Problem \#385}
\label{prob:385}

\noindent\textbf{Status:} OPEN \quad|\quad \textbf{Prize:} none \quad|\quad \textbf{Updated:} 2025-08-31

\noindent\textbf{Tags:} number theory

\medskip

\noindent\textbf{Statement:}

Let\[F(n) = \max_{\substack{m&lt;n\\ m\textrm{ composite}}} m+p(m),\]where $p(m)$ is the least prime divisor of $m$. Is it true that $F(n)&gt;n$ for all sufficiently large $n$? Does $F(n)-n\to \infty$ as $n\to\infty$?

\noindent\rule{\linewidth}{0.4pt}


%% ============================================================
\section{Problem \#386}
\label{prob:386}

\noindent\textbf{Status:} OPEN \quad|\quad \textbf{Prize:} none \quad|\quad \textbf{Updated:} 2025-08-31

\noindent\textbf{Tags:} number theory, binomial coefficients

\medskip

\noindent\textbf{Statement:}

Let $2\leq k\leq n-2$. Can $\binom{n}{k}$ be the product of consecutive primes infinitely often? For example\[\binom{21}{2}=2\cdot 3\cdot 5\cdot 7.\]

\noindent\rule{\linewidth}{0.4pt}


%% ============================================================
\section{Problem \#387}
\label{prob:387}

\noindent\textbf{Status:} OPEN \quad|\quad \textbf{Prize:} none \quad|\quad \textbf{Updated:} 2025-08-31

\noindent\textbf{Tags:} number theory, binomial coefficients

\medskip

\noindent\textbf{Statement:}

Is there an absolute constant $c&#62;0$ such that, for all $1\leq k&#60; n$, the binomial coefficient $\binom{n}{k}$ has a divisor in $(cn,n]$?

\noindent\rule{\linewidth}{0.4pt}


%% ============================================================
\section{Problem \#388}
\label{prob:388}

\noindent\textbf{Status:} OPEN \quad|\quad \textbf{Prize:} none \quad|\quad \textbf{Updated:} 2025-08-31

\noindent\textbf{Tags:} number theory

\medskip

\noindent\textbf{Statement:}

Can one classify all solutions of\[\prod_{1\leq i\leq k_1}(m_1+i)=\prod_{1\leq j\leq k_2}(m_2+j)\]where $k_1,k_2&#62;3$ and $m_1+k_1\leq m_2$? Are there only finitely many solutions?

\noindent\rule{\linewidth}{0.4pt}


%% ============================================================
\section{Problem \#389}
\label{prob:389}

\noindent\textbf{Status:} OPEN \quad|\quad \textbf{Prize:} none \quad|\quad \textbf{Updated:} 2025-08-31

\noindent\textbf{Tags:} number theory

\medskip

\noindent\textbf{Statement:}

Is it true that for every $n\geq 1$ there is a $k$ such that\[n(n+1)\cdots(n+k-1)\mid (n+k)\cdots (n+2k-1)?\]

\noindent\rule{\linewidth}{0.4pt}


%% ============================================================
\section{Problem \#390}
\label{prob:390}

\noindent\textbf{Status:} OPEN \quad|\quad \textbf{Prize:} none \quad|\quad \textbf{Updated:} 2025-08-31

\noindent\textbf{Tags:} number theory, factorials

\medskip

\noindent\textbf{Statement:}

Let $f(n)$ be the minimal $m$ such that\[n! = a_1\cdots a_k\]with $n&#60; a_1&#60;\cdots &#60;a_k=m$. Is there (and what is it) a constant $c$ such that\[f(n)-2n \sim c\frac{n}{\log n}?\]

\noindent\rule{\linewidth}{0.4pt}


%% ============================================================
\section{Problem \#393}
\label{prob:393}

\noindent\textbf{Status:} OPEN \quad|\quad \textbf{Prize:} none \quad|\quad \textbf{Updated:} 2025-08-31

\noindent\textbf{Tags:} number theory, factorials

\medskip

\noindent\textbf{Statement:}

Let $f(n)$ denote the minimal $m\geq 1$ such that\[n! = a_1\cdots a_t\]with $a_1&#60;\cdots &#60;a_t=a_1+m$. What is the behaviour of $f(n)$?

\noindent\rule{\linewidth}{0.4pt}


%% ============================================================
\section{Problem \#394}
\label{prob:394}

\noindent\textbf{Status:} OPEN \quad|\quad \textbf{Prize:} none \quad|\quad \textbf{Updated:} 2025-10-28

\noindent\textbf{Tags:} number theory

\medskip

\noindent\textbf{Statement:}

Let $t_k(n)$ denote the least $m$ such that\[n\mid m(m+1)(m+2)\cdots (m+k-1).\]Is it true that\[\sum_{n\leq x}t_2(n)\ll \frac{x^2}{(\log x)^c}\]for some $c&#62;0$? Is it true that, for $k\geq 2$,\[\sum_{n\leq x}t_{k+1}(n) =o\left(\sum_{n\leq x}t_k(n)\right)?\]

\noindent\rule{\linewidth}{0.4pt}


%% ============================================================
\section{Problem \#396}
\label{prob:396}

\noindent\textbf{Status:} OPEN \quad|\quad \textbf{Prize:} none \quad|\quad \textbf{Updated:} 2025-08-31

\noindent\textbf{Tags:} number theory, binomial coefficients

\medskip

\noindent\textbf{Statement:}

Is it true that for every $k$ there exists $n$ such that\[\prod_{0\leq i\leq k}(n-i) \mid \binom{2n}{n}?\]

\noindent\rule{\linewidth}{0.4pt}


%% ============================================================
\section{Problem \#400}
\label{prob:400}

\noindent\textbf{Status:} OPEN \quad|\quad \textbf{Prize:} none \quad|\quad \textbf{Updated:} 2025-08-31

\noindent\textbf{Tags:} number theory, factorials

\medskip

\noindent\textbf{Statement:}

For any $k\geq 2$ let $g_k(n)$ denote the maximum value of\[(a_1+\cdots+a_k)-n\]where $a_1,\ldots,a_k$ are integers such that $a_1!\cdots a_k! \mid n!$. Can one show that\[\sum_{n\leq x}g_k(n) \sim c_k x\log x\]for some constant $c_k$? Is it true that there is a constant $c_k$ such that for almost all $n&#60;x$ we have\[g_k(n)=c_k\log x+o(\log x)?\]

\noindent\rule{\linewidth}{0.4pt}


%% ============================================================
\section{Problem \#404}
\label{prob:404}

\noindent\textbf{Status:} OPEN \quad|\quad \textbf{Prize:} none \quad|\quad \textbf{Updated:} 2025-08-31

\noindent\textbf{Tags:} number theory, factorials

\medskip

\noindent\textbf{Statement:}

For which integers $a\geq 1$ and primes $p$ is there a finite upper bound on those $k$ such that there are $a=a_1&#60;\cdots&#60;a_n$ with\[p^k \mid (a_1!+\cdots+a_n!)?\]If $f(a,p)$ is the greatest such $k$, how does this function behave? Is there a prime $p$ and an infinite sequence $a_1&#60;a_2&#60;\cdots$ such that if $p^{m_k}$ is the highest power of $p$ dividing $\sum_{i\leq k}a_i!$ then $m_k\to \infty$?

\noindent\rule{\linewidth}{0.4pt}


%% ============================================================
\section{Problem \#406}
\label{prob:406}

\noindent\textbf{Status:} OPEN \quad|\quad \textbf{Prize:} none \quad|\quad \textbf{Updated:} 2025-08-31

\noindent\textbf{Tags:} number theory, base representations

\medskip

\noindent\textbf{Statement:}

Is it true that there are only finitely many powers of $2$ which have only the digits $0$ and $1$ when written in base $3$?

\noindent\rule{\linewidth}{0.4pt}


%% ============================================================
\section{Problem \#408}
\label{prob:408}

\noindent\textbf{Status:} OPEN \quad|\quad \textbf{Prize:} none \quad|\quad \textbf{Updated:} 2025-08-31

\noindent\textbf{Tags:} number theory, iterated functions

\medskip

\noindent\textbf{Statement:}

Let $\phi(n)$ be the Euler totient function and $\phi_k(n)$ be the iterated $\phi$ function, so that $\phi_1(n)=\phi(n)$ and $\phi_k(n)=\phi(\phi_{k-1}(n))$. Let\[f(n) = \min \{ k : \phi_k(n)=1\}.\]Does $f(n)/\log n$ have a distribution function? Is $f(n)/\log n$ almost always constant? What can be said about the largest prime factor of $\phi_k(n)$ when, say, $k=\log\log n$?

\noindent\rule{\linewidth}{0.4pt}


%% ============================================================
\section{Problem \#409}
\label{prob:409}

\noindent\textbf{Status:} OPEN \quad|\quad \textbf{Prize:} none \quad|\quad \textbf{Updated:} 2025-08-31

\noindent\textbf{Tags:} number theory, iterated functions

\medskip

\noindent\textbf{Statement:}

How many iterations of $n\mapsto \phi(n)+1$ are needed before a prime is reached? Can infinitely many $n$ reach the same prime? What is the density of $n$ which reach any fixed prime?

\noindent\rule{\linewidth}{0.4pt}


%% ============================================================
\section{Problem \#410}
\label{prob:410}

\noindent\textbf{Status:} OPEN \quad|\quad \textbf{Prize:} none \quad|\quad \textbf{Updated:} 2025-08-31

\noindent\textbf{Tags:} number theory, iterated functions

\medskip

\noindent\textbf{Statement:}

Let $\sigma_1(n)=\sigma(n)$, the sum of divisors function, and $\sigma_k(n)=\sigma(\sigma_{k-1}(n))$. Is it true that for all $n\geq 2$\[\lim_{k\to \infty} \sigma_k(n)^{1/k}=\infty?\]

\noindent\rule{\linewidth}{0.4pt}


%% ============================================================
\section{Problem \#411}
\label{prob:411}

\noindent\textbf{Status:} OPEN \quad|\quad \textbf{Prize:} none \quad|\quad \textbf{Updated:} 2025-08-31

\noindent\textbf{Tags:} number theory, iterated functions

\medskip

\noindent\textbf{Statement:}

Let $g_1=g(n)=n+\phi(n)$ and $g_k(n)=g(g_{k-1}(n))$. For which $n$ and $r$ is it true that $g_{k+r}(n)=2g_k(n)$ for all large $k$?

\noindent\rule{\linewidth}{0.4pt}


%% ============================================================
\section{Problem \#412}
\label{prob:412}

\noindent\textbf{Status:} OPEN \quad|\quad \textbf{Prize:} none \quad|\quad \textbf{Updated:} 2025-08-31

\noindent\textbf{Tags:} number theory, iterated functions

\medskip

\noindent\textbf{Statement:}

Let $\sigma_1(n)=\sigma(n)$, the sum of divisors function, and $\sigma_k(n)=\sigma(\sigma_{k-1}(n))$. Is it true that, for every $m,n\geq 2$, there exist some $i,j$ such that $\sigma_i(m)=\sigma_j(n)$?

\noindent\rule{\linewidth}{0.4pt}


%% ============================================================
\section{Problem \#413}
\label{prob:413}

\noindent\textbf{Status:} OPEN \quad|\quad \textbf{Prize:} none \quad|\quad \textbf{Updated:} 2025-08-31

\noindent\textbf{Tags:} number theory, iterated functions

\medskip

\noindent\textbf{Statement:}

Let $\omega(n)$ count the number of distinct primes dividing $n$. Are there infinitely many $n$ such that, for all $m&lt;n$, we have $m+\omega(m) \leq n$? Can one show that there exists an $\epsilon&gt;0$ such that there are infinitely many $n$ where $m+\epsilon \omega(m)\leq n$ for all $m&#60;n$?

\noindent\rule{\linewidth}{0.4pt}


%% ============================================================
\section{Problem \#414}
\label{prob:414}

\noindent\textbf{Status:} OPEN \quad|\quad \textbf{Prize:} none \quad|\quad \textbf{Updated:} 2025-08-31

\noindent\textbf{Tags:} number theory, iterated functions

\medskip

\noindent\textbf{Statement:}

Let $h_1(n)=h(n)=n+\tau(n)$ (where $\tau(n)$ counts the number of divisors of $n$) and $h_k(n)=h(h_{k-1}(n))$. Is it true, for any $m,n$, there exist $i$ and $j$ such that $h_i(m)=h_j(n)$?

\noindent\rule{\linewidth}{0.4pt}


%% ============================================================
\section{Problem \#415}
\label{prob:415}

\noindent\textbf{Status:} OPEN \quad|\quad \textbf{Prize:} none \quad|\quad \textbf{Updated:} 2025-08-31

\noindent\textbf{Tags:} number theory

\medskip

\noindent\textbf{Statement:}

For any $n$ let $F(n)$ be the largest $k$ such that any of the $k!$ possible ordering patterns appears in some sequence of $\phi(m+1),\ldots,\phi(m+k)$ with $m+k\leq n$. Is it true that\[F(n)=(c+o(1))\log\log\log n\]for some constant $c$? Is the first pattern which fails to appear always\[\phi(m+1)&#62;\phi(m+2)&#62;\cdots &#62;\phi(m+k)?\]Is it true that the 'natural' ordering which mimics what happens to $\phi(1),\ldots,\phi(k)$ is the most likely to appear?

\noindent\rule{\linewidth}{0.4pt}


%% ============================================================
\section{Problem \#416}
\label{prob:416}

\noindent\textbf{Status:} OPEN \quad|\quad \textbf{Prize:} none \quad|\quad \textbf{Updated:} 2025-08-31

\noindent\textbf{Tags:} number theory

\medskip

\noindent\textbf{Statement:}

Let $V(x)$ count the number of $n\leq x$ such that $\phi(m)=n$ is solvable. Does $V(2x)/V(x)\to 2$? Is there an asymptotic formula for $V(x)$?

\noindent\rule{\linewidth}{0.4pt}


%% ============================================================
\section{Problem \#417}
\label{prob:417}

\noindent\textbf{Status:} OPEN \quad|\quad \textbf{Prize:} none \quad|\quad \textbf{Updated:} 2025-08-31

\noindent\textbf{Tags:} number theory

\medskip

\noindent\textbf{Statement:}

Let\[V'(x)=\#\{\phi(m) : 1\leq m\leq x\}\]and\[V(x)=\#\{\phi(m) \leq x : 1\leq m\}.\]Does $\lim V(x)/V'(x)$ exist? Is it $&#62;1$?

\noindent\rule{\linewidth}{0.4pt}


%% ============================================================
\section{Problem \#420}
\label{prob:420}

\noindent\textbf{Status:} OPEN \quad|\quad \textbf{Prize:} none \quad|\quad \textbf{Updated:} 2025-08-31

\noindent\textbf{Tags:} number theory

\medskip

\noindent\textbf{Statement:}

If $\tau(n)$ counts the number of divisors of $n$ then let\[F(f,n)=\frac{\tau((n+\lfloor f(n)\rfloor)!)}{\tau(n!)}.\]Is it true that\[\lim_{n\to \infty}F((\log n)^C,n)=\infty\]for large $C$? Is it true that $F(\log n,n)$ is everywhere dense in $(1,\infty)$? More generally, if $f(n)\leq \log n$ is a monotonic function such that $f(n)\to \infty$ as $n\to \infty$, then is $F(f,n)$ everywhere dense?

\noindent\rule{\linewidth}{0.4pt}


%% ============================================================
\section{Problem \#421}
\label{prob:421}

\noindent\textbf{Status:} OPEN \quad|\quad \textbf{Prize:} none \quad|\quad \textbf{Updated:} 2025-08-31

\noindent\textbf{Tags:} number theory

\medskip

\noindent\textbf{Statement:}

Is there a sequence $1\leq d_1&#60;d_2&#60;\cdots$ with density $1$ such that all products $\prod_{u\leq i\leq v}d_i$ are distinct?

\noindent\rule{\linewidth}{0.4pt}


%% ============================================================
\section{Problem \#422}
\label{prob:422}

\noindent\textbf{Status:} OPEN \quad|\quad \textbf{Prize:} none \quad|\quad \textbf{Updated:} 2025-08-31

\noindent\textbf{Tags:} number theory

\medskip

\noindent\textbf{Statement:}

Let $f(1)=f(2)=1$ and for $n&#62;2$\[f(n) = f(n-f(n-1))+f(n-f(n-2)).\]Does $f(n)$ miss infinitely many integers? What is its behaviour?

\noindent\rule{\linewidth}{0.4pt}


%% ============================================================
\section{Problem \#423}
\label{prob:423}

\noindent\textbf{Status:} OPEN \quad|\quad \textbf{Prize:} none \quad|\quad \textbf{Updated:} 2025-08-31

\noindent\textbf{Tags:} number theory

\medskip

\noindent\textbf{Statement:}

Let $a_1=1$ and $a_2=2$ and for $k\geq 3$ choose $a_k$ to be the least integer $&#62;a_{k-1}$ which is the sum of at least two consecutive terms of the sequence. What is the asymptotic behaviour of this sequence?

\noindent\rule{\linewidth}{0.4pt}


%% ============================================================
\section{Problem \#424}
\label{prob:424}

\noindent\textbf{Status:} OPEN \quad|\quad \textbf{Prize:} none \quad|\quad \textbf{Updated:} 2025-08-31

\noindent\textbf{Tags:} number theory

\medskip

\noindent\textbf{Statement:}

Let $a_1=2$ and $a_2=3$ and continue the sequence by appending to $a_1,\ldots,a_n$ all possible values of $a_ia_j-1$ with $i\neq j$. Is it true that the set of integers which eventually appear has positive density?

\noindent\rule{\linewidth}{0.4pt}


%% ============================================================
\section{Problem \#425}
\label{prob:425}

\noindent\textbf{Status:} OPEN \quad|\quad \textbf{Prize:} none \quad|\quad \textbf{Updated:} 2025-08-31

\noindent\textbf{Tags:} number theory, sidon sets

\medskip

\noindent\textbf{Statement:}

Let $F(n)$ be the maximum possible size of a subset $A\subseteq\{1,\ldots,N\}$ such that the products $ab$ are distinct for all $a&#60;b$. Is there a constant $c$ such that\[F(n)=\pi(n)+(c+o(1))n^{3/4}(\log n)^{-3/2}?\]If $A\subseteq \{1,\ldots,n\}$ is such that all products $a_1\cdots a_r$ are distinct for $a_1&#60;\cdots &#60;a_r$ then is it true that\[\lvert A\rvert \leq \pi(n)+O(n^{\frac{r+1}{2r}})?\]

\noindent\rule{\linewidth}{0.4pt}


%% ============================================================
\section{Problem \#428}
\label{prob:428}

\noindent\textbf{Status:} OPEN \quad|\quad \textbf{Prize:} none \quad|\quad \textbf{Updated:} 2025-08-31

\noindent\textbf{Tags:} number theory, primes

\medskip

\noindent\textbf{Statement:}

Is there a set $A\subseteq \mathbb{N}$ such that, for infinitely many $n$, all of $n-a$ are prime for all $a\in A$ with $0&lt;a&lt;n$ and\[\liminf\frac{\lvert A\cap [1,x]\rvert}{\pi(x)}&gt;0?\]

\noindent\rule{\linewidth}{0.4pt}


%% ============================================================
\section{Problem \#430}
\label{prob:430}

\noindent\textbf{Status:} OPEN \quad|\quad \textbf{Prize:} none \quad|\quad \textbf{Updated:} 2025-08-31

\noindent\textbf{Tags:} number theory

\medskip

\noindent\textbf{Statement:}

Fix some integer $n$ and define a decreasing sequence in $[1,n)$ by $a_1=n-1$ and, for $k\geq 2$, letting $a_k$ be the greatest integer in $[1,a_{k-1})$ such that all of the prime factors of $a_k$ are $&#62;n-a_k$. Is it true that, for sufficiently large $n$, not all of this sequence can be prime?

\noindent\rule{\linewidth}{0.4pt}


%% ============================================================
\section{Problem \#431}
\label{prob:431}

\noindent\textbf{Status:} OPEN \quad|\quad \textbf{Prize:} none \quad|\quad \textbf{Updated:} 2025-08-31

\noindent\textbf{Tags:} number theory, primes

\medskip
\noindent\textit{inverse Goldbach problem}


\noindent\textbf{Statement:}

Are there two infinite sets $A$ and $B$ such that $A+B$ agrees with the set of prime numbers up to finitely many exceptions?

\noindent\rule{\linewidth}{0.4pt}


%% ============================================================
\section{Problem \#432}
\label{prob:432}

\noindent\textbf{Status:} OPEN \quad|\quad \textbf{Prize:} none \quad|\quad \textbf{Updated:} 2025-08-31

\noindent\textbf{Tags:} number theory

\medskip

\noindent\textbf{Statement:}

Let $A,B\subseteq \mathbb{N}$ be two infinite sets. How dense can $A+B$ be if all elements of $A+B$ are pairwise relatively prime?

\noindent\rule{\linewidth}{0.4pt}


%% ============================================================
\section{Problem \#436}
\label{prob:436}

\noindent\textbf{Status:} OPEN \quad|\quad \textbf{Prize:} none \quad|\quad \textbf{Updated:} 2025-08-31

\noindent\textbf{Tags:} number theory

\medskip

\noindent\textbf{Statement:}

If $p$ is a prime and $k,m\geq 2$ then let $r(k,m,p)$ be the minimal $r$ such that $r,r+1,\ldots,r+m-1$ are all $k$th power residues modulo $p$. Let\[\Lambda(k,m)=\limsup_{p\to \infty} r(k,m,p).\]Is it true that $\Lambda(k,2)$ is finite for all $k$? Is $\Lambda(k,3)$ finite for all odd $k$? How large are they?

\noindent\rule{\linewidth}{0.4pt}


%% ============================================================
\section{Problem \#445}
\label{prob:445}

\noindent\textbf{Status:} OPEN \quad|\quad \textbf{Prize:} none \quad|\quad \textbf{Updated:} 2025-08-31

\noindent\textbf{Tags:} number theory

\medskip

\noindent\textbf{Statement:}

Is it true that, for any $c&#62;1/2$, if $p$ is a sufficiently large prime then, for any $n\geq 0$, there exist $a,b\in(n,n+p^c)$ such that $ab\equiv 1\pmod{p}$?

\noindent\rule{\linewidth}{0.4pt}


%% ============================================================
\section{Problem \#450}
\label{prob:450}

\noindent\textbf{Status:} OPEN \quad|\quad \textbf{Prize:} none \quad|\quad \textbf{Updated:} 2025-08-31

\noindent\textbf{Tags:} number theory, divisors

\medskip

\noindent\textbf{Statement:}

How large must $y=y(\epsilon,n)$ be such that the number of integers in $(x,x+y)$ with a divisor in $(n,2n)$ is at most $\epsilon y$?

\noindent\rule{\linewidth}{0.4pt}


%% ============================================================
\section{Problem \#451}
\label{prob:451}

\noindent\textbf{Status:} OPEN \quad|\quad \textbf{Prize:} none \quad|\quad \textbf{Updated:} 2025-08-31

\noindent\textbf{Tags:} number theory

\medskip

\noindent\textbf{Statement:}

Estimate $n_k$, the smallest integer $&#62;2k$ such that $\prod_{1\leq i\leq k}(n_k-i)$ has no prime factor in $(k,2k)$.

\noindent\rule{\linewidth}{0.4pt}


%% ============================================================
\section{Problem \#452}
\label{prob:452}

\noindent\textbf{Status:} OPEN \quad|\quad \textbf{Prize:} none \quad|\quad \textbf{Updated:} 2025-08-31

\noindent\textbf{Tags:} number theory

\medskip

\noindent\textbf{Statement:}

Let $\omega(n)$ count the number of distinct prime factors of $n$. What is the size of the largest interval $I\subseteq [x,2x]$ such that $\omega(n)&#62;\log\log n$ for all $n\in I$?

\noindent\rule{\linewidth}{0.4pt}


%% ============================================================
\section{Problem \#454}
\label{prob:454}

\noindent\textbf{Status:} OPEN \quad|\quad \textbf{Prize:} none \quad|\quad \textbf{Updated:} 2025-08-31

\noindent\textbf{Tags:} number theory, primes

\medskip

\noindent\textbf{Statement:}

Let\[f(n) = \min_{i&#60;n} (p_{n+i}+p_{n-i}),\]where $p_k$ is the $k$th prime. Is it true that\[\limsup_n (f(n)-2p_n)=\infty?\]

\noindent\rule{\linewidth}{0.4pt}


%% ============================================================
\section{Problem \#455}
\label{prob:455}

\noindent\textbf{Status:} OPEN \quad|\quad \textbf{Prize:} none \quad|\quad \textbf{Updated:} 2025-08-31

\noindent\textbf{Tags:} number theory

\medskip

\noindent\textbf{Statement:}

Let $q_1&#60;q_2&#60;\cdots$ be a sequence of primes such that\[q_{n+1}-q_n\geq q_n-q_{n-1}.\]Must\[\lim_n \frac{q_n}{n^2}=\infty?\]

\noindent\rule{\linewidth}{0.4pt}


%% ============================================================
\section{Problem \#456}
\label{prob:456}

\noindent\textbf{Status:} OPEN \quad|\quad \textbf{Prize:} none \quad|\quad \textbf{Updated:} 2025-08-31

\noindent\textbf{Tags:} number theory

\medskip

\noindent\textbf{Statement:}

Let $p_n$ be the smallest prime $\equiv 1\pmod{n}$ and let $m_n$ be the smallest integer such that $n\mid \phi(m_n)$. Is it true that $m_n&#60;p_n$ for almost all $n$? Does $p_n/m_n\to \infty$ for almost all $n$? Are there infinitely many primes $p$ such that $p-1$ is the only $n$ for which $m_n=p$?

\noindent\rule{\linewidth}{0.4pt}


%% ============================================================
\section{Problem \#460}
\label{prob:460}

\noindent\textbf{Status:} OPEN \quad|\quad \textbf{Prize:} none \quad|\quad \textbf{Updated:} 2025-08-31

\noindent\textbf{Tags:} number theory

\medskip
\noindent\textit{ambiguous statement}


\noindent\textbf{Statement:}

Let $a_0=0$ and $a_1=1$, and in general define $a_k$ to be the least integer $&#62;a_{k-1}$ for which $(n-a_k,n-a_i)=1$ for all $0\leq i&#60;k$. Does\[\sum_{0&#60;a_i&#60; n}\frac{1}{a_i}\to \infty\]as $n\to \infty$? What about if we restrict the sum to those $i$ such that $n-a_j$ is divisible by some prime $\leq a_j$, or the complement of such $i$?

\noindent\rule{\linewidth}{0.4pt}


%% ============================================================
\section{Problem \#461}
\label{prob:461}

\noindent\textbf{Status:} OPEN \quad|\quad \textbf{Prize:} none \quad|\quad \textbf{Updated:} 2025-08-31

\noindent\textbf{Tags:} number theory, primes

\medskip

\noindent\textbf{Statement:}

Let $s_t(n)$ be the $t$-smooth component of $n$ - that is, the product of all primes $p$ (with multiplicity) dividing $n$ such that $p&#60;t$. Let $f(n,t)$ count the number of distinct possible values for $s_t(m)$ for $m\in [n+1,n+t]$. Is it true that\[f(n,t)\gg t\](uniformly, for all $t$ and $n$)?

\noindent\rule{\linewidth}{0.4pt}


%% ============================================================
\section{Problem \#462}
\label{prob:462}

\noindent\textbf{Status:} OPEN \quad|\quad \textbf{Prize:} none \quad|\quad \textbf{Updated:} 2025-08-31

\noindent\textbf{Tags:} number theory, primes

\medskip

\noindent\textbf{Statement:}

Let $p(n)$ denote the least prime factor of $n$. There is a constant $c&#62;0$ such that\[\sum_{\substack{n&lt;x\\ n\textrm{ not prime}}}\frac{p(n)}{n}\sim c\frac{x^{1/2}}{(\log x)^2}.\]Is it true that there exists a constant $C&gt;0$ such that\[\sum_{x\leq n\leq x+Cx^{1/2}(\log x)^2}\frac{p(n)}{n} \gg 1\]for all large $x$?

\noindent\rule{\linewidth}{0.4pt}


%% ============================================================
\section{Problem \#463}
\label{prob:463}

\noindent\textbf{Status:} OPEN \quad|\quad \textbf{Prize:} none \quad|\quad \textbf{Updated:} 2025-08-31

\noindent\textbf{Tags:} number theory, primes

\medskip

\noindent\textbf{Statement:}

Is there a function $f$ with $f(n)\to \infty$ as $n\to \infty$ such that, for all large $n$, there is a composite number $m$ such that\[n+f(n)&#60;m&#60;n+p(m)?\](Here $p(m)$ is the least prime factor of $m$.)

\noindent\rule{\linewidth}{0.4pt}


%% ============================================================
\section{Problem \#467}
\label{prob:467}

\noindent\textbf{Status:} OPEN \quad|\quad \textbf{Prize:} none \quad|\quad \textbf{Updated:} 2025-08-31

\noindent\textbf{Tags:} number theory

\medskip
\noindent\textit{ambiguous statement}


\noindent\textbf{Statement:}

Prove the following for all large $x$: there is a choice of congruence classes $a_p$ for all primes $p\leq x$ and a decomposition $\{p\leq x\}=A\sqcup B$ into two non-empty sets such that, for all $n&#60;x$, there exist some $p\in A$ and $q\in B$ such that $n\equiv a_p\pmod{p}$ and $n\equiv a_q\pmod{q}$.

\noindent\rule{\linewidth}{0.4pt}


%% ============================================================
\section{Problem \#468}
\label{prob:468}

\noindent\textbf{Status:} OPEN \quad|\quad \textbf{Prize:} none \quad|\quad \textbf{Updated:} 2025-08-31

\noindent\textbf{Tags:} number theory, divisors

\medskip

\noindent\textbf{Statement:}

For any $n$ let $D_n$ be the set of sums of the shape $d_1,d_1+d_2,d_1+d_2+d_3,\ldots$ where $1&#60;d_1&#60;d_2&#60;\cdots$ are the divisors of $n$. What is the size of $D_n\backslash \cup_{m&#60;n}D_m$? If $f(N)$ is the minimal $n$ such that $N\in D_n$ then is it true that $f(N)=o(N)$? Perhaps just for almost all $N$?

\noindent\rule{\linewidth}{0.4pt}


%% ============================================================
\section{Problem \#469}
\label{prob:469}

\noindent\textbf{Status:} OPEN \quad|\quad \textbf{Prize:} none \quad|\quad \textbf{Updated:} 2025-08-31

\noindent\textbf{Tags:} number theory, divisors

\medskip

\noindent\textbf{Statement:}

Let $A$ be the set of all $n$ such that $n=d_1+\cdots+d_k$ with $d_i$ distinct proper divisors of $n$, but this is not true for any $m\mid n$ with $m&#60;n$. Does\[\sum_{n\in A}\frac{1}{n}\]converge?

\noindent\rule{\linewidth}{0.4pt}


%% ============================================================
\section{Problem \#470}
\label{prob:470}

\noindent\textbf{Status:} OPEN \quad|\quad \textbf{Prize:} \$10 \quad|\quad \textbf{Updated:} 2025-08-31

\noindent\textbf{Tags:} number theory, divisors

\medskip

\noindent\textbf{Statement:}

Call $n$ weird if $\sigma(n)\geq 2n$ and $n$ is not pseudoperfect, that is, it is not the sum of any set of its divisors. Are there any odd weird numbers? Are there infinitely many primitive weird numbers, i.e. those such that no proper divisor of $n$ is weird?

\noindent\rule{\linewidth}{0.4pt}


%% ============================================================
\section{Problem \#472}
\label{prob:472}

\noindent\textbf{Status:} OPEN \quad|\quad \textbf{Prize:} none \quad|\quad \textbf{Updated:} 2025-08-31

\noindent\textbf{Tags:} number theory

\medskip

\noindent\textbf{Statement:}

Given some initial finite sequence of primes $q_1&#60;\cdots&#60;q_m$ extend it so that $q_{n+1}$ is the smallest prime of the form $q_n+q_i-1$ for $n\geq m$. Is there an initial starting sequence so that the resulting sequence is infinite?

\noindent\rule{\linewidth}{0.4pt}


%% ============================================================
\section{Problem \#477}
\label{prob:477}

\noindent\textbf{Status:} OPEN \quad|\quad \textbf{Prize:} none \quad|\quad \textbf{Updated:} 2025-08-31

\noindent\textbf{Tags:} number theory

\medskip

\noindent\textbf{Statement:}

Is there a polynomial $f:\mathbb{Z}\to \mathbb{Z}$ of degree at least $2$ and a set $A\subset \mathbb{Z}$ such that for any $n\in \mathbb{Z}$ there is exactly one $a\in A$ and $b\in \{ f(k) : k\in\mathbb{Z}\}$ such that $n=a+b$?

\noindent\rule{\linewidth}{0.4pt}


%% ============================================================
\section{Problem \#478}
\label{prob:478}

\noindent\textbf{Status:} OPEN \quad|\quad \textbf{Prize:} none \quad|\quad \textbf{Updated:} 2025-08-31

\noindent\textbf{Tags:} number theory, factorials

\medskip

\noindent\textbf{Statement:}

Let $p$ be a prime and\[A_p = \{ k! \pmod{p} : 1\leq k&#60;p\}.\]Is it true that\[\lvert A_p\rvert \sim (1-\tfrac{1}{e})p?\]

\noindent\rule{\linewidth}{0.4pt}


%% ============================================================
\section{Problem \#479}
\label{prob:479}

\noindent\textbf{Status:} OPEN \quad|\quad \textbf{Prize:} none \quad|\quad \textbf{Updated:} 2025-08-31

\noindent\textbf{Tags:} number theory

\medskip

\noindent\textbf{Statement:}

Is it true that, for all $k\neq 1$, there are infinitely many $n$ such that $2^n\equiv k\pmod{n}$?

\noindent\rule{\linewidth}{0.4pt}


%% ============================================================
\section{Problem \#483}
\label{prob:483}

\noindent\textbf{Status:} OPEN \quad|\quad \textbf{Prize:} none \quad|\quad \textbf{Updated:} 2025-08-31

\noindent\textbf{Tags:} number theory, additive combinatorics, ramsey theory

\medskip

\noindent\textbf{Statement:}

Let $f(k)$ be the minimal $N$ such that if $\{1,\ldots,N\}$ is $k$-coloured then there is a monochromatic solution to $a+b=c$. Estimate $f(k)$. In particular, is it true that $f(k) &lt; c^k$ for some constant $c&gt;0$?

\noindent\rule{\linewidth}{0.4pt}


%% ============================================================
\section{Problem \#486}
\label{prob:486}

\noindent\textbf{Status:} OPEN \quad|\quad \textbf{Prize:} none \quad|\quad \textbf{Updated:} 2025-08-31

\noindent\textbf{Tags:} number theory, primitive sets

\medskip

\noindent\textbf{Statement:}

Let $A\subseteq \mathbb{N}$, and for each $n\in A$ choose some $X_n\subseteq \mathbb{Z}/n\mathbb{Z}$. Let\[B = \{ m\in \mathbb{N} : m\not\in X_n\pmod{n}\textrm{ for all }n\in A\textrm{ with }m&#62;n\}.\]Must $B$ have a logarithmic density, i.e. is it true that\[\lim_{x\to \infty} \frac{1}{\log x}\sum_{\substack{m\in B\\ m&#60;x}}\frac{1}{m}\]exists?

\noindent\rule{\linewidth}{0.4pt}


%% ============================================================
\section{Problem \#489}
\label{prob:489}

\noindent\textbf{Status:} OPEN \quad|\quad \textbf{Prize:} none \quad|\quad \textbf{Updated:} 2025-08-31

\noindent\textbf{Tags:} number theory

\medskip

\noindent\textbf{Statement:}

Let $A\subseteq \mathbb{N}$ be a set such that $\lvert A\cap [1,x]\rvert=o(x^{1/2})$. Let\[B=\{ n\geq 1 : a\nmid n\textrm{ for all }a\in A\}.\]If $B=\{b_1&#60;b_2&#60;\cdots\}$ then is it true that\[\lim \frac{1}{x}\sum_{b_i&#60;x}(b_{i+1}-b_i)^2\]exists (and is finite)?

\noindent\rule{\linewidth}{0.4pt}


%% ============================================================
\section{Problem \#495}
\label{prob:495}

\noindent\textbf{Status:} OPEN \quad|\quad \textbf{Prize:} none \quad|\quad \textbf{Updated:} 2025-08-31

\noindent\textbf{Tags:} diophantine approximation, number theory

\medskip
\noindent\textit{Littlewood conjecture}


\noindent\textbf{Statement:}

Let $\alpha,\beta \in \mathbb{R}$. Is it true that\[\liminf_{n\to \infty} n \| n\alpha \| \| n\beta\| =0\]where $\|x\|$ is the distance from $x$ to the nearest integer?

\noindent\rule{\linewidth}{0.4pt}


%% ============================================================
\section{Problem \#500}
\label{prob:500}

\noindent\textbf{Status:} OPEN \quad|\quad \textbf{Prize:} \$500 \quad|\quad \textbf{Updated:} 2025-08-31

\noindent\textbf{Tags:} graph theory, hypergraphs, turan number

\medskip

\noindent\textbf{Statement:}

What is $\mathrm{ex}_3(n,K_4^3)$? That is, the largest number of $3$-edges which can placed on $n$ vertices so that there exists no $K_4^3$, a set of 4 vertices which is covered by all 4 possible $3$-edges.

\noindent\rule{\linewidth}{0.4pt}


%% ============================================================
\section{Problem \#501}
\label{prob:501}

\noindent\textbf{Status:} OPEN \quad|\quad \textbf{Prize:} none \quad|\quad \textbf{Updated:} 2025-08-31

\noindent\textbf{Tags:} combinatorics, set theory

\medskip

\noindent\textbf{Statement:}

For every $x\in\mathbb{R}$ let $A_x\subset \mathbb{R}$ be a bounded set with outer measure $&#60;1$. Must there exist an infinite independent set, that is, some infinite $X\subseteq \mathbb{R}$ such that $x\not\in A_y$ for all $x\neq y\in X$? If the sets $A_x$ are closed and have measure $&#60;1$, then must there exist an independent set of size $3$?

\noindent\rule{\linewidth}{0.4pt}


%% ============================================================
\section{Problem \#503}
\label{prob:503}

\noindent\textbf{Status:} OPEN \quad|\quad \textbf{Prize:} none \quad|\quad \textbf{Updated:} 2025-08-31

\noindent\textbf{Tags:} geometry, distances

\medskip

\noindent\textbf{Statement:}

What is the size of the largest $A\subseteq \mathbb{R}^d$ such that every three points from $A$ determine an isosceles triangle? That is, for any three points $x,y,z$ from $A$, at least two of the distances $\lvert x-y\rvert,\lvert y-z\rvert,\lvert x-z\rvert$ are equal.

\noindent\rule{\linewidth}{0.4pt}


%% ============================================================
\section{Problem \#507}
\label{prob:507}

\noindent\textbf{Status:} OPEN \quad|\quad \textbf{Prize:} none \quad|\quad \textbf{Updated:} 2025-08-31

\noindent\textbf{Tags:} geometry

\medskip
\noindent\textit{Heilbronn's triangle problem}


\noindent\textbf{Statement:}

Let $\alpha(n)$ be such that every set of $n$ points in the unit disk contains three points which determine a triangle of area at most $\alpha(n)$. Estimate $\alpha(n)$.

\noindent\rule{\linewidth}{0.4pt}


%% ============================================================
\section{Problem \#508}
\label{prob:508}

\noindent\textbf{Status:} OPEN \quad|\quad \textbf{Prize:} none \quad|\quad \textbf{Updated:} 2025-08-31

\noindent\textbf{Tags:} geometry, ramsey theory

\medskip
\noindent\textit{Hadwiger-Nelson problem}


\noindent\textbf{Statement:}

What is the chromatic number of the plane? That is, what is the smallest number of colours required to colour $\mathbb{R}^2$ such that no two points of the same colour are distance $1$ apart?

\noindent\rule{\linewidth}{0.4pt}


%% ============================================================
\section{Problem \#509}
\label{prob:509}

\noindent\textbf{Status:} OPEN \quad|\quad \textbf{Prize:} none \quad|\quad \textbf{Updated:} 2025-08-31

\noindent\textbf{Tags:} analysis

\medskip

\noindent\textbf{Statement:}

Let $f(z)\in\mathbb{C}[z]$ be a monic non-constant polynomial. Can the set\[\{ z\in \mathbb{C} : \lvert f(z)\rvert \leq 1\}\]be covered by a set of circles the sum of whose radii is $\leq 2$?

\noindent\rule{\linewidth}{0.4pt}


%% ============================================================
\section{Problem \#510}
\label{prob:510}

\noindent\textbf{Status:} OPEN \quad|\quad \textbf{Prize:} none \quad|\quad \textbf{Updated:} 2025-08-31

\noindent\textbf{Tags:} analysis

\medskip
\noindent\textit{Chowla's cosine problem}


\noindent\textbf{Statement:}

If $A\subset \mathbb{Z}$ is a finite set of size $N$ then is there some absolute constant $c&#62;0$ and $\theta$ such that\[\sum_{n\in A}\cos(n\theta) &#60; -cN^{1/2}?\]

\noindent\rule{\linewidth}{0.4pt}


%% ============================================================
\section{Problem \#513}
\label{prob:513}

\noindent\textbf{Status:} OPEN \quad|\quad \textbf{Prize:} none \quad|\quad \textbf{Updated:} 2025-08-31

\noindent\textbf{Tags:} analysis

\medskip

\noindent\textbf{Statement:}

Let $f=\sum_{n=0}^\infty a_nz^n$ be a transcendental entire function. What is the greatest possible value of\[\liminf_{r\to \infty} \frac{\max_n\lvert a_nr^n\rvert}{\max_{\lvert z\rvert=r}\lvert f(z)\rvert}?\]

\noindent\rule{\linewidth}{0.4pt}


%% ============================================================
\section{Problem \#514}
\label{prob:514}

\noindent\textbf{Status:} OPEN \quad|\quad \textbf{Prize:} none \quad|\quad \textbf{Updated:} 2025-08-31

\noindent\textbf{Tags:} analysis

\medskip

\noindent\textbf{Statement:}

Let $f(z)$ be an entire transcendental function. Does there exist a path $L$ so that, for every $n$,\[\lvert f(z)/z^n\rvert \to \infty\]as $z\to \infty$ along $L$? Can the length of this path be estimated in terms of $M(r)=\max_{\lvert z\rvert=r}\lvert f(z)\rvert$? Does there exist a path along which $\lvert f(z)\rvert$ tends to $\infty$ faster than a fixed function of $M(r)$ (such that $M(r)^\epsilon$)?

\noindent\rule{\linewidth}{0.4pt}


%% ============================================================
\section{Problem \#517}
\label{prob:517}

\noindent\textbf{Status:} OPEN \quad|\quad \textbf{Prize:} none \quad|\quad \textbf{Updated:} 2025-08-31

\noindent\textbf{Tags:} analysis

\medskip

\noindent\textbf{Statement:}

Let $f(z)=\sum_{k=1}^\infty a_kz^{n_k}$ be an entire function (with $a_k\neq 0$ for all $k\geq 1$). Is it true that if $n_k/k\to \infty$ then $f(z)$ assumes every value infinitely often?

\noindent\rule{\linewidth}{0.4pt}


%% ============================================================
\section{Problem \#520}
\label{prob:520}

\noindent\textbf{Status:} OPEN \quad|\quad \textbf{Prize:} none \quad|\quad \textbf{Updated:} 2025-08-31

\noindent\textbf{Tags:} number theory, probability

\medskip

\noindent\textbf{Statement:}

Let $f$ be a Rademacher multiplicative function: a random $\{-1,0,1\}$-valued multiplicative function, where for each prime $p$ we independently choose $f(p)\in \{-1,1\}$ uniformly at random, and for square-free integers $n$ we extend $f(p_1\cdots p_r)=f(p_1)\cdots f(p_r)$ (and $f(n)=0$ if $n$ is not squarefree). Does there exist some constant $c&#62;0$ such that, almost surely,\[\limsup_{N\to \infty}\frac{\sum_{m\leq N}f(m)}{\sqrt{N\log\log N}}=c?\]

\noindent\rule{\linewidth}{0.4pt}


%% ============================================================
\section{Problem \#521}
\label{prob:521}

\noindent\textbf{Status:} OPEN \quad|\quad \textbf{Prize:} none \quad|\quad \textbf{Updated:} 2025-08-31

\noindent\textbf{Tags:} analysis, polynomials, probability

\medskip

\noindent\textbf{Statement:}

Let $(\epsilon_k)_{k\geq 0}$ be independently uniformly chosen at random from $\{-1,1\}$. If $R_n$ counts the number of real roots of $f_n(z)=\sum_{0\leq k\leq n}\epsilon_k z^k$ then is it true that, almost surely,\[\lim_{n\to \infty}\frac{R_n}{\log n}=\frac{2}{\pi}?\]

\noindent\rule{\linewidth}{0.4pt}


%% ============================================================
\section{Problem \#522}
\label{prob:522}

\noindent\textbf{Status:} OPEN \quad|\quad \textbf{Prize:} none \quad|\quad \textbf{Updated:} 2025-12-08

\noindent\textbf{Tags:} analysis, polynomials, probability

\medskip

\noindent\textbf{Statement:}

Let $f(z)=\sum_{0\leq k\leq n} \epsilon_k z^k$ be a random polynomial, where $\epsilon_k\in \{-1,1\}$ independently uniformly at random for $0\leq k\leq n$. Is it true that, if $R_n$ is the number of roots of $f(z)$ in $\{ z\in \mathbb{C} : \lvert z\rvert \leq 1\}$, then\[\frac{R_n}{n/2}\to 1\]almost surely?

\noindent\rule{\linewidth}{0.4pt}


%% ============================================================
\section{Problem \#524}
\label{prob:524}

\noindent\textbf{Status:} OPEN \quad|\quad \textbf{Prize:} none \quad|\quad \textbf{Updated:} 2025-08-31

\noindent\textbf{Tags:} analysis, probability, polynomials

\medskip

\noindent\textbf{Statement:}

For any $t\in (0,1)$ let $t=\sum_{k=1}^\infty \epsilon_k(t)2^{-k}$ (where $\epsilon_k(t)\in \{0,1\}$). What is the correct order of magnitude (for almost all $t\in(0,1)$) for\[M_n(t)=\max_{x\in [-1,1]}\left\lvert \sum_{k\leq n}(-1)^{\epsilon_k(t)}x^k\right\rvert?\]

\noindent\rule{\linewidth}{0.4pt}


%% ============================================================
\section{Problem \#528}
\label{prob:528}

\noindent\textbf{Status:} OPEN \quad|\quad \textbf{Prize:} none \quad|\quad \textbf{Updated:} 2025-08-31

\noindent\textbf{Tags:} geometry

\medskip

\noindent\textbf{Statement:}

Let $f(n,k)$ count the number of self-avoiding walks of $n$ steps (beginning at the origin) in $\mathbb{Z}^k$ (i.e. those walks which do not intersect themselves). Determine\[C_k=\lim_{n\to\infty}f(n,k)^{1/n}.\]

\noindent\rule{\linewidth}{0.4pt}


%% ============================================================
\section{Problem \#529}
\label{prob:529}

\noindent\textbf{Status:} OPEN \quad|\quad \textbf{Prize:} none \quad|\quad \textbf{Updated:} 2025-08-31

\noindent\textbf{Tags:} geometry, probability

\medskip

\noindent\textbf{Statement:}

Let $d_k(n)$ be the expected distance from the origin after taking $n$ random steps from the origin in $\mathbb{Z}^k$ (conditional on no self intersections) - that is, a self-avoiding walk . Is it true that\[\lim_{n\to \infty}\frac{d_2(n)}{n^{1/2}}= \infty?\]Is it true that\[d_k(n)\ll n^{1/2}\]for $k\geq 3$?

\noindent\rule{\linewidth}{0.4pt}


%% ============================================================
\section{Problem \#530}
\label{prob:530}

\noindent\textbf{Status:} OPEN \quad|\quad \textbf{Prize:} none \quad|\quad \textbf{Updated:} 2025-08-31

\noindent\textbf{Tags:} number theory, sidon sets

\medskip

\noindent\textbf{Statement:}

Let $\ell(N)$ be maximal such that in any finite set $A\subset \mathbb{R}$ of size $N$ there exists a Sidon subset $S$ of size $\ell(N)$ (i.e. the only solutions to $a+b=c+d$ in $S$ are the trivial ones). Determine the order of $\ell(N)$. In particular, is it true that $\ell(N)\sim N^{1/2}$?

\noindent\rule{\linewidth}{0.4pt}


%% ============================================================
\section{Problem \#531}
\label{prob:531}

\noindent\textbf{Status:} OPEN \quad|\quad \textbf{Prize:} none \quad|\quad \textbf{Updated:} 2025-08-31

\noindent\textbf{Tags:} number theory, ramsey theory

\medskip

\noindent\textbf{Statement:}

Let $F(k)$ be the minimal $N$ such that if we two-colour $\{1,\ldots,N\}$ there is a set $A$ of size $k$ such that all subset sums $\sum_{a\in S}a$ (for $\emptyset\neq S\subseteq A$) are monochromatic. Estimate $F(k)$.

\noindent\rule{\linewidth}{0.4pt}


%% ============================================================
\section{Problem \#535}
\label{prob:535}

\noindent\textbf{Status:} OPEN \quad|\quad \textbf{Prize:} none \quad|\quad \textbf{Updated:} 2025-08-31

\noindent\textbf{Tags:} number theory

\medskip

\noindent\textbf{Statement:}

Let $r\geq 3$, and let $f_r(N)$ denote the size of the largest subset of $\{1,\ldots,N\}$ such that no subset of size $r$ has the same pairwise greatest common divisor between all elements. Estimate $f_r(N)$.

\noindent\rule{\linewidth}{0.4pt}


%% ============================================================
\section{Problem \#536}
\label{prob:536}

\noindent\textbf{Status:} OPEN \quad|\quad \textbf{Prize:} none \quad|\quad \textbf{Updated:} 2025-08-31

\noindent\textbf{Tags:} number theory

\medskip

\noindent\textbf{Statement:}

Let $f(N)$ be the largest size of $A\subseteq \{1,\ldots,N\}$ with the property that there are no distinct $a,b,c\in A$ such that\[[a,b]=[b,c]=[a,c],\]where $[a,b]$ denotes the least common multiple. Estimate $f(N)$ - in particular, is it true that $f(N)=o(N)$?

\noindent\rule{\linewidth}{0.4pt}


%% ============================================================
\section{Problem \#538}
\label{prob:538}

\noindent\textbf{Status:} OPEN \quad|\quad \textbf{Prize:} none \quad|\quad \textbf{Updated:} 2025-08-31

\noindent\textbf{Tags:} number theory

\medskip

\noindent\textbf{Statement:}

Let $r\geq 2$ and suppose that $A\subseteq\{1,\ldots,N\}$ is such that, for any $m$, there are at most $r$ solutions to $m=pa$ where $p$ is prime and $a\in A$. Give the best possible upper bound for\[\sum_{n\in A}\frac{1}{n}.\]

\noindent\rule{\linewidth}{0.4pt}


%% ============================================================
\section{Problem \#539}
\label{prob:539}

\noindent\textbf{Status:} OPEN \quad|\quad \textbf{Prize:} none \quad|\quad \textbf{Updated:} 2025-08-31

\noindent\textbf{Tags:} number theory

\medskip

\noindent\textbf{Statement:}

Let $h(n)$ be such that, for any set $A\subseteq \mathbb{N}$ of size $n$, the set\[\left\{ \frac{a}{(a,b)}: a,b\in A\right\}\]has size at least $h(n)$. Estimate $h(n)$.

\noindent\rule{\linewidth}{0.4pt}


%% ============================================================
\section{Problem \#544}
\label{prob:544}

\noindent\textbf{Status:} OPEN \quad|\quad \textbf{Prize:} none \quad|\quad \textbf{Updated:} 2025-08-31

\noindent\textbf{Tags:} graph theory, ramsey theory

\medskip

\noindent\textbf{Statement:}

Show that\[R(3,k+1)-R(3,k)\to\infty\]as $k\to \infty$. Similarly, prove or disprove that\[R(3,k+1)-R(3,k)=o(k).\]

\noindent\rule{\linewidth}{0.4pt}


%% ============================================================
\section{Problem \#545}
\label{prob:545}

\noindent\textbf{Status:} OPEN \quad|\quad \textbf{Prize:} none \quad|\quad \textbf{Updated:} 2025-12-02

\noindent\textbf{Tags:} graph theory, ramsey theory

\medskip

\noindent\textbf{Statement:}

Let $G$ be a graph with $m$ edges and no isolated vertices. Is the Ramsey number $R(G)$ maximised when $G$ is 'as complete as possible'? That is, if $m=\binom{n}{2}+t$ edges with $0\leq t&#60;n$ then is\[R(G)\leq R(H),\]where $H$ is the graph formed by connecting a new vertex to $t$ of the vertices of $K_n$?

\noindent\rule{\linewidth}{0.4pt}


%% ============================================================
\section{Problem \#550}
\label{prob:550}

\noindent\textbf{Status:} OPEN \quad|\quad \textbf{Prize:} none \quad|\quad \textbf{Updated:} 2025-08-31

\noindent\textbf{Tags:} graph theory, ramsey theory

\medskip

\noindent\textbf{Statement:}

Let $m_1\leq\cdots\leq m_k$ and $n$ be sufficiently large. If $T$ is a tree on $n$ vertices and $G$ is the complete multipartite graph with vertex class sizes $m_1,\ldots,m_k$ then prove that\[R(T,G)\leq (\chi(G)-1)(R(T,K_{m_1,m_2})-1)+m_1.\]

\noindent\rule{\linewidth}{0.4pt}


%% ============================================================
\section{Problem \#552}
\label{prob:552}

\noindent\textbf{Status:} OPEN \quad|\quad \textbf{Prize:} none \quad|\quad \textbf{Updated:} 2025-08-31

\noindent\textbf{Tags:} graph theory, ramsey theory

\medskip

\noindent\textbf{Statement:}

Determine the Ramsey number\[R(C_4,S_n),\]where $S_n=K_{1,n}$ is the star on $n+1$ vertices. In particular, is it true that, for any $c&#62;0$, there are infinitely many $n$ such that\[R(C_4,S_n)\leq n+\sqrt{n}-c?\]

\noindent\rule{\linewidth}{0.4pt}


%% ============================================================
\section{Problem \#554}
\label{prob:554}

\noindent\textbf{Status:} OPEN \quad|\quad \textbf{Prize:} none \quad|\quad \textbf{Updated:} 2025-08-31

\noindent\textbf{Tags:} graph theory, ramsey theory

\medskip

\noindent\textbf{Statement:}

Let $R_k(G)$ denote the minimal $m$ such that if the edges of $K_m$ are $k$-coloured then there is a monochromatic copy of $G$. Show that\[\lim_{k\to \infty}\frac{R_k(C_{2n+1})}{R_k(K_3)}=0\]for any $n\geq 2$.

\noindent\rule{\linewidth}{0.4pt}


%% ============================================================
\section{Problem \#555}
\label{prob:555}

\noindent\textbf{Status:} OPEN \quad|\quad \textbf{Prize:} none \quad|\quad \textbf{Updated:} 2025-08-31

\noindent\textbf{Tags:} graph theory, ramsey theory

\medskip

\noindent\textbf{Statement:}

Let $R_k(G)$ denote the minimal $m$ such that if the edges of $K_m$ are $k$-coloured then there is a monochromatic copy of $G$. Determine the value of\[R_k(C_{2n}).\]

\noindent\rule{\linewidth}{0.4pt}


%% ============================================================
\section{Problem \#557}
\label{prob:557}

\noindent\textbf{Status:} OPEN \quad|\quad \textbf{Prize:} none \quad|\quad \textbf{Updated:} 2025-08-31

\noindent\textbf{Tags:} graph theory, ramsey theory

\medskip

\noindent\textbf{Statement:}

Let $R_k(G)$ denote the minimal $m$ such that if the edges of $K_m$ are $k$-coloured then there is a monochromatic copy of $G$. Is it true that\[R_k(T)\leq kn+O(1)\]for any tree $T$ on $n$ vertices?

\noindent\rule{\linewidth}{0.4pt}


%% ============================================================
\section{Problem \#558}
\label{prob:558}

\noindent\textbf{Status:} OPEN \quad|\quad \textbf{Prize:} none \quad|\quad \textbf{Updated:} 2025-08-31

\noindent\textbf{Tags:} graph theory, ramsey theory

\medskip

\noindent\textbf{Statement:}

Let $R_k(G)$ denote the minimal $m$ such that if the edges of $K_m$ are $k$-coloured then there is a monochromatic copy of $G$. Determine\[R_k(K_{s,t})\]where $K_{s,t}$ is the complete bipartite graph with $s$ vertices in one component and $t$ in the other.

\noindent\rule{\linewidth}{0.4pt}


%% ============================================================
\section{Problem \#560}
\label{prob:560}

\noindent\textbf{Status:} OPEN \quad|\quad \textbf{Prize:} none \quad|\quad \textbf{Updated:} 2025-08-31

\noindent\textbf{Tags:} graph theory, ramsey theory

\medskip

\noindent\textbf{Statement:}

Let $\hat{R}(G)$ denote the size Ramsey number, the minimal number of edges $m$ such that there is a graph $H$ with $m$ edges such that in any $2$-colouring of the edges of $H$ there is a monochromatic copy of $G$. Determine\[\hat{R}(K_{n,n}),\]where $K_{n,n}$ is the complete bipartite graph with $n$ vertices in each component.

\noindent\rule{\linewidth}{0.4pt}


%% ============================================================
\section{Problem \#561}
\label{prob:561}

\noindent\textbf{Status:} OPEN \quad|\quad \textbf{Prize:} none \quad|\quad \textbf{Updated:} 2025-08-31

\noindent\textbf{Tags:} graph theory, ramsey theory

\medskip

\noindent\textbf{Statement:}

Let $\hat{R}(G)$ denote the size Ramsey number, the minimal number of edges $m$ such that there is a graph $H$ with $m$ edges such that in any $2$-colouring of the edges of $H$ there is a monochromatic copy of $G$. Let $F_1$ and $F_2$ be the union of stars. More precisely, let $F_1=\cup_{i\leq s} K_{1,n_i}$ and $F_2=\cup_{j\leq t} K_{1,m_j}$ with $n_1\geq \cdots \geq n_s\geq 1$ and $m_1\geq \cdots \geq m_t\geq 1$. Prove that\[\hat{R}(F_1,F_2) = \sum_{2\leq k\leq s+t}l_k\]where\[l_k=\max\{n_i+m_j-1 : i+j=k\}.\]

\noindent\rule{\linewidth}{0.4pt}


%% ============================================================
\section{Problem \#562}
\label{prob:562}

\noindent\textbf{Status:} OPEN \quad|\quad \textbf{Prize:} none \quad|\quad \textbf{Updated:} 2025-08-31

\noindent\textbf{Tags:} graph theory, ramsey theory, hypergraphs

\medskip

\noindent\textbf{Statement:}

Let $R_r(n)$ denote the $r$-uniform hypergraph Ramsey number: the minimal $m$ such that if we $2$-colour all edges of the complete $r$-uniform hypergraph on $m$ vertices then there must be some monochromatic copy of the complete $r$-uniform hypergraph on $n$ vertices. Prove that, for $r\geq 3$,\[\log_{r-1} R_r(n) \asymp_r n,\]where $\log_{r-1}$ denotes the $(r-1)$-fold iterated logarithm. That is, does $R_r(n)$ grow like\[2^{2^{\cdots n}}\]where the tower of exponentials has height $r-1$?

\noindent\rule{\linewidth}{0.4pt}


%% ============================================================
\section{Problem \#563}
\label{prob:563}

\noindent\textbf{Status:} OPEN \quad|\quad \textbf{Prize:} none \quad|\quad \textbf{Updated:} 2025-08-31

\noindent\textbf{Tags:} graph theory, ramsey theory, hypergraphs

\medskip

\noindent\textbf{Statement:}

Let $F(n,\alpha)$ denote the smallest $m$ such that there exists a $2$-colouring of the edges of $K_n$ so that every $X\subseteq [n]$ with $\lvert X\rvert\geq m$ contains more than $\alpha \binom{\lvert X\rvert}{2}$ many edges of each colour. Prove that, for every $0\leq \alpha&#60; 1/2$,\[F(n,\alpha)\sim c_\alpha\log n\]for some constant $c_\alpha$ depending only on $\alpha$.

\noindent\rule{\linewidth}{0.4pt}


%% ============================================================
\section{Problem \#564}
\label{prob:564}

\noindent\textbf{Status:} OPEN \quad|\quad \textbf{Prize:} \$500 \quad|\quad \textbf{Updated:} 2025-08-31

\noindent\textbf{Tags:} graph theory, ramsey theory, hypergraphs

\medskip

\noindent\textbf{Statement:}

Let $R_3(n)$ be the minimal $m$ such that if the edges of the $3$-uniform hypergraph on $m$ vertices are $2$-coloured then there is a monochromatic copy of the complete $3$-uniform hypergraph on $n$ vertices. Is there some constant $c&#62;0$ such that\[R_3(n) \geq 2^{2^{cn}}?\]

\noindent\rule{\linewidth}{0.4pt}


%% ============================================================
\section{Problem \#566}
\label{prob:566}

\noindent\textbf{Status:} OPEN \quad|\quad \textbf{Prize:} none \quad|\quad \textbf{Updated:} 2025-08-31

\noindent\textbf{Tags:} graph theory, ramsey theory

\medskip

\noindent\textbf{Statement:}

Let $G$ be such that any subgraph on $k$ vertices has at most $2k-3$ edges. Is it true that, if $H$ has $m$ edges and no isolated vertices, then\[R(G,H)\ll m?\]

\noindent\rule{\linewidth}{0.4pt}


%% ============================================================
\section{Problem \#567}
\label{prob:567}

\noindent\textbf{Status:} OPEN \quad|\quad \textbf{Prize:} none \quad|\quad \textbf{Updated:} 2025-08-31

\noindent\textbf{Tags:} graph theory, ramsey theory

\medskip

\noindent\textbf{Statement:}

Let $G$ be either $Q_3$ or $K_{3,3}$ or $H_5$ (the last formed by adding two vertex-disjoint chords to $C_5$). Is it true that, if $H$ has $m$ edges and no isolated vertices, then\[R(G,H)\ll m?\]

\noindent\rule{\linewidth}{0.4pt}


%% ============================================================
\section{Problem \#568}
\label{prob:568}

\noindent\textbf{Status:} OPEN \quad|\quad \textbf{Prize:} none \quad|\quad \textbf{Updated:} 2025-08-31

\noindent\textbf{Tags:} graph theory, ramsey theory

\medskip

\noindent\textbf{Statement:}

Let $G$ be a graph such that $R(G,T_n)\ll n$ for any tree $T_n$ on $n$ vertices and $R(G,K_n)\ll n^2$. Is it true that, for any $H$ with $m$ edges and no isolated vertices,\[R(G,H)\ll m?\]

\noindent\rule{\linewidth}{0.4pt}


%% ============================================================
\section{Problem \#569}
\label{prob:569}

\noindent\textbf{Status:} OPEN \quad|\quad \textbf{Prize:} none \quad|\quad \textbf{Updated:} 2025-08-31

\noindent\textbf{Tags:} graph theory, ramsey theory

\medskip

\noindent\textbf{Statement:}

Let $k\geq 1$. What is the best possible $c_k$ such that\[R(C_{2k+1},H)\leq c_k m\]for any graph $H$ on $m$ edges without isolated vertices?

\noindent\rule{\linewidth}{0.4pt}


%% ============================================================
\section{Problem \#571}
\label{prob:571}

\noindent\textbf{Status:} OPEN \quad|\quad \textbf{Prize:} none \quad|\quad \textbf{Updated:} 2025-08-31

\noindent\textbf{Tags:} graph theory, turan number

\medskip

\noindent\textbf{Statement:}

Show that for any rational $\alpha \in [1,2)$ there exists a bipartite graph $G$ such that\[\mathrm{ex}(n;G)\asymp n^{\alpha}.\]

\noindent\rule{\linewidth}{0.4pt}


%% ============================================================
\section{Problem \#572}
\label{prob:572}

\noindent\textbf{Status:} OPEN \quad|\quad \textbf{Prize:} none \quad|\quad \textbf{Updated:} 2025-08-31

\noindent\textbf{Tags:} graph theory, turan number

\medskip

\noindent\textbf{Statement:}

Show that for $k\geq 3$\[\mathrm{ex}(n;C_{2k})\gg n^{1+\frac{1}{k}}.\]

\noindent\rule{\linewidth}{0.4pt}


%% ============================================================
\section{Problem \#573}
\label{prob:573}

\noindent\textbf{Status:} OPEN \quad|\quad \textbf{Prize:} none \quad|\quad \textbf{Updated:} 2025-08-31

\noindent\textbf{Tags:} graph theory, turan number

\medskip

\noindent\textbf{Statement:}

Is it true that\[\mathrm{ex}(n;\{C_3,C_4\})\sim (n/2)^{3/2}?\]

\noindent\rule{\linewidth}{0.4pt}


%% ============================================================
\section{Problem \#575}
\label{prob:575}

\noindent\textbf{Status:} OPEN \quad|\quad \textbf{Prize:} none \quad|\quad \textbf{Updated:} 2025-08-31

\noindent\textbf{Tags:} graph theory, turan number

\medskip

\noindent\textbf{Statement:}

If $\mathcal{F}$ is a finite set of finite graphs then $\mathrm{ex}(n;\mathcal{F})$ is the maximum number of edges a graph on $n$ vertices can have without containing any subgraphs from $\mathcal{F}$. Note that it is trivial that $\mathrm{ex}(n;\mathcal{F})\leq \mathrm{ex}(n;G)$ for every $G\in\mathcal{F}$. Is it true that, for every $\mathcal{F}$, if there is a bipartite graph in $\mathcal{F}$ then there exists some bipartite $G\in\mathcal{F}$ such that\[\mathrm{ex}(n;G)\ll_{\mathcal{F}}\mathrm{ex}(n;\mathcal{F})?\]

\noindent\rule{\linewidth}{0.4pt}


%% ============================================================
\section{Problem \#576}
\label{prob:576}

\noindent\textbf{Status:} OPEN \quad|\quad \textbf{Prize:} none \quad|\quad \textbf{Updated:} 2025-08-31

\noindent\textbf{Tags:} graph theory, turan number

\medskip

\noindent\textbf{Statement:}

Let $Q_k$ be the $k$-dimensional hypercube graph (so that $Q_k$ has $2^k$ vertices and $k2^{k-1}$ edges). Determine the behaviour of\[\mathrm{ex}(n;Q_k).\]

\noindent\rule{\linewidth}{0.4pt}


%% ============================================================
\section{Problem \#579}
\label{prob:579}

\noindent\textbf{Status:} OPEN \quad|\quad \textbf{Prize:} none \quad|\quad \textbf{Updated:} 2025-08-31

\noindent\textbf{Tags:} graph theory, turan number

\medskip

\noindent\textbf{Statement:}

Let $\delta&#62;0$. If $n$ is sufficiently large and $G$ is a graph on $n$ vertices with no $K_{2,2,2}$ and at least $\delta n^2$ edges then $G$ contains an independent set of size $\gg_\delta n$.

\noindent\rule{\linewidth}{0.4pt}


%% ============================================================
\section{Problem \#584}
\label{prob:584}

\noindent\textbf{Status:} OPEN \quad|\quad \textbf{Prize:} none \quad|\quad \textbf{Updated:} 2025-08-31

\noindent\textbf{Tags:} graph theory, cycles

\medskip

\noindent\textbf{Statement:}

Let $G$ be a graph with $n$ vertices and $\delta n^{2}$ edges. Are there subgraphs $H_1,H_2\subseteq G$ such that {UL} {LI}$H_1$ has $\gg \delta^3n^2$ edges and every two edges in $H_1$ are contained in a cycle of length at most $6$, and furthermore if two edges share a vertex they are on a cycle of length $4$, and {LI}$H_2$ has $\gg \delta^2n^2$ edges and every two edges in $H_2$ are contained in a cycle of length at most $8$. {/UL}

\noindent\rule{\linewidth}{0.4pt}


%% ============================================================
\section{Problem \#585}
\label{prob:585}

\noindent\textbf{Status:} OPEN \quad|\quad \textbf{Prize:} none \quad|\quad \textbf{Updated:} 2025-08-31

\noindent\textbf{Tags:} graph theory, cycles

\medskip

\noindent\textbf{Statement:}

What is the maximum number of edges that a graph on $n$ vertices can have if it does not contain two edge-disjoint cycles with the same vertex set?

\noindent\rule{\linewidth}{0.4pt}


%% ============================================================
\section{Problem \#588}
\label{prob:588}

\noindent\textbf{Status:} OPEN \quad|\quad \textbf{Prize:} \$100 \quad|\quad \textbf{Updated:} 2025-08-31

\noindent\textbf{Tags:} geometry

\medskip

\noindent\textbf{Statement:}

Let $f_k(n)$ be minimal such that if $n$ points in $\mathbb{R}^2$ have no $k+1$ points on a line then there must be at most $f_k(n)$ many lines containing at least $k$ points. Is it true that\[f_k(n)=o(n^2)\]for $k\geq 4$?

\noindent\rule{\linewidth}{0.4pt}


%% ============================================================
\section{Problem \#589}
\label{prob:589}

\noindent\textbf{Status:} OPEN \quad|\quad \textbf{Prize:} none \quad|\quad \textbf{Updated:} 2025-08-31

\noindent\textbf{Tags:} geometry

\medskip

\noindent\textbf{Statement:}

Let $g(n)$ be maximal such that in any set of $n$ points in $\mathbb{R}^2$ with no four points on a line there exists a subset on $g(n)$ points with no three points on a line. Estimate $g(n)$.

\noindent\rule{\linewidth}{0.4pt}


%% ============================================================
\section{Problem \#592}
\label{prob:592}

\noindent\textbf{Status:} OPEN \quad|\quad \textbf{Prize:} \$1000 \quad|\quad \textbf{Updated:} 2025-08-31

\noindent\textbf{Tags:} set theory, ramsey theory

\medskip

\noindent\textbf{Statement:}

Determine which countable ordinals $\beta$ have the property that, if $\alpha=\omega^{^\beta}$, then in any red/blue colouring of the edges of $K_\alpha$ there is either a red $K_\alpha$ or a blue $K_3$.

\noindent\rule{\linewidth}{0.4pt}


%% ============================================================
\section{Problem \#593}
\label{prob:593}

\noindent\textbf{Status:} OPEN \quad|\quad \textbf{Prize:} \$500 \quad|\quad \textbf{Updated:} 2025-08-31

\noindent\textbf{Tags:} set theory, graph theory, hypergraphs, chromatic number

\medskip

\noindent\textbf{Statement:}

Characterize those finite 3-uniform hypergraphs which appear in every 3-uniform hypergraph of chromatic number $&#62;\aleph_0$.

\noindent\rule{\linewidth}{0.4pt}


%% ============================================================
\section{Problem \#595}
\label{prob:595}

\noindent\textbf{Status:} OPEN \quad|\quad \textbf{Prize:} \$250 \quad|\quad \textbf{Updated:} 2025-08-31

\noindent\textbf{Tags:} graph theory, set theory

\medskip

\noindent\textbf{Statement:}

Is there an infinite graph $G$ which contains no $K_4$ and is not the union of countably many triangle-free graphs?

\noindent\rule{\linewidth}{0.4pt}


%% ============================================================
\section{Problem \#596}
\label{prob:596}

\noindent\textbf{Status:} OPEN \quad|\quad \textbf{Prize:} none \quad|\quad \textbf{Updated:} 2025-08-31

\noindent\textbf{Tags:} graph theory, ramsey theory, set theory

\medskip

\noindent\textbf{Statement:}

For which graphs $G_1,G_2$ is it true that {UL} {LI} for every $n\geq 1$ there is a graph $H$ without a $G_1$ but if the edges of $H$ are $n$-coloured then there is a monochromatic copy of $G_2$, and yet{/LI} {LI} for every graph $H$ without a $G_1$ there is an $\aleph_0$-colouring of the edges of $H$ without a monochromatic $G_2$. {/UL}

\noindent\rule{\linewidth}{0.4pt}


%% ============================================================
\section{Problem \#597}
\label{prob:597}

\noindent\textbf{Status:} OPEN \quad|\quad \textbf{Prize:} none \quad|\quad \textbf{Updated:} 2025-08-31

\noindent\textbf{Tags:} graph theory, ramsey theory, set theory

\medskip

\noindent\textbf{Statement:}

Let $G$ be a graph on at most $\aleph_1$ vertices which contains no $K_4$ and no $K_{\aleph_0,\aleph_0}$ (the complete bipartite graph with $\aleph_0$ vertices in each class). Is it true that\[\omega_1^2 \to (\omega_1\omega, G)^2?\]What about finite $G$?

\noindent\rule{\linewidth}{0.4pt}


%% ============================================================
\section{Problem \#598}
\label{prob:598}

\noindent\textbf{Status:} OPEN \quad|\quad \textbf{Prize:} none \quad|\quad \textbf{Updated:} 2025-08-31

\noindent\textbf{Tags:} set theory, ramsey theory

\medskip

\noindent\textbf{Statement:}

Let $m$ be an infinite cardinal and $\kappa$ be the successor cardinal of $2^{\aleph_0}$. Can one colour the countable subsets of $m$ using $\kappa$ many colours so that every $X\subseteq m$ with $\lvert X\rvert=\kappa$ contains subsets of all possible colours?

\noindent\rule{\linewidth}{0.4pt}


%% ============================================================
\section{Problem \#600}
\label{prob:600}

\noindent\textbf{Status:} OPEN \quad|\quad \textbf{Prize:} none \quad|\quad \textbf{Updated:} 2025-08-31

\noindent\textbf{Tags:} graph theory

\medskip

\noindent\textbf{Statement:}

Let $e(n,r)$ be minimal such that every graph on $n$ vertices with at least $e(n,r)$ edges, each edge contained in at least one triangle, must have an edge contained in at least $r$ triangles. Let $r\geq 2$. Is it true that\[e(n,r+1)-e(n,r)\to \infty\]as $n\to \infty$? Is it true that\[\frac{e(n,r+1)}{e(n,r)}\to 1\]as $n\to \infty$?

\noindent\rule{\linewidth}{0.4pt}


%% ============================================================
\section{Problem \#601}
\label{prob:601}

\noindent\textbf{Status:} OPEN \quad|\quad \textbf{Prize:} \$500 \quad|\quad \textbf{Updated:} 2025-08-31

\noindent\textbf{Tags:} graph theory, set theory

\medskip

\noindent\textbf{Statement:}

For which limit ordinals $\alpha$ is it true that if $G$ is a graph with vertex set $\alpha$ then $G$ must have either an infinite path or independent set on a set of vertices with order type $\alpha$?

\noindent\rule{\linewidth}{0.4pt}


%% ============================================================
\section{Problem \#602}
\label{prob:602}

\noindent\textbf{Status:} OPEN \quad|\quad \textbf{Prize:} none \quad|\quad \textbf{Updated:} 2025-08-31

\noindent\textbf{Tags:} combinatorics, set theory

\medskip
\noindent\textit{Property B}


\noindent\textbf{Statement:}

Let $(A_i)$ be a family of sets with $\lvert A_i\rvert=\aleph_0$ for all $i$, such that for any $i\neq j$ we have $\lvert A_i\cap A_j\rvert$ finite and $\neq 1$. Is there a $2$-colouring of $\cup A_i$ such that no $A_i$ is monochromatic?

\noindent\rule{\linewidth}{0.4pt}


%% ============================================================
\section{Problem \#603}
\label{prob:603}

\noindent\textbf{Status:} OPEN \quad|\quad \textbf{Prize:} none \quad|\quad \textbf{Updated:} 2025-08-31

\noindent\textbf{Tags:} combinatorics, set theory

\medskip

\noindent\textbf{Statement:}

Let $(A_i)$ be a family of countably infinite sets such that $\lvert A_i\cap A_j\rvert \neq 2$ for all $i\neq j$. Find the smallest cardinal $C$ such that $\cup A_i$ can always be coloured with at most $C$ colours so that no $A_i$ is monochromatic.

\noindent\rule{\linewidth}{0.4pt}


%% ============================================================
\section{Problem \#604}
\label{prob:604}

\noindent\textbf{Status:} OPEN \quad|\quad \textbf{Prize:} \$500 \quad|\quad \textbf{Updated:} 2025-08-31

\noindent\textbf{Tags:} geometry, distances

\medskip
\noindent\textit{pinned distance problem}


\noindent\textbf{Statement:}

Given $n$ distinct points $A\subset\mathbb{R}^2$ must there be a point $x\in A$ such that\[\#\{ d(x,y) : y \in A\} \gg n^{1-o(1)}?\]Or even $\gg n/\sqrt{\log n}$?

\noindent\rule{\linewidth}{0.4pt}


%% ============================================================
\section{Problem \#609}
\label{prob:609}

\noindent\textbf{Status:} OPEN \quad|\quad \textbf{Prize:} none \quad|\quad \textbf{Updated:} 2025-08-31

\noindent\textbf{Tags:} graph theory, ramsey theory

\medskip

\noindent\textbf{Statement:}

Let $f(n)$ be the minimal $m$ such that if the edges of $K_{2^n+1}$ are coloured with $n$ colours then there must be a monochromatic odd cycle of length at most $m$. Estimate $f(n)$.

\noindent\rule{\linewidth}{0.4pt}


%% ============================================================
\section{Problem \#611}
\label{prob:611}

\noindent\textbf{Status:} OPEN \quad|\quad \textbf{Prize:} none \quad|\quad \textbf{Updated:} 2025-08-31

\noindent\textbf{Tags:} graph theory

\medskip

\noindent\textbf{Statement:}

For a graph $G$ let $\tau(G)$ denote the minimal number of vertices that include at least one from each maximal clique of $G$ (sometimes called the clique transversal number). Is it true that if all maximal cliques in $G$ have at least $cn$ vertices then $\tau(G)=o_c(n)$? Similarly, estimate for $c&#62;0$ the minimal $k_c(n)$ such that if every maximal clique in $G$ has at least $k_c(n)$ vertices then $\tau(G)&#60;(1-c)n$.

\noindent\rule{\linewidth}{0.4pt}


%% ============================================================
\section{Problem \#612}
\label{prob:612}

\noindent\textbf{Status:} OPEN \quad|\quad \textbf{Prize:} none \quad|\quad \textbf{Updated:} 2025-08-31

\noindent\textbf{Tags:} graph theory

\medskip

\noindent\textbf{Statement:}

Let $G$ be a connected graph with $n$ vertices, minimum degree $d$, and diameter $D$. Show if that $G$ contains no $K_{2r}$ and $(r-1)(3r+2)\mid d$ then\[D\leq \frac{2(r-1)(3r+2)}{2r^2-1}\frac{n}{d}+O(1),\]and if $G$ contains no $K_{2r+1}$ and $3r-1 \mid d$ then\[D\leq \frac{3r-1}{r}\frac{n}{d}+O(1).\]

\noindent\rule{\linewidth}{0.4pt}


%% ============================================================
\section{Problem \#614}
\label{prob:614}

\noindent\textbf{Status:} OPEN \quad|\quad \textbf{Prize:} none \quad|\quad \textbf{Updated:} 2025-08-31

\noindent\textbf{Tags:} graph theory

\medskip

\noindent\textbf{Statement:}

Let $f(n,k)$ be minimal such that there is a graph with $n$ vertices and $f(n,k)$ edges where every set of $k+2$ vertices induces a subgraph with maximum degree at least $k$. Determine $f(n,k)$.

\noindent\rule{\linewidth}{0.4pt}


%% ============================================================
\section{Problem \#616}
\label{prob:616}

\noindent\textbf{Status:} OPEN \quad|\quad \textbf{Prize:} none \quad|\quad \textbf{Updated:} 2025-08-31

\noindent\textbf{Tags:} graph theory

\medskip

\noindent\textbf{Statement:}

Let $r\geq 3$. For an $r$-uniform hypergraph $G$ let $\tau(G)$ denote the covering number (or transversal number), the minimum size of a set of vertices which includes at least one from each edge in $G$. Determine the best possible $t$ such that, if $G$ is an $r$-uniform hypergraph $G$ where every subgraph $G'$ on at most $3r-3$ vertices has $\tau(G')\leq 1$, we have $\tau(G)\leq t$.

\noindent\rule{\linewidth}{0.4pt}


%% ============================================================
\section{Problem \#619}
\label{prob:619}

\noindent\textbf{Status:} OPEN \quad|\quad \textbf{Prize:} none \quad|\quad \textbf{Updated:} 2025-08-31

\noindent\textbf{Tags:} graph theory

\medskip

\noindent\textbf{Statement:}

For a triangle-free graph $G$ let $h_r(G)$ be the smallest number of edges that need to be added to $G$ so that it has diameter $r$ (while preserving the property of being triangle-free). Is it true that there exists a constant $c&#62;0$ such that if $G$ is a connected graph on $n$ vertices then $h_4(G)&#60;(1-c)n$?

\noindent\rule{\linewidth}{0.4pt}


%% ============================================================
\section{Problem \#620}
\label{prob:620}

\noindent\textbf{Status:} OPEN \quad|\quad \textbf{Prize:} none \quad|\quad \textbf{Updated:} 2025-08-31

\noindent\textbf{Tags:} graph theory

\medskip
\noindent\textit{Erdős-Rogers problem}


\noindent\textbf{Statement:}

If $G$ is a graph on $n$ vertices without a $K_4$ then how large a triangle-free induced subgraph must $G$ contain?

\noindent\rule{\linewidth}{0.4pt}


%% ============================================================
\section{Problem \#623}
\label{prob:623}

\noindent\textbf{Status:} OPEN \quad|\quad \textbf{Prize:} none \quad|\quad \textbf{Updated:} 2025-08-31

\noindent\textbf{Tags:} set theory

\medskip

\noindent\textbf{Statement:}

Let $X$ be a set of cardinality $\aleph_\omega$ and $f$ be a function from the finite subsets of $X$ to $X$ such that $f(A)\not\in A$ for all $A$. Must there exist an infinite $Y\subseteq X$ that is independent - that is, for all finite $B\subset Y$ we have $f(B)\not\in Y$?

\noindent\rule{\linewidth}{0.4pt}


%% ============================================================
\section{Problem \#624}
\label{prob:624}

\noindent\textbf{Status:} OPEN \quad|\quad \textbf{Prize:} none \quad|\quad \textbf{Updated:} 2025-08-31

\noindent\textbf{Tags:} combinatorics

\medskip

\noindent\textbf{Statement:}

Let $X$ be a finite set of size $n$ and $H(n)$ be such that there is a function $f:\{A : A\subseteq X\}\to X$ so that for every $Y\subseteq X$ with $\lvert Y\rvert \geq H(n)$ we have\[\{ f(A) : A\subseteq Y\}=X.\]Prove that\[H(n)-\log_2 n \to \infty.\]

\noindent\rule{\linewidth}{0.4pt}


%% ============================================================
\section{Problem \#625}
\label{prob:625}

\noindent\textbf{Status:} OPEN \quad|\quad \textbf{Prize:} \$1000 \quad|\quad \textbf{Updated:} 2025-08-31

\noindent\textbf{Tags:} graph theory, chromatic number

\medskip

\noindent\textbf{Statement:}

The cochromatic number of $G$, denoted by $\zeta(G)$, is the minimum number of colours needed to colour the vertices of $G$ such that each colour class induces either a complete graph or empty graph. Let $\chi(G)$ denote the chromatic number. If $G$ is a random graph with $n$ vertices and each edge included independently with probability $1/2$ then is it true that almost surely\[\chi(G) - \zeta(G) \to \infty\]as $n\to \infty$?

\noindent\rule{\linewidth}{0.4pt}


%% ============================================================
\section{Problem \#626}
\label{prob:626}

\noindent\textbf{Status:} OPEN \quad|\quad \textbf{Prize:} none \quad|\quad \textbf{Updated:} 2025-08-31

\noindent\textbf{Tags:} graph theory, chromatic number, cycles

\medskip

\noindent\textbf{Statement:}

Let $k\geq 4$ and $g_k(n)$ denote the largest $m$ such that there is a graph on $n$ vertices with chromatic number $k$ and girth $&#62;m$ (i.e. contains no cycle of length $\leq m$). Does\[\lim_{n\to \infty}\frac{g_k(n)}{\log n}\]exist? Conversely, if $h^{(m)}(n)$ is the maximal chromatic number of a graph on $n$ vertices with girth $&#62;m$ then does\[\lim_{n\to \infty}\frac{\log h^{(m)}(n)}{\log n}\]exist, and what is its value?

\noindent\rule{\linewidth}{0.4pt}


%% ============================================================
\section{Problem \#627}
\label{prob:627}

\noindent\textbf{Status:} OPEN \quad|\quad \textbf{Prize:} none \quad|\quad \textbf{Updated:} 2025-08-31

\noindent\textbf{Tags:} graph theory, chromatic number

\medskip

\noindent\textbf{Statement:}

Let $\omega(G)$ denote the clique number of $G$ and $\chi(G)$ the chromatic number. If $f(n)$ is the maximum value of $\chi(G)/\omega(G)$, as $G$ ranges over all graphs on $n$ vertices, then does\[\lim_{n\to\infty}\frac{f(n)}{n/(\log_2n)^2}\]exist?

\noindent\rule{\linewidth}{0.4pt}


%% ============================================================
\section{Problem \#629}
\label{prob:629}

\noindent\textbf{Status:} OPEN \quad|\quad \textbf{Prize:} none \quad|\quad \textbf{Updated:} 2025-08-31

\noindent\textbf{Tags:} graph theory, chromatic number

\medskip

\noindent\textbf{Statement:}

The list chromatic number $\chi_L(G)$ is defined to be the minimal $k$ such that for any assignment of a list of $k$ colours to each vertex of $G$ (perhaps different lists for different vertices) a colouring of each vertex by a colour on its list can be chosen such that adjacent vertices receive distinct colours. Determine the minimal number of vertices $n(k)$ of a bipartite graph $G$ such that $\chi_L(G)&#62;k$.

\noindent\rule{\linewidth}{0.4pt}


%% ============================================================
\section{Problem \#634}
\label{prob:634}

\noindent\textbf{Status:} OPEN \quad|\quad \textbf{Prize:} \$25 \quad|\quad \textbf{Updated:} 2025-08-31

\noindent\textbf{Tags:} geometry

\medskip

\noindent\textbf{Statement:}

Find all $n$ such that there is at least one triangle which can be cut into $n$ congruent triangles.

\noindent\rule{\linewidth}{0.4pt}


%% ============================================================
\section{Problem \#635}
\label{prob:635}

\noindent\textbf{Status:} OPEN \quad|\quad \textbf{Prize:} none \quad|\quad \textbf{Updated:} 2026-01-30

\noindent\textbf{Tags:} number theory

\medskip

\noindent\textbf{Statement:}

Let $t\geq 1$ and $A\subseteq \{1,\ldots,N\}$ be such that whenever $a,b\in A$ with $b-a\geq t$ we have $b-a\nmid b$. How large can $\lvert A\rvert$ be? Is it true that\[\lvert A\rvert \leq \left(\frac{1}{2}+o_t(1)\right)N?\]

\noindent\rule{\linewidth}{0.4pt}


%% ============================================================
\section{Problem \#638}
\label{prob:638}

\noindent\textbf{Status:} OPEN \quad|\quad \textbf{Prize:} none \quad|\quad \textbf{Updated:} 2025-08-31

\noindent\textbf{Tags:} graph theory, ramsey theory

\medskip

\noindent\textbf{Statement:}

Let $S$ be a family of finite graphs such that for every $n$ there is some $G_n\in S$ such that if the edges of $G_n$ are coloured with $n$ colours then there is a monochromatic triangle. Is it true that for every infinite cardinal $\aleph$ there is a graph $G$ of which every finite subgraph is in $S$ and if the edges of $G$ are coloured with $\aleph$ many colours then there is a monochromatic triangle.

\noindent\rule{\linewidth}{0.4pt}


%% ============================================================
\section{Problem \#640}
\label{prob:640}

\noindent\textbf{Status:} OPEN \quad|\quad \textbf{Prize:} none \quad|\quad \textbf{Updated:} 2025-08-31

\noindent\textbf{Tags:} graph theory, chromatic number

\medskip

\noindent\textbf{Statement:}

Let $k\geq 3$. Does there exist some $f(k)$ such that if a graph $G$ has chromatic number $\geq f(k)$ then $G$ must contain some odd cycle whose vertices span a graph of chromatic number $\geq k$?

\noindent\rule{\linewidth}{0.4pt}


%% ============================================================
\section{Problem \#642}
\label{prob:642}

\noindent\textbf{Status:} OPEN \quad|\quad \textbf{Prize:} none \quad|\quad \textbf{Updated:} 2025-08-31

\noindent\textbf{Tags:} graph theory, cycles

\medskip

\noindent\textbf{Statement:}

Let $f(n)$ be the maximal number of edges in a graph on $n$ vertices such that all cycles have more vertices than chords. Is it true that $f(n)\ll n$?

\noindent\rule{\linewidth}{0.4pt}


%% ============================================================
\section{Problem \#643}
\label{prob:643}

\noindent\textbf{Status:} OPEN \quad|\quad \textbf{Prize:} none \quad|\quad \textbf{Updated:} 2025-08-31

\noindent\textbf{Tags:} graph theory, hypergraphs

\medskip

\noindent\textbf{Statement:}

Let $f(n;t)$ be minimal such that if a $t$-uniform hypergraph on $n$ vertices contains at least $f(n;t)$ edges then there must be four edges $A,B,C,D$ such that\[A\cup B= C\cup D\]and\[A\cap B=C\cap D=\emptyset.\]Estimate $f(n;t)$ - in particular, is it true that for $t\geq 3$\[f(n;t)=(1+o(1))\binom{n}{t-1}?\]

\noindent\rule{\linewidth}{0.4pt}


%% ============================================================
\section{Problem \#644}
\label{prob:644}

\noindent\textbf{Status:} OPEN \quad|\quad \textbf{Prize:} none \quad|\quad \textbf{Updated:} 2025-08-31

\noindent\textbf{Tags:} combinatorics

\medskip

\noindent\textbf{Statement:}

Let $f(k,r)$ be minimal such that if $A_1,A_2,\ldots$ is a family of sets, all of size $k$, such that for every collection of $r$ of the $A_is$ there is some pair $\{x,y\}$ which intersects all of the $A_j$, then there is some set of size $f(k,r)$ which intersects all of the sets $A_i$. Is it true that\[f(k,7)=(1+o(1))\frac{3}{4}k?\]Is it true that for any $r\geq 3$ there exists some constant $c_r$ such that\[f(k,r)=(1+o(1))c_rk?\]

\noindent\rule{\linewidth}{0.4pt}


%% ============================================================
\section{Problem \#653}
\label{prob:653}

\noindent\textbf{Status:} OPEN \quad|\quad \textbf{Prize:} none \quad|\quad \textbf{Updated:} 2025-08-31

\noindent\textbf{Tags:} geometry, distances

\medskip

\noindent\textbf{Statement:}

Let $x_1,\ldots,x_n\in \mathbb{R}^2$ and let $R(x_i)=\#\{ \lvert x_j-x_i\rvert : j\neq i\}$, where the points are ordered such that\[R(x_1)\leq \cdots \leq R(x_n).\]Let $g(n)$ be the maximum number of distinct values the $R(x_i)$ can take. Is it true that $g(n) \geq (1-o(1))n$?

\noindent\rule{\linewidth}{0.4pt}


%% ============================================================
\section{Problem \#654}
\label{prob:654}

\noindent\textbf{Status:} OPEN \quad|\quad \textbf{Prize:} none \quad|\quad \textbf{Updated:} 2025-08-31

\noindent\textbf{Tags:} geometry, distances

\medskip

\noindent\textbf{Statement:}

Let $f(n)$ be such that, given any $x_1,\ldots,x_n\in \mathbb{R}^2$ with no four points on a circle, there exists some $x_i$ with at least $f(n)$ many distinct distances to other $x_j$. Estimate $f(n)$ - in particular, is it true that\[f(n)&#62;(1-o(1))n?\]Or at least\[f(n) &#62; (1/3+c)n\]for some $c&#62;0$, for all large $n$?

\noindent\rule{\linewidth}{0.4pt}


%% ============================================================
\section{Problem \#655}
\label{prob:655}

\noindent\textbf{Status:} OPEN \quad|\quad \textbf{Prize:} none \quad|\quad \textbf{Updated:} 2025-08-31

\noindent\textbf{Tags:} geometry, distances

\medskip
\noindent\textit{ambiguous statement}


\noindent\textbf{Statement:}

Let $x_1,\ldots,x_n\in \mathbb{R}^2$ be such that no circle whose centre is one of the $x_i$ contains three other points. Are there at least\[(1+c)\frac{n}{2}\]distinct distances determined between the $x_i$, for some constant $c&#62;0$ and all $n$ sufficiently large?

\noindent\rule{\linewidth}{0.4pt}


%% ============================================================
\section{Problem \#657}
\label{prob:657}

\noindent\textbf{Status:} OPEN \quad|\quad \textbf{Prize:} none \quad|\quad \textbf{Updated:} 2025-08-31

\noindent\textbf{Tags:} geometry, distances

\medskip

\noindent\textbf{Statement:}

Is it true that if $A\subset \mathbb{R}^2$ is a set of $n$ points such that every subset of $3$ points determines $3$ distinct distances (i.e. $A$ has no isosceles triangles) then $A$ must determine at least $f(n)n$ distinct distances, for some $f(n)\to \infty$?

\noindent\rule{\linewidth}{0.4pt}


%% ============================================================
\section{Problem \#660}
\label{prob:660}

\noindent\textbf{Status:} OPEN \quad|\quad \textbf{Prize:} none \quad|\quad \textbf{Updated:} 2025-08-31

\noindent\textbf{Tags:} geometry, distances, convex

\medskip
\noindent\textit{ambiguous statement}


\noindent\textbf{Statement:}

Let $x_1,\ldots,x_n\in \mathbb{R}^3$ be the vertices of a convex polyhedron. Are there at least\[(1-o(1))\frac{n}{2}\]many distinct distances between the $x_i$?

\noindent\rule{\linewidth}{0.4pt}


%% ============================================================
\section{Problem \#661}
\label{prob:661}

\noindent\textbf{Status:} OPEN \quad|\quad \textbf{Prize:} \$50 \quad|\quad \textbf{Updated:} 2025-08-31

\noindent\textbf{Tags:} geometry, distances

\medskip

\noindent\textbf{Statement:}

Are there, for all large $n$, some points $x_1,\ldots,x_n,y_1,\ldots,y_n\in \mathbb{R}^2$ such that the number of distinct distances $d(x_i,y_j)$ is\[o\left(\frac{n}{\sqrt{\log n}}\right)?\]

\noindent\rule{\linewidth}{0.4pt}


%% ============================================================
\section{Problem \#662}
\label{prob:662}

\noindent\textbf{Status:} OPEN \quad|\quad \textbf{Prize:} none \quad|\quad \textbf{Updated:} 2025-08-31

\noindent\textbf{Tags:} geometry, distances

\medskip
\noindent\textit{ambiguous statement}


\noindent\textbf{Statement:}

Consider the triangular lattice with minimal distance between two points $1$. Denote by $f(t)$ the number of distances from any points $\leq t$. For example $f(1)=6$, $f(\sqrt{3})=12$, and $f(3)=18$. Let $x_1,\ldots,x_n\in \mathbb{R}^2$ be such that $d(x_i,x_j)\geq 1$ for all $i\neq j$. Is it true that, provided $n$ is sufficiently large depending on $t$, the number of distances $d(x_i,x_j)\leq t$ is less than or equal to $f(t)$ with equality perhaps only for the triangular lattice? In particular, is it true that the number of distances $\leq \sqrt{3}-\epsilon$ is less than $1$?

\noindent\rule{\linewidth}{0.4pt}


%% ============================================================
\section{Problem \#663}
\label{prob:663}

\noindent\textbf{Status:} OPEN \quad|\quad \textbf{Prize:} none \quad|\quad \textbf{Updated:} 2025-08-31

\noindent\textbf{Tags:} number theory

\medskip

\noindent\textbf{Statement:}

Let $k\geq 2$ and $q(n,k)$ denote the least prime which does not divide $\prod_{1\leq i\leq k}(n+i)$. Is it true that, if $k$ is fixed and $n$ is sufficiently large, we have\[q(n,k)&#60;(1+o(1))\log n?\]

\noindent\rule{\linewidth}{0.4pt}


%% ============================================================
\section{Problem \#665}
\label{prob:665}

\noindent\textbf{Status:} OPEN \quad|\quad \textbf{Prize:} none \quad|\quad \textbf{Updated:} 2025-08-31

\noindent\textbf{Tags:} combinatorics

\medskip

\noindent\textbf{Statement:}

A pairwise balanced design for $\{1,\ldots,n\}$ is a collection of sets $A_1,\ldots,A_m\subseteq \{1,\ldots,n\}$ such that $2\leq \lvert A_i\rvert &lt;n$ and every pair of distinct elements $x,y\in \{1,\ldots,n\}$ is contained in exactly one $A_i$. Is there a constant $C&gt;0$ and, for all large $n$, a pairwise balanced design such that\[\lvert A_i\rvert &#62; n^{1/2}-C\]for all $1\leq i\leq m$?

\noindent\rule{\linewidth}{0.4pt}


%% ============================================================
\section{Problem \#667}
\label{prob:667}

\noindent\textbf{Status:} OPEN \quad|\quad \textbf{Prize:} none \quad|\quad \textbf{Updated:} 2025-08-31

\noindent\textbf{Tags:} graph theory, ramsey theory

\medskip

\noindent\textbf{Statement:}

Let $p,q\geq 1$ be fixed integers. We define $H(n)=H(N;p,q)$ to be the largest $m$ such that any graph on $n$ vertices where every set of $p$ vertices spans at least $q$ edges must contain a complete graph on $m$ vertices. Is\[c(p,q)=\liminf \frac{\log H(n)}{\log n}\]a strictly increasing function of $q$ for $1\leq q\leq \binom{p-1}{2}+1$?

\noindent\rule{\linewidth}{0.4pt}


%% ============================================================
\section{Problem \#668}
\label{prob:668}

\noindent\textbf{Status:} OPEN \quad|\quad \textbf{Prize:} none \quad|\quad \textbf{Updated:} 2025-08-31

\noindent\textbf{Tags:} geometry, distances

\medskip

\noindent\textbf{Statement:}

Is it true that the number of incongruent sets of $n$ points in $\mathbb{R}^2$ which maximise the number of unit distances tends to infinity as $n\to\infty$? Is it always $&#62;1$ for $n&#62;3$?

\noindent\rule{\linewidth}{0.4pt}


%% ============================================================
\section{Problem \#669}
\label{prob:669}

\noindent\textbf{Status:} OPEN \quad|\quad \textbf{Prize:} none \quad|\quad \textbf{Updated:} 2025-08-31

\noindent\textbf{Tags:} geometry

\medskip
\noindent\textit{orchard problems}


\noindent\textbf{Statement:}

Let $F_k(n)$ be minimal such that for any $n$ points in $\mathbb{R}^2$ there exist at most $F_k(n)$ many distinct lines passing through at least $k$ of the points, and $f_k(n)$ similarly but with lines passing through exactly $k$ points. Estimate $f_k(n)$ and $F_k(n)$ - in particular, determine $\lim F_k(n)/n^2$ and $\lim f_k(n)/n^2$.

\noindent\rule{\linewidth}{0.4pt}


%% ============================================================
\section{Problem \#670}
\label{prob:670}

\noindent\textbf{Status:} OPEN \quad|\quad \textbf{Prize:} none \quad|\quad \textbf{Updated:} 2026-04-16

\noindent\textbf{Tags:} geometry, distances

\medskip

\noindent\textbf{Statement:}

Let $A\subseteq \mathbb{R}^d$ be a set of $n$ points such that all pairwise distances differ by at least $1$. Is the diameter of $A$ at least $(1+o(1))n^2$?

\noindent\rule{\linewidth}{0.4pt}


%% ============================================================
\section{Problem \#671}
\label{prob:671}

\noindent\textbf{Status:} OPEN \quad|\quad \textbf{Prize:} \$250 \quad|\quad \textbf{Updated:} 2025-08-31

\noindent\textbf{Tags:} analysis

\medskip

\noindent\textbf{Statement:}

Given $a_{i}^n\in [-1,1]$ for all $1\leq i\leq n&#60;\infty$ we define $p_{i}^n$ as the unique polynomial of degree $n-1$ such that $p_{i}^n(a_{i}^n)=1$ and $p_{i}^n(a_{i'}^n)=0$ if $1\leq i'\leq n$ with $i\neq i'$. We similarly define\[\mathcal{L}^nf(x) = \sum_{1\leq i\leq n}f(a_i^n)p_i^n(x),\]the unique polynomial of degree $n-1$ which agrees with $f$ on $a_i^n$ for $1\leq i\leq n$ (that is, the sequence of Lagrange interpolation polynomials). Is there such a sequence of $a_i^n$ such that for every continuous $f:[-1,1]\to \mathbb{R}$ there exists some $x\in [-1,1]$ where\[\limsup_{n\to \infty} \sum_{1\leq i\leq n}\lvert p_{i}^n(x)\rvert=\infty\]and yet\[\mathcal{L}^nf(x) \to f(x)?\]Is there such a sequence such that\[\limsup_{n\to \infty} \sum_{1\leq i\leq n}\lvert p_{i}^n(x)\rvert=\infty\]for every $x\in [-1,1]$ and yet for every continuous $f:[-1,1]\to \mathbb{R}$ there exists $x\in [-1,1]$ with\[\mathcal{L}^nf(x) \to f(x)?\]

\noindent\rule{\linewidth}{0.4pt}


%% ============================================================
\section{Problem \#675}
\label{prob:675}

\noindent\textbf{Status:} OPEN \quad|\quad \textbf{Prize:} none \quad|\quad \textbf{Updated:} 2025-08-31

\noindent\textbf{Tags:} number theory

\medskip

\noindent\textbf{Statement:}

We say that $A\subset \mathbb{N}$ has the translation property if, for every $n$, there exists some integer $t_n\geq 1$ such that, for all $1\leq a\leq n$,\[a\in A\quad\textrm{ if and only if }\quad a+t_n\in A.\]{UL} {LI}Does the set of the sums of two squares have the translation property?{/LI} {LI}If we partition all primes into $P\sqcup Q$, such that each set contains $\gg x/\log x$ many primes $\leq x$ for all large $x$, then can the set of integers only divisible by primes from $P$ have the translation property?{/LI} {LI}If $A$ is the set of squarefree numbers then how fast does the minimal such $t_n$ grow? Is it true that $t_n&#62;\exp(n^c)$ for some constant $c&#62;0$?{/LI} {/UL}

\noindent\rule{\linewidth}{0.4pt}


%% ============================================================
\section{Problem \#676}
\label{prob:676}

\noindent\textbf{Status:} OPEN \quad|\quad \textbf{Prize:} none \quad|\quad \textbf{Updated:} 2025-08-31

\noindent\textbf{Tags:} number theory

\medskip

\noindent\textbf{Statement:}

Is every sufficiently large integer of the form\[ap^2+b\]for some prime $p$ and integer $a\geq 1$ and $0\leq b&#60;p$?

\noindent\rule{\linewidth}{0.4pt}


%% ============================================================
\section{Problem \#677}
\label{prob:677}

\noindent\textbf{Status:} OPEN \quad|\quad \textbf{Prize:} none \quad|\quad \textbf{Updated:} 2025-08-31

\noindent\textbf{Tags:} number theory

\medskip

\noindent\textbf{Statement:}

Let $M(n,k)=[n+1,\ldots,n+k]$ be the least common multiple of $\{n+1,\ldots,n+k\}$. Is it true that for all $m\geq n+k$\[M(n,k) \neq M(m,k)?\]

\noindent\rule{\linewidth}{0.4pt}


%% ============================================================
\section{Problem \#679}
\label{prob:679}

\noindent\textbf{Status:} OPEN \quad|\quad \textbf{Prize:} none \quad|\quad \textbf{Updated:} 2025-08-31

\noindent\textbf{Tags:} number theory

\medskip

\noindent\textbf{Statement:}

Let $\epsilon&#62;0$ and $\omega(n)$ count the number of distinct prime factors of $n$. Are there infinitely many values of $n$ such that\[\omega(n-k) &#60; (1+\epsilon)\frac{\log k}{\log\log k}\]for all $k&#60;n$ which are sufficiently large depending on $\epsilon$ only? Can one show the stronger version with\[\omega(n-k) &#60; \frac{\log k}{\log\log k}+O(1)\]is false?

\noindent\rule{\linewidth}{0.4pt}


%% ============================================================
\section{Problem \#680}
\label{prob:680}

\noindent\textbf{Status:} OPEN \quad|\quad \textbf{Prize:} none \quad|\quad \textbf{Updated:} 2025-08-31

\noindent\textbf{Tags:} number theory, primes

\medskip

\noindent\textbf{Statement:}

Is it true that, for all sufficiently large $n$, there exists some $k$ such that\[p(n+k)&#62;k^2+1,\]where $p(m)$ denotes the least prime factor of $m$? Can one prove this is false if we replace $k^2+1$ by $e^{(1+\epsilon)\sqrt{k}}+C_\epsilon$, for all $\epsilon&#62;0$, where $C_\epsilon&#62;0$ is some constant?

\noindent\rule{\linewidth}{0.4pt}


%% ============================================================
\section{Problem \#681}
\label{prob:681}

\noindent\textbf{Status:} OPEN \quad|\quad \textbf{Prize:} none \quad|\quad \textbf{Updated:} 2025-08-31

\noindent\textbf{Tags:} number theory, primes

\medskip

\noindent\textbf{Statement:}

Is it true that for all large $n$ there exists $k$ such that $n+k$ is composite and\[p(n+k)&#62;k^2,\]where $p(m)$ is the least prime factor of $m$?

\noindent\rule{\linewidth}{0.4pt}


%% ============================================================
\section{Problem \#683}
\label{prob:683}

\noindent\textbf{Status:} OPEN \quad|\quad \textbf{Prize:} none \quad|\quad \textbf{Updated:} 2025-09-04

\noindent\textbf{Tags:} number theory, primes, binomial coefficients

\medskip

\noindent\textbf{Statement:}

Is it true that for every $1\leq k\leq n$ the largest prime divisor of $\binom{n}{k}$, say $P(\binom{n}{k})$, satisfies\[P\left(\binom{n}{k}\right)\geq \min(n-k+1, k^{1+c})\]for some constant $c&#62;0$?

\noindent\rule{\linewidth}{0.4pt}


%% ============================================================
\section{Problem \#684}
\label{prob:684}

\noindent\textbf{Status:} OPEN \quad|\quad \textbf{Prize:} none \quad|\quad \textbf{Updated:} 2025-08-31

\noindent\textbf{Tags:} number theory, primes, binomial coefficients

\medskip

\noindent\textbf{Statement:}

For $0\leq k\leq n$ write\[\binom{n}{k} = uv\]where the only primes dividing $u$ are in $[2,k]$ and the only primes dividing $v$ are in $(k,n]$. Let $f(n)$ be the smallest $k$ such that $u&#62;n^2$. Give bounds for $f(n)$.

\noindent\rule{\linewidth}{0.4pt}


%% ============================================================
\section{Problem \#685}
\label{prob:685}

\noindent\textbf{Status:} OPEN \quad|\quad \textbf{Prize:} none \quad|\quad \textbf{Updated:} 2025-08-31

\noindent\textbf{Tags:} number theory, primes, binomial coefficients

\medskip

\noindent\textbf{Statement:}

Let $\epsilon&#62;0$ and $n$ be large depending on $\epsilon$. Is it true that for all $n^\epsilon&#60;k\leq n^{1-\epsilon}$ the number of distinct prime divisors of $\binom{n}{k}$ is\[(1+o(1))k\sum_{k&#60;p&#60;n}\frac{1}{p}?\]Or perhaps even when $k \geq (\log n)^c$?

\noindent\rule{\linewidth}{0.4pt}


%% ============================================================
\section{Problem \#686}
\label{prob:686}

\noindent\textbf{Status:} OPEN \quad|\quad \textbf{Prize:} none \quad|\quad \textbf{Updated:} 2025-08-31

\noindent\textbf{Tags:} number theory

\medskip

\noindent\textbf{Statement:}

Can every integer $N\geq 2$ be written as\[N=\frac{\prod_{1\leq i\leq k}(m+i)}{\prod_{1\leq i\leq k}(n+i)}\]for some $k\geq 2$ and $m\geq n+k$?

\noindent\rule{\linewidth}{0.4pt}


%% ============================================================
\section{Problem \#687}
\label{prob:687}

\noindent\textbf{Status:} OPEN \quad|\quad \textbf{Prize:} \$1000 \quad|\quad \textbf{Updated:} 2025-08-31

\noindent\textbf{Tags:} number theory

\medskip

\noindent\textbf{Statement:}

Let $Y(x)$ be the maximal $y$ such that there exists a choice of congruence classes $a_p$ for all primes $p\leq x$ such that every integer in $[1,y]$ is congruent to at least one of the $a_p\pmod{p}$. Give good estimates for $Y(x)$. In particular, can one prove that $Y(x)=o(x^2)$ or even $Y(x)\ll x^{1+o(1)}$?

\noindent\rule{\linewidth}{0.4pt}


%% ============================================================
\section{Problem \#688}
\label{prob:688}

\noindent\textbf{Status:} OPEN \quad|\quad \textbf{Prize:} none \quad|\quad \textbf{Updated:} 2025-08-31

\noindent\textbf{Tags:} number theory

\medskip

\noindent\textbf{Statement:}

Define $\epsilon_n$ to be maximal such that there exists some choice of congruence class $a_p$ for all primes $n^{\epsilon_n}&#60;p\leq n$ such that every integer in $[1,n]$ satisfies at least one of the congruences $\equiv a_p\pmod{p}$. Estimate $\epsilon_n$ - in particular is it true that $\epsilon_n=o(1)$?

\noindent\rule{\linewidth}{0.4pt}


%% ============================================================
\section{Problem \#689}
\label{prob:689}

\noindent\textbf{Status:} OPEN \quad|\quad \textbf{Prize:} none \quad|\quad \textbf{Updated:} 2025-08-31

\noindent\textbf{Tags:} number theory

\medskip

\noindent\textbf{Statement:}

Let $n$ be sufficiently large. Is there some choice of congruence class $a_p$ for all primes $2\leq p\leq n$ such that every integer in $[1,n]$ satisfies at least two of the congruences $\equiv a_p\pmod{p}$?

\noindent\rule{\linewidth}{0.4pt}


%% ============================================================
\section{Problem \#691}
\label{prob:691}

\noindent\textbf{Status:} OPEN \quad|\quad \textbf{Prize:} none \quad|\quad \textbf{Updated:} 2025-08-31

\noindent\textbf{Tags:} number theory

\medskip

\noindent\textbf{Statement:}

Given $A\subseteq \mathbb{N}$ let $M_A=\{ n \geq 1 : a\mid n\textrm{ for some }a\in A\}$ be the set of multiples of $A$. Find a necessary and sufficient condition on $A$ for $M_A$ to have density $1$.

\noindent\rule{\linewidth}{0.4pt}


%% ============================================================
\section{Problem \#693}
\label{prob:693}

\noindent\textbf{Status:} OPEN \quad|\quad \textbf{Prize:} none \quad|\quad \textbf{Updated:} 2025-08-31

\noindent\textbf{Tags:} number theory, divisors

\medskip

\noindent\textbf{Statement:}

Let $k\geq 2$ and $n$ be sufficiently large depending on $k$. Let $A=\{a_1&#60;a_2&#60;\cdots \}$ be the set of those integers in $[n,n^k]$ which have a divisor in $(n,2n)$. Estimate\[\max_{i} a_{i+1}-a_i.\]Is this $\leq (\log n)^{O(1)}$?

\noindent\rule{\linewidth}{0.4pt}


%% ============================================================
\section{Problem \#695}
\label{prob:695}

\noindent\textbf{Status:} OPEN \quad|\quad \textbf{Prize:} none \quad|\quad \textbf{Updated:} 2025-08-31

\noindent\textbf{Tags:} number theory

\medskip

\noindent\textbf{Statement:}

Let $p_1&#60;p_2&#60;\cdots$ be a sequence of primes such that $p_{i+1}\equiv 1\pmod{p_i}$. Is it true that\[\lim_k p_k^{1/k}=\infty?\]Does there exist such a sequence with\[p_k \leq \exp(k(\log k)^{1+o(1)})?\]

\noindent\rule{\linewidth}{0.4pt}


%% ============================================================
\section{Problem \#700}
\label{prob:700}

\noindent\textbf{Status:} OPEN \quad|\quad \textbf{Prize:} none \quad|\quad \textbf{Updated:} 2025-08-31

\noindent\textbf{Tags:} number theory, binomial coefficients

\medskip

\noindent\textbf{Statement:}

Let\[f(n)=\min_{1&lt;k\leq n/2}\textrm{gcd}\left(n,\binom{n}{k}\right).\]{UL} {LI}Characterise those composite $n$ such that $f(n)=n/P(n)$, where $P(n)$ is the largest prime dividing $n$.{/LI} {LI}Are there infinitely many composite $n$ such that $f(n)&gt;n^{1/2}$?{/LI} {LI} Is it true that, for every composite $n$,\[f(n) \ll_A \frac{n}{(\log n)^A}\]for every $A&#62;0$?{/LI} {/UL}

\noindent\rule{\linewidth}{0.4pt}


%% ============================================================
\section{Problem \#701}
\label{prob:701}

\noindent\textbf{Status:} OPEN \quad|\quad \textbf{Prize:} none \quad|\quad \textbf{Updated:} 2025-08-31

\noindent\textbf{Tags:} combinatorics, intersecting family

\medskip

\noindent\textbf{Statement:}

Let $\mathcal{F}$ be a family of sets closed under taking subsets (i.e. if $B\subseteq A\in\mathcal{F}$ then $B\in \mathcal{F}$). There exists some element $x$ such that whenever $\mathcal{F}'\subseteq \mathcal{F}$ is an intersecting subfamily we have\[\lvert \mathcal{F}'\rvert \leq \lvert \{ A\in \mathcal{F} : x\in A\}\rvert.\]

\noindent\rule{\linewidth}{0.4pt}


%% ============================================================
\section{Problem \#704}
\label{prob:704}

\noindent\textbf{Status:} OPEN \quad|\quad \textbf{Prize:} none \quad|\quad \textbf{Updated:} 2025-08-31

\noindent\textbf{Tags:} graph theory, geometry, chromatic number

\medskip

\noindent\textbf{Statement:}

Let $G_n$ be the unit distance graph in $\mathbb{R}^n$, with two vertices joined by an edge if and only if the distance between them is $1$. Estimate the chromatic number $\chi(G_n)$. Does it grow exponentially in $n$? Does\[\lim_{n\to \infty}\chi(G_n)^{1/n}\]exist?

\noindent\rule{\linewidth}{0.4pt}


%% ============================================================
\section{Problem \#706}
\label{prob:706}

\noindent\textbf{Status:} OPEN \quad|\quad \textbf{Prize:} none \quad|\quad \textbf{Updated:} 2025-08-31

\noindent\textbf{Tags:} graph theory, chromatic number

\medskip

\noindent\textbf{Statement:}

Let $L(r)$ be such that if $G$ is a graph formed by taking a finite set of points $P$ in $\mathbb{R}^2$ and some set $A\subset (0,\infty)$ of size $r$, where the vertex set is $P$ and there is an edge between two points if and only if their distance is a member of $A$, then $\chi(G)\leq L(r)$. Estimate $L(r)$. In particular, is it true that $L(r)\leq r^{O(1)}$?

\noindent\rule{\linewidth}{0.4pt}


%% ============================================================
\section{Problem \#708}
\label{prob:708}

\noindent\textbf{Status:} OPEN \quad|\quad \textbf{Prize:} \$100 \quad|\quad \textbf{Updated:} 2025-08-31

\noindent\textbf{Tags:} number theory

\medskip

\noindent\textbf{Statement:}

Let $g(n)$ be minimal such that for any $A\subseteq [2,\infty)\cap \mathbb{N}$ with $\lvert A\rvert =n$ and any set $I$ of $\max(A)$ consecutive integers there exists some $B\subseteq I$ with $\lvert B\rvert=g(n)$ such that\[\prod_{a\in A} a \mid \prod_{b\in B}b.\]Is it true that\[g(n) \leq (2+o(1))n?\]Or perhaps even $g(n)\leq 2n$?

\noindent\rule{\linewidth}{0.4pt}


%% ============================================================
\section{Problem \#709}
\label{prob:709}

\noindent\textbf{Status:} OPEN \quad|\quad \textbf{Prize:} none \quad|\quad \textbf{Updated:} 2025-08-31

\noindent\textbf{Tags:} number theory

\medskip

\noindent\textbf{Statement:}

Let $f(n)$ be minimal such that, for any $A=\{a_1,\ldots,a_n\}\subseteq [2,\infty)\cap\mathbb{N}$ of size $n$, in any interval $I$ of $f(n)\max(A)$ consecutive integers there exist distinct $x_1,\ldots,x_n\in I$ such that $a_i\mid x_i$. Obtain good bounds for $f(n)$, or even an asymptotic formula.

\noindent\rule{\linewidth}{0.4pt}


%% ============================================================
\section{Problem \#710}
\label{prob:710}

\noindent\textbf{Status:} OPEN \quad|\quad \textbf{Prize:} ₹2000 \quad|\quad \textbf{Updated:} 2025-08-31

\noindent\textbf{Tags:} number theory

\medskip

\noindent\textbf{Statement:}

Let $f(n)$ be minimal such that in $(n,n+f(n))$ there exist distinct integers $a_1,\ldots,a_n$ such that $k\mid a_k$ for all $1\leq k\leq n$. Obtain an asymptotic formula for $f(n)$.

\noindent\rule{\linewidth}{0.4pt}


%% ============================================================
\section{Problem \#711}
\label{prob:711}

\noindent\textbf{Status:} OPEN \quad|\quad \textbf{Prize:} ₹1000 \quad|\quad \textbf{Updated:} 2025-08-31

\noindent\textbf{Tags:} number theory

\medskip

\noindent\textbf{Statement:}

Let $f(n,m)$ be minimal such that in $(m,m+f(n,m))$ there exist distinct integers $a_1,\ldots,a_n$ such that $k\mid a_k$ for all $1\leq k\leq n$. Prove that\[\max_m f(n,m) \leq n^{1+o(1)}\]and that\[\max_m (f(n,m)-f(n,n))\to \infty.\]

\noindent\rule{\linewidth}{0.4pt}


%% ============================================================
\section{Problem \#712}
\label{prob:712}

\noindent\textbf{Status:} OPEN \quad|\quad \textbf{Prize:} \$500 \quad|\quad \textbf{Updated:} 2025-08-31

\noindent\textbf{Tags:} graph theory, turan number, hypergraphs

\medskip

\noindent\textbf{Statement:}

Determine, for any $k&#62;r&#62;2$, the value of\[\frac{\mathrm{ex}_r(n,K_k^r)}{\binom{n}{r}},\]where $\mathrm{ex}_r(n,K_k^r)$ is the largest number of $r$-edges which can placed on $n$ vertices so that there exists no set of $k$ vertices which is covered by all $\binom{k}{r}$ possible $r$-edges.

\noindent\rule{\linewidth}{0.4pt}


%% ============================================================
\section{Problem \#713}
\label{prob:713}

\noindent\textbf{Status:} OPEN \quad|\quad \textbf{Prize:} \$500 \quad|\quad \textbf{Updated:} 2025-08-31

\noindent\textbf{Tags:} graph theory, turan number

\medskip

\noindent\textbf{Statement:}

Is it true that, for every bipartite graph $G$, there exists some $\alpha\in [1,2)$ and $c&#62;0$ such that\[\mathrm{ex}(n;G)\sim cn^\alpha?\]Must $\alpha$ be rational?

\noindent\rule{\linewidth}{0.4pt}


%% ============================================================
\section{Problem \#714}
\label{prob:714}

\noindent\textbf{Status:} OPEN \quad|\quad \textbf{Prize:} none \quad|\quad \textbf{Updated:} 2025-08-31

\noindent\textbf{Tags:} graph theory, turan number

\medskip

\noindent\textbf{Statement:}

Is it true that\[\mathrm{ex}(n; K_{r,r}) \gg n^{2-1/r}?\]

\noindent\rule{\linewidth}{0.4pt}


%% ============================================================
\section{Problem \#719}
\label{prob:719}

\noindent\textbf{Status:} OPEN \quad|\quad \textbf{Prize:} none \quad|\quad \textbf{Updated:} 2025-08-31

\noindent\textbf{Tags:} graph theory, hypergraphs

\medskip

\noindent\textbf{Statement:}

Let $\mathrm{ex}_r(n;K_{r+1}^r)$ be the maximum number of $r$-edges that can be placed on $n$ vertices without forming a $K_{r+1}^r$ (the $r$-uniform complete graph on $r+1$ vertices). Is every $r$-hypergraph $G$ on $n$ vertices the union of at most $\mathrm{ex}_{r}(n;K_{r+1}^r)$ many copies of $K_r^r$ and $K_{r+1}^r$, no two of which share a $K_r^r$?

\noindent\rule{\linewidth}{0.4pt}


%% ============================================================
\section{Problem \#724}
\label{prob:724}

\noindent\textbf{Status:} OPEN \quad|\quad \textbf{Prize:} none \quad|\quad \textbf{Updated:} 2025-08-31

\noindent\textbf{Tags:} combinatorics

\medskip

\noindent\textbf{Statement:}

Let $f(n)$ be the maximum number of mutually orthogonal Latin squares of order $n$. Is it true that\[f(n) \gg n^{1/2}?\]

\noindent\rule{\linewidth}{0.4pt}


%% ============================================================
\section{Problem \#725}
\label{prob:725}

\noindent\textbf{Status:} OPEN \quad|\quad \textbf{Prize:} none \quad|\quad \textbf{Updated:} 2025-08-31

\noindent\textbf{Tags:} combinatorics

\medskip

\noindent\textbf{Statement:}

Give an asymptotic formula for the number of $k\times n$ Latin rectangles .

\noindent\rule{\linewidth}{0.4pt}


%% ============================================================
\section{Problem \#726}
\label{prob:726}

\noindent\textbf{Status:} OPEN \quad|\quad \textbf{Prize:} none \quad|\quad \textbf{Updated:} 2025-08-31

\noindent\textbf{Tags:} number theory

\medskip

\noindent\textbf{Statement:}

As $n\to \infty$ ranges over integers\[\sum_{p\leq n}1_{n\in (p/2,p)\pmod{p}}\frac{1}{p}\sim \frac{\log\log n}{2}.\]

\noindent\rule{\linewidth}{0.4pt}


%% ============================================================
\section{Problem \#727}
\label{prob:727}

\noindent\textbf{Status:} OPEN \quad|\quad \textbf{Prize:} none \quad|\quad \textbf{Updated:} 2025-08-31

\noindent\textbf{Tags:} number theory, factorials

\medskip

\noindent\textbf{Statement:}

Let $k\geq 2$. Does\[(n+k)!^2 \mid (2n)!\]for infinitely many $n$?

\noindent\rule{\linewidth}{0.4pt}


%% ============================================================
\section{Problem \#730}
\label{prob:730}

\noindent\textbf{Status:} OPEN \quad|\quad \textbf{Prize:} none \quad|\quad \textbf{Updated:} 2025-08-31

\noindent\textbf{Tags:} number theory, binomial coefficients, base representations

\medskip

\noindent\textbf{Statement:}

Are there infinitely many pairs of integers $n\neq m$ such that $\binom{2n}{n}$ and $\binom{2m}{m}$ have the same set of prime divisors?

\noindent\rule{\linewidth}{0.4pt}


%% ============================================================
\section{Problem \#731}
\label{prob:731}

\noindent\textbf{Status:} OPEN \quad|\quad \textbf{Prize:} none \quad|\quad \textbf{Updated:} 2025-08-31

\noindent\textbf{Tags:} number theory, binomial coefficients

\medskip

\noindent\textbf{Statement:}

Find some reasonable function $f(n)$ such that, for almost all integers $n$, the least integer $m$ such that $m\nmid \binom{2n}{n}$ satisfies\[m\sim f(n).\]

\noindent\rule{\linewidth}{0.4pt}


%% ============================================================
\section{Problem \#734}
\label{prob:734}

\noindent\textbf{Status:} OPEN \quad|\quad \textbf{Prize:} none \quad|\quad \textbf{Updated:} 2025-08-31

\noindent\textbf{Tags:} combinatorics

\medskip

\noindent\textbf{Statement:}

Find, for all large $n$, a non-trivial pairwise balanced block design $A_1,\ldots,A_m\subseteq \{1,\ldots,n\}$ such that, for all $t$, there are $O(n^{1/2})$ many $i$ such that $\lvert A_i\rvert=t$.

\noindent\rule{\linewidth}{0.4pt}


%% ============================================================
\section{Problem \#738}
\label{prob:738}

\noindent\textbf{Status:} OPEN \quad|\quad \textbf{Prize:} none \quad|\quad \textbf{Updated:} 2025-08-31

\noindent\textbf{Tags:} graph theory, chromatic number

\medskip

\noindent\textbf{Statement:}

If $G$ has infinite chromatic number and is triangle-free (contains no $K_3$) then must $G$ contain every tree as an induced subgraph?

\noindent\rule{\linewidth}{0.4pt}


%% ============================================================
\section{Problem \#740}
\label{prob:740}

\noindent\textbf{Status:} OPEN \quad|\quad \textbf{Prize:} none \quad|\quad \textbf{Updated:} 2025-08-31

\noindent\textbf{Tags:} graph theory, chromatic number

\medskip

\noindent\textbf{Statement:}

Let $\mathfrak{m}$ be an infinite cardinal and $G$ be a graph with chromatic number $\mathfrak{m}$. Let $r\geq 1$. Must $G$ contain a subgraph of chromatic number $\mathfrak{m}$ which does not contain any odd cycle of length $\leq r$?

\noindent\rule{\linewidth}{0.4pt}


%% ============================================================
\section{Problem \#749}
\label{prob:749}

\noindent\textbf{Status:} OPEN \quad|\quad \textbf{Prize:} none \quad|\quad \textbf{Updated:} 2025-08-31

\noindent\textbf{Tags:} additive combinatorics

\medskip

\noindent\textbf{Statement:}

Let $\epsilon&#62;0$. Does there exist $A\subseteq \mathbb{N}$ such that the lower density of $A+A$ is at least $1-\epsilon$ and yet $1_A\ast 1_A(n) \ll_\epsilon 1$ for all $n$?

\noindent\rule{\linewidth}{0.4pt}


%% ============================================================
\section{Problem \#750}
\label{prob:750}

\noindent\textbf{Status:} OPEN \quad|\quad \textbf{Prize:} none \quad|\quad \textbf{Updated:} 2025-08-31

\noindent\textbf{Tags:} graph theory, chromatic number

\medskip

\noindent\textbf{Statement:}

Let $f(m)$ be some function such that $f(m)\to \infty$ as $m\to \infty$. Does there exist a graph $G$ of infinite chromatic number such that every subgraph on $m$ vertices contains an independent set of size at least $\frac{m}{2}-f(m)$?

\noindent\rule{\linewidth}{0.4pt}


%% ============================================================
\section{Problem \#757}
\label{prob:757}

\noindent\textbf{Status:} OPEN \quad|\quad \textbf{Prize:} none \quad|\quad \textbf{Updated:} 2025-08-31

\noindent\textbf{Tags:} geometry, distances, sidon sets

\medskip

\noindent\textbf{Statement:}

Let $A\subset \mathbb{R}$ be a set of size $n$ such that every subset $B\subseteq A$ with $\lvert B\rvert =4$ has $\lvert B-B\rvert\geq 11$. Find the best constant $c&#62;0$ such that $A$ must always contain a Sidon set of size $\geq cn$.

\noindent\rule{\linewidth}{0.4pt}


%% ============================================================
\section{Problem \#761}
\label{prob:761}

\noindent\textbf{Status:} OPEN \quad|\quad \textbf{Prize:} none \quad|\quad \textbf{Updated:} 2025-08-31

\noindent\textbf{Tags:} graph theory, chromatic number

\medskip

\noindent\textbf{Statement:}

The cochromatic number of $G$, denoted by $\zeta(G)$, is the minimum number of colours needed to colour the vertices of $G$ such that each colour class induces either a complete graph or empty graph. The dichromatic number of $G$, denoted by $\delta(G)$, is the minimum number $k$ of colours required such that, in any orientation of the edges of $G$, there is a $k$-colouring of the vertices of $G$ such that there are no monochromatic oriented cycles. Must a graph with large chromatic number have large dichromatic number? Must a graph with large cochromatic number contain a graph with large dichromatic number?

\noindent\rule{\linewidth}{0.4pt}


%% ============================================================
\section{Problem \#766}
\label{prob:766}

\noindent\textbf{Status:} OPEN \quad|\quad \textbf{Prize:} none \quad|\quad \textbf{Updated:} 2025-08-31

\noindent\textbf{Tags:} graph theory, turan number

\medskip

\noindent\textbf{Statement:}

Let $f(n;k,l)=\min \mathrm{ex}(n;G)$, where $G$ ranges over all graphs with $k$ vertices and $l$ edges. Give good estimates for $f(n;k,l)$ in the range $k&#60;l\leq k^2/4$. For fixed $k$ and large $n$ is $f(n;k,l)$ a strictly monotone function of $l$?

\noindent\rule{\linewidth}{0.4pt}


%% ============================================================
\section{Problem \#768}
\label{prob:768}

\noindent\textbf{Status:} OPEN \quad|\quad \textbf{Prize:} none \quad|\quad \textbf{Updated:} 2025-08-31

\noindent\textbf{Tags:} number theory

\medskip

\noindent\textbf{Statement:}

Let $A\subset\mathbb{N}$ be the set of $n$ such that for every prime $p\mid n$ there exists some $d\mid n$ with $d&#62;1$ such that $d\equiv 1\pmod{p}$. Is it true that there exists some constant $c&#62;0$ such that for all large $N$\[\frac{\lvert A\cap [1,N]\rvert}{N}=\exp(-(c+o(1))\sqrt{\log N}\log\log N).\]

\noindent\rule{\linewidth}{0.4pt}


%% ============================================================
\section{Problem \#769}
\label{prob:769}

\noindent\textbf{Status:} OPEN \quad|\quad \textbf{Prize:} none \quad|\quad \textbf{Updated:} 2025-08-31

\noindent\textbf{Tags:} number theory, geometry

\medskip

\noindent\textbf{Statement:}

Let $c(n)$ be minimal such that if $k\geq c(n)$ then the $n$-dimensional unit cube can be decomposed into $k$ homothetic $n$-dimensional cubes. Give good bounds for $c(n)$ - in particular, is it true that $c(n) \gg n^n$?

\noindent\rule{\linewidth}{0.4pt}


%% ============================================================
\section{Problem \#770}
\label{prob:770}

\noindent\textbf{Status:} OPEN \quad|\quad \textbf{Prize:} none \quad|\quad \textbf{Updated:} 2025-08-31

\noindent\textbf{Tags:} number theory

\medskip

\noindent\textbf{Statement:}

Let $h(n)$ be minimal such that $2^n-1,3^n-1,\ldots,h(n)^n-1$ are mutually coprime. Does, for every prime $p$, the density $\delta_p$ of integers with $h(n)=p$ exist? Does $\liminf h(n)=\infty$? Is it true that if $p$ is the greatest prime such that $p-1\mid n$ and $p&#62;n^\epsilon$ then $h(n)=p$?

\noindent\rule{\linewidth}{0.4pt}


%% ============================================================
\section{Problem \#773}
\label{prob:773}

\noindent\textbf{Status:} OPEN \quad|\quad \textbf{Prize:} none \quad|\quad \textbf{Updated:} 2025-08-31

\noindent\textbf{Tags:} number theory, sidon sets, squares

\medskip

\noindent\textbf{Statement:}

What is the size of the largest Sidon subset $A\subseteq\{1,2^2,\ldots,N^2\}$? Is it $N^{1-o(1)}$?

\noindent\rule{\linewidth}{0.4pt}


%% ============================================================
\section{Problem \#774}
\label{prob:774}

\noindent\textbf{Status:} OPEN \quad|\quad \textbf{Prize:} none \quad|\quad \textbf{Updated:} 2025-08-31

\noindent\textbf{Tags:} number theory

\medskip

\noindent\textbf{Statement:}

We call $A\subset \mathbb{N}$ dissociated if $\sum_{n\in X}n\neq \sum_{m\in Y}m$ for all finite $X,Y\subset A$ with $X\neq Y$. Let $A\subset \mathbb{N}$ be an infinite set. We call $A$ proportionately dissociated if every finite $B\subset A$ contains a dissociated set of size $\gg \lvert B\rvert$. Is every proportionately dissociated set the union of a finite number of dissociated sets?

\noindent\rule{\linewidth}{0.4pt}


%% ============================================================
\section{Problem \#776}
\label{prob:776}

\noindent\textbf{Status:} OPEN \quad|\quad \textbf{Prize:} none \quad|\quad \textbf{Updated:} 2025-08-31

\noindent\textbf{Tags:} combinatorics

\medskip

\noindent\textbf{Statement:}

Let $r\geq 2$ and $A_1,\ldots,A_m\subseteq \{1,\ldots,n\}$ be such that $A_i\not\subseteq A_j$ for all $i\neq j$ and for any $t$ if there exists some $i$ with $\lvert A_i\rvert=t$ then there must exist at least $r$ sets of that size. How large must $n$ be (as a function of $r$) to ensure that there is such a family which achieves $n-3$ distinct sizes of sets?

\noindent\rule{\linewidth}{0.4pt}


%% ============================================================
\section{Problem \#778}
\label{prob:778}

\noindent\textbf{Status:} OPEN \quad|\quad \textbf{Prize:} none \quad|\quad \textbf{Updated:} 2025-08-31

\noindent\textbf{Tags:} graph theory

\medskip

\noindent\textbf{Statement:}

Alice and Bob play a game on the edges of $K_n$, alternating colouring edges by red (Alice) and blue (Bob). Alice goes first, and wins if at the end the largest red clique is larger than any of the blue cliques. Does Bob have a winning strategy for $n\geq 3$? (Erd\H{o}s believed the answer is yes.) If we change the game so that Bob colours two edges after each edge that Alice colours, but now require Bob's largest clique to be strictly larger than Alice's, then does Bob have a winning strategy for $n&#62;3$? Finally, consider the game when Alice wins if the maximum degree of the red subgraph is larger than the maximum degree of the blue subgraph. Who wins?

\noindent\rule{\linewidth}{0.4pt}


%% ============================================================
\section{Problem \#782}
\label{prob:782}

\noindent\textbf{Status:} OPEN \quad|\quad \textbf{Prize:} none \quad|\quad \textbf{Updated:} 2025-08-31

\noindent\textbf{Tags:} number theory

\medskip

\noindent\textbf{Statement:}

Do the squares contain arbitrarily long quasi-progressions? That is, does there exist some constant $C&#62;0$ such that, for any $k$, the squares contain a sequence $x_1,\ldots,x_k$ where, for some $d$ and all $1\leq i&#60;k$,\[x_i+d\leq x_{i+1}\leq x_i+d+C.\]Do the squares contain arbitrarily large cubes\[a+\left\{ \sum_i \epsilon_ib_i : \epsilon_i\in \{0,1\}\right\}?\]

\noindent\rule{\linewidth}{0.4pt}


%% ============================================================
\section{Problem \#786}
\label{prob:786}

\noindent\textbf{Status:} OPEN \quad|\quad \textbf{Prize:} none \quad|\quad \textbf{Updated:} 2025-08-31

\noindent\textbf{Tags:} number theory

\medskip

\noindent\textbf{Statement:}

Let $\epsilon&#62;0$. Is there some set $A\subset \mathbb{N}$ of density $&#62;1-\epsilon$ such that $a_1\cdots a_r=b_1\cdots b_s$ with $a_i,b_j\in A$ can only hold when $r=s$? Similarly, can one always find a set $A\subset\{1,\ldots,N\}$ with this property of size $\geq (1-o(1))N$?

\noindent\rule{\linewidth}{0.4pt}


%% ============================================================
\section{Problem \#787}
\label{prob:787}

\noindent\textbf{Status:} OPEN \quad|\quad \textbf{Prize:} none \quad|\quad \textbf{Updated:} 2025-08-31

\noindent\textbf{Tags:} additive combinatorics

\medskip

\noindent\textbf{Statement:}

Let $g(n)$ be maximal such that given any set $A\subset \mathbb{R}$ with $\lvert A\rvert=n$ there exists some $B\subseteq A$ of size $\lvert B\rvert\geq g(n)$ such that $b_1+b_2\not\in A$ for all $b_1\neq b_2\in B$. Estimate $g(n)$.

\noindent\rule{\linewidth}{0.4pt}


%% ============================================================
\section{Problem \#788}
\label{prob:788}

\noindent\textbf{Status:} OPEN \quad|\quad \textbf{Prize:} none \quad|\quad \textbf{Updated:} 2025-08-31

\noindent\textbf{Tags:} additive combinatorics

\medskip

\noindent\textbf{Statement:}

Let $f(n)$ be maximal such that if $B\subset (2n,4n)\cap \mathbb{N}$ there exists some $C\subset (n,2n)\cap \mathbb{N}$ such that $c_1+c_2\not\in B$ for all $c_1\neq c_2\in C$ and $\lvert C\rvert+\lvert B\rvert \geq f(n)$. Estimate $f(n)$. In particular is it true that $f(n)\leq n^{1/2+o(1)}$?

\noindent\rule{\linewidth}{0.4pt}


%% ============================================================
\section{Problem \#789}
\label{prob:789}

\noindent\textbf{Status:} OPEN \quad|\quad \textbf{Prize:} none \quad|\quad \textbf{Updated:} 2025-08-31

\noindent\textbf{Tags:} additive combinatorics

\medskip

\noindent\textbf{Statement:}

Let $h(n)$ be maximal such that if $A\subseteq \mathbb{Z}$ with $\lvert A\rvert=n$ then there is $B\subseteq A$ with $\lvert B\rvert \geq h(n)$ such that if $a_1+\cdots+a_r=b_1+\cdots+b_s$ with $a_i,b_i\in B$ then $r=s$. Estimate $h(n)$.

\noindent\rule{\linewidth}{0.4pt}


%% ============================================================
\section{Problem \#790}
\label{prob:790}

\noindent\textbf{Status:} OPEN \quad|\quad \textbf{Prize:} none \quad|\quad \textbf{Updated:} 2025-08-31

\noindent\textbf{Tags:} additive combinatorics

\medskip

\noindent\textbf{Statement:}

Let $l(n)$ be maximal such that if $A\subset\mathbb{Z}$ with $\lvert A\rvert=n$ then there exists a sum-free $B\subseteq A$ with $\lvert B\rvert \geq l(n)$ - that is, $B$ is such that there are no solutions to\[a_1=a_2+\cdots+a_r\]with $a_i\in B$ all distinct. Estimate $l(n)$. In particular, is it true that $l(n)n^{-1/2}\to \infty$? Is it true that $l(n)&lt; n^{1-c}$ for some $c&gt;0$?

\noindent\rule{\linewidth}{0.4pt}


%% ============================================================
\section{Problem \#791}
\label{prob:791}

\noindent\textbf{Status:} OPEN \quad|\quad \textbf{Prize:} none \quad|\quad \textbf{Updated:} 2025-08-31

\noindent\textbf{Tags:} additive combinatorics

\medskip

\noindent\textbf{Statement:}

Let $g(n)$ be minimal such that there exists $A\subseteq \{0,\ldots,n\}$ of size $g(n)$ with $\{0,\ldots,n\}\subseteq A+A$. Estimate $g(n)$. In particular is it true that $g(n)\sim 2n^{1/2}$?

\noindent\rule{\linewidth}{0.4pt}


%% ============================================================
\section{Problem \#792}
\label{prob:792}

\noindent\textbf{Status:} OPEN \quad|\quad \textbf{Prize:} none \quad|\quad \textbf{Updated:} 2025-08-31

\noindent\textbf{Tags:} additive combinatorics

\medskip

\noindent\textbf{Statement:}

Let $f(n)$ be maximal such that in any $A\subset \mathbb{Z}$ with $\lvert A\rvert=n$ there exists some sum-free subset $B\subseteq A$ with $\lvert B\rvert \geq f(n)$, so that there are no solutions to\[a+b=c\]with $a,b,c\in B$. Estimate $f(n)$.

\noindent\rule{\linewidth}{0.4pt}


%% ============================================================
\section{Problem \#793}
\label{prob:793}

\noindent\textbf{Status:} OPEN \quad|\quad \textbf{Prize:} none \quad|\quad \textbf{Updated:} 2025-08-31

\noindent\textbf{Tags:} number theory

\medskip

\noindent\textbf{Statement:}

Let $F(n)$ be the maximum possible size of a subset $A\subseteq\{1,\ldots,n\}$ such that $a\nmid bc$ whenever $a,b,c\in A$ with $a\neq b$ and $a\neq c$. Is there a constant $C$ such that\[F(n)=\pi(n)+(C+o(1))n^{2/3}(\log n)^{-2}?\]

\noindent\rule{\linewidth}{0.4pt}


%% ============================================================
\section{Problem \#796}
\label{prob:796}

\noindent\textbf{Status:} OPEN \quad|\quad \textbf{Prize:} none \quad|\quad \textbf{Updated:} 2025-08-31

\noindent\textbf{Tags:} number theory

\medskip

\noindent\textbf{Statement:}

Let $k\geq 2$ and let $g_k(n)$ be the largest possible size of $A\subseteq \{1,\ldots,n\}$ such that every $m$ has $&#60;k$ solutions to $m=a_1a_2$ with $a_1&#60;a_2\in A$. Is it true that\[g_3(n)=\frac{\log\log n}{\log n}n+(c+o(1))\frac{n}{\log n}\]for some constant $c$?

\noindent\rule{\linewidth}{0.4pt}


%% ============================================================
\section{Problem \#802}
\label{prob:802}

\noindent\textbf{Status:} OPEN \quad|\quad \textbf{Prize:} none \quad|\quad \textbf{Updated:} 2025-08-31

\noindent\textbf{Tags:} graph theory

\medskip

\noindent\textbf{Statement:}

Is it true that any $K_r$-free graph on $n$ vertices with average degree $t$ contains an independent set on\[\gg_r \frac{\log t}{t}n\]many vertices?

\noindent\rule{\linewidth}{0.4pt}


%% ============================================================
\section{Problem \#805}
\label{prob:805}

\noindent\textbf{Status:} OPEN \quad|\quad \textbf{Prize:} none \quad|\quad \textbf{Updated:} 2025-08-31

\noindent\textbf{Tags:} graph theory

\medskip

\noindent\textbf{Statement:}

For which functions $g(n)$ with $n&#62;g(n)\geq (\log n)^2$ is there a graph on $n$ vertices in which every induced subgraph on $g(n)$ vertices contains a clique of size $\geq \log n$ and an independent set of size $\geq \log n$? In particular, is there such a graph for $g(n)=(\log n)^3$?

\noindent\rule{\linewidth}{0.4pt}


%% ============================================================
\section{Problem \#809}
\label{prob:809}

\noindent\textbf{Status:} OPEN \quad|\quad \textbf{Prize:} none \quad|\quad \textbf{Updated:} 2025-08-31

\noindent\textbf{Tags:} graph theory, ramsey theory

\medskip

\noindent\textbf{Statement:}

Define the anti-Ramsey number $\chi_S(n,e,G)$ as the smallest $r$ such that there is a graph with $n$ vertices and $e$ edges with an $r$-colouring of its edges in which every copy of $G$ has entirely distinct edge colours. Is it true that, for all $k\geq 3$,\[\chi_S(n, \lfloor n^2/4\rfloor+1,C_{2k+1})\sim n^2/8?\]

\noindent\rule{\linewidth}{0.4pt}


%% ============================================================
\section{Problem \#810}
\label{prob:810}

\noindent\textbf{Status:} OPEN \quad|\quad \textbf{Prize:} none \quad|\quad \textbf{Updated:} 2025-08-31

\noindent\textbf{Tags:} graph theory, ramsey theory

\medskip

\noindent\textbf{Statement:}

Does there exist some $\epsilon&#62;0$ such that, for all sufficiently large $n$, there exists a graph $G$ on $n$ vertices with at least $\epsilon n^2$ many edges such that the edges can be coloured with $n$ colours so that every $C_4$ receives $4$ distinct colours?

\noindent\rule{\linewidth}{0.4pt}


%% ============================================================
\section{Problem \#811}
\label{prob:811}

\noindent\textbf{Status:} OPEN \quad|\quad \textbf{Prize:} none \quad|\quad \textbf{Updated:} 2025-08-31

\noindent\textbf{Tags:} graph theory, ramsey theory

\medskip

\noindent\textbf{Statement:}

Suppose $n\equiv 1\pmod{m}$. We say that an edge-colouring of $K_n$ using $m$ colours is balanced if every vertex sees exactly $\lfloor n/m\rfloor$ many edges of each colours. For which graphs $G$ is it true that, if $m=e(G)$, for all large $n\equiv 1\pmod{m}$, every balanced edge-colouring of $K_n$ with $m$ colours contains a rainbow copy of $G$? (That is, a subgraph isomorphic to $G$ where each edge receives a different colour.)

\noindent\rule{\linewidth}{0.4pt}


%% ============================================================
\section{Problem \#812}
\label{prob:812}

\noindent\textbf{Status:} OPEN \quad|\quad \textbf{Prize:} none \quad|\quad \textbf{Updated:} 2025-08-31

\noindent\textbf{Tags:} graph theory, ramsey theory

\medskip

\noindent\textbf{Statement:}

Is it true that\[\frac{R(n+1)}{R(n)}\geq 1+c\]for some constant $c&#62;0$, for all large $n$? Is it true that\[R(n+1)-R(n) \gg n^2?\]

\noindent\rule{\linewidth}{0.4pt}


%% ============================================================
\section{Problem \#813}
\label{prob:813}

\noindent\textbf{Status:} OPEN \quad|\quad \textbf{Prize:} none \quad|\quad \textbf{Updated:} 2025-08-31

\noindent\textbf{Tags:} graph theory

\medskip

\noindent\textbf{Statement:}

Let $h(n)$ be minimal such that every graph on $n$ vertices where every set of $7$ vertices contains a triangle (a copy of $K_3$) must contain a clique on at least $h(n)$ vertices. Estimate $h(n)$ - in particular, do there exist constants $c_1,c_2&#62;0$ such that\[n^{1/3+c_1}\ll h(n) \ll n^{1/2-c_2}?\]

\noindent\rule{\linewidth}{0.4pt}


%% ============================================================
\section{Problem \#817}
\label{prob:817}

\noindent\textbf{Status:} OPEN \quad|\quad \textbf{Prize:} none \quad|\quad \textbf{Updated:} 2025-08-31

\noindent\textbf{Tags:} additive combinatorics

\medskip

\noindent\textbf{Statement:}

Let $k\geq 3$ and define $g_k(n)$ to be the minimal $N$ such that $\{1,\ldots,N\}$ contains some $A$ of size $\lvert A\rvert=n$ such that\[\langle A\rangle = \left\{\sum_{a\in A}\epsilon_aa: \epsilon_a\in \{0,1\}\right\}\]contains no non-trivial $k$-term arithmetic progression. Estimate $g_k(n)$. In particular, is it true that\[g_3(n) \gg 3^n?\]

\noindent\rule{\linewidth}{0.4pt}


%% ============================================================
\section{Problem \#819}
\label{prob:819}

\noindent\textbf{Status:} OPEN \quad|\quad \textbf{Prize:} none \quad|\quad \textbf{Updated:} 2025-08-31

\noindent\textbf{Tags:} additive combinatorics

\medskip

\noindent\textbf{Statement:}

Let $f(N)$ be maximal such that there exists $A\subseteq \{1,\ldots,N\}$ with $\lvert A\rvert=\lfloor N^{1/2}\rfloor$ such that $\lvert (A+A)\cap [1,N]\rvert=f(N)$. Estimate $f(N)$.

\noindent\rule{\linewidth}{0.4pt}


%% ============================================================
\section{Problem \#820}
\label{prob:820}

\noindent\textbf{Status:} OPEN \quad|\quad \textbf{Prize:} none \quad|\quad \textbf{Updated:} 2025-08-31

\noindent\textbf{Tags:} number theory

\medskip

\noindent\textbf{Statement:}

Let $H(n)$ be the smallest integer $l$ such that there exist $k&lt;l$ with $(k^n-1,l^n-1)=1$. Is it true that $H(n)=3$ infinitely often? (That is, $(2^n-1,3^n-1)=1$ infinitely often?) Estimate $H(n)$. Is it true that there exists some constant $c&gt;0$ such that, for all $\epsilon&#62;0$,\[H(n) &#62; \exp(n^{(c-\epsilon)/\log\log n})\]for infinitely many $n$ and\[H(n) &#60; \exp(n^{(c+\epsilon)/\log\log n})\]for all large enough $n$? Does a similar upper bound hold for the smallest $k$ such that $(k^n-1,2^n-1)=1$?

\noindent\rule{\linewidth}{0.4pt}


%% ============================================================
\section{Problem \#821}
\label{prob:821}

\noindent\textbf{Status:} OPEN \quad|\quad \textbf{Prize:} none \quad|\quad \textbf{Updated:} 2025-08-31

\noindent\textbf{Tags:} number theory

\medskip

\noindent\textbf{Statement:}

Let $g(n)$ count the number of $m$ such that $\phi(m)=n$. Is it true that, for every $\epsilon&#62;0$, there exist infinitely many $n$ such that\[g(n) &#62; n^{1-\epsilon}?\]

\noindent\rule{\linewidth}{0.4pt}


%% ============================================================
\section{Problem \#824}
\label{prob:824}

\noindent\textbf{Status:} OPEN \quad|\quad \textbf{Prize:} none \quad|\quad \textbf{Updated:} 2025-08-31

\noindent\textbf{Tags:} number theory

\medskip

\noindent\textbf{Statement:}

Let $h(x)$ count the number of integers $1\leq a&lt;b&lt;x$ such that $(a,b)=1$ and $\sigma(a)=\sigma(b)$, where $\sigma$ is the sum of divisors function. Is it true that $h(x)&gt;x^{2-o(1)}$?

\noindent\rule{\linewidth}{0.4pt}


%% ============================================================
\section{Problem \#826}
\label{prob:826}

\noindent\textbf{Status:} OPEN \quad|\quad \textbf{Prize:} none \quad|\quad \textbf{Updated:} 2025-08-31

\noindent\textbf{Tags:} number theory

\medskip

\noindent\textbf{Statement:}

Are there infinitely many $n$ such that, for all $k\geq 1$,\[\tau(n+k)\ll k?\]

\noindent\rule{\linewidth}{0.4pt}


%% ============================================================
\section{Problem \#827}
\label{prob:827}

\noindent\textbf{Status:} OPEN \quad|\quad \textbf{Prize:} none \quad|\quad \textbf{Updated:} 2025-08-31

\noindent\textbf{Tags:} geometry

\medskip

\noindent\textbf{Statement:}

Let $n_k$ be minimal such that if $n_k$ points in $\mathbb{R}^2$ are in general position then there exists a subset of $k$ points such that all $\binom{k}{3}$ triples determine circles of different radii. Determine $n_k$.

\noindent\rule{\linewidth}{0.4pt}


%% ============================================================
\section{Problem \#828}
\label{prob:828}

\noindent\textbf{Status:} OPEN \quad|\quad \textbf{Prize:} none \quad|\quad \textbf{Updated:} 2025-08-31

\noindent\textbf{Tags:} number theory

\medskip

\noindent\textbf{Statement:}

Is it true that, for any $a\in\mathbb{Z}$, there are infinitely many $n$ such that\[\phi(n) \mid n+a?\]

\noindent\rule{\linewidth}{0.4pt}


%% ============================================================
\section{Problem \#829}
\label{prob:829}

\noindent\textbf{Status:} OPEN \quad|\quad \textbf{Prize:} none \quad|\quad \textbf{Updated:} 2025-08-31

\noindent\textbf{Tags:} number theory

\medskip

\noindent\textbf{Statement:}

Let $A\subset\mathbb{N}$ be the set of cubes. Is it true that\[1_A\ast 1_A(n) \ll (\log n)^{O(1)}?\]

\noindent\rule{\linewidth}{0.4pt}


%% ============================================================
\section{Problem \#830}
\label{prob:830}

\noindent\textbf{Status:} OPEN \quad|\quad \textbf{Prize:} none \quad|\quad \textbf{Updated:} 2025-08-31

\noindent\textbf{Tags:} number theory

\medskip

\noindent\textbf{Statement:}

We say that $a,b\in \mathbb{N}$ are an amicable pair if $\sigma(a)=\sigma(b)=a+b$. Are there infinitely many amicable pairs? If $A(x)$ counts the number of amicable $1\leq a\leq b\leq x$ then is it true that\[A(x)&gt;x^{1-o(1)}?\]

\noindent\rule{\linewidth}{0.4pt}


%% ============================================================
\section{Problem \#831}
\label{prob:831}

\noindent\textbf{Status:} OPEN \quad|\quad \textbf{Prize:} none \quad|\quad \textbf{Updated:} 2025-08-31

\noindent\textbf{Tags:} geometry

\medskip

\noindent\textbf{Statement:}

Let $h(n)$ be maximal such that in any $n$ points in $\mathbb{R}^2$ (with no three on a line and no four on a circle) there are at least $h(n)$ many circles of different radii passing through three points. Estimate $h(n)$.

\noindent\rule{\linewidth}{0.4pt}


%% ============================================================
\section{Problem \#836}
\label{prob:836}

\noindent\textbf{Status:} OPEN \quad|\quad \textbf{Prize:} none \quad|\quad \textbf{Updated:} 2025-08-31

\noindent\textbf{Tags:} graph theory, hypergraphs, chromatic number

\medskip

\noindent\textbf{Statement:}

Let $r\geq 2$ and $G$ be a $r$-uniform hypergraph with chromatic number $3$ (that is, there is a $3$-colouring of the vertices of $G$ such that no edge is monochromatic). Suppose any two edges of $G$ have a non-empty intersection. Must $G$ contain $O(r^2)$ many vertices? Must there be two edges which meet in $\gg r$ many vertices?

\noindent\rule{\linewidth}{0.4pt}


%% ============================================================
\section{Problem \#837}
\label{prob:837}

\noindent\textbf{Status:} OPEN \quad|\quad \textbf{Prize:} none \quad|\quad \textbf{Updated:} 2025-08-31

\noindent\textbf{Tags:} graph theory, hypergraphs

\medskip

\noindent\textbf{Statement:}

Let $k\geq 2$ and $A_k\subseteq [0,1]$ be the set of $\alpha$ such that there exists some $\beta(\alpha)&#62;\alpha$ with the property that, if $G_1,G_2,\ldots$ is a sequence of $k$-uniform hypergraphs with\[\liminf \frac{e(G_n)}{\binom{\lvert G_n\rvert}{k}} &#62;\alpha\]then there exist subgraphs $H_n\subseteq G_n$ such that $\lvert H_n\rvert \to \infty$ and\[\liminf \frac{e(H_n)}{\binom{\lvert H_n\rvert}{k}} &#62;\beta,\]and further that this property does not necessarily hold if $&#62;\alpha$ is replaced by $\geq \alpha$. What is $A_3$?

\noindent\rule{\linewidth}{0.4pt}


%% ============================================================
\section{Problem \#838}
\label{prob:838}

\noindent\textbf{Status:} OPEN \quad|\quad \textbf{Prize:} none \quad|\quad \textbf{Updated:} 2025-08-31

\noindent\textbf{Tags:} geometry, convex

\medskip

\noindent\textbf{Statement:}

Let $f(n)$ be maximal such that any $n$ points in $\mathbb{R}^2$, with no three on a line, determine at least $f(n)$ different convex subsets. Estimate $f(n)$ - in particular, does there exist a constant $c$ such that\[\lim \frac{\log f(n)}{(\log n)^2}=c?\]

\noindent\rule{\linewidth}{0.4pt}


%% ============================================================
\section{Problem \#839}
\label{prob:839}

\noindent\textbf{Status:} OPEN \quad|\quad \textbf{Prize:} none \quad|\quad \textbf{Updated:} 2025-08-31

\noindent\textbf{Tags:} number theory

\medskip

\noindent\textbf{Statement:}

Let $1\leq a_1&#60;a_2&#60;\cdots$ be a sequence of integers such that no $a_i$ is the sum of consecutive $a_j$ for $j&#60;i$. Is it true that\[\limsup \frac{a_n}{n}=\infty?\]Or even\[\lim \frac{1}{\log x}\sum_{a_n&#60;x}\frac{1}{a_n}=0?\]

\noindent\rule{\linewidth}{0.4pt}


%% ============================================================
\section{Problem \#840}
\label{prob:840}

\noindent\textbf{Status:} OPEN \quad|\quad \textbf{Prize:} none \quad|\quad \textbf{Updated:} 2025-08-31

\noindent\textbf{Tags:} additive combinatorics, sidon sets

\medskip

\noindent\textbf{Statement:}

Let $f(N)$ be the size of the largest quasi-Sidon subset $A\subset\{1,\ldots,N\}$, where we say that $A$ is quasi-Sidon if\[\lvert A+A\rvert=(1+o(1))\binom{\lvert A\rvert}{2}.\]How does $f(N)$ grow?

\noindent\rule{\linewidth}{0.4pt}


%% ============================================================
\section{Problem \#849}
\label{prob:849}

\noindent\textbf{Status:} OPEN \quad|\quad \textbf{Prize:} none \quad|\quad \textbf{Updated:} 2025-08-31

\noindent\textbf{Tags:} number theory, binomial coefficients

\medskip
\noindent\textit{Singmaster's conjecture}


\noindent\textbf{Statement:}

Is it true that, for every integer $t\geq 1$, there is some integer $a$ such that\[\binom{n}{k}=a\](with $1\leq k\leq n/2$) has exactly $t$ solutions?

\noindent\rule{\linewidth}{0.4pt}


%% ============================================================
\section{Problem \#850}
\label{prob:850}

\noindent\textbf{Status:} OPEN \quad|\quad \textbf{Prize:} none \quad|\quad \textbf{Updated:} 2025-08-31

\noindent\textbf{Tags:} number theory, primes

\medskip
\noindent\textit{Erdős-Woods conjecture}


\noindent\textbf{Statement:}

Can there exist two distinct integers $x$ and $y$ such that $x,y$ have the same prime factors, $x+1,y+1$ have the same prime factors, and $x+2,y+2$ also have the same prime factors?

\noindent\rule{\linewidth}{0.4pt}


%% ============================================================
\section{Problem \#852}
\label{prob:852}

\noindent\textbf{Status:} OPEN \quad|\quad \textbf{Prize:} none \quad|\quad \textbf{Updated:} 2025-08-31

\noindent\textbf{Tags:} number theory, primes

\medskip

\noindent\textbf{Statement:}

Let $d_n=p_{n+1}-p_n$, where $p_n$ is the $n$th prime. Let $h(x)$ be maximal such that for some $n&lt;x$ the numbers $d_n,d_{n+1},\ldots,d_{n+h(x)-1}$ are all distinct. Estimate $h(x)$. In particular, is it true that\[h(x) &gt;(\log x)^c\]for some constant $c&#62;0$, and\[h(x)=o(\log x)?\]

\noindent\rule{\linewidth}{0.4pt}


%% ============================================================
\section{Problem \#853}
\label{prob:853}

\noindent\textbf{Status:} OPEN \quad|\quad \textbf{Prize:} none \quad|\quad \textbf{Updated:} 2025-08-31

\noindent\textbf{Tags:} number theory, primes

\medskip

\noindent\textbf{Statement:}

Let $d_n=p_{n+1}-p_n$, where $p_n$ is the $n$th prime. Let $r(x)$ be the smallest even integer $t$ such that $d_n=t$ has no solutions for $n\leq x$. Is it true that $r(x)\to \infty$? Or even $r(x)/\log x \to \infty$?

\noindent\rule{\linewidth}{0.4pt}


%% ============================================================
\section{Problem \#854}
\label{prob:854}

\noindent\textbf{Status:} OPEN \quad|\quad \textbf{Prize:} none \quad|\quad \textbf{Updated:} 2025-08-31

\noindent\textbf{Tags:} number theory

\medskip

\noindent\textbf{Statement:}

Let $n_k$ denote the $k$th primorial, i.e. the product of the first $k$ primes. If $1=a_1&#60;a_2&#60;\cdots a_{\phi(n_k)}=n_k-1$ is the sequence of integers coprime to $n_k$, then estimate the smallest even integer not of the form $a_{i+1}-a_i$. Are there\[\gg \max_i (a_{i+1}-a_i)\]many even integers of the form $a_{j+1}-a_j$?

\noindent\rule{\linewidth}{0.4pt}


%% ============================================================
\section{Problem \#855}
\label{prob:855}

\noindent\textbf{Status:} OPEN \quad|\quad \textbf{Prize:} none \quad|\quad \textbf{Updated:} 2026-03-14

\noindent\textbf{Tags:} number theory, primes

\medskip
\noindent\textit{second Hardy-Littlewood conjecture}


\noindent\textbf{Statement:}

If $\pi(x)$ counts the number of primes in $[1,x]$ then is it true that (for large $x$ and $y$)\[\pi(x+y) \leq \pi(x)+\pi(y)?\]

\noindent\rule{\linewidth}{0.4pt}


%% ============================================================
\section{Problem \#856}
\label{prob:856}

\noindent\textbf{Status:} OPEN \quad|\quad \textbf{Prize:} none \quad|\quad \textbf{Updated:} 2025-08-31

\noindent\textbf{Tags:} number theory

\medskip

\noindent\textbf{Statement:}

Let $k\geq 3$ and $f_k(N)$ be the maximum value of $\sum_{n\in A}\frac{1}{n}$, where $A$ ranges over all subsets of $\{1,\ldots,N\}$ which contain no subset of size $k$ with the same pairwise least common multiple. Estimate $f_k(N)$.

\noindent\rule{\linewidth}{0.4pt}


%% ============================================================
\section{Problem \#857}
\label{prob:857}

\noindent\textbf{Status:} OPEN \quad|\quad \textbf{Prize:} none \quad|\quad \textbf{Updated:} 2025-08-31

\noindent\textbf{Tags:} combinatorics

\medskip
\noindent\textit{weak sunflower problem}


\noindent\textbf{Statement:}

Let $m=m(n,k)$ be minimal such that in any collection of sets $A_1,\ldots,A_m\subseteq \{1,\ldots,n\}$ there must exist a sunflower of size $k$ - that is, some collection of $k$ of the $A_i$ which pairwise have the same intersection. Estimate $m(n,k)$, or even better, give an asymptotic formula.

\noindent\rule{\linewidth}{0.4pt}


%% ============================================================
\section{Problem \#859}
\label{prob:859}

\noindent\textbf{Status:} OPEN \quad|\quad \textbf{Prize:} none \quad|\quad \textbf{Updated:} 2025-08-31

\noindent\textbf{Tags:} number theory, divisors

\medskip

\noindent\textbf{Statement:}

Let $t\geq 1$ and let $d_t$ be the density of the set of integers $n\in\mathbb{N}$ for which $t$ can be represented as the sum of distinct divisors of $n$. Do there exist constants $c_1,c_2&#62;0$ such that\[d_t \sim \frac{c_1}{(\log t)^{c_2}}\]as $t\to \infty$?

\noindent\rule{\linewidth}{0.4pt}


%% ============================================================
\section{Problem \#860}
\label{prob:860}

\noindent\textbf{Status:} OPEN \quad|\quad \textbf{Prize:} none \quad|\quad \textbf{Updated:} 2025-08-31

\noindent\textbf{Tags:} number theory, primes

\medskip

\noindent\textbf{Statement:}

Let $h(n)$ be such that, for any $m\geq 1$, in the interval $(m,m+h(n))$ there exist distinct integers $a_i$ for $1\leq i\leq \pi(n)$ such that $p_i\mid a_i$, where $p_i$ denotes the $i$th prime. Estimate $h(n)$.

\noindent\rule{\linewidth}{0.4pt}


%% ============================================================
\section{Problem \#864}
\label{prob:864}

\noindent\textbf{Status:} OPEN \quad|\quad \textbf{Prize:} none \quad|\quad \textbf{Updated:} 2025-08-31

\noindent\textbf{Tags:} number theory, sidon sets, additive combinatorics

\medskip

\noindent\textbf{Statement:}

Let $A\subseteq \{1,\ldots N\}$ be a set such that there exists at most one $n$ with more than one solution to $n=a+b$ (with $a\leq b\in A$). Estimate the maximal possible size of $\lvert A\rvert$ - in particular, is it true that\[\lvert A\rvert \leq (1+o(1))\frac{2}{\sqrt{3}}N^{1/2}?\]

\noindent\rule{\linewidth}{0.4pt}


%% ============================================================
\section{Problem \#865}
\label{prob:865}

\noindent\textbf{Status:} OPEN \quad|\quad \textbf{Prize:} none \quad|\quad \textbf{Updated:} 2025-08-31

\noindent\textbf{Tags:} number theory, additive combinatorics

\medskip

\noindent\textbf{Statement:}

There exists a constant $C&#62;0$ such that, for all large $N$, if $A\subseteq \{1,\ldots,N\}$ has size at least $\frac{5}{8}N+C$ then there are distinct $a,b,c\in A$ such that $a+b,a+c,b+c\in A$.

\noindent\rule{\linewidth}{0.4pt}


%% ============================================================
\section{Problem \#866}
\label{prob:866}

\noindent\textbf{Status:} OPEN \quad|\quad \textbf{Prize:} none \quad|\quad \textbf{Updated:} 2025-08-31

\noindent\textbf{Tags:} number theory, additive combinatorics

\medskip

\noindent\textbf{Statement:}

Let $k\geq 3$ and $g_k(N)$ be minimal such that if $A\subseteq \{1,\ldots,2N\}$ has $\lvert A\rvert \geq N+g_k(N)$ then there exist integers $b_1,\ldots,b_k$ such that all $\binom{k}{2}$ pairwise sums are in $A$ (but the $b_i$ themselves need not be in $A$). Estimate $g_k(N)$.

\noindent\rule{\linewidth}{0.4pt}


%% ============================================================
\section{Problem \#870}
\label{prob:870}

\noindent\textbf{Status:} OPEN \quad|\quad \textbf{Prize:} none \quad|\quad \textbf{Updated:} 2025-08-31

\noindent\textbf{Tags:} number theory, additive basis

\medskip

\noindent\textbf{Statement:}

Let $k\geq 3$ and $A$ be an additive basis of order $k$. Does there exist a constant $c=c(k)&#62;0$ such that if $r(n)\geq c\log n$ for all large $n$ then $A$ must contain a minimal basis of order $k$? (Here $r(n)$ counts the number of representations of $n$ as the sum of at most $k$ elements from $A$.)

\noindent\rule{\linewidth}{0.4pt}


%% ============================================================
\section{Problem \#872}
\label{prob:872}

\noindent\textbf{Status:} OPEN \quad|\quad \textbf{Prize:} none \quad|\quad \textbf{Updated:} 2025-08-31

\noindent\textbf{Tags:} number theory, primitive sets

\medskip

\noindent\textbf{Statement:}

Consider the two-player game in which players alternately choose integers from $\{2,3,\ldots,n\}$ to be included in some set $A$ (the same set for both players) such that no $a\mid b$ for $a\neq b\in A$. The game ends when no legal move is possible. One player wants the game to last as long as possible, the other wants the game to end quickly. How long can the game be guaranteed to last for? At least $\epsilon n$ moves? (For $\epsilon&#62;0$ and $n$ sufficiently large.) At least $(1-\epsilon)\frac{n}{2}$ moves?

\noindent\rule{\linewidth}{0.4pt}


%% ============================================================
\section{Problem \#873}
\label{prob:873}

\noindent\textbf{Status:} OPEN \quad|\quad \textbf{Prize:} none \quad|\quad \textbf{Updated:} 2025-08-31

\noindent\textbf{Tags:} number theory

\medskip

\noindent\textbf{Statement:}

Let $A=\{a_1&lt;a_2&lt;\cdots\}\subseteq \mathbb{N}$ and let $F(A,X,k)$ count the number of $i$ such that\[[a_i,a_{i+1},\ldots,a_{i+k-1}] &lt; X,\]where the left-hand side is the least common multiple. Is it true that, for every $\epsilon &gt;0$, there exists some $k$ such that\[F(A,X,k)&#60;X^\epsilon?\]

\noindent\rule{\linewidth}{0.4pt}


%% ============================================================
\section{Problem \#875}
\label{prob:875}

\noindent\textbf{Status:} OPEN \quad|\quad \textbf{Prize:} none \quad|\quad \textbf{Updated:} 2025-08-31

\noindent\textbf{Tags:} additive combinatorics

\medskip

\noindent\textbf{Statement:}

Let $A=\{a_1&#60;a_2&#60;\cdots\}\subset \mathbb{N}$ be an infinite set such that the sets\[S_r = \{ a_1+\cdots +a_r : a_1&#60;\cdots&#60;a_r\in A\}\]are disjoint for distinct $r\geq 1$. How fast can such a sequence grow? How small can $a_{n+1}-a_n$ be? In particular, for which $c$ is it possible that $a_{n+1}-a_n\leq n^{c}$?

\noindent\rule{\linewidth}{0.4pt}


%% ============================================================
\section{Problem \#876}
\label{prob:876}

\noindent\textbf{Status:} OPEN \quad|\quad \textbf{Prize:} none \quad|\quad \textbf{Updated:} 2025-08-31

\noindent\textbf{Tags:} additive combinatorics

\medskip

\noindent\textbf{Statement:}

Let $A=\{a_1&#60;a_2&#60;\cdots\}\subset \mathbb{N}$ be an infinite sum-free set - that is, there are no solutions to\[a=b_1+\cdots+b_r\]with $b_1&#60;\cdots&#60;b_r&#60;a\in A$. How small can $a_{n+1}-a_n$ be? Is it possible that $a_{n+1}-a_n&#60;n$?

\noindent\rule{\linewidth}{0.4pt}


%% ============================================================
\section{Problem \#878}
\label{prob:878}

\noindent\textbf{Status:} OPEN \quad|\quad \textbf{Prize:} none \quad|\quad \textbf{Updated:} 2025-08-31

\noindent\textbf{Tags:} number theory

\medskip

\noindent\textbf{Statement:}

If $n=\prod_{1\leq i\leq t} p_i^{k_i}$ is the factorisation of $n$ into distinct primes then let\[f(n)=\sum p_i^{\ell_i},\]where $\ell_i$ is chosen such that $n\in [p_i^{\ell_i},p_i^{\ell_i+1})$. Furthermore, let\[F(n)=\max \sum_{i} a_i\]where the maximum is taken over all distinct $a_1,\ldots,a_k\leq n$ such that $(a_i,a_j)=1$ for $i\neq j$ and all prime factors of each $a_i$ are prime factors of $n$. Is it true that, for almost all $n$,\[f(n)=o(n\log\log n)\]and\[F(n) \gg n\log\log n?\]Is it true that\[\max_{n\leq x}f(n)\sim \frac{x\log x}{\log\log x}?\]Is it true that (for all $x$, or perhaps just for all large $x$)\[\max_{n\leq x}f(n)=\max_{n\leq x}F(n)?\]Find an asymptotic formula for the number of $n&#60;x$ such that $f(n)=F(n)$. Find an asymptotic formula for\[H(x)=\sum_{n&#60;x}\frac{f(n)}{n}.\]Is it true that\[H(x) \ll x\log\log\log\log x?\]

\noindent\rule{\linewidth}{0.4pt}


%% ============================================================
\section{Problem \#879}
\label{prob:879}

\noindent\textbf{Status:} OPEN \quad|\quad \textbf{Prize:} none \quad|\quad \textbf{Updated:} 2025-08-31

\noindent\textbf{Tags:} number theory

\medskip

\noindent\textbf{Statement:}

Call a set $S\subseteq \{1,\ldots,n\}$ admissible if $(a,b)=1$ for all $a\neq b\in S$. Let\[G(n) = \max_{S\subseteq \{1,\ldots,n\}} \sum_{a\in S}a\]and\[H(n)=\sum_{p&lt;n}p+ n\pi(n^{1/2}).\]Is it true that\[G(n) &gt;H(n)-n^{1+o(1)}?\]Is it true that, for every $k\geq 2$, if $n$ is sufficiently large then the admissible set which maximises $G(n)$ contains at least one integer with at least $k$ prime factors?

\noindent\rule{\linewidth}{0.4pt}


%% ============================================================
\section{Problem \#881}
\label{prob:881}

\noindent\textbf{Status:} OPEN \quad|\quad \textbf{Prize:} none \quad|\quad \textbf{Updated:} 2025-08-31

\noindent\textbf{Tags:} number theory, additive basis

\medskip

\noindent\textbf{Statement:}

Let $A\subset\mathbb{N}$ be an additive basis of order $k$ which is minimal, in the sense that if $B\subset A$ is any infinite set then $A\backslash B$ is not a basis of order $k$. Must there exist an infinite $B\subset A$ such that $A\backslash B$ is a basis of order $k+1$?

\noindent\rule{\linewidth}{0.4pt}


%% ============================================================
\section{Problem \#883}
\label{prob:883}

\noindent\textbf{Status:} OPEN \quad|\quad \textbf{Prize:} none \quad|\quad \textbf{Updated:} 2025-08-31

\noindent\textbf{Tags:} number theory, graph theory

\medskip

\noindent\textbf{Statement:}

For $A\subseteq \{1,\ldots,n\}$ let $G(A)$ be the graph with vertex set $A$, where two integers are joined by an edge if they are coprime. Is it true that if\[\lvert A\rvert &#62;\lfloor\tfrac{n}{2}\rfloor+\lfloor\tfrac{n}{3}\rfloor-\lfloor\tfrac{n}{6}\rfloor\]then $G(A)$ contains all odd cycles of length $\leq \frac{n}{3}+1$? Is it true that, for every $\ell\geq 1$, if $n$ is sufficiently large and\[\lvert A\rvert &#62;\lfloor\tfrac{n}{2}\rfloor+\lfloor\tfrac{n}{3}\rfloor-\lfloor\tfrac{n}{6}\rfloor\]then $G(A)$ must contain a complete $(1,\ell,\ell)$ triparite graph on $2\ell+1$ vertices?

\noindent\rule{\linewidth}{0.4pt}


%% ============================================================
\section{Problem \#885}
\label{prob:885}

\noindent\textbf{Status:} OPEN \quad|\quad \textbf{Prize:} none \quad|\quad \textbf{Updated:} 2025-08-31

\noindent\textbf{Tags:} number theory, divisors

\medskip

\noindent\textbf{Statement:}

For integer $n\geq 1$ we define the factor difference set of $n$ by\[D(n) = \{\lvert a-b\rvert : n=ab\}.\]Is it true that, for every $k\geq 1$, there exist integers $N_1&#60;\cdots&#60;N_k$ such that\[\lvert \cap_i D(N_i)\rvert \geq k?\]

\noindent\rule{\linewidth}{0.4pt}


%% ============================================================
\section{Problem \#886}
\label{prob:886}

\noindent\textbf{Status:} OPEN \quad|\quad \textbf{Prize:} none \quad|\quad \textbf{Updated:} 2025-08-31

\noindent\textbf{Tags:} number theory, divisors

\medskip

\noindent\textbf{Statement:}

Let $\epsilon&#62;0$. Is it true that, for all large $n$, the number of divisors of $n$ in $(n^{1/2},n^{1/2}+n^{1/2-\epsilon})$ is $O_\epsilon(1)$?

\noindent\rule{\linewidth}{0.4pt}


%% ============================================================
\section{Problem \#887}
\label{prob:887}

\noindent\textbf{Status:} OPEN \quad|\quad \textbf{Prize:} none \quad|\quad \textbf{Updated:} 2025-08-31

\noindent\textbf{Tags:} number theory, divisors

\medskip

\noindent\textbf{Statement:}

Is there an absolute constant $K$ such that, for every $C&#62;0$, if $n$ is sufficiently large then $n$ has at most $K$ divisors in $(n^{1/2},n^{1/2}+C n^{1/4})$.

\noindent\rule{\linewidth}{0.4pt}


%% ============================================================
\section{Problem \#889}
\label{prob:889}

\noindent\textbf{Status:} OPEN \quad|\quad \textbf{Prize:} none \quad|\quad \textbf{Updated:} 2025-08-31

\noindent\textbf{Tags:} number theory

\medskip

\noindent\textbf{Statement:}

For $k\geq 0$ and $n\geq 1$ let $v(n,k)$ count the prime factors of $n+k$ which do not divide $n+i$ for $0\leq i&lt;k$. Equivalently, $v(n,k)$ counts the number of prime factors of $n+k$ which are $&gt;k$. Is it true that\[v_0(n)=\max_{k\geq 0}v(n,k)\to \infty\]as $n\to \infty$?

\noindent\rule{\linewidth}{0.4pt}


%% ============================================================
\section{Problem \#890}
\label{prob:890}

\noindent\textbf{Status:} OPEN \quad|\quad \textbf{Prize:} none \quad|\quad \textbf{Updated:} 2025-08-31

\noindent\textbf{Tags:} number theory, primes

\medskip

\noindent\textbf{Statement:}

If $\omega_k(n)$ counts the number of distinct prime factors of $n$ which are $&#62;k$, then is it true that, for every $k\geq 1$,\[\liminf_{n\to \infty}\sum_{0\leq i&#60;k}\omega_k(n+i)\leq k?\]Is it true that\[\limsup_{n\to \infty}\left(\sum_{0\leq i&#60;k}\omega(n+i)\right) \frac{\log\log n}{\log n}=1,\]where $\omega$ counts the number of distinct prime factors without restriction?

\noindent\rule{\linewidth}{0.4pt}


%% ============================================================
\section{Problem \#891}
\label{prob:891}

\noindent\textbf{Status:} OPEN \quad|\quad \textbf{Prize:} none \quad|\quad \textbf{Updated:} 2025-08-31

\noindent\textbf{Tags:} number theory

\medskip

\noindent\textbf{Statement:}

Let $2=p_1&lt;p_2&lt;\cdots$ be the primes and $k\geq 2$. Is it true that, for all sufficiently large $n$, there must exist an integer in $[n,n+p_1\cdots p_k)$ with $&gt;k$ many prime factors?

\noindent\rule{\linewidth}{0.4pt}


%% ============================================================
\section{Problem \#892}
\label{prob:892}

\noindent\textbf{Status:} OPEN \quad|\quad \textbf{Prize:} none \quad|\quad \textbf{Updated:} 2025-08-31

\noindent\textbf{Tags:} number theory, primitive sets

\medskip

\noindent\textbf{Statement:}

Is there a necessary and sufficient condition for a sequence of integers $b_1&#60;b_2&#60;\cdots$ that ensures there exists a primitive sequence $a_1&#60;a_2&#60;\cdots$ (i.e. no element divides another) with $a_n \ll b_n$ for all $n$? In particular, is this always possible if there are no non-trivial solutions to $(b_i,b_j)=b_k$? Similarly, find necessary and sufficient conditions on a sequence $n_1&#60;n_2&#60;\cdots$ that ensure there exists a primitive set $A$ such that\[\lvert A\cap [1,2^{n_i}]\rvert \gg 2^{n_i}\]for every $i$.

\noindent\rule{\linewidth}{0.4pt}


%% ============================================================
\section{Problem \#893}
\label{prob:893}

\noindent\textbf{Status:} OPEN \quad|\quad \textbf{Prize:} none \quad|\quad \textbf{Updated:} 2025-08-31

\noindent\textbf{Tags:} number theory, divisors

\medskip

\noindent\textbf{Statement:}

If $\tau(n)$ counts the divisors of $n$ then let\[f(n)=\sum_{1\leq k\leq n}\tau(2^k-1).\]Does $f(2n)/f(n)$ tend to a limit?

\noindent\rule{\linewidth}{0.4pt}


%% ============================================================
\section{Problem \#901}
\label{prob:901}

\noindent\textbf{Status:} OPEN \quad|\quad \textbf{Prize:} none \quad|\quad \textbf{Updated:} 2025-08-31

\noindent\textbf{Tags:} combinatorics, hypergraphs

\medskip

\noindent\textbf{Statement:}

Let $m(n)$ be minimal such that there is an $n$-uniform hypergraph with $m(n)$ edges which is $3$-chromatic. Estimate $m(n)$.

\noindent\rule{\linewidth}{0.4pt}


%% ============================================================
\section{Problem \#902}
\label{prob:902}

\noindent\textbf{Status:} OPEN \quad|\quad \textbf{Prize:} none \quad|\quad \textbf{Updated:} 2025-08-31

\noindent\textbf{Tags:} graph theory

\medskip

\noindent\textbf{Statement:}

Let $f(n)$ be minimal such that there is a tournament (a complete directed graph) on $f(n)$ vertices such that every set of $n$ vertices is dominated by at least one other vertex. Estimate $f(n)$.

\noindent\rule{\linewidth}{0.4pt}


%% ============================================================
\section{Problem \#906}
\label{prob:906}

\noindent\textbf{Status:} OPEN \quad|\quad \textbf{Prize:} none \quad|\quad \textbf{Updated:} 2025-08-31

\noindent\textbf{Tags:} analysis, iterated functions

\medskip

\noindent\textbf{Statement:}

Is there an entire non-zero function $f:\mathbb{C}\to \mathbb{C}$ such that, for any infinite sequence $n_1&#60;n_2&#60;\cdots$, the set\[\{ z: f^{(n_k)}(z)=0 \textrm{ for some }k\geq 1\}\]is everywhere dense?

\noindent\rule{\linewidth}{0.4pt}


%% ============================================================
\section{Problem \#911}
\label{prob:911}

\noindent\textbf{Status:} OPEN \quad|\quad \textbf{Prize:} none \quad|\quad \textbf{Updated:} 2025-08-31

\noindent\textbf{Tags:} graph theory, ramsey theory

\medskip

\noindent\textbf{Statement:}

Let $\hat{R}(G)$ denote the size Ramsey number, the minimal number of edges $m$ such that there is a graph $H$ with $m$ edges that is Ramsey for $G$. Is there a function $f$ such that $f(x)/x\to \infty$ as $x\to \infty$ such that, for all large $C$, if $G$ is a graph with $n$ vertices and $e\geq Cn$ edges then\[\hat{R}(G) &#62; f(C) e?\]

\noindent\rule{\linewidth}{0.4pt}


%% ============================================================
\section{Problem \#912}
\label{prob:912}

\noindent\textbf{Status:} OPEN \quad|\quad \textbf{Prize:} none \quad|\quad \textbf{Updated:} 2025-08-31

\noindent\textbf{Tags:} number theory, factorials

\medskip

\noindent\textbf{Statement:}

If\[n! = \prod_i p_i^{k_i}\]is the factorisation into distinct primes then let $h(n)$ count the number of distinct exponents $k_i$. Prove that there exists some $c&#62;0$ such that\[h(n) \sim c \left(\frac{n}{\log n}\right)^{1/2}\]as $n\to \infty$.

\noindent\rule{\linewidth}{0.4pt}


%% ============================================================
\section{Problem \#913}
\label{prob:913}

\noindent\textbf{Status:} OPEN \quad|\quad \textbf{Prize:} none \quad|\quad \textbf{Updated:} 2025-08-31

\noindent\textbf{Tags:} number theory

\medskip

\noindent\textbf{Statement:}

Are there infinitely many $n$ such that if\[n(n+1) = \prod_i p_i^{k_i}\]is the factorisation into distinct primes then all exponents $k_i$ are distinct?

\noindent\rule{\linewidth}{0.4pt}


%% ============================================================
\section{Problem \#917}
\label{prob:917}

\noindent\textbf{Status:} OPEN \quad|\quad \textbf{Prize:} none \quad|\quad \textbf{Updated:} 2025-08-31

\noindent\textbf{Tags:} graph theory, chromatic number

\medskip

\noindent\textbf{Statement:}

Let $k\geq 4$ and $f_k(n)$ be the largest number of edges in a graph on $n$ vertices which has chromatic number $k$ and is critical (i.e. deleting any edge reduces the chromatic number). Is it true that\[f_k(n) \gg_k n^2?\]Is it true that\[f_6(n)\sim n^2/4?\]More generally, is it true that, for $k\geq 6$,\[f_k(n) \sim \frac{1}{2}\left(1-\frac{1}{\lfloor k/3\rfloor}\right)n^2?\]

\noindent\rule{\linewidth}{0.4pt}


%% ============================================================
\section{Problem \#918}
\label{prob:918}

\noindent\textbf{Status:} OPEN \quad|\quad \textbf{Prize:} none \quad|\quad \textbf{Updated:} 2025-08-31

\noindent\textbf{Tags:} graph theory, chromatic number

\medskip

\noindent\textbf{Statement:}

Is there a graph with $\aleph_2$ vertices and chromatic number $\aleph_2$ such that every subgraph on $\aleph_1$ vertices has chromatic number $\leq\aleph_0$? Is there a graph with $\aleph_{\omega+1}$ vertices and chromatic number $\aleph_1$ such that every subgraph on $\aleph_\omega$ vertices has chromatic number $\leq\aleph_0$?

\noindent\rule{\linewidth}{0.4pt}


%% ============================================================
\section{Problem \#919}
\label{prob:919}

\noindent\textbf{Status:} OPEN \quad|\quad \textbf{Prize:} none \quad|\quad \textbf{Updated:} 2025-08-31

\noindent\textbf{Tags:} graph theory, chromatic number

\medskip

\noindent\textbf{Statement:}

Is there a graph $G$ with vertex set $\omega_2^2$ and chromatic number $\aleph_2$ such that every subgraph whose vertices have a lesser type has chromatic number $\leq \aleph_0$? What if instead we ask for $G$ to have chromatic number $\aleph_1$?

\noindent\rule{\linewidth}{0.4pt}


%% ============================================================
\section{Problem \#920}
\label{prob:920}

\noindent\textbf{Status:} OPEN \quad|\quad \textbf{Prize:} none \quad|\quad \textbf{Updated:} 2025-08-31

\noindent\textbf{Tags:} graph theory, chromatic number

\medskip

\noindent\textbf{Statement:}

Let $f_k(n)$ be the maximum possible chromatic number of a graph with $n$ vertices which contains no $K_k$. Is it true that, for $k\geq 4$,\[f_k(n) \gg \frac{n^{1-\frac{1}{k-1}}}{(\log n)^{c_k}}\]for some constant $c_k&#62;0$?

\noindent\rule{\linewidth}{0.4pt}


%% ============================================================
\section{Problem \#928}
\label{prob:928}

\noindent\textbf{Status:} OPEN \quad|\quad \textbf{Prize:} none \quad|\quad \textbf{Updated:} 2025-09-04

\noindent\textbf{Tags:} number theory

\medskip

\noindent\textbf{Statement:}

Let $\alpha,\beta\in (0,1)$ and let $P(n)$ denote the largest prime divisor of $n$. Does the density of integers $n$ such that $P(n)&#60;n^{\alpha}$ and $P(n+1)&#60;(n+1)^\beta$ exist?

\noindent\rule{\linewidth}{0.4pt}


%% ============================================================
\section{Problem \#929}
\label{prob:929}

\noindent\textbf{Status:} OPEN \quad|\quad \textbf{Prize:} none \quad|\quad \textbf{Updated:} 2025-08-31

\noindent\textbf{Tags:} number theory

\medskip

\noindent\textbf{Statement:}

Let $k\geq 2$ be large and let $S(k)$ be the minimal $x$ such that there is a positive density set of $n$ where\[n+1,n+2,\ldots,n+k\]are all divisible by primes $\leq x$. Estimate $S(k)$ - in particular, is it true that $S(k)\geq k^{1-o(1)}$?

\noindent\rule{\linewidth}{0.4pt}


%% ============================================================
\section{Problem \#930}
\label{prob:930}

\noindent\textbf{Status:} OPEN \quad|\quad \textbf{Prize:} none \quad|\quad \textbf{Updated:} 2025-08-31

\noindent\textbf{Tags:} number theory

\medskip

\noindent\textbf{Statement:}

Is it true that, for every $r$, there is a $k$ such that if $I_1,\ldots,I_r$ are disjoint intervals of consecutive integers, all of length at least $k$, then\[\prod_{1\leq i\leq r}\prod_{m\in I_i}m\]is not a perfect power?

\noindent\rule{\linewidth}{0.4pt}


%% ============================================================
\section{Problem \#931}
\label{prob:931}

\noindent\textbf{Status:} OPEN \quad|\quad \textbf{Prize:} none \quad|\quad \textbf{Updated:} 2025-08-31

\noindent\textbf{Tags:} number theory

\medskip

\noindent\textbf{Statement:}

Let $k_1\geq k_2\geq 3$. Are there only finitely many $n_2\geq n_1+k_1$ such that\[\prod_{1\leq i\leq k_1}(n_1+i)\textrm{ and }\prod_{1\leq j\leq k_2}(n_2+j)\]have the same prime factors?

\noindent\rule{\linewidth}{0.4pt}


%% ============================================================
\section{Problem \#932}
\label{prob:932}

\noindent\textbf{Status:} OPEN \quad|\quad \textbf{Prize:} none \quad|\quad \textbf{Updated:} 2025-08-31

\noindent\textbf{Tags:} number theory

\medskip

\noindent\textbf{Statement:}

Let $p_k$ denote the $k$th prime. For infinitely many $r$ there are at least two integers $p_r&#60;n&#60;p_{r+1}$ all of whose prime factors are $&#60;p_{r+1}-p_r$.

\noindent\rule{\linewidth}{0.4pt}


%% ============================================================
\section{Problem \#933}
\label{prob:933}

\noindent\textbf{Status:} OPEN \quad|\quad \textbf{Prize:} none \quad|\quad \textbf{Updated:} 2025-08-31

\noindent\textbf{Tags:} number theory

\medskip

\noindent\textbf{Statement:}

If $n(n+1)=2^k3^lm$, where $(m,6)=1$, then is it true that\[\limsup_{n\to \infty} \frac{2^k3^l}{n\log n}=\infty?\]

\noindent\rule{\linewidth}{0.4pt}


%% ============================================================
\section{Problem \#934}
\label{prob:934}

\noindent\textbf{Status:} OPEN \quad|\quad \textbf{Prize:} none \quad|\quad \textbf{Updated:} 2025-08-31

\noindent\textbf{Tags:} graph theory

\medskip

\noindent\textbf{Statement:}

Let $h_t(d)$ be minimal such that every graph $G$ with $h_t(d)$ edges and maximal degree $\leq d$ contains two edges whose shortest path between them has length $\geq t$. Estimate $h_t(d)$.

\noindent\rule{\linewidth}{0.4pt}


%% ============================================================
\section{Problem \#935}
\label{prob:935}

\noindent\textbf{Status:} OPEN \quad|\quad \textbf{Prize:} none \quad|\quad \textbf{Updated:} 2025-09-04

\noindent\textbf{Tags:} number theory

\medskip

\noindent\textbf{Statement:}

For any integer $n=\prod p^{k_p}$ let $Q_2(n)$ be the powerful part of $n$, so that\[Q_2(n) = \prod_{\substack{p\\ k_p\geq 2}}p^{k_p}.\]Is it true that, for every $\epsilon&#62;0$ and $\ell\geq 1$, if $n$ is sufficiently large then\[Q_2(n(n+1)\cdots(n+\ell))&#60;n^{2+\epsilon}?\]If $\ell\geq 2$ then is\[\limsup_{n\to \infty}\frac{Q_2(n(n+1)\cdots(n+\ell))}{n^2}\]infinite? If $\ell\geq 2$ then is\[\lim_{n\to \infty}\frac{Q_2(n(n+1)\cdots(n+\ell))}{n^{\ell+1}}=0?\]

\noindent\rule{\linewidth}{0.4pt}


%% ============================================================
\section{Problem \#936}
\label{prob:936}

\noindent\textbf{Status:} OPEN \quad|\quad \textbf{Prize:} none \quad|\quad \textbf{Updated:} 2025-08-31

\noindent\textbf{Tags:} number theory

\medskip

\noindent\textbf{Statement:}

Are\[2^n\pm 1\]and\[n!\pm 1\]powerful (i.e. if $p\mid m$ then $p^2\mid m$) for only finitely many $n$?

\noindent\rule{\linewidth}{0.4pt}


%% ============================================================
\section{Problem \#938}
\label{prob:938}

\noindent\textbf{Status:} OPEN \quad|\quad \textbf{Prize:} none \quad|\quad \textbf{Updated:} 2025-08-31

\noindent\textbf{Tags:} number theory

\medskip

\noindent\textbf{Statement:}

Let $A=\{n_1&#60;n_2&#60;\cdots\}$ be the sequence of powerful numbers (if $p\mid n$ then $p^2\mid n$). Are there only finitely many three-term progressions of consecutive terms $n_k,n_{k+1},n_{k+2}$?

\noindent\rule{\linewidth}{0.4pt}


%% ============================================================
\section{Problem \#939}
\label{prob:939}

\noindent\textbf{Status:} OPEN \quad|\quad \textbf{Prize:} none \quad|\quad \textbf{Updated:} 2025-08-31

\noindent\textbf{Tags:} number theory

\medskip

\noindent\textbf{Statement:}

Let $r\geq 2$. An $r$-powerful number $n$ is one such that if $p\mid n$ then $p^r\mid n$. If $r\geq 4$ then can the sum of $r-2$ coprime $r$-powerful numbers ever be itself $r$-powerful? Are there at most finitely many such solutions? Are there infinitely many triples of coprime $3$-powerful numbers $a,b,c$ such that $a+b=c$?

\noindent\rule{\linewidth}{0.4pt}


%% ============================================================
\section{Problem \#940}
\label{prob:940}

\noindent\textbf{Status:} OPEN \quad|\quad \textbf{Prize:} none \quad|\quad \textbf{Updated:} 2025-08-31

\noindent\textbf{Tags:} number theory

\medskip

\noindent\textbf{Statement:}

Let $r\geq 3$. A number $n$ is $r$-powerful if for every prime $p$ which divides $n$ we have $p^r\mid n$. Are there infinitely many integers which are not the sum of at most $r$ many $r$-powerful numbers? Does the set of integers which are the sum of at most $r$ $r$-powerful numbers have density $0$?

\noindent\rule{\linewidth}{0.4pt}


%% ============================================================
\section{Problem \#942}
\label{prob:942}

\noindent\textbf{Status:} OPEN \quad|\quad \textbf{Prize:} none \quad|\quad \textbf{Updated:} 2025-08-31

\noindent\textbf{Tags:} number theory

\medskip

\noindent\textbf{Statement:}

Let $h(n)$ count the number of powerful (if $p\mid m$ then $p^2\mid m$) integers in $[n^2,(n+1)^2)$. Estimate $h(n)$. In particular is there some constant $c&#62;0$ such that\[h(n) &lt; (\log n)^{c+o(1)}\]and, for infinitely many $n$,\[h(n) &gt;(\log n)^{c-o(1)}?\]

\noindent\rule{\linewidth}{0.4pt}


%% ============================================================
\section{Problem \#943}
\label{prob:943}

\noindent\textbf{Status:} OPEN \quad|\quad \textbf{Prize:} none \quad|\quad \textbf{Updated:} 2025-08-31

\noindent\textbf{Tags:} number theory

\medskip

\noindent\textbf{Statement:}

Let $A$ be the set of powerful numbers (if $p\mid n$ then $p^2\mid n$). Is it true that\[1_A\ast 1_A(n)=n^{o(1)}\]for every $n$?

\noindent\rule{\linewidth}{0.4pt}


%% ============================================================
\section{Problem \#944}
\label{prob:944}

\noindent\textbf{Status:} OPEN \quad|\quad \textbf{Prize:} none \quad|\quad \textbf{Updated:} 2025-08-31

\noindent\textbf{Tags:} graph theory, chromatic number

\medskip

\noindent\textbf{Statement:}

A critical vertex, edge, or set of edges, is one whose deletion lowers the chromatic number. Let $k\geq 4$ and $r\geq 1$. Must there exist a graph $G$ with chromatic number $k$ such that every vertex is critical, yet every critical set of edges has size $&#62;r$?

\noindent\rule{\linewidth}{0.4pt}


%% ============================================================
\section{Problem \#945}
\label{prob:945}

\noindent\textbf{Status:} OPEN \quad|\quad \textbf{Prize:} none \quad|\quad \textbf{Updated:} 2025-08-31

\noindent\textbf{Tags:} number theory, divisors

\medskip

\noindent\textbf{Statement:}

Let $F(x)$ be the maximal $k$ such that there exist $n+1,\ldots,n+k\leq x$ with $\tau(n+1),\ldots,\tau(n+k)$ all distinct (where $\tau(m)$ counts the divisors of $m$). Estimate $F(x)$. In particular, is it true that\[F(x) \leq (\log x)^{O(1)}?\]In other words, is there a constant $C&#62;0$ such that, for all large $x$, every interval $[x,x+(\log x)^C]$ contains two integers with the same number of divisors?

\noindent\rule{\linewidth}{0.4pt}


%% ============================================================
\section{Problem \#948}
\label{prob:948}

\noindent\textbf{Status:} OPEN \quad|\quad \textbf{Prize:} none \quad|\quad \textbf{Updated:} 2025-08-31

\noindent\textbf{Tags:} number theory, ramsey theory

\medskip

\noindent\textbf{Statement:}

Is there a function $f(n)$ and a $k$ such that in any $k$-colouring of the integers there exists a sequence $a_1&#60;\cdots$ such that $a_n&#60;f(n)$ for infinitely many $n$ and the set\[\left\{ \sum_{i\in S}a_i : \textrm{finite }S\right\}\]does not contain all colours?

\noindent\rule{\linewidth}{0.4pt}


%% ============================================================
\section{Problem \#949}
\label{prob:949}

\noindent\textbf{Status:} OPEN \quad|\quad \textbf{Prize:} none \quad|\quad \textbf{Updated:} 2025-08-31

\noindent\textbf{Tags:} ramsey theory

\medskip

\noindent\textbf{Statement:}

Let $S\subset \mathbb{R}$ be a set containing no solutions to $a+b=c$. Must there be a set $A\subseteq \mathbb{R}\backslash S$ of cardinality continuum such that $A+A\subseteq \mathbb{R}\backslash S$?

\noindent\rule{\linewidth}{0.4pt}


%% ============================================================
\section{Problem \#950}
\label{prob:950}

\noindent\textbf{Status:} OPEN \quad|\quad \textbf{Prize:} none \quad|\quad \textbf{Updated:} 2025-08-31

\noindent\textbf{Tags:} number theory, primes

\medskip

\noindent\textbf{Statement:}

Let\[f(n) = \sum_{p&#60;n}\frac{1}{n-p}.\]Is it true that\[\liminf f(n)=1\]and\[\limsup f(n)=\infty?\]Is it true that $f(n)=o(\log\log n)$ for all $n$?

\noindent\rule{\linewidth}{0.4pt}


%% ============================================================
\section{Problem \#951}
\label{prob:951}

\noindent\textbf{Status:} OPEN \quad|\quad \textbf{Prize:} none \quad|\quad \textbf{Updated:} 2025-08-31

\noindent\textbf{Tags:} number theory

\medskip

\noindent\textbf{Statement:}

Let $1&#60;a_1&#60;\cdots$ be a sequence of real numbers such that\[\left\lvert \prod_i a_i^{k_i}-\prod_j a_j^{\ell_j}\right\rvert \geq 1\]for every distinct pair of non-negative finitely supported integer tuples $k_i,\ell_j\geq 0$. Is it true that\[\#\{ a_i \leq x\} \leq \pi(x)?\]

\noindent\rule{\linewidth}{0.4pt}


%% ============================================================
\section{Problem \#952}
\label{prob:952}

\noindent\textbf{Status:} OPEN \quad|\quad \textbf{Prize:} none \quad|\quad \textbf{Updated:} 2025-08-31

\noindent\textbf{Tags:} number theory

\medskip
\noindent\textit{Gaussian moat problem}


\noindent\textbf{Statement:}

Is there an infinite sequence of distinct Gaussian primes $x_1,x_2,\ldots$ such that\[\lvert x_{n+1}-x_n\rvert \ll 1?\]

\noindent\rule{\linewidth}{0.4pt}


%% ============================================================
\section{Problem \#953}
\label{prob:953}

\noindent\textbf{Status:} OPEN \quad|\quad \textbf{Prize:} none \quad|\quad \textbf{Updated:} 2025-08-31

\noindent\textbf{Tags:} geometry, distances

\medskip

\noindent\textbf{Statement:}

Let $A\subset \{ x\in \mathbb{R}^2 : \lvert x\rvert &#60;r\}$ be a measurable set with no integer distances, that is, such that $\lvert a-b\rvert \not\in \mathbb{Z}$ for any distinct $a,b\in A$. How large can the measure of $A$ be?

\noindent\rule{\linewidth}{0.4pt}


%% ============================================================
\section{Problem \#954}
\label{prob:954}

\noindent\textbf{Status:} OPEN \quad|\quad \textbf{Prize:} none \quad|\quad \textbf{Updated:} 2025-08-31

\noindent\textbf{Tags:} number theory

\medskip

\noindent\textbf{Statement:}

Let $0=a_0&#60;a_1&#60;a_2&#60;\cdots$ be the sequence of integers defined by $a_0=0$ and $a_1=1$, and $a_{k+1}$ is the smallest integer $n$ for which the number of solutions to $a_i+a_j \leq n$ (with $0\leq i\leq j\leq k$ and $j\geq 1$) is $&#60;n$. Is the number of solutions to $a_i+a_j \leq x$ equal to $x+O(x^{1/4+o(1)})$?

\noindent\rule{\linewidth}{0.4pt}


%% ============================================================
\section{Problem \#955}
\label{prob:955}

\noindent\textbf{Status:} OPEN \quad|\quad \textbf{Prize:} none \quad|\quad \textbf{Updated:} 2025-08-31

\noindent\textbf{Tags:} number theory

\medskip

\noindent\textbf{Statement:}

Let\[s(n)=\sigma(n)-n=\sum_{\substack{d\mid n\\ d&#60;n}}d\]be the sum of proper divisors function. If $A\subset \mathbb{N}$ has density $0$ then $s^{-1}(A)$ must also have density $0$.

\noindent\rule{\linewidth}{0.4pt}


%% ============================================================
\section{Problem \#956}
\label{prob:956}

\noindent\textbf{Status:} OPEN \quad|\quad \textbf{Prize:} none \quad|\quad \textbf{Updated:} 2025-08-31

\noindent\textbf{Tags:} geometry, distances, convex

\medskip

\noindent\textbf{Statement:}

If $C,D\subseteq \mathbb{R}^2$ then the distance between $C$ and $D$ is defined by\[\delta(C,D)=\inf_{\substack{c\in C\\ d\in D}}\| c-d\|.\]Let $h(n)$ be the maximal number of unit distances between disjoint convex translates. That is, the maximal $m$ such that there is a compact convex set $C\subset \mathbb{R}^2$ and a set $X$ of size $n$ such that all $(C+x)_{x\in X}$ are disjoint and there are $m$ pairs $x_1,x_2\in X$ such that\[\delta(C+x_1,C+x_2)=1.\]Determine $h(n)$ - in particular, prove that there exists a constant $c&#62;0$ such that $h(n)&#62;n^{1+c}$ for all large $n$.

\noindent\rule{\linewidth}{0.4pt}


%% ============================================================
\section{Problem \#959}
\label{prob:959}

\noindent\textbf{Status:} OPEN \quad|\quad \textbf{Prize:} none \quad|\quad \textbf{Updated:} 2025-08-31

\noindent\textbf{Tags:} geometry, distances

\medskip

\noindent\textbf{Statement:}

Let $A\subset \mathbb{R}^2$ be a set of size $n$ and let $\{d_1,\ldots,d_k\}$ be the set of distinct distances determined by $A$. Let $f(d)$ be the number of times the distance $d$ is determined, and suppose the $d_i$ are ordered such that\[f(d_1)\geq f(d_2)\geq \cdots \geq f(d_k).\]Estimate\[\max (f(d_1)-f(d_2)),\]where the maximum is taken over all $A$ of size $n$.

\noindent\rule{\linewidth}{0.4pt}


%% ============================================================
\section{Problem \#961}
\label{prob:961}

\noindent\textbf{Status:} OPEN \quad|\quad \textbf{Prize:} none \quad|\quad \textbf{Updated:} 2025-08-31

\noindent\textbf{Tags:} number theory

\medskip

\noindent\textbf{Statement:}

Let $f(k)$ be the minimal $n$ such that every set of $n$ consecutive integers $&#62;k$ contains an integer divisible by a prime $&#62;k$. Estimate $f(k)$.

\noindent\rule{\linewidth}{0.4pt}


%% ============================================================
\section{Problem \#962}
\label{prob:962}

\noindent\textbf{Status:} OPEN \quad|\quad \textbf{Prize:} none \quad|\quad \textbf{Updated:} 2025-08-31

\noindent\textbf{Tags:} number theory

\medskip

\noindent\textbf{Statement:}

Let $k(n)$ be the maximal $k$ such that there exists $m\leq n$ such that each of the integers\[m+1,\ldots,m+k\]are divisible by at least one prime $&#62;k$. Estimate $k(n)$ - in particular, is it true that\[\log k(n) \leq (\log n)^{1/2+o(1)}?\]

\noindent\rule{\linewidth}{0.4pt}


%% ============================================================
\section{Problem \#963}
\label{prob:963}

\noindent\textbf{Status:} OPEN \quad|\quad \textbf{Prize:} none \quad|\quad \textbf{Updated:} 2025-08-31

\noindent\textbf{Tags:} number theory

\medskip

\noindent\textbf{Statement:}

Let $f(n)$ be the maximal $k$ such that in any set $A\subset \mathbb{R}$ of size $n$ there is a subset $B\subseteq A$ of size $\lvert B\rvert\geq k$ which is dissociated that is, the sums $\sum_{b\in S}b$ are distinct for all $S\subseteq B$. Estimate $f(n)$ - in particular, is it true that\[f(n)\geq \lfloor \log_2 n\rfloor?\]

\noindent\rule{\linewidth}{0.4pt}


%% ============================================================
\section{Problem \#968}
\label{prob:968}

\noindent\textbf{Status:} OPEN \quad|\quad \textbf{Prize:} none \quad|\quad \textbf{Updated:} 2025-08-31

\noindent\textbf{Tags:} number theory

\medskip

\noindent\textbf{Statement:}

Let $u_n=p_n/n$, where $p_n$ is the $n$th prime. Does the set of $n$ such that $u_n&#60;u_{n+1}$ have positive density?

\noindent\rule{\linewidth}{0.4pt}


%% ============================================================
\section{Problem \#969}
\label{prob:969}

\noindent\textbf{Status:} OPEN \quad|\quad \textbf{Prize:} none \quad|\quad \textbf{Updated:} 2025-08-31

\noindent\textbf{Tags:} number theory

\medskip

\noindent\textbf{Statement:}

Let $Q(x)$ count the number of squarefree integers in $[1,x]$. Determine the order of magnitude in the error term in the asymptotic\[Q(x)=\frac{6}{\pi^2}x+E(x).\]

\noindent\rule{\linewidth}{0.4pt}


%% ============================================================
\section{Problem \#970}
\label{prob:970}

\noindent\textbf{Status:} OPEN \quad|\quad \textbf{Prize:} none \quad|\quad \textbf{Updated:} 2025-08-31

\noindent\textbf{Tags:} number theory

\medskip
\noindent\textit{Jacobsthal's function}


\noindent\textbf{Statement:}

Let $h(k)$ be Jacobsthal's function, defined to as the minimal $m$ such that, if $n$ has at most $k$ prime factors, then in any set of $m$ consecutive integers there exists an integer coprime to $n$. Determine the order of magnitude of $h(k)$. In particular, is it true that\[h(k) \ll k^2?\]

\noindent\rule{\linewidth}{0.4pt}


%% ============================================================
\section{Problem \#971}
\label{prob:971}

\noindent\textbf{Status:} OPEN \quad|\quad \textbf{Prize:} none \quad|\quad \textbf{Updated:} 2025-08-31

\noindent\textbf{Tags:} number theory

\medskip

\noindent\textbf{Statement:}

Let $p(a,d)$ be the least prime congruent to $a\pmod{d}$. Does there exist a constant $c&#62;0$ such that, for all large $d$,\[p(a,d) &#62; (1+c)\phi(d)\log d\]for $\gg \phi(d)$ many values of $a$?

\noindent\rule{\linewidth}{0.4pt}


%% ============================================================
\section{Problem \#972}
\label{prob:972}

\noindent\textbf{Status:} OPEN \quad|\quad \textbf{Prize:} none \quad|\quad \textbf{Updated:} 2025-08-31

\noindent\textbf{Tags:} number theory

\medskip

\noindent\textbf{Statement:}

Let $\alpha&#62;1$ be irrational. Are there infinitely many primes $p$ such that $\lfloor p\alpha\rfloor$ is also prime?

\noindent\rule{\linewidth}{0.4pt}


%% ============================================================
\section{Problem \#973}
\label{prob:973}

\noindent\textbf{Status:} OPEN \quad|\quad \textbf{Prize:} none \quad|\quad \textbf{Updated:} 2025-08-31

\noindent\textbf{Tags:} analysis

\medskip

\noindent\textbf{Statement:}

Does there exist a constant $C&#62;1$ such that, for every $n\geq 2$, there exists a sequence $z_i\in \mathbb{C}$ with $z_1=1$ and $\lvert z_i\rvert \geq 1$ for all $1\leq i\leq n$ with\[\max_{2\leq k\leq n+1}\left\lvert \sum_{1\leq i\leq n}z_i^k\right\rvert &#60; C^{-n}?\]

\noindent\rule{\linewidth}{0.4pt}


%% ============================================================
\section{Problem \#975}
\label{prob:975}

\noindent\textbf{Status:} OPEN \quad|\quad \textbf{Prize:} none \quad|\quad \textbf{Updated:} 2025-08-31

\noindent\textbf{Tags:} number theory, divisors

\medskip

\noindent\textbf{Statement:}

Let $f\in \mathbb{Z}[x]$ be an irreducible non-constant polynomial such that $f(n)\geq 1$ for all large $n\in\mathbb{N}$. Does there exist a constant $c=c(f)&#62;0$ such that\[\sum_{n\leq X} \tau(f(n))\sim cX\log X,\]where $\tau$ is the divisor function?

\noindent\rule{\linewidth}{0.4pt}


%% ============================================================
\section{Problem \#976}
\label{prob:976}

\noindent\textbf{Status:} OPEN \quad|\quad \textbf{Prize:} none \quad|\quad \textbf{Updated:} 2025-08-31

\noindent\textbf{Tags:} number theory

\medskip

\noindent\textbf{Statement:}

Let $f\in \mathbb{Z}[x]$ be an irreducible polynomial of degree $d\geq 2$. Let $F_f(n)$ be maximal such that there exists $1\leq m\leq n$ with $f(m)$ is divisible by a prime $\geq F_f(n)$. Equivalently, $F_f(n)$ is the greatest prime divisor of\[\prod_{1\leq m\leq n}f(m).\]Estimate $F_f(n)$. In particular, is it true that $F_f(n)\gg n^{1+c}$ for some constant $c&#62;0$? Or even $\gg n^d$?

\noindent\rule{\linewidth}{0.4pt}


%% ============================================================
\section{Problem \#978}
\label{prob:978}

\noindent\textbf{Status:} OPEN \quad|\quad \textbf{Prize:} none \quad|\quad \textbf{Updated:} 2025-08-31

\noindent\textbf{Tags:} number theory

\medskip

\noindent\textbf{Statement:}

Let $f\in \mathbb{Z}[x]$ be an irreducible polynomial of degree $k&#62;2$ (and suppose that $k\neq 2^l$ for any $l\geq 1$) such that the leading coefficient of $f$ is positive. Does the set of integers $n\geq 1$ for which $f(n)$ is $(k-1)$-power-free have positive density? If $k&#62;3$, and for all primes $p$ there exists $n$ such that $p^{k-2}\nmid f(n)$, then are there infinitely many $n$ for which $f(n)$ is $(k-2)$-power-free? In particular, does\[n^4+2\]represent infinitely many squarefree numbers?

\noindent\rule{\linewidth}{0.4pt}


%% ============================================================
\section{Problem \#979}
\label{prob:979}

\noindent\textbf{Status:} OPEN \quad|\quad \textbf{Prize:} none \quad|\quad \textbf{Updated:} 2025-08-31

\noindent\textbf{Tags:} number theory

\medskip

\noindent\textbf{Statement:}

Let $k\geq 2$, and let $f_k(n)$ count the number of solutions to\[n=p_1^k+\cdots+p_k^k,\]where the $p_i$ are prime numbers. Is it true that $\limsup f_k(n)=\infty$?

\noindent\rule{\linewidth}{0.4pt}


%% ============================================================
\section{Problem \#983}
\label{prob:983}

\noindent\textbf{Status:} OPEN \quad|\quad \textbf{Prize:} none \quad|\quad \textbf{Updated:} 2025-08-31

\noindent\textbf{Tags:} number theory

\medskip

\noindent\textbf{Statement:}

Let $n\geq 2$ and $\pi(n)&lt;k\leq n$. Let $f(k,n)$ be the smallest integer $r$ such that in any $A\subseteq \{1,\ldots,n\}$ of size $\lvert A\rvert=k$ there exist primes $p_1,\ldots,p_r$ such that $&gt;r$ many $a\in A$ are only divisible by primes from $\{p_1,\ldots,p_r\}$. Is it true that\[2\pi(n^{1/2})-f(\pi(n)+1,n)\to \infty\]as $n\to \infty$? In general, estimate $f(k,n)$, particularly when $\pi(n)+1&#60;k=o(n)$.

\noindent\rule{\linewidth}{0.4pt}


%% ============================================================
\section{Problem \#985}
\label{prob:985}

\noindent\textbf{Status:} OPEN \quad|\quad \textbf{Prize:} none \quad|\quad \textbf{Updated:} 2025-08-31

\noindent\textbf{Tags:} number theory

\medskip

\noindent\textbf{Statement:}

Is it true that, for every prime $p$, there is a prime $q&#60;p$ which is a primitive root modulo $p$?

\noindent\rule{\linewidth}{0.4pt}


%% ============================================================
\section{Problem \#986}
\label{prob:986}

\noindent\textbf{Status:} OPEN \quad|\quad \textbf{Prize:} none \quad|\quad \textbf{Updated:} 2025-08-31

\noindent\textbf{Tags:} graph theory, ramsey theory

\medskip

\noindent\textbf{Statement:}

For any fixed $s\geq 3$,\[R(s,k) \gg \frac{k^{s-1}}{(\log k)^c}\]for some constant $c=c(s)&#62;0$.

\noindent\rule{\linewidth}{0.4pt}


%% ============================================================
\section{Problem \#995}
\label{prob:995}

\noindent\textbf{Status:} OPEN \quad|\quad \textbf{Prize:} none \quad|\quad \textbf{Updated:} 2025-09-07

\noindent\textbf{Tags:} analysis, discrepancy

\medskip

\noindent\textbf{Statement:}

Let $n_1&#60;n_2&#60;\cdots$ be a lacunary sequence of integers and $f\in L^2([0,1])$. Estimate the growth of, for almost all $\alpha$,\[\sum_{1\leq k\leq N}f(\{ \alpha n_k\}).\]For example, is it true that, for almost all $\alpha$,\[\sum_{1\leq k\leq N}f(\{ \alpha n_k\})=o(N\sqrt{\log\log N})?\]

\noindent\rule{\linewidth}{0.4pt}


%% ============================================================
\section{Problem \#996}
\label{prob:996}

\noindent\textbf{Status:} OPEN \quad|\quad \textbf{Prize:} none \quad|\quad \textbf{Updated:} 2025-09-07

\noindent\textbf{Tags:} analysis

\medskip

\noindent\textbf{Statement:}

Let $n_1&lt;n_2&lt;\cdots$ be a lacunary sequence of integers, and let $f\in L^2([0,1])$. Let $f_n$ be the $n$th partial sum of the Fourier series of $f(x)$. Is there an absolute constant $C&gt;0$ such that, if\[\| f-f_n\|_2 \ll \frac{1}{(\log\log\log n)^{C}}\]then\[\lim_{N\to\infty}\frac{1}{N}\sum_{k\leq N}f(\{\alpha n_k\})=\int_0^1 f(x)\mathrm{d}x\]for almost every $\alpha$?

\noindent\rule{\linewidth}{0.4pt}


%% ============================================================
\section{Problem \#1002}
\label{prob:1002}

\noindent\textbf{Status:} OPEN \quad|\quad \textbf{Prize:} none \quad|\quad \textbf{Updated:} 2025-09-07

\noindent\textbf{Tags:} analysis, diophantine approximation

\medskip

\noindent\textbf{Statement:}

For any $0&#60;\alpha&#60;1$, let\[f(\alpha,n)=\frac{1}{\log n}\sum_{1\leq k\leq n}(\tfrac{1}{2}-\{ \alpha k\}).\]Does $f(\alpha,n)$ have an asymptotic distribution function? In other words, is there a non-decreasing function $g$ such that $g(-\infty)=0$, $g(\infty)=1$, and\[\lim_{n\to \infty}\lvert \{ \alpha\in (0,1): f(\alpha,n)\leq c\}\rvert=g(c)?\]

\noindent\rule{\linewidth}{0.4pt}


%% ============================================================
\section{Problem \#1003}
\label{prob:1003}

\noindent\textbf{Status:} OPEN \quad|\quad \textbf{Prize:} none \quad|\quad \textbf{Updated:} 2025-09-08

\noindent\textbf{Tags:} number theory

\medskip

\noindent\textbf{Statement:}

Are there infinitely many solutions to $\phi(n)=\phi(n+1)$, where $\phi$ is the Euler totient function?

\noindent\rule{\linewidth}{0.4pt}


%% ============================================================
\section{Problem \#1004}
\label{prob:1004}

\noindent\textbf{Status:} OPEN \quad|\quad \textbf{Prize:} none \quad|\quad \textbf{Updated:} 2025-09-07

\noindent\textbf{Tags:} number theory

\medskip

\noindent\textbf{Statement:}

Let $c&#62;0$. If $x$ is sufficiently large then does there exist $n\leq x$ such that the values of $\phi(n+k)$ are all distinct for $1\leq k\leq (\log x)^c$, where $\phi$ is the Euler totient function?

\noindent\rule{\linewidth}{0.4pt}


%% ============================================================
\section{Problem \#1005}
\label{prob:1005}

\noindent\textbf{Status:} OPEN \quad|\quad \textbf{Prize:} none \quad|\quad \textbf{Updated:} 2025-09-09

\noindent\textbf{Tags:} number theory

\medskip

\noindent\textbf{Statement:}

Let $\frac{a_1}{b_1},\frac{a_2}{b_2},\ldots$ be the Farey fractions of order $n\geq 4$. Let $f(n)$ be the largest integer such that if $1\leq k&lt;l\leq k+f(n)$ then $\frac{a_k}{b_k}$ and $\frac{a_l}{b_l}$ are similarly ordered - in other words,\[(a_k-a_l)(b_k-b_l)\geq 0.\]Estimate $f(n)$ - in particular, is there a constant $c&gt;0$ such that $f(n)=(c+o(1))n$ for all large $n$?

\noindent\rule{\linewidth}{0.4pt}


%% ============================================================
\section{Problem \#1011}
\label{prob:1011}

\noindent\textbf{Status:} OPEN \quad|\quad \textbf{Prize:} none \quad|\quad \textbf{Updated:} 2025-09-10

\noindent\textbf{Tags:} graph theory

\medskip

\noindent\textbf{Statement:}

Let $f_r(n)$ be minimal such that every graph on $n$ vertices with $\geq f_r(n)$ edges and chromatic number $\geq r$ contains a triangle. Determine $f_r(n)$.

\noindent\rule{\linewidth}{0.4pt}


%% ============================================================
\section{Problem \#1013}
\label{prob:1013}

\noindent\textbf{Status:} OPEN \quad|\quad \textbf{Prize:} none \quad|\quad \textbf{Updated:} 2025-09-10

\noindent\textbf{Tags:} graph theory

\medskip

\noindent\textbf{Statement:}

Let $h_3(k)$ be the minimal $n$ such that there exists a triangle-free graph on $n$ vertices with chromatic number $k$. Find an asymptotic for $h_3(k)$, and also prove\[\lim_{k\to \infty}\frac{h_3(k+1)}{h_3(k)}=1.\]

\noindent\rule{\linewidth}{0.4pt}


%% ============================================================
\section{Problem \#1016}
\label{prob:1016}

\noindent\textbf{Status:} OPEN \quad|\quad \textbf{Prize:} none \quad|\quad \textbf{Updated:} 2025-09-10

\noindent\textbf{Tags:} graph theory, cycles

\medskip
\noindent\textit{pancyclic graphs}


\noindent\textbf{Statement:}

Let $h(n)$ be minimal such that there is a graph on $n$ vertices with $n+h(n)$ edges which contains a cycle on $k$ vertices, for all $3\leq k\leq n$. Estimate $h(n)$. In particular, is it true that\[h(n) \geq \log_2n+\log_*n-O(1),\]where $\log_*n$ is the iterated logarithmic function?

\noindent\rule{\linewidth}{0.4pt}


%% ============================================================
\section{Problem \#1017}
\label{prob:1017}

\noindent\textbf{Status:} OPEN \quad|\quad \textbf{Prize:} none \quad|\quad \textbf{Updated:} 2025-09-12

\noindent\textbf{Tags:} graph theory

\medskip

\noindent\textbf{Statement:}

Let $f(n,k)$ be such that every graph on $n$ vertices and $k$ edges can be partitioned into at most $f(n,k)$ edge-disjoint complete graphs. Estimate $f(n,k)$ for $k&#62;n^2/4$.

\noindent\rule{\linewidth}{0.4pt}


%% ============================================================
\section{Problem \#1029}
\label{prob:1029}

\noindent\textbf{Status:} OPEN \quad|\quad \textbf{Prize:} \$100 \quad|\quad \textbf{Updated:} 2025-09-13

\noindent\textbf{Tags:} graph theory, ramsey theory

\medskip

\noindent\textbf{Statement:}

If $R(k)$ is the Ramsey number for $K_k$, the minimal $n$ such that every $2$-colouring of the edges of $K_n$ contains a monochromatic copy of $K_k$, then\[\frac{R(k)}{k2^{k/2}}\to \infty.\]

\noindent\rule{\linewidth}{0.4pt}


%% ============================================================
\section{Problem \#1030}
\label{prob:1030}

\noindent\textbf{Status:} OPEN \quad|\quad \textbf{Prize:} none \quad|\quad \textbf{Updated:} 2025-09-13

\noindent\textbf{Tags:} graph theory, ramsey theory

\medskip

\noindent\textbf{Statement:}

Let $R(k,l)$ be the usual Ramsey number: the smallest $n$ such that if the edges of $K_n$ are coloured red and blue then there exists either a red $K_k$ or a blue $K_l$. Prove the existence of some $c&#62;0$ such that\[\lim_{k\to \infty}\frac{R(k+1,k)}{R(k,k)}&#62; 1+c.\]

\noindent\rule{\linewidth}{0.4pt}


%% ============================================================
\section{Problem \#1032}
\label{prob:1032}

\noindent\textbf{Status:} OPEN \quad|\quad \textbf{Prize:} none \quad|\quad \textbf{Updated:} 2025-09-13

\noindent\textbf{Tags:} graph theory, chromatic number

\medskip

\noindent\textbf{Statement:}

We say that a graph is $4$-chromatic critical if it has chromatic number $4$, and removing any edge decreases the chromatic number to $3$. Is there, for arbitrarily large $n$, a $4$-chromatic critical graph on $n$ vertices with minimum degree $\gg n$?

\noindent\rule{\linewidth}{0.4pt}


%% ============================================================
\section{Problem \#1033}
\label{prob:1033}

\noindent\textbf{Status:} OPEN \quad|\quad \textbf{Prize:} none \quad|\quad \textbf{Updated:} 2025-12-12

\noindent\textbf{Tags:} graph theory

\medskip

\noindent\textbf{Statement:}

Let $h(n)$ be such that every graph on $n$ vertices with $&#62;n^2/4$ many edges contains a triangle whose vertices have degrees summing to at least $h(n)$. Estimate $h(n)$. In particular, is it true that\[h(n)\geq (2(\sqrt{3}-1)-o(1))n?\]

\noindent\rule{\linewidth}{0.4pt}


%% ============================================================
\section{Problem \#1035}
\label{prob:1035}

\noindent\textbf{Status:} OPEN \quad|\quad \textbf{Prize:} none \quad|\quad \textbf{Updated:} 2025-12-26

\noindent\textbf{Tags:} graph theory

\medskip

\noindent\textbf{Statement:}

Is there a constant $c&#62;0$ such that every graph on $2^n$ vertices with minimum degree $&#62;(1-c)2^n$ contains the $n$-dimensional hypercube $Q_n$?

\noindent\rule{\linewidth}{0.4pt}


%% ============================================================
\section{Problem \#1038}
\label{prob:1038}

\noindent\textbf{Status:} OPEN \quad|\quad \textbf{Prize:} none \quad|\quad \textbf{Updated:} 2025-09-15

\noindent\textbf{Tags:} analysis

\medskip

\noindent\textbf{Statement:}

Determine the infimum and supremum of\[\lvert \{ x\in \mathbb{R} : \lvert f(x)\rvert &#60; 1\}\rvert\]as $f\in \mathbb{R}[x]$ ranges over all non-constant monic polynomials, all of whose roots are real and in the interval $[-1,1]$.

\noindent\rule{\linewidth}{0.4pt}


%% ============================================================
\section{Problem \#1039}
\label{prob:1039}

\noindent\textbf{Status:} OPEN \quad|\quad \textbf{Prize:} none \quad|\quad \textbf{Updated:} 2025-09-15

\noindent\textbf{Tags:} analysis

\medskip

\noindent\textbf{Statement:}

Let $f(z)=\prod_{i=1}^n(z-z_i)\in \mathbb{C}[z]$ with $\lvert z_i\rvert \leq 1$ for all $i$. Let $\rho(f)$ be the radius of the largest disc which is contained in $\{z: \lvert f(z)\rvert&#60; 1\}$. Determine the behaviour of $\rho(f)$. In particular, is it always true that $\rho(f)\gg 1/n$?

\noindent\rule{\linewidth}{0.4pt}


%% ============================================================
\section{Problem \#1040}
\label{prob:1040}

\noindent\textbf{Status:} OPEN \quad|\quad \textbf{Prize:} none \quad|\quad \textbf{Updated:} 2025-09-15

\noindent\textbf{Tags:} analysis

\medskip

\noindent\textbf{Statement:}

Let $F\subseteq \mathbb{C}$ be a closed infinite set, and let $\mu(F)$ be the infimum of\[\lvert \{ z: \lvert f(z)\rvert &#60; 1\}\rvert,\]as $f$ ranges over all polynomials of the shape $\prod (z-z_i)$ with $z_i\in F$. Is $\mu(F)$ determined by the transfinite diameter of $F$? In particular, is $\mu(F)=0$ whenever the transfinite diameter of $F$ is $\geq 1$?

\noindent\rule{\linewidth}{0.4pt}


%% ============================================================
\section{Problem \#1045}
\label{prob:1045}

\noindent\textbf{Status:} OPEN \quad|\quad \textbf{Prize:} none \quad|\quad \textbf{Updated:} 2026-03-14

\noindent\textbf{Tags:} analysis

\medskip

\noindent\textbf{Statement:}

Let $z_1,\ldots,z_n\in \mathbb{C}$ with $\lvert z_i-z_j\rvert\leq 2$ for all $i,j$, and\[\Delta(z_1,\ldots,z_n)=\prod_{i\neq j}\lvert z_i-z_j\rvert.\]What is the maximum possible value of $\Delta$? Is it maximised by taking the $z_i$ to be the vertices of a regular polygon?

\noindent\rule{\linewidth}{0.4pt}


%% ============================================================
\section{Problem \#1049}
\label{prob:1049}

\noindent\textbf{Status:} OPEN \quad|\quad \textbf{Prize:} none \quad|\quad \textbf{Updated:} 2025-09-28

\noindent\textbf{Tags:} irrationality

\medskip

\noindent\textbf{Statement:}

Let $t&#62;1$ be a rational number. Is\[\sum_{n=1}^\infty\frac{1}{t^n-1}=\sum_{n=1}^\infty \frac{\tau(n)}{t^n}\]irrational, where $\tau(n)$ counts the divisors of $n$?

\noindent\rule{\linewidth}{0.4pt}


%% ============================================================
\section{Problem \#1052}
\label{prob:1052}

\noindent\textbf{Status:} OPEN \quad|\quad \textbf{Prize:} \$10 \quad|\quad \textbf{Updated:} 2025-09-28

\noindent\textbf{Tags:} number theory

\medskip
\noindent\textit{unitary perfect numbers}


\noindent\textbf{Statement:}

A unitary divisor of $n$ is $d\mid n$ such that $(d,n/d)=1$. A number $n\geq 1$ is a unitary perfect number if it is the sum of its unitary divisors (aside from $n$ itself). Are there only finite many unitary perfect numbers?

\noindent\rule{\linewidth}{0.4pt}


%% ============================================================
\section{Problem \#1053}
\label{prob:1053}

\noindent\textbf{Status:} OPEN \quad|\quad \textbf{Prize:} none \quad|\quad \textbf{Updated:} 2025-09-28

\noindent\textbf{Tags:} number theory

\medskip
\noindent\textit{multiply perfect numbers}


\noindent\textbf{Statement:}

Call a number $k$-perfect if $\sigma(n)=kn$, where $\sigma(n)$ is the sum of the divisors of $n$. Must $k=o(\log\log n)$?

\noindent\rule{\linewidth}{0.4pt}


%% ============================================================
\section{Problem \#1054}
\label{prob:1054}

\noindent\textbf{Status:} OPEN \quad|\quad \textbf{Prize:} none \quad|\quad \textbf{Updated:} 2025-09-28

\noindent\textbf{Tags:} number theory, divisors

\medskip

\noindent\textbf{Statement:}

Let $f(n)$ be the minimal integer $m$ such that $n$ is the sum of the $k$ smallest divisors of $m$ for some $k\geq 1$. Is it true that $f(n)=o(n)$? Or is this true only for almost all $n$, and $\limsup f(n)/n=\infty$?

\noindent\rule{\linewidth}{0.4pt}


%% ============================================================
\section{Problem \#1055}
\label{prob:1055}

\noindent\textbf{Status:} OPEN \quad|\quad \textbf{Prize:} none \quad|\quad \textbf{Updated:} 2025-09-28

\noindent\textbf{Tags:} number theory, primes

\medskip

\noindent\textbf{Statement:}

A prime $p$ is in class $1$ if the only prime divisors of $p+1$ are $2$ or $3$. In general, a prime $p$ is in class $r$ if every prime factor of $p+1$ is in some class $\leq r-1$, with equality for at least one prime factor. Are there infinitely many primes in each class? If $p_r$ is the least prime in class $r$, then how does $p_r^{1/r}$ behave?

\noindent\rule{\linewidth}{0.4pt}


%% ============================================================
\section{Problem \#1056}
\label{prob:1056}

\noindent\textbf{Status:} OPEN \quad|\quad \textbf{Prize:} none \quad|\quad \textbf{Updated:} 2025-09-28

\noindent\textbf{Tags:} number theory

\medskip

\noindent\textbf{Statement:}

Let $k\geq 2$. Does there exist a prime $p$ and consecutive intervals $I_1,\ldots,I_k$ such that\[\prod_{n\in I_i}n \equiv 1\pmod{p}\]for all $1\leq i\leq k$?

\noindent\rule{\linewidth}{0.4pt}


%% ============================================================
\section{Problem \#1057}
\label{prob:1057}

\noindent\textbf{Status:} OPEN \quad|\quad \textbf{Prize:} none \quad|\quad \textbf{Updated:} 2025-09-28

\noindent\textbf{Tags:} number theory

\medskip
\noindent\textit{Carmichael numbers}


\noindent\textbf{Statement:}

Let $C(x)$ count the number of Carmichael numbers in the interval $[1,x]$. Is it true that $C(x)=x^{1-o(1)}$?

\noindent\rule{\linewidth}{0.4pt}


%% ============================================================
\section{Problem \#1059}
\label{prob:1059}

\noindent\textbf{Status:} OPEN \quad|\quad \textbf{Prize:} none \quad|\quad \textbf{Updated:} 2025-09-28

\noindent\textbf{Tags:} number theory, primes

\medskip

\noindent\textbf{Statement:}

Are there infinitely many primes $p$ such that $p-k!$ is composite for each $k$ such that $1\leq k!&#60;p$?

\noindent\rule{\linewidth}{0.4pt}


%% ============================================================
\section{Problem \#1060}
\label{prob:1060}

\noindent\textbf{Status:} OPEN \quad|\quad \textbf{Prize:} none \quad|\quad \textbf{Updated:} 2025-09-28

\noindent\textbf{Tags:} number theory

\medskip

\noindent\textbf{Statement:}

Let $f(n)$ count the number of solutions to $k\sigma(k)=n$, where $\sigma(k)$ is the sum of divisors of $k$. Is it true that $f(n)\leq n^{o(\frac{1}{\log\log n})}$? Perhaps even $\leq (\log n)^{O(1)}$?

\noindent\rule{\linewidth}{0.4pt}


%% ============================================================
\section{Problem \#1061}
\label{prob:1061}

\noindent\textbf{Status:} OPEN \quad|\quad \textbf{Prize:} none \quad|\quad \textbf{Updated:} 2025-09-28

\noindent\textbf{Tags:} number theory

\medskip

\noindent\textbf{Statement:}

How many solutions are there to\[\sigma(a)+\sigma(b)=\sigma(a+b)\]with $a+b\leq x$, where $\sigma$ is the sum of divisors function? Is it $\sim cx$ for some constant $c&#62;0$?

\noindent\rule{\linewidth}{0.4pt}


%% ============================================================
\section{Problem \#1062}
\label{prob:1062}

\noindent\textbf{Status:} OPEN \quad|\quad \textbf{Prize:} none \quad|\quad \textbf{Updated:} 2025-09-28

\noindent\textbf{Tags:} number theory

\medskip

\noindent\textbf{Statement:}

Let $f(n)$ be the size of the largest subset $A\subseteq \{1,\ldots,n\}$ such that there are no three distinct elements $a,b,c\in A$ such that $a\mid b$ and $a\mid c$. How large can $f(n)$ be? Is $\lim f(n)/n$ irrational?

\noindent\rule{\linewidth}{0.4pt}


%% ============================================================
\section{Problem \#1063}
\label{prob:1063}

\noindent\textbf{Status:} OPEN \quad|\quad \textbf{Prize:} none \quad|\quad \textbf{Updated:} 2025-10-01

\noindent\textbf{Tags:} number theory

\medskip

\noindent\textbf{Statement:}

Let $k\geq 2$ and define $n_k\geq 2k$ to be the least value of $n$ such that $n-i$ divides $\binom{n}{k}$ for all but one $0\leq i&#60;k$. Estimate $n_k$.

\noindent\rule{\linewidth}{0.4pt}


%% ============================================================
\section{Problem \#1065}
\label{prob:1065}

\noindent\textbf{Status:} OPEN \quad|\quad \textbf{Prize:} none \quad|\quad \textbf{Updated:} 2025-10-01

\noindent\textbf{Tags:} number theory

\medskip

\noindent\textbf{Statement:}

Are there infinitely many primes $p$ such that $p=2^kq+1$ for some prime $q$ and $k\geq 0$? Or $p=2^k3^lq+1$?

\noindent\rule{\linewidth}{0.4pt}


%% ============================================================
\section{Problem \#1066}
\label{prob:1066}

\noindent\textbf{Status:} OPEN \quad|\quad \textbf{Prize:} none \quad|\quad \textbf{Updated:} 2025-10-01

\noindent\textbf{Tags:} graph theory, planar graphs

\medskip

\noindent\textbf{Statement:}

Let $G$ be a graph given by $n$ points in $\mathbb{R}^2$, where any two distinct points are at least distance $1$ apart, and we draw an edge between two points if they are distance $1$ apart. Let $g(n)$ be maximal such that any such graph always has an independent set on at least $g(n)$ vertices. Estimate $g(n)$, or perhaps $\lim \frac{g(n)}{n}$.

\noindent\rule{\linewidth}{0.4pt}


%% ============================================================
\section{Problem \#1068}
\label{prob:1068}

\noindent\textbf{Status:} OPEN \quad|\quad \textbf{Prize:} none \quad|\quad \textbf{Updated:} 2025-10-01

\noindent\textbf{Tags:} graph theory, set theory, chromatic number

\medskip

\noindent\textbf{Statement:}

Does every graph with chromatic number $\aleph_1$ contain a countable subgraph which is infinitely vertex-connected?

\noindent\rule{\linewidth}{0.4pt}


%% ============================================================
\section{Problem \#1070}
\label{prob:1070}

\noindent\textbf{Status:} OPEN \quad|\quad \textbf{Prize:} none \quad|\quad \textbf{Updated:} 2025-10-05

\noindent\textbf{Tags:} geometry

\medskip

\noindent\textbf{Statement:}

Let $f(n)$ be maximal such that, given any $n$ points in $\mathbb{R}^2$, there exist $f(n)$ points such that no two are distance $1$ apart. Estimate $f(n)$. In particular, is it true that $f(n)\geq n/4$?

\noindent\rule{\linewidth}{0.4pt}


%% ============================================================
\section{Problem \#1072}
\label{prob:1072}

\noindent\textbf{Status:} OPEN \quad|\quad \textbf{Prize:} none \quad|\quad \textbf{Updated:} 2025-10-05

\noindent\textbf{Tags:} number theory

\medskip

\noindent\textbf{Statement:}

For any prime $p$, let $f(p)$ be the least integer such that $f(p)!+1\equiv 0\pmod{p}$. Is it true that there are infinitely many $p$ for which $f(p)=p-1$? Is it true that $f(p)/p\to 0$ for almost all $p$?

\noindent\rule{\linewidth}{0.4pt}


%% ============================================================
\section{Problem \#1073}
\label{prob:1073}

\noindent\textbf{Status:} OPEN \quad|\quad \textbf{Prize:} none \quad|\quad \textbf{Updated:} 2025-10-05

\noindent\textbf{Tags:} number theory

\medskip

\noindent\textbf{Statement:}

Let $A(x)$ count the number of composite $u&#60;x$ such that $n!+1\equiv 0\pmod{u}$ for some $n$. Is it true that $A(x)\leq x^{o(1)}$?

\noindent\rule{\linewidth}{0.4pt}


%% ============================================================
\section{Problem \#1074}
\label{prob:1074}

\noindent\textbf{Status:} OPEN \quad|\quad \textbf{Prize:} none \quad|\quad \textbf{Updated:} 2025-10-05

\noindent\textbf{Tags:} number theory

\medskip

\noindent\textbf{Statement:}

Let $S$ be the set of all $m\geq 1$ such that there exists a prime $p\not\equiv 1\pmod{m}$ such that $m!+1\equiv 0\pmod{p}$. Does\[\lim \frac{\lvert S\cap [1,x]\rvert}{x}\]exist? What is it? Similarly, if $P$ is the set of all primes $p$ such that there exists an $m$ with $p\not\equiv 1\pmod{m}$ such that $m!+1\equiv 0\pmod{p}$, then does\[\lim \frac{\lvert P\cap [1,x]\rvert}{\pi(x)}\]exist? What is it?

\noindent\rule{\linewidth}{0.4pt}


%% ============================================================
\section{Problem \#1075}
\label{prob:1075}

\noindent\textbf{Status:} OPEN \quad|\quad \textbf{Prize:} none \quad|\quad \textbf{Updated:} 2025-10-05

\noindent\textbf{Tags:} hypergraphs

\medskip

\noindent\textbf{Statement:}

Let $r\geq 3$. There exists $c_r&#62;r^{-r}$ such that, for any $\epsilon&#62;0$, if $n$ is sufficiently large, the following holds. Any $r$-uniform hypergraph on $n$ vertices with at least $(1+\epsilon)(n/r)^r$ many edges contains a subgraph on $m$ vertices with at least $c_rm^r$ edges, where $m=m(n)\to \infty$ as $n\to \infty$.

\noindent\rule{\linewidth}{0.4pt}


%% ============================================================
\section{Problem \#1083}
\label{prob:1083}

\noindent\textbf{Status:} OPEN \quad|\quad \textbf{Prize:} none \quad|\quad \textbf{Updated:} 2025-10-17

\noindent\textbf{Tags:} geometry, distances

\medskip

\noindent\textbf{Statement:}

Let $d\geq 3$, and let $f_d(n)$ be the minimal $m$ such that every set of $n$ points in $\mathbb{R}^d$ determines at least $m$ distinct distances. Estimate $f_d(n)$ - in particular, is it true that\[f_d(n)=n^{\frac{2}{d}-o(1)}?\]

\noindent\rule{\linewidth}{0.4pt}


%% ============================================================
\section{Problem \#1084}
\label{prob:1084}

\noindent\textbf{Status:} OPEN \quad|\quad \textbf{Prize:} none \quad|\quad \textbf{Updated:} 2025-10-17

\noindent\textbf{Tags:} geometry, distances

\medskip
\noindent\textit{contact number problem}


\noindent\textbf{Statement:}

Let $f_d(n)$ be minimal such that in any collection of $n$ points in $\mathbb{R}^d$, all of distance at least $1$ apart, there are at most $f_d(n)$ many pairs of points which are distance $1$ apart. Estimate $f_d(n)$.

\noindent\rule{\linewidth}{0.4pt}


%% ============================================================
\section{Problem \#1085}
\label{prob:1085}

\noindent\textbf{Status:} OPEN \quad|\quad \textbf{Prize:} none \quad|\quad \textbf{Updated:} 2025-10-17

\noindent\textbf{Tags:} geometry, distances

\medskip

\noindent\textbf{Statement:}

Let $f_d(n)$ be minimal such that, in any set of $n$ points in $\mathbb{R}^d$, there exist at most $f_d(n)$ pairs of points which distance $1$ apart. Estimate $f_d(n)$.

\noindent\rule{\linewidth}{0.4pt}


%% ============================================================
\section{Problem \#1086}
\label{prob:1086}

\noindent\textbf{Status:} OPEN \quad|\quad \textbf{Prize:} none \quad|\quad \textbf{Updated:} 2025-10-17

\noindent\textbf{Tags:} geometry, distances

\medskip

\noindent\textbf{Statement:}

Let $g(n)$ be minimal such that any set of $n$ points in $\mathbb{R}^2$ contains the vertices of at most $g(n)$ many triangles with the same area. Estimate $g(n)$.

\noindent\rule{\linewidth}{0.4pt}


%% ============================================================
\section{Problem \#1087}
\label{prob:1087}

\noindent\textbf{Status:} OPEN \quad|\quad \textbf{Prize:} none \quad|\quad \textbf{Updated:} 2025-10-17

\noindent\textbf{Tags:} geometry, distances

\medskip

\noindent\textbf{Statement:}

Let $f(n)$ be minimal such that every set of $n$ points in $\mathbb{R}^2$ contains at most $f(n)$ many sets of four points which are 'degenerate' in the sense that some pair are the same distance apart. Estimate $f(n)$ - in particular, is it true that $f(n)\leq n^{3+o(1)}$?

\noindent\rule{\linewidth}{0.4pt}


%% ============================================================
\section{Problem \#1088}
\label{prob:1088}

\noindent\textbf{Status:} OPEN \quad|\quad \textbf{Prize:} none \quad|\quad \textbf{Updated:} 2025-10-17

\noindent\textbf{Tags:} geometry

\medskip

\noindent\textbf{Statement:}

Let $f_d(n)$ be the minimal $m$ such that any set of $m$ points in $\mathbb{R}^d$ contains a set of $n$ points such that any two determined distances are distinct. Estimate $f_d(n)$. In particular, is it true that, for fixed $n\geq 3$,\[f_d(n)=2^{o(d)}?\]

\noindent\rule{\linewidth}{0.4pt}


%% ============================================================
\section{Problem \#1093}
\label{prob:1093}

\noindent\textbf{Status:} OPEN \quad|\quad \textbf{Prize:} none \quad|\quad \textbf{Updated:} 2025-10-18

\noindent\textbf{Tags:} number theory, binomial coefficients

\medskip

\noindent\textbf{Statement:}

For $n\geq 2k$ we define the deficiency of $\binom{n}{k}$ as follows. If $\binom{n}{k}$ is divisible by a prime $p\leq k$ then the deficiency is undefined. Otherwise, the deficiency is the number of $0\leq i&lt;k$ such that $n-i$ is $k$-smooth, that is, divisible only by primes $\leq k$. Are there infinitely many binomial coefficients with deficiency $1$? Are there only finitely many with deficiency $&gt;1$?

\noindent\rule{\linewidth}{0.4pt}


%% ============================================================
\section{Problem \#1094}
\label{prob:1094}

\noindent\textbf{Status:} OPEN \quad|\quad \textbf{Prize:} none \quad|\quad \textbf{Updated:} 2025-10-18

\noindent\textbf{Tags:} number theory, binomial coefficients

\medskip

\noindent\textbf{Statement:}

For all $n\geq 2k$ the least prime factor of $\binom{n}{k}$ is $\leq \max(n/k,k)$, with only finitely many exceptions.

\noindent\rule{\linewidth}{0.4pt}


%% ============================================================
\section{Problem \#1095}
\label{prob:1095}

\noindent\textbf{Status:} OPEN \quad|\quad \textbf{Prize:} none \quad|\quad \textbf{Updated:} 2025-10-18

\noindent\textbf{Tags:} number theory, binomial coefficients

\medskip

\noindent\textbf{Statement:}

Let $g(k)&#62;k+1$ be the smallest $n$ such that all prime factors of $\binom{n}{k}$ are $&#62;k$. Estimate $g(k)$.

\noindent\rule{\linewidth}{0.4pt}


%% ============================================================
\section{Problem \#1097}
\label{prob:1097}

\noindent\textbf{Status:} OPEN \quad|\quad \textbf{Prize:} none \quad|\quad \textbf{Updated:} 2025-10-18

\noindent\textbf{Tags:} number theory, additive combinatorics

\medskip

\noindent\textbf{Statement:}

Let $A$ be a set of $n$ integers. How many distinct $d$ can occur as the common difference of a three-term arithmetic progression in $A$? In particular, are there always $O(n^{3/2})$ many such $d$?

\noindent\rule{\linewidth}{0.4pt}


%% ============================================================
\section{Problem \#1100}
\label{prob:1100}

\noindent\textbf{Status:} OPEN \quad|\quad \textbf{Prize:} none \quad|\quad \textbf{Updated:} 2025-10-19

\noindent\textbf{Tags:} number theory, divisors

\medskip

\noindent\textbf{Statement:}

If $1=d_1&#60;\cdots&#60;d_{\tau(n)}=n$ are the divisors of $n$, then let $\tau_\perp(n)$ count the number of $i$ for which $(d_i,d_{i+1})=1$. Is it true that $\tau_\perp(n)/\omega(n)\to \infty$ for almost all $n$? Is it true that\[\tau_\perp(n)&#60; \exp((\log n)^{o(1)})\]for all $n$? Let\[g(k) = \max_{\omega(n)=k}\tau_\perp(n),\]where $\omega(n)$ counts the number of distinct prime divisors of $n$, and $n$ is restricted to squarefree integers. Determine the growth of $g(k)$.

\noindent\rule{\linewidth}{0.4pt}


%% ============================================================
\section{Problem \#1101}
\label{prob:1101}

\noindent\textbf{Status:} OPEN \quad|\quad \textbf{Prize:} none \quad|\quad \textbf{Updated:} 2025-10-19

\noindent\textbf{Tags:} number theory

\medskip

\noindent\textbf{Statement:}

If $u=\{u_1&lt;u_2&lt;\cdots\}$ is a sequence of integers such that $(u_i,u_j)=1$ for all $i\neq j$ and $\sum \frac{1}{u_i}&lt;\infty$ then let $\{a_1&lt;a_2&lt;\cdots\}$ be the sequence of integers which are not divisible by any of the $u_i$. For any $x$ define $t_x$ by\[u_1\cdots u_{t_x}\leq x&lt; u_1\cdots u_{t_x}u_{t_x+1}.\]We call such a sequence $u_i$ good if, for all $\epsilon&gt;0$, if $x$ is sufficiently large then\[\max_{a_k&#60;x} (a_{k+1}-a_k) &#60; (1+\epsilon)t_x \prod_{i}\left(1-\frac{1}{u_i}\right)^{-1}.\]Is there a good sequence such that $u_n&#60; n^{O(1)}$? Is there a good sequence such that $u_n\leq e^{o(n)}$?

\noindent\rule{\linewidth}{0.4pt}


%% ============================================================
\section{Problem \#1103}
\label{prob:1103}

\noindent\textbf{Status:} OPEN \quad|\quad \textbf{Prize:} none \quad|\quad \textbf{Updated:} 2025-10-19

\noindent\textbf{Tags:} number theory

\medskip

\noindent\textbf{Statement:}

Let $A$ be an infinite sequence of integers such that every $n\in A+A$ is squarefree. How fast must $A$ grow?

\noindent\rule{\linewidth}{0.4pt}


%% ============================================================
\section{Problem \#1104}
\label{prob:1104}

\noindent\textbf{Status:} OPEN \quad|\quad \textbf{Prize:} none \quad|\quad \textbf{Updated:} 2025-10-26

\noindent\textbf{Tags:} graph theory, chromatic number

\medskip

\noindent\textbf{Statement:}

Let $f(n)$ be the maximum possible chromatic number of a triangle-free graph on $n$ vertices. Estimate $f(n)$.

\noindent\rule{\linewidth}{0.4pt}


%% ============================================================
\section{Problem \#1106}
\label{prob:1106}

\noindent\textbf{Status:} OPEN \quad|\quad \textbf{Prize:} none \quad|\quad \textbf{Updated:} 2025-11-17

\noindent\textbf{Tags:} number theory

\medskip

\noindent\textbf{Statement:}

Let $p(n)$ denote the partition function of $n$ and let $F(n)$ count the number of distinct prime factors of\[\prod_{1\leq k\leq n}p(k).\]Does $F(n)\to \infty$ with $n$? Is $F(n)&gt;n$ for all sufficiently large $n$?

\noindent\rule{\linewidth}{0.4pt}


%% ============================================================
\section{Problem \#1107}
\label{prob:1107}

\noindent\textbf{Status:} OPEN \quad|\quad \textbf{Prize:} none \quad|\quad \textbf{Updated:} 2025-11-17

\noindent\textbf{Tags:} number theory, powerful

\medskip

\noindent\textbf{Statement:}

Let $r\geq 2$. A number $n$ is $r$-powerful if for every prime $p$ which divides $n$ we have $p^r\mid n$. Is every large integer the sum of at most $r+1$ many $r$-powerful numbers?

\noindent\rule{\linewidth}{0.4pt}


%% ============================================================
\section{Problem \#1108}
\label{prob:1108}

\noindent\textbf{Status:} OPEN \quad|\quad \textbf{Prize:} none \quad|\quad \textbf{Updated:} 2025-11-17

\noindent\textbf{Tags:} number theory, factorials

\medskip

\noindent\textbf{Statement:}

Let\[A = \left\{ \sum_{n\in S}n! : S\subset \mathbb{N}\textrm{ finite}\right\}.\]If $k\geq 2$, then does $A$ contain only finitely many $k$th powers? Does it contain only finitely many powerful numbers?

\noindent\rule{\linewidth}{0.4pt}


%% ============================================================
\section{Problem \#1109}
\label{prob:1109}

\noindent\textbf{Status:} OPEN \quad|\quad \textbf{Prize:} none \quad|\quad \textbf{Updated:} 2025-12-03

\noindent\textbf{Tags:} number theory

\medskip

\noindent\textbf{Statement:}

Let $f(N)$ be the size of the largest subset $A\subseteq \{1,\ldots,N\}$ such that every $n\in A+A$ is squarefree. Estimate $f(N)$. In particular, is it true that $f(N)\leq N^{o(1)}$, or even $f(N) \leq (\log N)^{O(1)}$?

\noindent\rule{\linewidth}{0.4pt}


%% ============================================================
\section{Problem \#1110}
\label{prob:1110}

\noindent\textbf{Status:} OPEN \quad|\quad \textbf{Prize:} none \quad|\quad \textbf{Updated:} 2025-12-07

\noindent\textbf{Tags:} number theory

\medskip

\noindent\textbf{Statement:}

Let $p&#62;q\geq 2$ be two coprime integers. We call $n$ representable if it is the sum of integers of the form $p^kq^l$, none of which divide each other. If $\{p,q\}\neq \{2,3\}$ then what can be said about the density of non-representable numbers? Are there infinitely many coprime non-representable numbers?

\noindent\rule{\linewidth}{0.4pt}


%% ============================================================
\section{Problem \#1111}
\label{prob:1111}

\noindent\textbf{Status:} OPEN \quad|\quad \textbf{Prize:} none \quad|\quad \textbf{Updated:} 2025-12-07

\noindent\textbf{Tags:} graph theory

\medskip

\noindent\textbf{Statement:}

If $G$ is a finite graph and $A,B$ are disjoint sets of vertices then we call $A,B$ anticomplete if there are no edges between $A$ and $B$. If $t,c\geq 1$ then there exists $d\geq 1$ such that if $\chi(G)\geq d$ and $\omega(G)&#60;t$ then there are anticomplete sets $A,B$ with $\chi(A)\geq \chi(B)\geq c$.

\noindent\rule{\linewidth}{0.4pt}


%% ============================================================
\section{Problem \#1112}
\label{prob:1112}

\noindent\textbf{Status:} OPEN \quad|\quad \textbf{Prize:} none \quad|\quad \textbf{Updated:} 2025-12-28

\noindent\textbf{Tags:} additive combinatorics

\medskip

\noindent\textbf{Statement:}

Let $1\leq d_1&#60;d_2$ and $k\geq 3$. Does there exist an integer $r$ such that if $B=\{b_1&#60;\cdots\}$ is a lacunary sequence of positive integers with $b_{i+1}\geq rb_i$ then there exists a sequence of positive integers $A=\{a_1&#60;\cdots\}$ such that\[d_1\leq a_{i+1}-a_i\leq d_2\]for all $i\geq 1$ and $(kA)\cap B=\emptyset$, where $kA$ is the $k$-fold sumset?

\noindent\rule{\linewidth}{0.4pt}


%% ============================================================
\section{Problem \#1113}
\label{prob:1113}

\noindent\textbf{Status:} OPEN \quad|\quad \textbf{Prize:} none \quad|\quad \textbf{Updated:} 2025-12-28

\noindent\textbf{Tags:} number theory, covering systems

\medskip
\noindent\textit{Sierpinski numbers}


\noindent\textbf{Statement:}

A positive odd integer $m$ such that none of $2^km+1$ are prime for $k\geq 0$ is called a Sierpinski number . We say that a set of primes $P$ is a covering set for $m$ if every $2^km+1$ is divisible by some $p\in P$. Are there Sierpinski numbers with no finite covering set of primes?

\noindent\rule{\linewidth}{0.4pt}


%% ============================================================
\section{Problem \#1117}
\label{prob:1117}

\noindent\textbf{Status:} OPEN \quad|\quad \textbf{Prize:} none \quad|\quad \textbf{Updated:} 2025-12-29

\noindent\textbf{Tags:} analysis

\medskip

\noindent\textbf{Statement:}

Let $f(z)$ be an entire function which is not a monomial. Let $\nu(r)$ count the number of $z$ with $\lvert z\rvert=r$ such that $\lvert f(z)\rvert=\max_{\lvert z\rvert=r}\lvert f(z)\rvert$. (This is a finite quantity if $f$ is not a monomial.) Is it possible for\[\limsup \nu(r)=\infty?\]Is it possible for\[\liminf \nu(r)=\infty?\]

\noindent\rule{\linewidth}{0.4pt}


%% ============================================================
\section{Problem \#1120}
\label{prob:1120}

\noindent\textbf{Status:} OPEN \quad|\quad \textbf{Prize:} none \quad|\quad \textbf{Updated:} 2025-12-29

\noindent\textbf{Tags:} analysis

\medskip

\noindent\textbf{Statement:}

Let $f\in \mathbb{C}[z]$ be a monic polynomial of degree $n$, all of whose roots satisfy $\lvert z\rvert\leq 1$. Let\[E= \{ z : \lvert f(z)\rvert \leq 1\}.\]What is the shortest length of a path in $E$ joining $z=0$ to $\lvert z\rvert =1$?

\noindent\rule{\linewidth}{0.4pt}


%% ============================================================
\section{Problem \#1122}
\label{prob:1122}

\noindent\textbf{Status:} OPEN \quad|\quad \textbf{Prize:} none \quad|\quad \textbf{Updated:} 2025-12-30

\noindent\textbf{Tags:} number theory

\medskip

\noindent\textbf{Statement:}

Let $f:\mathbb{N}\to \mathbb{R}$ be an additive function (i.e. $f(ab)=f(a)+f(b)$ whenever $(a,b)=1$). Let\[A=\{ n \geq 1: f(n+1)&#60; f(n)\}.\]If $\lvert A\cap [1,X]\rvert =o(X)$ then must $f(n)=c\log n$ for some $c\in \mathbb{R}$?

\noindent\rule{\linewidth}{0.4pt}


%% ============================================================
\section{Problem \#1131}
\label{prob:1131}

\noindent\textbf{Status:} OPEN \quad|\quad \textbf{Prize:} none \quad|\quad \textbf{Updated:} 2026-01-01

\noindent\textbf{Tags:} analysis, polynomials

\medskip

\noindent\textbf{Statement:}

For $x_1,\ldots,x_n\in [-1,1]$ let\[l_k(x)=\frac{\prod_{i\neq k}(x-x_i)}{\prod_{i\neq k}(x_k-x_i)},\]which are such that $l_k(x_k)=1$ and $l_k(x_i)=0$ for $i\neq k$. What is the minimal value of\[I(x_1,\ldots,x_n)=\int_{-1}^1 \sum_k \lvert l_k(x)\rvert^2\mathrm{d}x?\]In particular, is it true that\[\min I =2-(1+o(1))\frac{1}{n}?\]

\noindent\rule{\linewidth}{0.4pt}


%% ============================================================
\section{Problem \#1132}
\label{prob:1132}

\noindent\textbf{Status:} OPEN \quad|\quad \textbf{Prize:} none \quad|\quad \textbf{Updated:} 2026-01-01

\noindent\textbf{Tags:} analysis, polynomials

\medskip

\noindent\textbf{Statement:}

For $x_1,\ldots,x_n\in [-1,1]$ let\[l_k(x)=\frac{\prod_{i\neq k}(x-x_i)}{\prod_{i\neq k}(x_k-x_i)},\]which are such that $l_k(x_k)=1$ and $l_k(x_i)=0$ for $i\neq k$. Let $x_1,x_2,\ldots\in [-1,1]$ be an infinite sequence, and let\[L_n(x) = \sum_{1\leq k\leq n}\lvert l_k(x)\rvert,\]where each $l_k(x)$ is defined above with respect to $x_1,\ldots,x_n$. Must there exist $x\in (-1,1)$ such that\[L_n(x) &#62;\frac{2}{\pi}\log n-O(1)\]for infinitely many $n$? Is it true that\[\limsup_{n\to \infty}\frac{L_n(x)}{\log n}\geq \frac{2}{\pi}\]for almost all $x\in (-1,1)$?

\noindent\rule{\linewidth}{0.4pt}


%% ============================================================
\section{Problem \#1133}
\label{prob:1133}

\noindent\textbf{Status:} OPEN \quad|\quad \textbf{Prize:} none \quad|\quad \textbf{Updated:} 2026-01-01

\noindent\textbf{Tags:} analysis, polynomials

\medskip

\noindent\textbf{Statement:}

Let $C&#62;0$. There exists $\epsilon&#62;0$ such that if $n$ is sufficiently large the following holds. For any $x_1,\ldots,x_n\in [-1,1]$ there exist $y_1,\ldots,y_n\in [-1,1]$ such that, if $P$ is a polynomial of degree $m&lt;(1+\epsilon)n$ with $P(x_i)=y_i$ for at least $(1-\epsilon)n$ many $1\leq i\leq n$, then\[\max_{x\in [-1,1]}\lvert P(x)\rvert &gt;C.\]

\noindent\rule{\linewidth}{0.4pt}


%% ============================================================
\section{Problem \#1135}
\label{prob:1135}

\noindent\textbf{Status:} OPEN \quad|\quad \textbf{Prize:} \$500 \quad|\quad \textbf{Updated:} 2026-01-11

\noindent\textbf{Tags:} number theory

\medskip
\noindent\textit{Collatz conjecture}


\noindent\textbf{Statement:}

Define $f:\mathbb{N}\to \mathbb{N}$ by $f(n)=n/2$ if $n$ is even and $f(n)=\frac{3n+1}{2}$ if $n$ is odd. Given any integer $m\geq 1$ does there exist $k\geq 1$ such that $f^{(k)}(m)=1$?

\noindent\rule{\linewidth}{0.4pt}


%% ============================================================
\section{Problem \#1137}
\label{prob:1137}

\noindent\textbf{Status:} OPEN \quad|\quad \textbf{Prize:} none \quad|\quad \textbf{Updated:} 2026-01-23

\noindent\textbf{Tags:} number theory, primes

\medskip

\noindent\textbf{Statement:}

Let $d_n=p_{n+1}-p_n$, where $p_n$ denotes the $n$th prime. Is it true that\[\frac{\max_{n&#60;x}d_{n}d_{n-1}}{(\max_{n&#60;x}d_n)^2}\to 0\]as $x\to \infty$?

\noindent\rule{\linewidth}{0.4pt}


%% ============================================================
\section{Problem \#1139}
\label{prob:1139}

\noindent\textbf{Status:} OPEN \quad|\quad \textbf{Prize:} none \quad|\quad \textbf{Updated:} 2026-01-23

\noindent\textbf{Tags:} number theory, primes

\medskip

\noindent\textbf{Statement:}

Let $1\leq u_1&#60;u_2&#60;\cdots$ be the sequence of integers with at most $2$ prime factors. Is it true that\[\limsup \frac{u_{k+1}-u_k}{\log k}=\infty?\]

\noindent\rule{\linewidth}{0.4pt}


%% ============================================================
\section{Problem \#1142}
\label{prob:1142}

\noindent\textbf{Status:} OPEN \quad|\quad \textbf{Prize:} none \quad|\quad \textbf{Updated:} 2026-01-23

\noindent\textbf{Tags:} number theory, primes

\medskip

\noindent\textbf{Statement:}

Are there infinitely many $n$ (or any $n&#62;105$) such that $n-2^k$ is prime for all $1&#60;2^k&#60;n$?

\noindent\rule{\linewidth}{0.4pt}


%% ============================================================
\section{Problem \#1143}
\label{prob:1143}

\noindent\textbf{Status:} OPEN \quad|\quad \textbf{Prize:} none \quad|\quad \textbf{Updated:} 2026-01-23

\noindent\textbf{Tags:} number theory, primes

\medskip

\noindent\textbf{Statement:}

Let $p_1&lt;\cdots&lt;p_u$ be primes and let $k\geq 1$. Let $F_k(p_1,\ldots,p_u)$ be such that every interval of $k$ positive integers contains at least $F_k(p_1,\ldots,p_u)$ multiples of at least one of the $p_i$. Estimate $F_k(p_1,\ldots,p_u)$, particularly in the range $k=\alpha p_u$ for constant $\alpha&gt;2$.

\noindent\rule{\linewidth}{0.4pt}


%% ============================================================
\section{Problem \#1144}
\label{prob:1144}

\noindent\textbf{Status:} OPEN \quad|\quad \textbf{Prize:} none \quad|\quad \textbf{Updated:} 2026-01-23

\noindent\textbf{Tags:} number theory, probability

\medskip

\noindent\textbf{Statement:}

Let $f$ be a random completely multiplicative function, where for each prime $p$ we independently choose $f(p)\in \{-1,1\}$ uniformly at random. Is it true that\[\limsup_{N\to \infty}\frac{\sum_{m\leq N}f(m)}{\sqrt{N}}=\infty\]with probability $1$?

\noindent\rule{\linewidth}{0.4pt}


%% ============================================================
\section{Problem \#1145}
\label{prob:1145}

\noindent\textbf{Status:} OPEN \quad|\quad \textbf{Prize:} none \quad|\quad \textbf{Updated:} 2026-01-23

\noindent\textbf{Tags:} additive combinatorics, additive basis

\medskip

\noindent\textbf{Statement:}

Let $A=\{1\leq a_1&#60;a_2&#60;\cdots\}$ and $B=\{1\leq b_1&#60;b_2&#60;\cdots\}$ be sets of integers with $a_n/b_n\to 1$. If $A+B$ contains all sufficiently large positive integers then is it true that $\limsup 1_A\ast 1_B(n)=\infty$?

\noindent\rule{\linewidth}{0.4pt}


%% ============================================================
\section{Problem \#1146}
\label{prob:1146}

\noindent\textbf{Status:} OPEN \quad|\quad \textbf{Prize:} none \quad|\quad \textbf{Updated:} 2026-01-23

\noindent\textbf{Tags:} number theory

\medskip

\noindent\textbf{Statement:}

We say that $A\subset \mathbb{N}$ is an essential component if $d_s(A+B)&gt;d_s(B)$ for every $B\subset \mathbb{N}$ with $0&lt;d_s(B)&lt;1$ where $d_s$ is the Schnirelmann density . Is $B=\{2^m3^n : m,n\geq 0\}$ an essential component?

\noindent\rule{\linewidth}{0.4pt}


%% ============================================================
\section{Problem \#1150}
\label{prob:1150}

\noindent\textbf{Status:} OPEN \quad|\quad \textbf{Prize:} none \quad|\quad \textbf{Updated:} 2026-01-23

\noindent\textbf{Tags:} analysis, polynomials

\medskip

\noindent\textbf{Statement:}

Does there exist a constant $c&#62;0$ such that, for all large $n$ and all polynomials $P$ of degree $n$ with coefficients $\pm 1$,\[\max_{\lvert z\rvert=1}\lvert P(z)\rvert &#62; (1+c)\sqrt{n}?\]

\noindent\rule{\linewidth}{0.4pt}


%% ============================================================
\section{Problem \#1151}
\label{prob:1151}

\noindent\textbf{Status:} OPEN \quad|\quad \textbf{Prize:} none \quad|\quad \textbf{Updated:} 2026-01-23

\noindent\textbf{Tags:} analysis, polynomials

\medskip

\noindent\textbf{Statement:}

Given $a_1,\ldots,a_n\in [-1,1]$ let\[\mathcal{L}^nf(x) = \sum_{1\leq i\leq n}f(a_i)\ell_i(x)\]be the unique polynomial of degree $n-1$ which agrees with $f$ on $a_i$ for $1\leq i\leq n$ (that is, the Lagrange interpolation polynomial). Let $a_i$ be the set of Chebyshev nodes . Prove that, for any closed $A\subseteq [-1,1]$, there exists a continuous function $f$ such that $A$ is the set of limit points of $\mathcal{L}^nf(x)$.

\noindent\rule{\linewidth}{0.4pt}


%% ============================================================
\section{Problem \#1152}
\label{prob:1152}

\noindent\textbf{Status:} OPEN \quad|\quad \textbf{Prize:} none \quad|\quad \textbf{Updated:} 2026-01-23

\noindent\textbf{Tags:} analysis, polynomials

\medskip

\noindent\textbf{Statement:}

For $n\geq 1$ fix some sequence of $n$ distinct numbers $x_{1n},\ldots,x_{nn}\in [-1,1]$. Let $\epsilon=\epsilon(n)\to 0$. Does there always exist a continuous function $f:[-1,1]\to \mathbb{R}$ such that if $p_n$ is a sequence of polynomials, with degrees $\deg p_n&#60;(1+\epsilon(n))n$, such that $p_n(x_{kn})=f(x_{kn})$ for all $1\leq k\leq n$, then $p_n(x)\not\to f(x)$ for almost all $x\in [-1,1]$?

\noindent\rule{\linewidth}{0.4pt}


%% ============================================================
\section{Problem \#1155}
\label{prob:1155}

\noindent\textbf{Status:} OPEN \quad|\quad \textbf{Prize:} none \quad|\quad \textbf{Updated:} 2026-01-23

\noindent\textbf{Tags:} graph theory

\medskip

\noindent\textbf{Statement:}

Construct a random graph on $n$ vertices in the following way: begin with the complete graph $K_n$. At each stage, choose uniformly a random triangle in the graph and delete all the edges of this triangle. Repeat until the graph is triangle-free. Describe the typical parameters and structure of such a graph. In particular, if $f(n)$ is the number of edges remaining, then is it true that\[\mathbb{E}f(n)\asymp n^{3/2}\]and that $f(n) \ll n^{3/2}$ almost surely?

\noindent\rule{\linewidth}{0.4pt}


%% ============================================================
\section{Problem \#1156}
\label{prob:1156}

\noindent\textbf{Status:} OPEN \quad|\quad \textbf{Prize:} none \quad|\quad \textbf{Updated:} 2026-01-23

\noindent\textbf{Tags:} graph theory, chromatic number

\medskip

\noindent\textbf{Statement:}

Let $G$ be a random graph on $n$ vertices, in which every edge is included independently with probability $1/2$. Is there some constant $C$ such that that chromatic number $\chi(G)$ is, almost surely, concentrated on at most $C$ values? Is it true that, if $\omega(n)\to \infty$ sufficiently slowly, then for every function $f(n)$\[\mathbb{P}(\lvert\chi(G)-f(n)\rvert&#60;\omega(n))&#60;1/2\]if $n$ is sufficiently large?

\noindent\rule{\linewidth}{0.4pt}


%% ============================================================
\section{Problem \#1157}
\label{prob:1157}

\noindent\textbf{Status:} OPEN \quad|\quad \textbf{Prize:} none \quad|\quad \textbf{Updated:} 2026-01-23

\noindent\textbf{Tags:} hypergraphs, turan number

\medskip

\noindent\textbf{Statement:}

Let $t,k,r\geq 2$. Let $\mathcal{F}$ be the family of all $r$-uniform hypergraphs with $k$ vertices and $s$ edges. Determine\[\mathrm{ex}_r(n,\mathcal{F}).\]

\noindent\rule{\linewidth}{0.4pt}


%% ============================================================
\section{Problem \#1158}
\label{prob:1158}

\noindent\textbf{Status:} OPEN \quad|\quad \textbf{Prize:} none \quad|\quad \textbf{Updated:} 2026-01-23

\noindent\textbf{Tags:} hypergraphs, turan number

\medskip

\noindent\textbf{Statement:}

Let $K_{t}(r)$ be the complete $t$-partite $t$-uniform hypergraph with $r$ vertices in each class. Is it true that\[\mathrm{ex}_t(n,K_t(r)) \geq n^{t-r^{1-t}-o(1)}\]for all $t,r$?

\noindent\rule{\linewidth}{0.4pt}


%% ============================================================
\section{Problem \#1159}
\label{prob:1159}

\noindent\textbf{Status:} OPEN \quad|\quad \textbf{Prize:} none \quad|\quad \textbf{Updated:} 2026-01-23

\noindent\textbf{Tags:} combinatorics

\medskip

\noindent\textbf{Statement:}

Determine whether there exists a constant $C&gt;1$ such that the following holds. Let $P$ be a finite projective plane . Must there exist a set of points $S$ such that $1\leq \lvert S\cap \ell\rvert \leq C$ for all lines $\ell$?

\noindent\rule{\linewidth}{0.4pt}


%% ============================================================
\section{Problem \#1160}
\label{prob:1160}

\noindent\textbf{Status:} OPEN \quad|\quad \textbf{Prize:} none \quad|\quad \textbf{Updated:} 2026-01-23

\noindent\textbf{Tags:} group theory

\medskip

\noindent\textbf{Statement:}

Let $g(n)$ denote the number of groups of order $n$. If $n\leq 2^m$ then $g(n)\leq g(2^m)$.

\noindent\rule{\linewidth}{0.4pt}


%% ============================================================
\section{Problem \#1162}
\label{prob:1162}

\noindent\textbf{Status:} OPEN \quad|\quad \textbf{Prize:} none \quad|\quad \textbf{Updated:} 2026-01-23

\noindent\textbf{Tags:} group theory

\medskip

\noindent\textbf{Statement:}

Give an asymptotic formula for the number of subgroups of $S_n$. Is there a statistical theorem on their order?

\noindent\rule{\linewidth}{0.4pt}


%% ============================================================
\section{Problem \#1163}
\label{prob:1163}

\noindent\textbf{Status:} OPEN \quad|\quad \textbf{Prize:} none \quad|\quad \textbf{Updated:} 2026-01-23

\noindent\textbf{Tags:} group theory

\medskip
\noindent\textit{ambiguous statement}


\noindent\textbf{Statement:}

Describe (by statistical means) the arithmetic structure of the orders of subgroups of $S_n$.

\noindent\rule{\linewidth}{0.4pt}


%% ============================================================
\section{Problem \#1167}
\label{prob:1167}

\noindent\textbf{Status:} OPEN \quad|\quad \textbf{Prize:} none \quad|\quad \textbf{Updated:} 2026-01-23

\noindent\textbf{Tags:} set theory, probability

\medskip

\noindent\textbf{Statement:}

Let $r\geq 2$ be finite and $\lambda$ be an infinite cardinal. Let $\kappa_\alpha$ be cardinals for all $\alpha&#60;\gamma$. Is it true that\[2^\lambda \to (\kappa_\alpha+1)_{\alpha&#60;\gamma}^{r+1}\]implies\[\lambda \to (\kappa_\alpha)_{\alpha&#60;\gamma}^{r}?\]Here $+$ means cardinal addition, so that $\kappa_\alpha+1=\kappa_\alpha$ if $\kappa_\alpha$ is infinite.

\noindent\rule{\linewidth}{0.4pt}


%% ============================================================
\section{Problem \#1168}
\label{prob:1168}

\noindent\textbf{Status:} OPEN \quad|\quad \textbf{Prize:} none \quad|\quad \textbf{Updated:} 2026-01-23

\noindent\textbf{Tags:} set theory, ramsey theory

\medskip

\noindent\textbf{Statement:}

Prove that\[\aleph_{\omega+1}\not\to (\aleph_{\omega+1}, 3,\ldots,3)_{\aleph_0}^2\]without assuming the generalised continuum hypothesis.

\noindent\rule{\linewidth}{0.4pt}


%% ============================================================
\section{Problem \#1170}
\label{prob:1170}

\noindent\textbf{Status:} OPEN \quad|\quad \textbf{Prize:} none \quad|\quad \textbf{Updated:} 2026-01-23

\noindent\textbf{Tags:} set theory, ramsey theory

\medskip

\noindent\textbf{Statement:}

Is it consistent that\[\omega_2\to (\alpha)_2^2\]for every $\alpha &#60;\omega_2$?

\noindent\rule{\linewidth}{0.4pt}


%% ============================================================
\section{Problem \#1171}
\label{prob:1171}

\noindent\textbf{Status:} OPEN \quad|\quad \textbf{Prize:} none \quad|\quad \textbf{Updated:} 2026-01-23

\noindent\textbf{Tags:} set theory, ramsey theory

\medskip

\noindent\textbf{Statement:}

Is it true that, for all finite $k&#60;\omega$,\[\omega_1^2\to (\omega_1\omega, 3,\ldots,3)_{k+1}^2?\]

\noindent\rule{\linewidth}{0.4pt}


%% ============================================================
\section{Problem \#1172}
\label{prob:1172}

\noindent\textbf{Status:} OPEN \quad|\quad \textbf{Prize:} none \quad|\quad \textbf{Updated:} 2026-01-23

\noindent\textbf{Tags:} set theory, ramsey theory

\medskip

\noindent\textbf{Statement:}

Establish whether the following are true assuming the generalised continuum hypothesis:\[\omega_3 \to (\omega_2,\omega_1+2)^2,\]\[\omega_3\to (\omega_2+\omega_1,\omega_2+\omega)^2,\]\[\omega_2\to (\omega_1^{\omega+2}+2, \omega_1+2)^2.\]Establish whether the following is consistent with the generalised continuum hypothesis:\[\omega_2\to (\omega_1+\omega)_2^2,\]or even $\omega_2 \to (\xi)_2^2$ for all $\xi&#60;\omega_2$.

\noindent\rule{\linewidth}{0.4pt}


%% ============================================================
\section{Problem \#1173}
\label{prob:1173}

\noindent\textbf{Status:} OPEN \quad|\quad \textbf{Prize:} none \quad|\quad \textbf{Updated:} 2026-01-23

\noindent\textbf{Tags:} set theory, combinatorics

\medskip

\noindent\textbf{Statement:}

Assume the generalised continuum hypothesis. Let\[f: \omega_{\omega+1}\to [\omega_{\omega+1}]^{\leq \aleph_\omega}\]be a set mapping such that\[\lvert f(\alpha)\cap f(\beta)\rvert &#60;\aleph_\omega\]for all $\alpha\neq \beta$. Does there exist a free set of cardinality $\aleph_{\omega+1}$?

\noindent\rule{\linewidth}{0.4pt}


%% ============================================================
\section{Problem \#1175}
\label{prob:1175}

\noindent\textbf{Status:} OPEN \quad|\quad \textbf{Prize:} none \quad|\quad \textbf{Updated:} 2026-01-23

\noindent\textbf{Tags:} set theory, chromatic number

\medskip

\noindent\textbf{Statement:}

Let $\kappa$ be an uncountable cardinal. Must there exist a cardinal $\lambda$ such that every graph with chromatic number $\lambda$ contains a triangle-free subgraph with chromatic number $\kappa$?

\noindent\rule{\linewidth}{0.4pt}


%% ============================================================
\section{Problem \#1177}
\label{prob:1177}

\noindent\textbf{Status:} OPEN \quad|\quad \textbf{Prize:} none \quad|\quad \textbf{Updated:} 2026-01-23

\noindent\textbf{Tags:} set theory, chromatic number, hypergraphs

\medskip

\noindent\textbf{Statement:}

Let $G$ be a finite $3$-uniform hypergraph, and let $F_G(\kappa)$ denote the collection of $3$-uniform hypergraphs with chromatic number $\kappa$ not containing $G$. If $F_G(\aleph_1)$ is not empty then there exists $X\in F_G(\aleph_1)$ of cardinality at most $2^{2^{\aleph_0}}$. If both $F_G(\aleph_1)$ and $F_H(\aleph_1)$ are non-empty then $F_G(\aleph_1)\cap F_H(\aleph_1)$ is non-empty. If $\kappa,\lambda$ are uncountable cardinals and $F_G(\kappa)$ is non-empty then $F_G(\lambda)$ is non-empty.

\noindent\rule{\linewidth}{0.4pt}


%% ============================================================
\section{Problem \#1178}
\label{prob:1178}

\noindent\textbf{Status:} OPEN \quad|\quad \textbf{Prize:} none \quad|\quad \textbf{Updated:} 2026-01-25

\noindent\textbf{Tags:} graph theory, hypergraphs

\medskip

\noindent\textbf{Statement:}

For $r\geq 3$ let $d_r(e)$ be the minimal $d$ such that\[\mathrm{ex}_r(n,\mathcal{F})=o(n^2),\]where $\mathcal{F}$ is the family of $r$-uniform hypergraphs on $d$ vertices with $e$ edges. Prove that\[d_r(e)=(r-2)e+3\]for all $r,e\geq 3$.

\noindent\rule{\linewidth}{0.4pt}


%% ============================================================
\section{Problem \#1181}
\label{prob:1181}

\noindent\textbf{Status:} OPEN \quad|\quad \textbf{Prize:} none \quad|\quad \textbf{Updated:} 2026-03-07

\noindent\textbf{Tags:} number theory

\medskip

\noindent\textbf{Statement:}

Let $q(n,k)$ denote the least prime which does not divide $\prod_{1\leq i\leq k}(n+i)$. Is it true that there exists some $c&#62;0$ such that, for all large $n$,\[q(n,\log n)&#60;(1-c)(\log n)^2?\]

\noindent\rule{\linewidth}{0.4pt}


%% ============================================================
\section{Problem \#1182}
\label{prob:1182}

\noindent\textbf{Status:} OPEN \quad|\quad \textbf{Prize:} none \quad|\quad \textbf{Updated:} 2026-03-07

\noindent\textbf{Tags:} graph theory, ramsey theory

\medskip

\noindent\textbf{Statement:}

Let $f(n)$ be maximal such that there is a connected graph $G$ with $n$ vertices and $f(n)$ edges such that\[R(K_3,G)= 2n-1.\]Let $F(n)$ be maximal such that every connected graph $G$ with $n$ vertices and $\leq F(n)$ edges has\[R(K_3,G)= 2n-1.\]Estimate $f(n)$ and $F(n)$. In particular, is it true that $F(n)/n\to \infty$?

\noindent\rule{\linewidth}{0.4pt}


%% ============================================================
\section{Problem \#1183}
\label{prob:1183}

\noindent\textbf{Status:} OPEN \quad|\quad \textbf{Prize:} none \quad|\quad \textbf{Updated:} 2026-03-07

\noindent\textbf{Tags:} combinatorics, ramsey theory

\medskip

\noindent\textbf{Statement:}

Let $f(n)$ be maximal such that in any $2$-colouring of the subsets of $\{1,\ldots,n\}$ there is always a monochromatic family of at least $f(n)$ sets which is closed under taking unions and intersections. Estimate $f(n)$. Let $F(n)$ be defined similarly, except that we only require the family be closed under taking unions. Estimate $F(n)$. In particular, is it true that $F(n)\geq n^{\omega(n)}$ for some $\omega(n)\to \infty$ as $n\to \infty$, and $F(n)&#60;(1+o(1))^n$?

\noindent\rule{\linewidth}{0.4pt}


%% ============================================================
\section{Problem \#1184}
\label{prob:1184}

\noindent\textbf{Status:} OPEN \quad|\quad \textbf{Prize:} none \quad|\quad \textbf{Updated:} 2026-04-04

\noindent\textbf{Tags:} number theory, primes

\medskip

\noindent\textbf{Statement:}

Let $f(n,k)$ count the number of $1\leq i\leq k$ such that $P(n+i)&gt;k$ (where $P(m)$ is the largest prime divisor of $m$). Is it true that, if $\alpha&gt;1$ is such that $n=k^{\alpha+o(1)}$, then\[f(n,k)=(1-\rho(\alpha)+o(1))k,\]where $\rho$ is the Dickman function ?

\noindent\rule{\linewidth}{0.4pt}


%% ============================================================
\section{Problem \#1186}
\label{prob:1186}

\noindent\textbf{Status:} OPEN \quad|\quad \textbf{Prize:} none \quad|\quad \textbf{Updated:} 2026-04-04

\noindent\textbf{Tags:} additive combinatorics, arithmetic progressions

\medskip

\noindent\textbf{Statement:}

Let $\delta_k$ be such that in any $2$-colouring of $\{1,\ldots,n\}$ there exist at least $(\delta_k+o(1))n^2$ many monochromatic $k$-term arithmetic progressions. Give reasonable bounds (or even an asymptotic formula) for $\delta_k$.

\noindent\rule{\linewidth}{0.4pt}


%% ============================================================
\section{Problem \#1188}
\label{prob:1188}

\noindent\textbf{Status:} OPEN \quad|\quad \textbf{Prize:} none \quad|\quad \textbf{Updated:} 2026-04-04

\noindent\textbf{Tags:} number theory, covering systems

\medskip

\noindent\textbf{Statement:}

Call a set of distinct integers $1&#60;n_1&#60;\cdots&#60;n_k$ with associated congruence classes $a_i\pmod{n_i}$ a distinct covering system if every integer satisfies at least one of these congruences. A minimal distinct covering system is one such that no proper subset forms a covering system. Let $F(x)$ count the number of minimal distinct covering systems with all moduli in $[1,x]$. Estimate $F(x)$.

\noindent\rule{\linewidth}{0.4pt}


%% ============================================================
\section{Problem \#1189}
\label{prob:1189}

\noindent\textbf{Status:} OPEN \quad|\quad \textbf{Prize:} none \quad|\quad \textbf{Updated:} 2026-04-04

\noindent\textbf{Tags:} number theory, covering systems

\medskip

\noindent\textbf{Statement:}

Call a set of distinct integers $1&lt;n_1&lt;\cdots&lt;n_k$ a covering set if there is a choice of $a_i\pmod{n_i}$ for $1\leq i\leq k$ such that every integer satisfies at least one of these congruences. A set is an irreducible covering set if no proper subset is a covering set. How many irreducible covering sets of size $k$ are there? What is the minimum and maximum that $n_k$ can be? Determine or estimate $\max \sum\frac{1}{n_i}$, where the maximum ranges over all irreducible covering sets of size $k$. Are there infinitely many $n$ such that the divisors of $n$ (which are $&gt;1$) form an irreducible covering set?

\noindent\rule{\linewidth}{0.4pt}


%% ============================================================
\section{Problem \#1191}
\label{prob:1191}

\noindent\textbf{Status:} OPEN \quad|\quad \textbf{Prize:} \$1000 \quad|\quad \textbf{Updated:} 2026-04-04

\noindent\textbf{Tags:} additive combinatorics, sidon sets

\medskip

\noindent\textbf{Statement:}

Let $A\subset\mathbb{N}$ be an infinite Sidon set. Is it true that\[\liminf_{x\to \infty} \frac{\lvert A\cap [1,x]\rvert}{x^{1/2}}(\log x)^{1/2}=0?\]Does there exist an infinite Sidon set $A$ such that\[\liminf_{x\to \infty} \frac{\lvert A\cap [1,x]\rvert}{x^{1/2}}(\log x)^c&#62;0\]for some $c&#62;0$?

\noindent\rule{\linewidth}{0.4pt}


%% ============================================================
\section{Problem \#1192}
\label{prob:1192}

\noindent\textbf{Status:} OPEN \quad|\quad \textbf{Prize:} none \quad|\quad \textbf{Updated:} 2026-04-04

\noindent\textbf{Tags:} additive combinatorics, additive basis

\medskip

\noindent\textbf{Statement:}

For $A\subset \mathbb{N}$ let $f_r(n)$ count the number of solutions to $n=a_1+\cdots+a_r$ with $a_i\in A$. Does there exist, for all $r\geq 2$, a basis $A$ of order $r$ (so that $f_r(n)&#62;0$ for all large $n$) such that\[\sum_{n\leq x}f_r(n)^2 \ll x\]for all $x$?

\noindent\rule{\linewidth}{0.4pt}


%% ============================================================
\section{Problem \#1194}
\label{prob:1194}

\noindent\textbf{Status:} OPEN \quad|\quad \textbf{Prize:} none \quad|\quad \textbf{Updated:} 2026-04-04

\noindent\textbf{Tags:} additive combinatorics, additive basis, sidon sets

\medskip

\noindent\textbf{Statement:}

Let $A\subset\mathbb{N}$ be such that every integer $n\geq 1$ can be written uniquely as $a_n-b_n$ for some $a_n,b_n\in A$. How fast must $a_n/n$ increase?

\noindent\rule{\linewidth}{0.4pt}


%% ============================================================
\section{Problem \#1199}
\label{prob:1199}

\noindent\textbf{Status:} OPEN \quad|\quad \textbf{Prize:} none \quad|\quad \textbf{Updated:} 2026-04-04

\noindent\textbf{Tags:} additive combinatorics, ramsey theory

\medskip

\noindent\textbf{Statement:}

Is it true that in any $2$-colouring of $\mathbb{N}$ there exists an infinite set $A$ such that all elements of $A+A$ are the same colour?

\noindent\rule{\linewidth}{0.4pt}


%% ============================================================
\section{Problem \#1200}
\label{prob:1200}

\noindent\textbf{Status:} OPEN \quad|\quad \textbf{Prize:} none \quad|\quad \textbf{Updated:} 2026-04-04

\noindent\textbf{Tags:} number theory, primes

\medskip

\noindent\textbf{Statement:}

There exists a constant $C$ such that for all large $x$ there is a collection of primes $p_1&#60;\ldots&#60;p_k&#60;x$ with $\sum\frac{1}{p_i}&#60;C$ together with a system of congruences $a_i\pmod{p_i}$ such that every integer $n&#60;x$ satisfies at least one of these congruences.

\noindent\rule{\linewidth}{0.4pt}


%% ============================================================
\section{Problem \#1201}
\label{prob:1201}

\noindent\textbf{Status:} OPEN \quad|\quad \textbf{Prize:} none \quad|\quad \textbf{Updated:} 2026-04-04

\noindent\textbf{Tags:} number theory, primes

\medskip

\noindent\textbf{Statement:}

Is it true that for every $\epsilon,\eta&#62;0$ there exists a $k$ such that the density of $n$ for which\[P(n(n+1)\cdots(n+k))&#62;n^{1-\epsilon}\]is at least $1-\eta$ (where $P(m)$ is the greatest prime divisor of $m$)?

\noindent\rule{\linewidth}{0.4pt}


%% ============================================================
\section{Problem \#1203}
\label{prob:1203}

\noindent\textbf{Status:} OPEN \quad|\quad \textbf{Prize:} none \quad|\quad \textbf{Updated:} 2026-04-04

\noindent\textbf{Tags:} number theory

\medskip

\noindent\textbf{Statement:}

If $\omega(n)$ counts the number of distinct prime divisors of $n$ then let\[F(n)=\max_k \omega(n+k)\frac{\log\log k}{\log k}.\]Prove that $F(n)\to \infty$ as $n\to \infty$.

\noindent\rule{\linewidth}{0.4pt}


%% ============================================================
\section{Problem \#1204}
\label{prob:1204}

\noindent\textbf{Status:} OPEN \quad|\quad \textbf{Prize:} none \quad|\quad \textbf{Updated:} 2026-04-04

\noindent\textbf{Tags:} number theory

\medskip

\noindent\textbf{Statement:}

We call a sequence of integers $0\leq a_1&#60;\cdots &#60;a_k$ admissible if it is missing at least one congruence class modulo every prime $p$. Let $A(k)=\min a_k$. Estimate $A(k)$ - in particular, is it true that\[A(k)\sim k\log k?\]Estimate\[B(k)=\min \frac{a_1+\cdots+a_k}{k}.\]

\noindent\rule{\linewidth}{0.4pt}


%% ============================================================
\section{Problem \#1206}
\label{prob:1206}

\noindent\textbf{Status:} OPEN \quad|\quad \textbf{Prize:} none \quad|\quad \textbf{Updated:} 2026-04-04

\noindent\textbf{Tags:} number theory, sidon sets

\medskip

\noindent\textbf{Statement:}

Does $\{1,2^3,\ldots,N^3\}$ contain a Sidon set of size $\gg N$? Is there an infinite set $A\subset \mathbb{N}$ of positive density such that $\{a^3 : a\in A\}$ is a Sidon set?

\noindent\rule{\linewidth}{0.4pt}


%% ============================================================
\section{Problem \#1207}
\label{prob:1207}

\noindent\textbf{Status:} OPEN \quad|\quad \textbf{Prize:} none \quad|\quad \textbf{Updated:} 2026-04-04

\noindent\textbf{Tags:} geometry, distances

\medskip

\noindent\textbf{Statement:}

Let $P_d(n)$ be such that in any set of $n$ points in $\mathbb{R}^d$ there exist at least $P_d(n)$ many points which do not contain an isosceles triangle. Estimate $P_d(n)$ - in particular, is it true that\[P_2(n)&lt;n^{1-c}\]for some constant $c&gt;0$?

\noindent\rule{\linewidth}{0.4pt}


%% ============================================================
\section{Problem \#1208}
\label{prob:1208}

\noindent\textbf{Status:} OPEN \quad|\quad \textbf{Prize:} none \quad|\quad \textbf{Updated:} 2026-04-04

\noindent\textbf{Tags:} geometry, distances

\medskip

\noindent\textbf{Statement:}

For $d\geq 2$ let $F_d(n)$ be minimal such that every set of $n$ points in $\mathbb{R}^d$ contains a set of $F_d(n)$ points with distinct distances. Estimate $F_d(n)$ for fixed $d$ as $n\to \infty$.

\noindent\rule{\linewidth}{0.4pt}


%% ============================================================
\section{Problem \#1209}
\label{prob:1209}

\noindent\textbf{Status:} OPEN \quad|\quad \textbf{Prize:} none \quad|\quad \textbf{Updated:} 2026-04-04

\noindent\textbf{Tags:} number theory

\medskip

\noindent\textbf{Statement:}

Let $A=\{a_1&#60;a_2&#60;\cdots\}$ be a sequence of integers which tends to infinity sufficiently fast. If there is an $n$ such that all $n+a_k$ are primes then must there exist infinitely many such $n$? What if we ask for $n+a_k$ to be squarefree instead of prime? Are there $n$ such that $n+2^{2^k}$ is always a prime (or always squarefree, or infinitely often a prime, or infinitely often squarefree)?

\noindent\rule{\linewidth}{0.4pt}


%% ============================================================
\section{Problem \#1210}
\label{prob:1210}

\noindent\textbf{Status:} OPEN \quad|\quad \textbf{Prize:} none \quad|\quad \textbf{Updated:} 2026-04-04

\noindent\textbf{Tags:} number theory

\medskip

\noindent\textbf{Statement:}

Let $A\subseteq [1,n)$ be a set of integers such that $(a,b)=1$ for all distinct $a,b\in A$. Is it true that\[\sum_{a\in A}\frac{1}{n-a}\leq \sum_{p&#60;n}\frac{1}{p}+O(1)?\]

\noindent\rule{\linewidth}{0.4pt}


%% ============================================================
\section{Problem \#1212}
\label{prob:1212}

\noindent\textbf{Status:} OPEN \quad|\quad \textbf{Prize:} none \quad|\quad \textbf{Updated:} 2026-04-04

\noindent\textbf{Tags:} number theory, primes

\medskip

\noindent\textbf{Statement:}

Let $G$ be the graph with vertex set those pairs $(x,y)\in \mathbb{N}^2$ with $\mathrm{gcd}(x,y)=1$, in which we join two vertices if the differ in only one coordinate, and there by $\pm 1$. Is there a path going to infinity on $G$, say $P$, such that for all $(x,y)\in P$ both $\min(x,y)&#62;1$ and at least one of $x$ or $y$ is composite?

\noindent\rule{\linewidth}{0.4pt}

\end{document}
